%%%%%%%% ICML 2025 EXAMPLE LATEX SUBMISSION FILE %%%%%%%%%%%%%%%%%

\documentclass{article}

% Recommended, but optional, packages for figures and better typesetting:
\usepackage{microtype}
\usepackage{graphicx}
\usepackage{subcaption}
\usepackage{booktabs} % for professional tables
% \usepackage{subcaption}
\usepackage{multirow}

% hyperref makes hyperlinks in the resulting PDF.
% If your build breaks (sometimes temporarily if a hyperlink spans a page)
% please comment out the following usepackage line and replace
% \usepackage{icml2025} with \usepackage[nohyperref]{icml2025} above.
\usepackage{hyperref}


% Attempt to make hyperref and algorithmic work together better:
\newcommand{\theHalgorithm}{\arabic{algorithm}}

% Use the following line for the initial blind version submitted for review:
\usepackage{icml2025}

% If accepted, instead use the following line for the camera-ready submission:
% \usepackage[accepted]{icml2025}

% For theorems and such
\usepackage{amsmath}
\usepackage{amssymb}
\usepackage{mathtools}
\usepackage{amsthm}

% if you use cleveref..
\usepackage[capitalize,noabbrev]{cleveref}

%%%%%%%%%%%%%%%%%%%%%%%%%%%%%%%%
% THEOREMS
%%%%%%%%%%%%%%%%%%%%%%%%%%%%%%%%
\theoremstyle{plain}
\newtheorem{theorem}{Theorem}[section]
\newtheorem{proposition}[theorem]{Proposition}
\newtheorem{lemma}[theorem]{Lemma}
\newtheorem{corollary}[theorem]{Corollary}
\theoremstyle{definition}
\newtheorem{definition}[theorem]{Definition}
\newtheorem{assumption}[theorem]{Assumption}
\theoremstyle{remark}
\newtheorem{remark}[theorem]{Remark}

\usepackage{caption}

% Todonotes is useful during development; simply uncomment the next line
%    and comment out the line below the next line to turn off comments
%\usepackage[disable,textsize=tiny]{todonotes}
\usepackage[textsize=tiny]{todonotes}

%\newcommand{\SM}[1]{{\color{orange}{\bf SM:} #1}}
%\newcommand{\YX}[1]{{\color{brown}{\bf YX:} #1}}
\newcommand{\JW}[1]{{\color{magenta}{\bf JW:} #1}}

% The \icmltitle you define below is probably too long as a header.
% Therefore, a short form for the running title is supplied here:
\icmltitlerunning{Adapter Subspace Sharpness Matters for LoRA-Based Continual Learning: A PAC-Bayesian Perspective}

\begin{document}

\twocolumn[
\icmltitle{
Investigating Sharpness in Low-rank Subspace for Continual Learning: A PAC-Bayes Perspective}
% 
% Adapter Subspace Sharpness Matters for LoRA-Based Continual Learning: A PAC-Bayesian Perspective}
% Revisit Sharpness in the view of PAC-Bayesian Bound under Lora-based Continual Learning}

% 
% It is OKAY to include author information, even for blind
% submissions: the style file will automatically remove it for you
% unless you've provided the [accepted] option to the icml2025
% package.

% List of affiliations: The first argument should be a (short)
% identifier you will use later to specify author affiliations
% Academic affiliations should list Department, University, City, Region, Country
% Industry affiliations should list Company, City, Region, Country

% You can specify symbols, otherwise they are numbered in order.
% Ideally, you should not use this facility. Affiliations will be numbered
% in order of appearance and this is the preferred way.
% \icmlsetsymbol{equal}{*}

\begin{icmlauthorlist}
% % \icmlauthor{Firstname1 Lastname1}{equal,yyy}
% % \icmlauthor{Firstname2 Lastname2}{equal,yyy,comp}
% % \icmlauthor{Firstname3 Lastname3}{comp}
% % \icmlauthor{Firstname4 Lastname4}{sch}
% % \icmlauthor{Firstname5 Lastname5}{yyy}
% % \icmlauthor{Firstname6 Lastname6}{sch,yyy,comp}
% % \icmlauthor{Firstname7 Lastname7}{comp}
% % %\icmlauthor{}{sch}
% % \icmlauthor{Firstname8 Lastname8}{sch}
% % \icmlauthor{Firstname8 Lastname8}{yyy,comp}
% % %\icmlauthor{}{sch}
% % %\icmlauthor{}{sch}
\end{icmlauthorlist}

% % \icmlaffiliation{yyy}{Department of XXX, University of YYY, Location, Country}
% % \icmlaffiliation{comp}{Company Name, Location, Country}
% % \icmlaffiliation{sch}{School of ZZZ, Institute of WWW, Location, Country}

\icmlcorrespondingauthor{Firstname1 Lastname1}{first1.last1@xxx.edu}
\icmlcorrespondingauthor{Firstname2 Lastname2}{first2.last2@www.uk}

% % % You may provide any keywords that you
% % % find helpful for describing your paper; these are used to populate
% % % the "keywords" metadata in the PDF but will not be shown in the document
% % \icmlkeywords{Machine Learning, ICML}

% % \vskip 0.3in
]

% this must go after the closing bracket ] following \twocolumn[ ...

% This command actually creates the footnote in the first column
% listing the affiliations and the copyright notice.
% The command takes one argument, which is text to display at the start of the footnote.
% The \icmlEqualContribution command is standard text for equal contribution.
% Remove it (just {}) if you do not need this facility.

\printAffiliationsAndNotice{}
%\printAffiliationsAndNotice{}  % leave blank if no need to mention equal contribution
% \printAffiliationsAndNotice{\icmlEqualContribution} % otherwise use the standard text.

\begin{abstract}
Low-Rank Adaptation (LoRA) offers a scalable route to Parameter Efficient Continual Learning (PECL) with large pre-trained models, yet shared adapter subspaces make it difficult to balance stability and plasticity. While training toward flatter solutions is known to improve generalization in full-model learning, its role {inside the adapter subspace} for CL has not been theoretically established. We address this gap by proposing \emph{adapter subspace sharpness} and deriving a sequential hierarchical PAC-Bayesian  bound that augments empirical risk with an adapter-subspace sharpness term and a hierarchical KL term that tracks posterior complexity across tasks. Guided by the bound, we study perturbation design in LoRA-based CL through three coupled questions: \emph{where} to perturb, \emph{how} to perturb, and how these choices interact with LoRA organization over the task sequence. Using a unified protocol across LoRA variants and perturbation type, we find that controlling sharpness within the adapter subspace improves accuracy over SGD and over full-parameter  perturbations in the  LoRA-Based CL regime. Our results provide a theory-to-practice account and design guidance for PECL.

% framework in which the generalization
% that augments empirical risk with an {adapter subspace sharpness} term and a {hierarchical KL} term reflecting posterior complexity.
 
% Guidedby the bound, we organize our study around three research questions on (i) perturbation scope, (ii) perturbation type, and (iii) interaction with LoRA organization. 


% Parameter efficient continual learning (PECL) adapts pretrained models to task sequences by training LoRA adapters while freezing the backbone. This concentrates cross task interference in the adapter update subspace and raises a geometric question: is flatness control in the full parameter space necessary, or is adapter subspace flatness sufficient?
% Parameter-efficient continual learning (PECL) adapts a pre-trained model to a sequence of tasks by freezing the backbone and updating lightweight modules such as LoRA adapters. 
% In this setting, extending sharpness-aware optimization is not straightforward: perturbations in the full parameter space necessarily include frozen backbone parameters that cannot be adjusted by subsequent tasks. 
% We thus focus on \emph{adapter-subspace sharpness}, defined as sensitivity to perturbations restricted to the adapter parameters that are allowed to be updated at the current task, explicitly excluding frozen components.
% % We therefore propose to measure and control flatness only within the trainable update region that remains reachable during continual adaptation, which we call update-region flatness. 
% We develop a sequential hierarchical PAC-Bayes analysis and show that, under the frozen-backbone constraint, both the sharpness-related term and the continual complexity reduce to quantities defined on task-admissible update directions, providing a principled resolution of the scope mismatch between PECL restricted and sharpness difinition.
% Guided by this perspective, we systematically study perturbation scope and perturbation type within the admissible region, together with LoRA organization across tasks, and derive  practical guidelines linking reduced curvature exposure and reduced cross-task overlap of admissible directions to improved PECL performance and stability.
% PECL updates only LoRA adapters while freezing the backbone, so continual adaptation occurs in a constrained update subspace. However, most flatness-based generalization arguments are phrased in the full parameter space and motivate full-space perturbations, which conflicts with the fixed-backbone regime. This tension leads to a concrete question: in PECL, is global flatness control necessary, or does adapter-subspace flatness suffice?
% We develop a sequential hierarchical PAC Bayes analysis for PECL and derive a bound with an explicit hyperposterior drift term that captures cross task interference. Under a frozen backbone and adapter supported posteriors, both the KL complexity and the sharpness dependent empirical term reduce exactly to adapter subspace quantities, identifying adapter flatness as the relevant notion of flatness in PECL. Experiments on continual learning benchmarks provide consistent evidence that sharpness aware optimization restricted to adapters improves performance, whereas full space perturbations that also affect the frozen backbone do not yield consistent gains and can degrade stability.


% Experiments on continual learning benchmarks suggest that enforcing flatness within the adapter subspace is a reliable strategy in PECL, and that extending perturbations into frozen backbone parameters is not necessary for the gains observed in our settings.

% Experiments on continual learning benchmarks confirm that adapter only sharpness aware optimization yields consistent gains, while full space perturbations that include the frozen backbone provide no systematic benefit and often harm stability.

% Parameter-efficient continual learning (PECL) adapts large pre-trained models to evolving task sequences by tuning lightweight adapters while freezing the backbone.
% In this setup, cross-task interference is concentrated within the shared adapter subspace. Conventional flatness-based approaches, such as Flat-LoRA, advocate for full-parameter perturbation to achieve global flatness—a requirement at odds with the fixed-backbone design of PECL.
% % In this setting, all cross-task interference is funneled into a shared trainable adapter subspace, traditional viewpoints (such as Flat-LoRA) suggest that full-parameter pertur-
% % bation is needed to pursue global flatness. This view conflicts with the core design of PECL, where backbone weights are intentionally kept fixed.
% To resolve this, we investigate whether flatness restricted to the adapter subspace alone can effectively balance stability and plasticity in PECL. Theoretically, we derive a PAC-Bayesian generalization bound for PECL under a frozen backbone and Gaussian posteriors over the adapter parameters. In this formulation, both the sharpness term and the KL divergence depend solely on the adapter subspace, establishing adapter flatness as the pertinent measure of flatness in PECL.
% Experiments on benchmarks demonstrate that sharpness-aware optimization confined to the adapter subspace consistently enhances performance, whereas full-parameter perturbations yield no systematic improvement. These findings provide a principled foundation for pursuing subspace-restricted flatness in PECL and inform geometry-aware designs for LoRA-based continual learning methods.

% We address this tension by asking whether flatness confined to the adapter subspace is sufficient to govern stability and plasticity. First, we derive a PAC-Bayesian generalization bound specialized to PECL with a frozen backbone and Gaussian posteriors supported on the adapter subspace. In this regime, both the sharpness term and the KL complexity depend only on the adapter subspace, which identifies adapter flatness as the relevant notion of flatness for PECL. 
% % We then relate factor-space flatness to the spectrum of a restricted Fisher information matrix. Finally, 
% Experiments on CL benchmarks show that subspace-restricted sharpness-aware optimization consistently improves continual learning performance, while full-parameter perturbations fail to provide systematic benefits. This provides a principled justification for flatness oriented optimisation confined to the adapter subspace and offers a conceptual basis for geometry aware designs of LoRA based continual learning methods. 


\end{abstract}


\section{Introduction}

%, as newly acquired knowledge may interfere with previously learned representations
%, which means the degradation of generalization on previously learned tasks after learning new ones
% Continual learning (CL) has become increasingly relevant for adapting large-scale pre-trained models to evolving tasks in dynamic environments~\cite{PARISI201954,10444954,a2923225ea7d94e6c86faff97e23c69be,10.1145/3735633}. 
Adapting large-scale pre-trained models to new tasks is now routine in modern machine learning~\cite{HAN2021225,bommasani2022opportunitiesrisksfoundationmodels}.
Continual learning (CL) studies sequential task learning while maintaining generalization on previously learned tasks~\cite{wangComprehensiveSurveyContinual2024,yangRecentAdvancesFoundation2024,shiContinualLearningLarge2024,zhouContinualLearningPreTrained2024}. Catastrophic forgetting is the deterioration of this past-task generalization, measured as increased test risk on earlier task distributions after learning new tasks~\cite{pmlr-v202-lin23f}.
% Continual learning (CL) studies tasks sequentially while maintaining generalization on previously learned tasks, where catastrophic forgetting refers to the degradation of generalization of this past-task generalization after learning new tasks.
% Continual learning (CL) studies catastrophic forgetting, which can be viewed as the degradation of generalization on previously learned tasks after learning new ones.
% how to 
% how to incorporate new tasks while preserving previously acquired knowledge, where catastrophic forgetting can be viewed as the degradation of generalization on previously learned tasks after learning new ones.
To reduce the computational and storage cost of full-model updates, parameter-efficient methods such as Low-Rank Adaptation (LoRA)\cite{hu2022lora} and Adapter\cite{pmlr-v97-houlsby19a} introduce lightweight task-specific modules while freezing the backbone weights~\cite{dingParameterefficientFinetuningLargescale2023,He_2025_CVPR,NEURIPS2024_b0bc711f,hanParameterEfficientFineTuningLarge2024}. Combining parameter-efficient adaptation with continual learning objectives yields parameter-efficient continual learning (PECL), which enables scalable sequential adaptation under restricted update subspaces but can still suffer from catastrophic forgetting~\cite{gaoUnifiedContinualLearning2023,qiaoLearnMoreBother2024}.




% Continual learning (CL) studies how generalization evolves over a task sequence, where catastrophic forgetting refers to the degradation of generalization on previously learned tasks after learning new ones~\cite{PARISI201954,10444954,a2923225ea7d94e6c86faff97e23c69be,10.1145/3735633}. While full fine-tuning remains the most expressive adaptation strategy, its application in sequential settings is computationally demanding and susceptible to catastrophic forgetting.~\cite{howard2018universallanguagemodelfinetuning,doi:10.1073/pnas.1611835114,pmlr-v234-zhai24a,luo2025empiricalstudycatastrophicforgetting}.
% To reduce the computational and storage cost of full-model updates, parameter efficient methods such as Low-Rank Adaptation (LoRA)~\cite{hu2022lora} introduce lightweight, task-specific modules while freezing the backbone weights~\cite{pmlr-v97-houlsby19a,dingParameterefficientFinetuningLargescale2023,He_2025_CVPR,NEURIPS2024_b0bc711f}. Combining parameter-efficient updates with continual learning objectives leads to parameter-efficient continual learning (PECL), which enables continual adaptation of large pre-trained models and recuce catastro, but remain susceptible to catastrophic forgetting~\cite{mccloskey1989catastrophic}.

% Adapting large-scale pre-trained models to new data is now routine in modern machine learning. Due to the scale of contemporary models, this adaptation is typically performed using parameter-efficient transfer learning methods~\cite{han2024parameter} such as low-rank adaptation (LoRA)~\cite{hu2022lora}, which restrict task-specific updates to low-rank adapter subspaces. However, parameter-efficient updates remain susceptible to catastrophic forgetting~\cite{mccloskey1989catastrophic}, as updates for new data can interfere with previously acquired knowledge.



% % Although PECL enables scalable adaptation, 
% However, despite its scalability, PECL remain susceptible to catastrophic forgetting~\cite{mccloskey1989catastrophic}.


% does not preclude continual-learning pathologies, parameter sharing can still induce catastrophic forgetting.

% However, parameter-efficient updates remain susceptible to catastrophic forgetting~\cite{mccloskey1989catastrophic}, as updates for new data can interfere with previously acquired knowledge.

% constrain generalization across the task sequence
% can still induce task interference and limit generalization.

% Parameter-efficient tuning methods such as Low-Rank Adaptation (LoRA)~\cite{hu2022lora}, address this challenge by introducing lightweight, task-specific modules while freezing the backbone weights~\cite{pmlr-v97-houlsby19a,dingParameterefficientFinetuningLargescale2023}. 


% Full finetuning across sequential tasks is not only computationally expensive, but also prone to catastrophic forgetting and degraded generalization.




% Recent studies suggest that flat-minima are often associated with improved generalization and robustness, as they are less sensitive to parameter perturbations and distribution shifts~\cite{wu2017understandinggeneralizationdeeplearning,chaudhariEntropySGDBiasingGradient2019,NEURIPS2020_518a38cc,NEURIPS2021_995f5e03,JMLRAnEmpiricalInvestigatio,yangRevisitingFlatnessAwareOptimization2025}.
% Recent studies have associated 
% flat minima with
% , as flat solutions are less sensitive to parameter perturbations and distribution shifts
Flat minima are widely associated with improved generalization and robustness~\cite{NIPS1994_01882513,wu2017understandinggeneralizationdeeplearning,chaudhariEntropySGDBiasingGradient2019,jiangFantasticGeneralizationMeasures2019,tsuzukuNormalizedFlatMinima2020,NEURIPS2021_995f5e03}, motivating Sharpness-aware minimization (SAM)~\cite{foretSharpnessAwareMinimizationEfficiently2021} and related methods~\citep{pmlr-v151-bisla22a,zhangGradientNormAware2023,NEURIPS2024_C-Flat,yangRevisitingFlatnessAwareOptimization2025,chen2026sharpnessflatnessdecompositionframework} to control sharpness under local perturbations in the full parameter space. Beyond single task setting, such flatness control has been argued to mitigate forgetting in CL~\cite{NEURIPS2020_518a38cc,JMLRAnEmpiricalInvestigatio,shiOvercomingCatastrophicForgetting2021,liRevisitingCatastrophicForgetting2024}.
% Beyond generalization, such flatness control has been argued to improve stability under sequential updates and help mitigate forgetting\cite{}.
% Sharpness-aware minimization (SAM)~\cite{foretSharpnessAwareMinimizationEfficiently2021} and related methods~\citep{pmlr-v151-bisla22a,zhangGradientNormAware2023,NEURIPS2024_C-Flat} formalize this connection by defining sharpness through the loss elevation in a local perturbation neighborhood and optimizing parameters to reduce the elevation with perturbation, thereby biasing training toward flatter solutions in full-model learning. 
% have demonstrated the benefits of encouraging flat minima  (i.e., promoting flatness or reducing sharpness; we use sharpness as a unified term to measure sensitivity to perturbations) in full-model training by optimize paramter with perturbation. 
% However these methods have not yet been extended to parameter-efficient updating techniques, where updates are confined to a low-rank, parameter-efficient region.
% In PECL, it remains unclear whether reducing sharpness within the adapter subspace is sufficient to ensure robust generalization across tasks.
In PECL, however, feasible adaptation is restricted to task-admissible update directions induced by the trainable adapters, while standard sharpness definitions are stated over full-parameter neighborhoods. This creates a scope mismatch between the perturbation geometry and the feasible update dynamics.
% However, these methods are not directly applicable to LoRA-based PECL, where adaptation over a task sequence is confined to a low-dimensional trainable adapter subspace.
% This restriction induces a scope mismatch: sharpness is typically defined via perturbations in the full parameter, while the learning dynamics can only evolve along adapter-induced update directions.
% while the learner can only move along the adapter-induced update directions.


Consequently, a central question in PECL is whether controlling sharpness only within the task-admissible update subspace is sufficient to ensure robust generalization across tasks.  This tension naturally motivates two PECL-specific design: \emph{where} perturbations should be applied and \emph{how} they should be instantiated, as illustrated in Figure~\ref{fig:scope_type1}.
Even in the single-task setting, prior work remains divided on perturbation scope, advocating either full-parameter perturbations or adapter-restricted perturbations~\citep{NEURIPS2024_4eb2c0ad,li2025flatloralowrankadaptationflat}.
% This tension motivates two PECL-specific design questions: \emph{where} perturbations should be applied and \emph{how} they should be instantiated as shown in Figure~\ref{fig:scope_type1}.
% More importantly, theoretical understanding of this question remains limited. 
% Despite its practical relevance, the above question lacks a principled theoretical explanation.
Despite its practical relevance, this question has not been characterized by a theory that links PECL-specific constraints to generalization.
The PAC-Bayesian framework~\cite{mcallesterPACBayesianStochasticModel2003,PAC-Bayesiansupervisedclassification2007,pentinaPACBayesianBoundLifelong2014,germainPACBayesianTheoryMeets2016,lotfiPACBayesCompressionBounds2022} provides a principled route to generalization guaranties via posterior complexity. However existing analyzes rarely address sequential continual learning and typically do not incorporate sharpness effects.
% provides a principled lens for analyzing generalization through posterior complexity, most existing analyses are not formulated for continual learning and do not incorporate the effects of sharpness.
% analyses rarely account for the role of sharpness or the adapter subspace constraints. 
% In particular, in PECL it remains poorly understood how controlling sharpness within the trainable adapter subspace effect with generalization across a sequence of tasks.
In particular, the interaction between adapter update subspace sharpness and generalization across tasks remains poorly understood.


% Even in single-task, existing work disagrees on the perturbation scope, contrasting full-parameter perturbations with perturbations restricted to the adapter~\citep{NEURIPS2024_4eb2c0ad,li2025flatloralowrankadaptationflat}.

% LoRA fine-tuning trainable
% Notably, even in single-task LoRA fine-tuning there is an unresolved debate on the appropriate perturbation scope: LoRA-SAM~\citep{NEURIPS2024_4eb2c0ad} advocates optimizing sharpness within the LoRA coordinates, whereas Flat-LoRA~\citep{li2025flatloralowrankadaptationflat} argues that sufficient flatness requires perturbing all model parameters. Yet how this scope choice interacts with the continual-learning trajectory and the stability–plasticity constraints in LoRA-based CL remains underexplored. 
% Consequently, a central open question in PECL is whether controlling sharpness restricted to the trainable adapter subspace is sufficient to achieve robust generalization over a task sequence. 
% Figure~\ref{fig:scope_type_main} makes this tension explicit and highlights two design questions that are specific to PECL: where perturbations should be defined (full space versus the task-admissible adapter subspace), and how the perturbation objective should be instantiated (randomized versus adversarial, zeroth versus first order).



% Consequently, a central open question in PECL is whether controlling sharpness only within the trainable adapter subspace suffices to yield cross task generalization.

% \begin{figure*}[thbp]
% \centering
%     % \begin{subfigure}{0.46\linewidth}
%     %     \centering
% \includegraphics[width=0.75\linewidth]{figures_TRML/shiyituxiao_forcut_re4.pdf}
%         % \subcaption{Perturbation Scope}
%         \label{fig: pertuabtion scope}
% %     \end{subfigure}
% %         \begin{subfigure}{0.455\linewidth}
% %         \centering
% % \includegraphics[width=1\linewidth]{figures_TRML/shiyituxiao_forcut_re2.pdf}
% %         % \subcaption{Perturbation type}
% %     \end{subfigure}
% %     \centering
%     \vspace{-8pt}
%     \caption{Perturbation design in frozen-backbone PECL. \textbf{(a) Perturbation scope:} Full-parameter sharpness introduces perturbations to both the frozen backbone parameters and the adapter parameters, whereas {adapter subspace sharpness} restricts perturbations to the trainable adapter coordinates at the current task.
%  % \textbf{(a) Perturbation range:} full-parameter sharpness inject perturbations into both the frozen backbone and adapters, whereas adapter sharpness restricted  perturbationsto the adapter coordinates that are trainable at the current task. 
% \textbf{(b) Perturbation type:} schematic comparison of randomized zeroth-order sharpness $\mathrm{RS}^{(0)}_S(W,\Sigma)$, adversarial zeroth-order sharpness $\mathrm{AS}^{(0)}_S(W,\rho)$, and adversarial first-order sharpness $\mathrm{AS}^{(1)}_S(W,\rho)$, illustrating their relative strength under comparable perturbation budgets.}
% \end{figure*}


{
\begin{figure*}[htbp]
    \centering
    \makebox[\linewidth]{%
    \begin{subfigure}[t]{0.42\linewidth}
        \raggedright
\includegraphics[height=3.8cm, keepaspectratio]{figures_TRML/shiyi_yanse.pdf}
        \subcaption{Scope mismatch: full vs.\ adapter subspace}
    \end{subfigure}
    \hfill
  \hspace{4pt}
    \begin{subfigure}[t]{0.55\linewidth}
        \raggedleft
\includegraphics[height=3.8cm, keepaspectratio]{figures_TRML/lossland_curv1d_combine_loss_title2.pdf}
       \subcaption{Type effects: Different subspace perturbations shape full-model flatness}
        \label{fig:perturb_type_landscape}
    \end{subfigure}
  }
    \vspace{-6pt}
    \caption{\textbf{Where and How to Perturb in LoRA-Based CL}.
    \textbf{(a) Where (scope)}: 
    % full-space neighborhoods probe unreachable directions, while adapter-subspace perturbations align with task-admissible updates
%     In LoRA-based PECL, the trainable degrees of freedom are confined to the task-admissible update subspace, and the  model class is the affine set $W_{\mathrm{froz}}+\mathcal W_{\Delta,t}\subseteq\mathbb R^d$.
% Full-space perturbations in the ambient parameter space $\mathbb R^d$ generally inject components outside $\mathcal W_{\Delta,t}$, which can create a scope mismatch between the perturbation neighborhood and the PECL update dynamics.
% In contrast, perturbations supported on $\mathcal W_{\Delta,t}$ align the sharpness regularization with feasible adaptation directions.
% In LoRA-based CL, adaptation is restricted to task-admissible update directions.
% In LoRA-based continual learning, only the LoRA adapters are updated.
LoRA-based continual learning restricts adaptation to a task-admissible set of update directions.
Full-space perturbations explore ambient directions that may lie outside the feasible update dynamics, whereas adapter subspace perturbations confine the scope to update consistent directions.
% Full perturbations change the whole model, whereas LoRA-restricted subspace perturbations vary only the adapter part and therefore match how the model actually updates
% Full-space perturbations in the full parameter space may inject variations outside these directions, which can create a scope mismatch between the perturbation neighborhood and the actual update dynamics.
% In contrast, LoRA-restricted perturbations align sharpness regularization with feasible adaptation directions.
    % In LoRA-Based CL, adaptation is confined to the task-admissible LoRA subspace. Therefore, full-space perturbations probe unreachable directions, while LoRA-restricted perturbations align sharpness with feasible updates.
    % the model can adapt only within the task-admissible LoRA update subspace. Consequently, full-space perturbation neighborhoods probe directions that are unreachable under the update dynamics, whereas restricting perturbations to the LoRA coordinates aligns sharpness with the feasible adaptation directions.
    % Final average accuracy across perturbation strengths $\rho$ under different LoRA organizations. Restricting perturbations to the task adapter subspace $\Delta W$ consistently improves over the SGD baseline, whereas applying perturbations in the full parameter space $W+\Delta W$ substantially degrades performance.
    % Final average accuracy under different perturbation strengths $\rho$ for LoRA organizations. Restricting perturbations to the task-admissible adapter subspace $\Delta W$ consistently outperforms the SGD baseline, whereas full-parameter perturbations $W+\Delta W$ substantially degrade performance.
    \textbf{(b) How (type)}:
    % Compared with SGD, subspace perturbations consistently improve performance. 
Different perturbation types shape the geometry of updates within the LoRA subspace, leading to different solutions and thus effective models with different full-space flatness.
    % The top panel evaluates loss under perturbations restricted to the LoRA updates, while the bottom panel evaluates loss under perturbations over the full parameter.
    % Subspace-aware objectives yield flatter profiles when evaluated within the LoRA coordinates, and are associated with reduced effective loss sensitivity of the resulting model on the current task.
    % Compared with SGD, subspace sharpness aware  methods lead to flatter loss profiles and better performance.
    % Even under a fixed LoRA scope, perturbation \emph{type} matters: different subspace geometries produce different $\Delta W_t$, and thus different full-model flatness of $W_0+\Delta W_t$ under full-parameter evaluation, with stronger adversarial types yielding flatter profiles and better performance than SGD.
    % Subspace-aware perturbations produce flatter loss profiles than SGD when evaluated within the LoRA coordinates, and are associated with a flatter effective solution of the full model and improved task performance.
    % Subspace-aware perturbations yield systematically flatter profiles than SGD, and this flattening in the LoRA coordinates translates into improved {task performance of the full model}
    }
    \label{fig:scope_type1}
\end{figure*}
}



We answer this question through a sequential hierarchical PAC-Bayesian framework for continual learning, and then specialize it to frozen-backbone PECL. Under this specialization, task-level priors and posteriors are supported on the task-admissible adapter update subspace, which forces both the sharpness-dependent empirical term and the continual complexity term to reduce to subspace-defined quantities.


% In this work, we propose a PAC-Bayesian framework for analyzing and improving CL under LoRA-based adaptation. 
In this work, we develop a PAC-Bayesian framework for continual learning. We derive a sequential hierarchical PAC-Bayes framework and specialize it to frozen-backbone PECL.
Under this specialization, task-level priors and posteriors are supported on the task-admissible adapter update subspace, and both the sharpness-dependent empirical term and the continual complexity term reduce to quantities defined on this subspace.
% both the sharpness-related empirical term and the continual complexity are reduced to quantities defined in the trainable adapter update subspace.
% We derive a generalization bound that incorporates empirical risk, sharpness-based perturbation risk, and a hierarchical KL divergence term reflecting posterior complexity.
We further conduct a systematic empirical study across perturbation strategies and LoRA organizations, providing consistent evidence that adapter-restricted sharpness regularization improves continual performance and robustness.
% We further conduct a systematic study  of the interaction between different LoRA configurations and various perturbation strategies, showing that sharpness-aware regularization confined to the adapter subspace improves continual learning across LoRA variants.


% Our main contributions can be summarized as:
% \begin{itemize}
% \item We develop a novel PAC-Bayesian generalization bound that explicitly incorporates sharpness within the adapter subspace, offering a formal explanation for sharpness-aware regularization in PECL.
% \item We propose a systematic analysis of how LoRA modularity and sharpness-based perturbations interact, showing that these design choices correspond to complementary mechanisms for controlling forgetting and enhancing generalization.

% \item We provide comprehensive empirical validation across benchmarks with varying task sequence lengths and inter-task similarity, showing that sharpness within the adapter subspace reliably predicts downstream performance.
% \end{itemize}



Our main contributions can be summarized as:
\begin{itemize}
\vspace{-8pt}
\item  \textbf{Theory}: We derive a sequential PAC-Bayesian bound for continual learning with a hierarchical prior, yielding a transition-induced KL drift term that quantifies cross-task interference.
We specialize the bound to frozen-backbone PECL in which both the perturbation term and the complexity term are supported on the trainable adapter subspace.
% by supporting task posteriors on the admissible adapter update directions.
% We propose a task-level sequential hierarchical PAC-Bayes template for continual learning, thereby making cross-task interference explicit.
% We then specialize this template to PECL and obtain a bound in which both the sharpness-dependent empirical term and the complexity term are supported on the trainable adapter subspace.
% We propose a task level sequential hierarchical PAC Bayes template that separates within task adaptation from a hyperposterior drift penalty, making cross task interference explicit across task transitions.
\vspace{-2pt}
\item \textbf{Principle}: The bound makes perturbation design in PECL explicitly {task-conditional} and decomposes it into two coupled choices: \emph{where} the bound suggests restricting to the current task's admissible update subspace rather than the full space, and \emph{how} to perturb, where perturbation that align with dominant high-curvature directions are expected to bias the training trajectory toward flatter regions.
% a perturbation geometry (randomized vs.\ adversarial; zeroth- vs.\ first-order) supported on $\mathcal W_{\Delta,t}$.
% Accordingly, within the same scope, increasingly adversarial objectives are expected to yield solutions with smaller subspace sharpness, as they concentrate perturbations toward dominant high-curvature directions and more effectively discourage sharp minima.
% Under a frozen backbone and adapter supported posteriors, we show that both the KL complexity and the sharpness dependent term are confined to the LoRA update subspace, providing a theoretical criterion for when adapter only flatness control is sufficient.
\vspace{-2pt}
\item \textbf{Evidence}: We systematically evaluate perturbation scope/type across LoRA organizations and robustness benchmarks, confirming that adapter-restricted sharpness regularization improves continual performance and robustness.
% We propose a systematic analysis and provide comprehensive empirical validation of how LoRA modularity and sharpness-based perturbations interact, showing that these design choices correspond to complementary mechanisms for controlling forgetting and enhancing generalization.
\end{itemize}




% \section{Introduction}

% Continual learning (CL) has become increasingly relevant for adapting large-scale pre-trained models to evolving tasks in dynamic environments~\cite{PARISI201954,10444954,a2923225ea7d94e6c86faff97e23c69be,10.1145/3735633}. While full fine-tuning remains the most expressive adaptation strategy, it is computationally demanding and highly susceptible to catastrophic forgetting, since newly acquired knowledge may interfere with previously learned representations~\cite{howard2018universallanguagemodelfinetuning,doi:10.1073/pnas.1611835114,pmlr-v234-zhai24a,luo2025empiricalstudycatastrophicforgetting}. Parameter-efficient continual learning (PECL)~\cite{pmlr-v97-houlsby19a,dingParameterefficientFinetuningLargescale2023,He_2025_CVPR,NEURIPS2024_b0bc711f} mitigates these issues by freezing a large pre-trained backbone and confining adaptation to lightweight modules. Low-rank adaptation (LoRA)~\cite{hu2022lora} is a prominent instance of this paradigm, where task-specific updates are routed through low-rank adapters while the backbone remains fixed.



% In PECL, freezing the backbone has a structural consequence that directly shapes the stability–plasticity trade-off. Since the backbone and previously learned adapters are kept fixed, cross-task interference is concentrated in a shared trainable subspace. Beyond computational efficiency, a central challenge is therefore to understand and control the local geometry induced by this adapter subspace,  because the shape of the loss landscape around the solution governs how gradients from new tasks interact with existing adapter parameters and thus mediates the stability–plasticity trade-off.

% Flat minima in the loss landscape are widely believed to promote generalization and robustness, as flatter solutions are less sensitive to parameter perturbations~\citep{FlatMinima1997,wu2017understandinggeneralizationdeeplearning,chaudhariEntropySGDBiasingGradient2019,NEURIPS2020_518a38cc,NEURIPS2021_995f5e03,JMLRAnEmpiricalInvestigatio,yangRevisitingFlatnessAwareOptimization2025}.
% Sharpness-aware minimization (SAM)~\cite{foretSharpnessAwareMinimizationEfficiently2021} and related methods~\citep{pmlr-v151-bisla22a,zhangGradientNormAware2023,NEURIPS2024_C-Flat} operationalize this principle by seeking parameters with low loss  under small perturbations. 
% % Sharpness-aware optimizers such as SAM~\citep{foretSharpnessAwareMinimizationEfficiently2021}, RWP~\citep{pmlr-v151-bisla22a}, GAM~\citep{zhangGradientNormAware2023} and C-FLAT~\cite{NEURIPS2024_C-Flat} instantiate this idea by explicitly penalizing increases of the loss under small parameter perturbations. 
% While these methods have proven effective in full fine-tuning, their behavior in the restricted optimization landscape of PECL remains under-explored. 
% % LoRA-SAM~\citep{NEURIPS2024_4eb2c0ad} provides evidence that limiting sharpness-aware updates to the LoRA coordinates already yields substantial implicit regularization, whereas Flat-LoRA~\citep{li2025flatloralowrankadaptationflat} argues that flatness restricted to the LoRA subspace may be insufficient and therefore advocates perturbing all model parameters. In the PECL regime that we consider, however, the core efficiency principle is to strictly freeze the pre-trained backbone and confine optimization to the adapter subspace $\mathcal W_\Delta$. This tension raises a central question:
% % While effective in full fine-tuning, their application in PECL is not straightforward. 
% Recent work presents conflicting views: LoRA-SAM~\citep{NEURIPS2024_4eb2c0ad} suggests that optimizing flatness within the LoRA subspace alone provides effective regularization, while Flat-LoRA~\citep{li2025flatloralowrankadaptationflat} argues that achieving sufficient flatness requires perturbing all model parameters, which implicitly includes the frozen backbone. The latter requirement, however, contradicts the fundamental PECL constraint of a fixed backbone, thus motivating our central question:


% % \emph{In PECL, is it necessary to pursue global flatness by perturbing all model parameters, or is controlling flatness solely within the shared adapter subspace sufficient to govern stability and plasticity?}

% \emph{In PECL with a frozen backbone, is controlling flatness within the adapter subspace sufficient to govern stability and plasticity, or is flatness control induced by full-parameter perturbation necessary?}




% An initial empirical observation informs our investigation. Applying SAM solely to the adapter parameters can yield solutions residing in broader, flatter basins of the full-model loss landscape (Fig.~\ref{fig:3d_loss_landscape}). This suggests that regularizing the geometry of the adapter subspace can implicitly shape the effective geometry experienced by the model. 
% % We therefore hypothesize that subspace-restricted flatness may be adequate to improve robustness and balance in PECL.


% % Our investigation is grounded in a simple empirical observation. When SAM is applied only to the low-rank adapter subspace, the loss landscape of the full model around the final solution already becomes wider and smoother than under SGD, even when probed in directions that extend beyond $\mathcal W_\Delta$ (Figure~\ref{fig:3d_loss_landscape}). This suggests that the geometry of the frozen backbone and that of the trainable adapters are not independent. Stabilizing the local curvature within the trainable subspace can implicitly regularize the effective geometry of the whole model. Motivated by this phenomenon, we hypothesize that subspace flatness, understood as geometric control confined to the trainable adapter subspace, may already suffice to substantially improve robustness and achieve a better balance between stability and plasticity in PECL.

% % To investigate this hypothesis, we develop a joint  theoretical and empirical analysis. Theoretically, we derive a PAC-Bayesian generalization bound specialized to PECL with a frozen backbone and Gaussian posteriors supported on the adapter subspace.
% % In this bound, both the sharpness and complexity terms depend exclusively on the adapter subspace $\mathcal{W}_\Delta$, formally identifying adapter flatness as the relevant notion of flatness for PECL. 
% % Our analysis departs from classical hierarchical PAC-Bayesian bounds in two structural ways that are crucial for LoRA-based PECL. First, we derive a task-level CL bound with a time-varying hyper-posterior adapted to the task filtration, which removes the independence assumption on task priors and yields a complexity term that separates task-wise KL divergences from a hyper-posterior drift term capturing cross-task interference. Second, we specialise this bound to PECL with a frozen backbone and a shared LoRA subspace by restricting the posterior support to $\mathcal W_\Delta$, so that both the KL and the empirical terms can be rewritten in terms of adapter-space quantities and an expected sharpness term defined only along $\mathcal W_\Delta$. This subspace-aware formulation is absent from prior PAC-Bayesian analyses and is what allows us to formally distinguish global versus adapter-subspace flatness in the PECL setting.

% To address this question, we develop a sequential hierarchical PAC Bayes analysis tailored to PECL with a frozen backbone. Unlike classical hierarchical bounds that assume fixed hyper distribution, our formulation allows the task level hyperposterior to evolve with the task index and introduces an explicit drift penalty that quantifies cross task interference. We then specialize the framework to LoRA based PECL and show that, under the frozen backbone constraint, both the KL complexity and the sharpness dependent empirical term reduce exactly to quantities defined on the LoRA update subspace.


% Empirically, our results support a single message: in frozen-backbone PECL, the geometry that matters is the trainable adapter subspace. Adapter-only sharpness-aware optimization consistently improves continual performance and robustness, while full-space perturbations that include the frozen backbone do not yield systematic benefits and often harm stability. 
% % The advantage remains stable across task lengths and adapter ranks, suggesting that subspace flatness is a structural driver of generalization in PECL.

% % Empirically, we compare subspace-restricted and full-parameter sharpness-aware optimization across vision CL benchmarks. 
% % Empirically, the results converge to a clear practical conclusion. Across ViT-based PECL benchmarks with a frozen backbone, sharpness-aware optimization restricted to the adapter subspace yields consistent improvements in continual learning performance and robustness, accompanied by reductions in sharpness diagnostics and more stable feature and prototype statistics. In contrast, applying flatness control via perturbations in the full parameter space, which necessarily injects noise into the frozen backbone, provides no systematic gains and often degrades stability. 
% % develop a joint theoretical and empirical analysis tailored to PECL. On the theory side, we derive a PAC Bayes generalization bound under a frozen backbone and posteriors supported on the adapter update subspace. The bound makes two constraints explicit. First, continual learning requires a sequential notion of knowledge accumulation, which we capture through a task level hyperposterior that can evolve with task index and yields a drift penalty across tasks. Second, in PECL all task dependent randomness lies in the adapter subspace, and we show that both the complexity term and the sharpness dependent empirical term reduce exactly to adapter space quantities. This establishes adapter subspace flatness, rather than global flatness over all parameters, as the relevant notion of flatness in PECL.
% % To investigate this hypothesis, we develop a joint theoretical and empirical analysis. On the theory side, we derive a PAC-Bayesian generalization bound tailored to parameter-efficient continual learning with a frozen backbone and Gaussian posteriors confined to the shared adapter subspace. In contrast to previous PAC-Bayesian bounds, our formulation (i) uses a task-level, time-varying hyper-posterior that relaxes the independence assumption on task priors and separates per-task divergence from a hyper-posterior drift term that measures cross-task interference, and (ii) specializes this bound to the LoRA setting so that both the complexity term and the sharpness-dependent term are expressed purely in adapter-space quantities. This subspace-aware view is absent from prior PAC-Bayesian analyses and is what allows us to formally distinguish global versus adapter-subspace flatness in the PECL regime.
% % We further connect flatness in the factor space of LoRA to the eigenspectrum of a restricted Fisher information matrix, linking optimization geometry to feature-space stability.


% % Our results show that flatness-aware optimization confined to the adapter subspace consistently improves performance, whereas perturbing all parameters fails to yield systematic gains.

% % Under this modeling assumption, the complexity and flatness contributions to the bound can be expressed entirely in terms of quantities restricted to update subspace.



% % framework that connects probabilistic generalization guarantees with geometric stability in LoRA-based PECL. On the theoretical side, we derive a PAC-Bayesian generalization bound specialized to PECL with a frozen backbone and Gaussian posteriors supported on the adapter subspace. Under this modelling assumption, both the sharpness term and the KL complexity that govern the bound can be written entirely in terms of quantities restricted to $\mathcal W_\Delta$ rather than to the full parameter space. We further clarify the mechanism linking weight-space flatness to feature-space robustness by showing that LoRA updates implement an abstract update subspace and that optimizing flatness in factor coordinates corresponds to suppressing dominant eigenvalues of the local Fisher information matrix, which stabilizes the learned representations. On the empirical side, we conduct a systematic study on vision benchmarks that compares flatness-aware optimization restricted to $\mathcal W_\Delta$ with full-parameter perturbations within the same PECL models.



% Our main contributions are
% \begin{itemize}
% \item \textbf{Sequential complexity decomposition for continual adaptation.} We propose a task level sequential hierarchical PAC Bayes template that separates within task adaptation from a hyperposterior drift penalty, making cross task interference explicit across task transitions.
% \item \textbf{Subspace restricted sharpness aware PAC Bayes bound for PECL.} Under a frozen backbone and adapter supported posteriors, we show that both the KL complexity and the sharpness dependent term are confined to the LoRA update subspace, providing a theoretical criterion for when adapter only flatness control is sufficient.
% \item \textbf{Empirical Evaluation.}  Across ImageNet R and Tiny ImageNet C/P with ViT backbones and multiple LoRA organizations, we find that sharpness aware optimization restricted to the adapter subspace yields consistent gains, while full parameter perturbations do not provide systematic benefits and can introduce additional instability.
% \end{itemize}

% % \section{Introduction: joost}
% % Adapting large-scale pre-trained models to new data is now routine in modern machine learning. Due to the scale of contemporary models, this adaptation is typically performed using parameter-efficient transfer learning methods~\cite{han2024parameter} such as low-rank adaptation (LoRA)~\cite{hu2022lora}, which restrict task-specific updates to low-rank adapter subspaces. However, parameter-efficient updates remain susceptible to catastrophic forgetting~\cite{mccloskey1989catastrophic}, as updates for new data can interfere with previously acquired knowledge. Continual learning (CL) provides a principled framework for mitigating this interference~\cite{doi:10.1073/pnas.1611835114,rebuffi2017icarl, de2021continual}. Combining parameter-efficient updates with continual learning objectives leads to parameter-efficient continual learning (PECL), which enables continual adaptation of large pre-trained models while mitigating catastrophic forgetting under constrained parameter budgets~\cite{dingParameterefficientFinetuningLargescale2023,He_2025_CVPR,NEURIPS2024_b0bc711f}.


% % % Better to move to low-rank already here
% % Another parallel strategy for mitigating forgetting is flat-minima optimization. Originally, flat minima in the loss landscape were shown to promote generalization and robustness~\citep{FlatMinima1997,wu2017understandinggeneralizationdeeplearning,chaudhariEntropySGDBiasingGradient2019,NEURIPS2020_518a38cc,NEURIPS2021_995f5e03,JMLRAnEmpiricalInvestigatio,yangRevisitingFlatnessAwareOptimization2025}. Sharpness-aware minimization (SAM))~\cite{foretSharpnessAwareMinimizationEfficiently2021} and related methods~\citep{pmlr-v151-bisla22a,zhangGradientNormAware2023,NEURIPS2024_C-Flat} operationalize this principle by seeking parameters that maintain low loss under small perturbations. Beyond improved generalization, flat minima have also been argued to reduce catastrophic forgetting as flatter solutions are less sensitive to parameter perturbations. Prior work derives that loss landscape flatness provides an upper bound on forgetting (CITE), and supporting empirical evidence has been reported in subsequent studies (CITE). While flat-minima optimization has proven effective in the context of full fine-tuning, its behavior under parameter-efficient updates, as considered in PECL, remains underexplored. In particular, it is unclear whether minimizing loss-landscape curvature within a restricted parameter subspace induces any meaningful control over the curvature of the full model.
% % %Recent work presents conflicting views: LoRA-SAM~\citep{NEURIPS2024_4eb2c0ad} suggests that optimizing flatness within the LoRA subspace alone provides effective regularization, while Flat-LoRA~\citep{li2025flatloralowrankadaptationflat} argues that achieving sufficient flatness requires perturbing all model parameters, which implicitly includes the frozen backbone. The latter requirement, however, contradicts the fundamental PECL constraint of a fixed backbone, thus motivating our central question:

% % Several works have suggested that over-parameterization is a key factor of sharp minima in deep networks (CITE). This raises the question of whether low-rank parameter updates, by severely restricting the optimization subspace, inherently preclude sharp minima, and thus render flat-minima optimization ineffective in parameter-efficient regimes. In this work, we investigate the relationship between curvature in sparse update subspaces and curvature in the full parameter space. \JW{Surprisingly, we find that the dominant curvature directions of the full updated model substantially align with the sparse update directions, indicating that flat-minima optimization can remain effective even under highly constrained parameter updates. In addition, we look at the effectiveness of two popular flat minima optimization methods and show that first-order minimization algorithms reduce principal curvature more significantly for sparse updates than for full model updates. We hypothesize that this is due ....  In extensive experiments, we verify that the improved flatness translates to considerable performance gains for continual learning.  }

% % %In addition, we develop a sequential hierarchical PAC Bayes analysis tailored to PECL. Unlike classical hierarchical bounds that assume fixed hyper distribution (CITE), our formulation allows the task level hyperposterior to evolve with the task index and introduces an explicit drift penalty that quantifies cross task interference. We then specialize the framework to LoRA based PECL and show that, under the frozen backbone constraint, both the KL complexity and the sharpness dependent empirical term reduce exactly to quantities defined on the LoRA update subspace. This founding is confirmed experimentally ???

% % In addition, we develop a sequential hierarchical PAC-Bayes analysis tailored to parameter-efficient continual learning (PECL). Unlike classical hierarchical PAC-Bayes bounds that assume a fixed hyperprior (CITE), our formulation allows the task-level hyperposterior to evolve across tasks and introduces an explicit drift penalty that quantifies cross-task interference. We further adapt this framework to LoRA-based PECL and show that, under the frozen-backbone constraint, both the KL complexity term and the sharpness-dependent empirical term reduce exactly to quantities defined on the LoRA update subspace. This result provides a principled theoretical criterion for when flatness control restricted to adapter parameters is sufficient.

% % %Empirically, our results support a single message: in frozen-backbone PECL, the geometry that matters is the trainable adapter subspace. Adapter-only sharpness-aware optimization consistently improves continual performance and robustness, while full-space perturbations that include the frozen backbone do not yield systematic benefits and often harm stability. 

% % Our main contributions are
% % \begin{itemize}
% % \item \textbf{First systematic investigation of flat-minima optimization in parameter-efficient network updates.} We show that parameter-efficient subspaces exhibit substantial overlap with the dominant curvature directions of the full model. This explains why flat-minima optimization remains effective even when restricted to low-rank adapter updates.
% % %Although flat-minima optimization methods operate on a restricted parameter subspace, we show that this subspace exhibits substantial overlap with the dominant curvature directions of the full model. As a result, flat-minima optimization remains effective at promoting flatter solutions even under parameter-efficient updates.
% % \item \textbf{Superiority of first-order flatness methods for parameter-efficient learning.} We demonstrate that first-order flatness-based methods significantly outperform zeroth-order alternatives in parameter-efficient regimes. We provide theoretical justification and empirically verify that these methods induce larger reductions in the principal curvature of the loss landscape.

% % \item \textbf{Subspace restricted sharpness-aware PAC Bayes bound for PECL.} Under a frozen backbone and adapter supported posteriors, we prove that both the KL complexity and the sharpness dependent term are confined to the LoRA subspace. This provides a theoretical criterion for when adapter-only control is sufficient to obtain flat minima.

% % \item \textbf{Empirical Evaluation.}  Across ImageNet-R and Tiny ImageNet C/P with ViT backbones and multiple LoRA organizations, we find that sharpness aware optimization restricted to the adapter subspace yields consistent gains, while full parameter perturbations do not provide systematic benefits and can introduce additional instability.
% % \end{itemize}

% % % \begin{itemize}
% % % \item \textbf{Subspace restricted PAC-Bayes theorem for PECL.} We derive a sharpness aware generalization bound in which both the KL complexity and the sharpness dependent term are confined to the adapter update subspace, providing a theoretical criterion for when adapter only flatness control is sufficient under a frozen backbone.
% % % \item \textbf{Sequential complexity decomposition.} We formulate a task level sequential hierarchical PAC Bayes template that separates within task adaptation from a hyperposterior drift penalty, making cross task interference explicit at the level of task transitions.
% % % \item \textbf{Empirical Evaluation.} Across ImageNet R and Tiny ImageNet C and P with ViT backbones and multiple LoRA organizations, we find that sharpness aware optimization restricted to the adapter subspace yields consistent gains, while full parameter perturbations do not provide systematic benefits and can introduce additional instability in PECL.
% % % \end{itemize}


% % % \begin{itemize}
% % % \item \textbf{PAC-Bayesian Analysis for PECL.} We provide a generalization bound for PECL where the sharpness and complexity are defined over the adapter subspace, offering a theoretical justification for subspace-restricted flatness optimization.
% % % \item \textbf{Geometry of Adapter Flatness.} We show that LoRA updates implement an abstract update subspace and that minimizing curvature in factor space is equivalent to compressing the Fisher spectrum along this subspace. This spectral effect stabilizes class prototypes and feature representations.
% % % \item \textbf{Empirical Evaluation.} On ImageNet R and Tiny ImageNet C/P with ViT B/16 backbones and three adapter organizations (SeqLoRA, IncLoRA, OLoRA), we find that sharpness aware optimization restricted to the adapter subspace consistently improves CL metrics and robustness, whereas full parameter perturbations do not yield systematic gains in the PECL regime and can introduce additional instability.
% % % \end{itemize}

% % \begin{figure}[tbp]
% % \centering
% % \begin{minipage}{1\linewidth}
% % \centering
% % \includegraphics[width=1\linewidth]{figures_TRML/2D_lossLandscape_task0_lora_opt_sam_vs_sgd_3d_compare_withxyz.pdf}
% % \caption{Visualization of the loss landscape for a ViT backbone trained with LoRA adapters, comparing SGD and SAM optimizers. Applying SAM to adapter parameters yields a flatter local landscape around the final solution.}
% % % The SAM optimizer not only discovers a wider and flatter basin but also achieves a lower loss value
% % \label{fig:3d_loss_landscape}
% % \end{minipage}
% % \end{figure}

% % % Our main contributions are as follows:
% % % \begin{itemize}
% % %     \item \textbf{Sharpness-aware PAC-Bayes for PECL.} We derive a PAC-Bayesian generalization bound tailored to PECL with a frozen backbone and adapter-subspace posteriors. In this regime, both the sharpness and the KL complexity that govern the bound are determined by the trainable adapter subspace $\mathcal W_\Delta$, which identifies adapter-subspace flatness as the relevant notion of flatness for PECL.
% % %     \item \textbf{Geometric link between adapter flatness and feature stability.} We clarify how LoRA updates realize an abstract update subspace and show that minimizing curvature in factor space is equivalent to suppressing dominant eigenvalues of the restricted Fisher information matrix. This spectral effect stabilizes feature representations and explains how subspace flatness shapes the global loss landscape.
% % %     \item \textbf{Empirical analysis of flatness strategies in PECL.} We provide systematic evidence on ImageNet-R, ImageNet-C and Tiny-ImageNet-P with ViT-B/16 adapters under three adapter-subspace sharing strategies (SeqLoRA, IncLoRA and OLoRA). Flatness-aware optimization confined to $\mathcal W_\Delta$ consistently improves continual learning metrics and feature stability compared with SGD, whereas full-parameter perturbations do not yield consistent gains in this regime and can introduce additional instability.
% % % \end{itemize}

% \section{Introduction before}

% Continual learning (CL) has become increasingly relevant for adapting large-scale pre-trained models to evolving tasks in dynamic environments~\cite{PARISI201954,10444954,a2923225ea7d94e6c86faff97e23c69be,10.1145/3735633}. While full fine-tuning remains the most expressive adaptation strategy, it is computationally demanding and highly susceptible to catastrophic forgetting, since newly acquired knowledge may interfere with previously learned representations~\cite{howard2018universallanguagemodelfinetuning,doi:10.1073/pnas.1611835114,pmlr-v234-zhai24a,luo2025empiricalstudycatastrophicforgetting}. Parameter-efficient continual learning (PECL)~\cite{pmlr-v97-houlsby19a,dingParameterefficientFinetuningLargescale2023,He_2025_CVPR,NEURIPS2024_b0bc711f} mitigates these issues by freezing a large pre-trained backbone and confining adaptation to lightweight modules. Low-rank adaptation (LoRA)~\cite{hu2022lora} is a prominent instance of this paradigm, where task-specific updates are routed through low-rank adapters while the backbone remains fixed.

% % \begin{figure}[tbp]
% % \centering
% % \begin{minipage}{1\linewidth}
% % \centering
% % \includegraphics[width=1\linewidth]{figures_TRML/2D_lossLandscape_task0_lora_opt_sam_vs_sgd_3d_compare_withxyz.pdf}
% % \caption{Visualization of the loss landscape for a ViT backbone trained with LoRA adapters, comparing SGD and SAM optimizers. Applying SAM to adapter parameters yields a flatter local landscape around the final solution.}
% % % The SAM optimizer not only discovers a wider and flatter basin but also achieves a lower loss value
% % \label{fig:3d_loss_landscape}
% % \end{minipage}
% % \end{figure}

% % In PECL, freezing the backbone has a structural consequence that directly shapes the stability–plasticity trade-off. Since the backbone and previously learned adapters are kept fixed, cross-task interference is concentrated in a shared trainable subspace. Beyond computational efficiency, a central challenge is therefore to understand and control the local geometry induced by this adapter subspace,  because the shape of the loss landscape around the solution governs how gradients from new tasks interact with existing adapter parameters and thus mediates the stability–plasticity trade-off.

% Flat minima in the loss landscape are widely believed to promote generalization and robustness, as flatter solutions are less sensitive to parameter perturbations~\citep{FlatMinima1997,wu2017understandinggeneralizationdeeplearning,chaudhariEntropySGDBiasingGradient2019,NEURIPS2020_518a38cc,NEURIPS2021_995f5e03,JMLRAnEmpiricalInvestigatio,yangRevisitingFlatnessAwareOptimization2025}.
% Sharpness-aware minimization (SAM)~\cite{foretSharpnessAwareMinimizationEfficiently2021} and related methods~\citep{pmlr-v151-bisla22a,zhangGradientNormAware2023,NEURIPS2024_C-Flat} operationalize this principle by seeking parameters with low loss  under small perturbations. 
% % Sharpness-aware optimizers such as SAM~\citep{foretSharpnessAwareMinimizationEfficiently2021}, RWP~\citep{pmlr-v151-bisla22a}, GAM~\citep{zhangGradientNormAware2023} and C-FLAT~\cite{NEURIPS2024_C-Flat} instantiate this idea by explicitly penalizing increases of the loss under small parameter perturbations. 
% % While these methods have proven effective in standard supervised learning, their behavior in the restricted optimization landscape of PECL remains under-explored. 
% While effective in standard supervised learning, how such flatness control should be formulated in parameter efficient continual learning (PECL) remains under explored.
% % LoRA-SAM~\citep{NEURIPS2024_4eb2c0ad} provides evidence that limiting sharpness-aware updates to the LoRA coordinates already yields substantial implicit regularization, whereas Flat-LoRA~\citep{li2025flatloralowrankadaptationflat} argues that flatness restricted to the LoRA subspace may be insufficient and therefore advocates perturbing all model parameters. In the PECL regime that we consider, however, the core efficiency principle is to strictly freeze the pre-trained backbone and confine optimization to the adapter subspace $\mathcal W_\Delta$. This tension raises a central question:
% % While effective in full fine-tuning, their application in PECL is not straightforward.


% % In CL, the notion of generalization extends beyond a single task: it is the aggregated test performance over a sequence of tasks, and forgetting manifests as a loss of this generalization on earlier tasks after training on later ones.
% % In continual learning, generalization is the joint test performance over all learned tasks, and forgetting corresponds to a degradation of this generalization on earlier tasks after learning later ones. 

% Unlike single task training where generalization is evaluated on one fixed distribution, continual learning evaluates generalization over a sequence of tasks, and forgetting is precisely the degradation of this generalization on earlier tasks after learning later ones \cite{}. And the definition of sharpness is alsorealterd to the data, but int the continual  learning, the data in the sequential task are changable.
% This sequential setting makes the geometry induced by the optimization constraints central, because only directions that remain reachable by future updates can influence how the model moves and thus how its test performance evolves across tasks.


% In standard fine tuning, flatness is  operationalized through perturbations in the full parameter space\cite{}. Although some study some the effective apply in this setting.

% However,in the PECL, the pre trained backbone is strictly frozen and all task adaptation is confined to an update subspace, so perturbations applied outside $\mathcal W_\Delta$ cannot be compensated by learning, since those directions are unreachable by subsequent task updates.
% % However, perturbations applied outside $\mathcal W_\Delta$ cannot be compensated by learning, since those directions are unreachable by subsequent task updates.
% This creates a structural mismatch between the usual global perturbation rationale and the PECL constraint.
% Recent findings reflect this tension: some works~\citep{NEURIPS2024_4eb2c0ad} report that sharpness control restricted to adapter coordinates is already sufficient, while others~\citep{li2025flatloralowrankadaptationflat} advocate full parameter perturbations that implicitly include the frozen backbone.
% This motivates our question:
% \emph{In PECL with a frozen backbone, where and how to apply perturabation to prompte  faltness, and what the relation ship between shrpness in the sequail task to the genereralization} 
% % \emph{In PECL with a frozen backbone, does generalization relevant flatness reside in the adapter subspace $\mathcal W_\Delta$, or must it be enforced by perturbing the full parameters?} 

% % \emph{In PECL with a frozen backbone, does generalization relevant flatness reside in the adapter subspace $\mathcal W_\Delta$, or must it be enforced by perturbing the full parameters?} 


% % In the PECL, the pre trained backbone is strictly frozen and all task adaptation is confined to an update subspace, for example the LoRA induced adapter subspace $\mathcal W_\Delta$.

% % the backbone and previously learned adapters are kept fixed, cross-task interference is concentrated in a shared trainable subspace.

% % In standard fine-tuning, ``flatness'' is typically operationalized through perturbations in the full parameter space.
% % In PECL, however, the backbone is frozen and optimization is restricted to an update subspace, e.g., the LoRA-induced adapter subspace $\mathcal W_\Delta$.
% % This creates a structural mismatch between the usual global-perturbation rationale and the PECL constraint. Recent findings reflect this tension: some works~\citep{NEURIPS2024_4eb2c0ad} report that subspace-only sharpness control is sufficient, while others~\citep{li2025flatloralowrankadaptationflat} advocate full-parameter perturbations that implicitly include the frozen backbone.
% % Recent work presents conflicting views: LoRA-SAM~\citep{NEURIPS2024_4eb2c0ad} suggests that optimizing flatness within the LoRA subspace alone provides effective regularization, while Flat-LoRA~\citep{li2025flatloralowrankadaptationflat} argues that achieving sufficient flatness requires perturbing all model parameters, which implicitly includes the frozen backbone. The latter requirement, however, contradicts the fundamental PECL constraint of a fixed backbone, 
% % \emph{In PECL, is it necessary to pursue global flatness by perturbing all model parameters, or is controlling flatness solely within the shared adapter subspace sufficient to govern stability and plasticity?}
% % \emph{In PECL with a frozen backbone, is controlling flatness within the adapter subspace sufficient to govern stability and plasticity, or is flatness control induced by full-parameter perturbation necessary?}
% % This motivating our central question:
% % \emph{In PECL with a frozen backbone, does generalization-relevant flatness reside in the adapter subspace $\mathcal W_\Delta$, or must it be enforced by perturbing the full parameter vector?}



% % An initial empirical observation informs our investigation. Applying SAM solely to the adapter parameters can yield solutions residing in broader, flatter basins of the full-model loss landscape (Fig.~\ref{fig:3d_loss_landscape}). This suggests that regularizing the geometry of the adapter subspace can implicitly shape the effective geometry experienced by the model. 
% % We therefore hypothesize that subspace-restricted flatness may be adequate to improve robustness and balance in PECL.


% % Our investigation is grounded in a simple empirical observation. When SAM is applied only to the low-rank adapter subspace, the loss landscape of the full model around the final solution already becomes wider and smoother than under SGD, even when probed in directions that extend beyond $\mathcal W_\Delta$ (Figure~\ref{fig:3d_loss_landscape}). This suggests that the geometry of the frozen backbone and that of the trainable adapters are not independent. Stabilizing the local curvature within the trainable subspace can implicitly regularize the effective geometry of the whole model. Motivated by this phenomenon, we hypothesize that subspace flatness, understood as geometric control confined to the trainable adapter subspace, may already suffice to substantially improve robustness and achieve a better balance between stability and plasticity in PECL.

% % To investigate this hypothesis, we develop a joint  theoretical and empirical analysis. Theoretically, we derive a PAC-Bayesian generalization bound specialized to PECL with a frozen backbone and Gaussian posteriors supported on the adapter subspace.
% % In this bound, both the sharpness and complexity terms depend exclusively on the adapter subspace $\mathcal{W}_\Delta$, formally identifying adapter flatness as the relevant notion of flatness for PECL. 
% % Our analysis departs from classical hierarchical PAC-Bayesian bounds in two structural ways that are crucial for LoRA-based PECL. First, we derive a task-level CL bound with a time-varying hyper-posterior adapted to the task filtration, which removes the independence assumption on task priors and yields a complexity term that separates task-wise KL divergences from a hyper-posterior drift term capturing cross-task interference. Second, we specialise this bound to PECL with a frozen backbone and a shared LoRA subspace by restricting the posterior support to $\mathcal W_\Delta$, so that both the KL and the empirical terms can be rewritten in terms of adapter-space quantities and an expected sharpness term defined only along $\mathcal W_\Delta$. This subspace-aware formulation is absent from prior PAC-Bayesian analyses and is what allows us to formally distinguish global versus adapter-subspace flatness in the PECL setting.

% % To address this question, we develop a sequential hierarchical PAC Bayes analysis tailored to PECL with a frozen backbone. Unlike  previous bounds that assume fixed hyper distribution, our formulation allows the task level hyperposterior to evolve with the task index and introduces an explicit drift penalty that quantifies cross task interference. We then specialize the framework to LoRA based PECL and show that, under the frozen backbone constraint, both the KL complexity and the sharpness dependent empirical term reduce exactly to quantities defined on the LoRA update subspace.

% \begin{figure*}[thbp]
% \centering
%     \begin{subfigure}{0.555\linewidth}
%         \centering
% \includegraphics[width=0.95\linewidth]{figures_TRML/perturbation_range2.pdf}
%         \subcaption{Perturbation Range}
%         \label{fig: pertuabtion scope}
%     \end{subfigure}
%         \begin{subfigure}{0.435\linewidth}
%         \centering
% \includegraphics[width=0.95\linewidth]{figures_TRML/perturbation_type1.pdf}
%         \subcaption{Perturbation type}
%     \end{subfigure}
%     \centering
% \end{figure*}

% To address this question, we develop a sequential hierarchical PAC-Bayes analysis for PECL under a frozen backbone. Our framework models continual adaptation as a task-indexed hyperposterior process and yields a pathwise bound whose complexity decomposes into (i) a within-task posterior divergence and (ii) an explicit hyperposterior drift term that quantifies cross-task interference across transitions.
% Specializing the bound to LoRA-based PECL, we show that the frozen-backbone constraint induces an exact reduction: both the task-wise KL complexity and the sharpness-dependent empirical term depend only on distributions supported on the LoRA update subspace $\mathcal W_{\Delta,t}$. This provides a theoretical criterion for flatness control in PECL: only perturbations that lie in the future-admissible update subspace can tighten the bound and are therefore relevant to continual generalization.
% Empirically, across class-incremental benchmarks with ViT backbones and multiple LoRA organizations, sharpness-aware optimization restricted to adapter parameters yields consistent gains, whereas full-parameter perturbations that include frozen backbone directions provide no systematic benefit and often reduce stability.

% % Empirically, our results support a single message: in frozen-backbone PECL, the geometry that matters is the trainable adapter subspace. Adapter-only sharpness-aware optimization consistently improves continual performance and robustness, while full-space perturbations that include the frozen backbone do not yield systematic benefits and often harm stability. 
% % The advantage remains stable across task lengths and adapter ranks, suggesting that subspace flatness is a structural driver of generalization in PECL.

% % Empirically, we compare subspace-restricted and full-parameter sharpness-aware optimization across vision CL benchmarks. 
% % Empirically, the results converge to a clear practical conclusion. Across ViT-based PECL benchmarks with a frozen backbone, sharpness-aware optimization restricted to the adapter subspace yields consistent improvements in continual learning performance and robustness, accompanied by reductions in sharpness diagnostics and more stable feature and prototype statistics. In contrast, applying flatness control via perturbations in the full parameter space, which necessarily injects noise into the frozen backbone, provides no systematic gains and often degrades stability. 
% % develop a joint theoretical and empirical analysis tailored to PECL. On the theory side, we derive a PAC Bayes generalization bound under a frozen backbone and posteriors supported on the adapter update subspace. The bound makes two constraints explicit. First, continual learning requires a sequential notion of knowledge accumulation, which we capture through a task level hyperposterior that can evolve with task index and yields a drift penalty across tasks. Second, in PECL all task dependent randomness lies in the adapter subspace, and we show that both the complexity term and the sharpness dependent empirical term reduce exactly to adapter space quantities. This establishes adapter subspace flatness, rather than global flatness over all parameters, as the relevant notion of flatness in PECL.
% % To investigate this hypothesis, we develop a joint theoretical and empirical analysis. On the theory side, we derive a PAC-Bayesian generalization bound tailored to parameter-efficient continual learning with a frozen backbone and Gaussian posteriors confined to the shared adapter subspace. In contrast to previous PAC-Bayesian bounds, our formulation (i) uses a task-level, time-varying hyper-posterior that relaxes the independence assumption on task priors and separates per-task divergence from a hyper-posterior drift term that measures cross-task interference, and (ii) specializes this bound to the LoRA setting so that both the complexity term and the sharpness-dependent term are expressed purely in adapter-space quantities. This subspace-aware view is absent from prior PAC-Bayesian analyses and is what allows us to formally distinguish global versus adapter-subspace flatness in the PECL regime.
% % We further connect flatness in the factor space of LoRA to the eigenspectrum of a restricted Fisher information matrix, linking optimization geometry to feature-space stability.


% % Our results show that flatness-aware optimization confined to the adapter subspace consistently improves performance, whereas perturbing all parameters fails to yield systematic gains.

% % Under this modeling assumption, the complexity and flatness contributions to the bound can be expressed entirely in terms of quantities restricted to update subspace.



% % framework that connects probabilistic generalization guarantees with geometric stability in LoRA-based PECL. On the theoretical side, we derive a PAC-Bayesian generalization bound specialized to PECL with a frozen backbone and Gaussian posteriors supported on the adapter subspace. Under this modelling assumption, both the sharpness term and the KL complexity that govern the bound can be written entirely in terms of quantities restricted to $\mathcal W_\Delta$ rather than to the full parameter space. We further clarify the mechanism linking weight-space flatness to feature-space robustness by showing that LoRA updates implement an abstract update subspace and that optimizing flatness in factor coordinates corresponds to suppressing dominant eigenvalues of the local Fisher information matrix, which stabilizes the learned representations. On the empirical side, we conduct a systematic study on vision benchmarks that compares flatness-aware optimization restricted to $\mathcal W_\Delta$ with full-parameter perturbations within the same PECL models.



% Our main contributions are
% \begin{itemize}
%     \item \textbf{A PAC-Bayes bound for continual learning}.We introduce a sequential hierarchical PAC-Bayes bound tailored to continual learning. The bound makes cross-task interference explicit through a hyperposterior drift term accumulated over task transitions.
%      \item \textbf{A PECL reduction that resolves the scope mismatch.} We specialize the bound to PECL with a frozen backbone and adapter-supported posteriors, and show that both the KL complexity and the sharpness-dependent empirical term reduce to the task-admissible update subspace $\mathcal W_{\Delta,t}$, identifying the geometry relevant for continual stability.
     
%      % We specialize the bound to PECL with a frozen backbone and adapter-supported posteriors. In this setting, both the KL complexity and the sharpness-dependent empirical term depend only on the task-admissible update subspace $\mathcal W_{\Delta,t}$, which identifies the relevant geometry for continual stability and explains why full-parameter perturbations can be misaligned with PECL.
%      \item \textbf{How to perturb and how to organize adapters in PECL.}
%      The bound yields concrete design rules: within $\mathcal W_{\Delta,t}$, choose perturbation covariances $\Sigma_t$ that emphasize high-curvature directions, and organize LoRA updates to reduce cross-task overlap of ${\mathcal W_{\Delta,t}}$ so that future adaptation is less amplified by inherited high-curvature directions.
     
%      % We use the bound to clarify how LoRA organization and perturbation  jointly control the empirical sharpness term and the KL budget.
%      % This yields concrete design principles for LoRA-based PECL, which we validate across multiple LoRA organizations and flatness-aware perturbation choices. (ii) choose LoRA organizations that reduce cross-task overlap of admissible updates, so that new-task adaptation is less amplified by high-curvature directions inherited from past tasks; and (iii) shape the perturbation covariance $\Sigma_t$ to emphasize high-curvature directions within $\mathcal W_{\Delta,t}$, since these directions dominate the sharpness-dependent empirical term under comparable KL and drift budgets.
% \end{itemize}


% % \begin{itemize}
% % \item \textbf{Sequential complexity decomposition for continual adaptation.} We propose a task level sequential hierarchical PAC Bayes template that separates within task adaptation from a hyperposterior drift penalty, making cross task interference explicit across task transitions.
% % \item \textbf{Subspace restricted sharpness aware PAC Bayes bound for PECL.} Under a frozen backbone and adapter supported posteriors, we show that both the KL complexity and the sharpness dependent term are confined to the LoRA update subspace, providing a theoretical criterion for when adapter only flatness control is sufficient.
% % \item \textbf{Empirical Evaluation.}  Across ImageNet R and Tiny ImageNet C/P with ViT backbones and multiple LoRA organizations, we find that sharpness aware optimization restricted to the adapter subspace yields consistent gains, while full parameter perturbations do not provide systematic benefits and can introduce additional instability.
% % \end{itemize}








% \section{Introduction}

% Continual learning (CL) has become increasingly relevant for adapting large-scale pre-trained models to evolving tasks in dynamic environments~\cite{PARISI201954,10444954,a2923225ea7d94e6c86faff97e23c69be,10.1145/3735633}. While full fine-tuning remains the most expressive adaptation strategy, it is computationally demanding and highly susceptible to catastrophic forgetting, since newly acquired knowledge may interfere with previously learned representations~\cite{howard2018universallanguagemodelfinetuning,doi:10.1073/pnas.1611835114,pmlr-v234-zhai24a,luo2025empiricalstudycatastrophicforgetting}. Parameter-efficient continual learning (PECL)~\cite{pmlr-v97-houlsby19a,dingParameterefficientFinetuningLargescale2023,He_2025_CVPR,NEURIPS2024_b0bc711f} mitigates these issues by freezing a large pre-trained backbone and confining adaptation to lightweight modules. Low-rank adaptation (LoRA)~\cite{hu2022lora} is a prominent instance of this paradigm, where task-specific updates are routed through low-rank adapters while the backbone remains fixed. It concentrates cross task interference into a shared trainable update subspace and reshapes the stability plasticity trade off. As a result, the local geometry induced by this subspace becomes a primary mechanism that governs how gradients from new tasks interact with previously learned adapter directions, thereby mediating forgetting which corresponds to the degradation of generalization for earlier task.




% % Continual learning (CL) studies how to learn a sequence of tasks while retaining performance on earlier ones. 
% % Unlike single-task training, Continual learning (CL) evaluates generalization over an evolving sequence of tasks, and forgetting corresponds to the degradation of earlier task performance after learning later ones. 
% % Forgetting arises when updates for new tasks interfere with previously acquired knowledge, making stability under repeated adaptation the central challenge. 
% %  Continual learning (CL) is increasingly relevant for large pre-trained models deployed in dynamic environments, where repeated full fine tuning is computationally costly and often unstable due to severe cross-task interference.
% % Parameter-efficient continual learning (PECL) mitigates the cost of adapting large pre-trained models by freezing the backbone and updating only lightweight modules. Low-rank adaptation (LoRA) is a prominent instance, where task-dependent changes are implemented by low-rank adapters while the backbone remains fixed. This design induces a structural constraint: at each task, optimization is restricted to the parameters designated as trainable, so cross-task interference is mediated by the local geometry along these task-admissible directions.


% % Flatness-aware optimization has improved generalization in standard supervised learning by penalizing sensitivity to small parameter perturbations. Sharpness-aware minimization (SAM) instantiates this idea by minimizing the worst-case loss increase in a neighborhood. In frozen-backbone PECL, however, extending sharpness-aware optimization is not straightforward: full-space perturbations inevitably include frozen backbone parameters that cannot be updated by subsequent task updates. Perturbing such unreachable parameters can therefore be misaligned with the update mechanism that governs stability and forgetting, creating a scope mismatch between global perturbation-based flatness control and PECL constrict.

% Flat minima in the loss landscape are widely associated with improved generalization and robustness in standard supervised learning, as flatter solutions are less sensitive to small parameter perturbations~\citep{FlatMinima1997,wu2017understandinggeneralizationdeeplearning,NEURIPS2020_518a38cc,NEURIPS2021_995f5e03,JMLRAnEmpiricalInvestigatio,yangRevisitingFlatnessAwareOptimization2025}.
% Sharpness-Aware Minimization (SAM)~\cite{foretSharpnessAwareMinimizationEfficiently2021} and related methods~\citep{pmlr-v151-bisla22a,zhangGradientNormAware2023,NEURIPS2024_C-Flat} operationalize this principle by penalizing the loss increase under small perturbations, thereby encouraging flat minima (i.e., promoting flatness or reducing sharpness; we use sharpness as the unified term to quantify perturbation sensitivity).
% % have demonstrated the benefits of encouraging flat minima  (i.e., promoting flatness or reducing sharpness; we use sharpness as the unified term in this work) in stand supervised learning.
% % Sharpness is  task-conditional in continual learning (CL) which proceeds over a sequence of tasks.
% % For each task $t$, it should be defined with respect to the task-specific risk and depends on the current task data distribution.

% In continual learning (CL), sharpness is inherently task conditional: for each task $t$, it is defined with respect to the task specific risk and depends on the current data distribution. 
% Under full fine tuning, sharpness is naturally defined in the full parameter space because future adaptation can proceed along arbitrary directions.
% However, extending sharpness related method  to parameter efficient continual learning (PECL) is not straightforward.



% % freezing a large pre-trained backbone and confining adaptation to lightweight modules. It has a structural consequence that directly shapes the stability--plasticity trade-off: since the backbone and previously learned adapters are kept fixed, cross-task interference is concentrated in a shared trainable update subspace. Controlling the local geometry induced by this subspace is therefore central, because the shape of the loss landscape around the current solution governs how gradients from new tasks interact with existing adapter directions and thereby mediates stability and forgetting.

% A central open question is where and how to apply perturbations for flatness control in the restricted optimization landscape of LoRA based PECL. Existing approaches offer conflicting prescriptions. LoRA-SAM~\citep{NEURIPS2024_4eb2c0ad} advocates optimizing sharpness only within the trainable LoRA coordinates, whereas Flat-LoRA~\citep{li2025flatloralowrankadaptationflat} argues that achieving sufficient flatness requires perturbing all model parameters, which implicitly includes the frozen backbone. Importantly, these conclusions are largely drawn in single task adaptation settings and are not evaluated under continual learning. In continual learning, sharpness is task conditional and the optimization landscape evolves across tasks, so the relevant question is not merely which perturbations reduce sensitivity around a fixed optimum, but which perturbations shape the stability plasticity trade off over task sequence. More broadly, existing flatness analyses do not provide a unified account of which notion of flatness is relevant under PECL constraints. This leaves the behavior of perturbation scope, perturbation geometry, and adapter organization in LoRA based continual learning largely unresolved, and motivates a unified theoretical and empirical analysis tailored to PECL in continual learning.






% % There is a argument on where should put perturbation. LoRA-SAM~\citep{NEURIPS2024_4eb2c0ad} suggests that optimizing flatness within the LoRA subspace alone provides effective regularization, while Flat-LoRA~\citep{li2025flatloralowrankadaptationflat} argues that achieving sufficient flatness requires perturbing all model parameters, which implicitly includes the frozen backbone. And which type of metohd should we apply is not sure. While sharpness methods have proven effective in full fine-tuning, their behavior in the restricted optimization landscape of PECL remains under-explored.


% % A central open question is where and how to apply perturbations in this restricted optimization landscape. 

% % A central open question is where and how to apply perturbations for flatness control in the restricted optimization landscape of LoRA based PECL. Existing approaches offer conflicting prescriptions.
% % LoRA-SAM~\citep{NEURIPS2024_4eb2c0ad} advocates optimizing sharpness only within the LoRA subspace, whereas Flat-LoRA~\citep{li2025flatloralowrankadaptationflat} argues that sufficient flatness requires perturbing all model parameters, which implicitly includes the frozen backbone. Despite the empirical success of sharpness aware methods in full fine tuning, their behavior, appropriate perturbation scope, and effective perturbation geometry in PECL remain under explored.  



% % More broadly, existing flatness analyses do not provide a unified account of which notion of flatness is relevant under PECL constraints, nor how adapter design reshapes it across tasks. In PECL, the relevant geometry is the sensitivity along directions that are admissible for optimization at the current task.





% % Freezing the backbone has a structural consequence that directly shapes the stability--plasticity trade-off: since the backbone and previously learned adapters are kept fixed, cross-task interference is concentrated in a shared trainable update subspace.
% % Controlling the local geometry induced by this subspace is therefore central, because the shape of the loss landscape around the current solution governs how gradients from new tasks interact with existing adapter directions and thereby mediates stability and forgetting. In this setting, extending sharpness-aware optimization is not straightforward.

% % While full fine-tuning is expressive, it is computationally demanding and often exacerbates catastrophic forgetting, namely the degradation of earlier task performance after learning later ones.


% % is a prominent instance, where task-specific updates are routed through low-rank adapters while the backbone remains fixed.




% % operationalize sharpness-aware learning by seeking parameters whose task-conditioned loss remains stable under small perturbations, where the perturbation neighborhood is typically defined in the full parameter space. This choice is matched to full fine-tuning, since future adaptation can move along arbitrary parameter directions.



% % Unlike single-task training, Continual learning (CL) learn a sequence of tasks and evaluates generalization over an evolving sequence of tasks.
% % In continual learning, this sensitivity is inherently \emph{task-conditional}: for each task $t$, sharpness is defined with respect to the task-specific risk, and thus depends on the current task data distribution.
% % Sharpness-Aware Minimization (SAM)~\cite{foretSharpnessAwareMinimizationEfficiently2021} and related methods~\citep{pmlr-v151-bisla22a,zhangGradientNormAware2023,NEURIPS2024_C-Flat} operationalize flatness-aware learning by seeking parameters whose task-conditioned loss remains low under small perturbations, where the neighborhood is typically taken in the full parameter space.
% % This formulation is well matched to full fine-tuning, since future adaptation can move along arbitrary parameter directions.



%  % While full fine-tuning remains the most expressive adaptation strategy, it is computationally demanding and highly susceptible to catastrophic forgetting which means the degradation of earlier task performance after learning later ones.
%  % Parameter-efficient continual learning (PECL)~\cite{pmlr-v97-houlsby19a,dingParameterefficientFinetuningLargescale2023,He_2025_CVPR,NEURIPS2024_b0bc711f} mitigates these issues by freezing a large pre-trained backbone and confining adaptation to lightweight modules . Low-rank adaptation (LoRA)~\cite{hu2022lora} is a prominent instance of this paradigm, where task-specific updates are routed through low-rank adapters while the backbone remains fixed. In PECL, freezing the backbone has a structural consequence that directly shapes the stability–plasticity trade-off. Since the backbone and previously learned adapters are kept fixed, cross-task interference is concentrated in a shared trainable subspace. Beyond computational efficiency, a central challenge is therefore to understand and control the local geometry induced by this adapter subspace,  because the shape of the loss landscape around the solution governs how gradients from new tasks interact with existing adapter parameters and thus mediates the stability–plasticity trade-off.

 

% %  however, extending sharpness-aware optimization is not straightforward: full-space perturbations inevitably include frozen backbone parameters that cannot be updated by subsequent task updates. While these methods have proven effective in full fine-tuning, their behavior in the restricted optimization landscape of PECL remains under-explored.Perturbing such unreachable parameters can therefore be misaligned with the update mechanism that governs stability and forgetting, creating a scope mismatch between global perturbation-based flatness control and PECL constrict.
% % More broadly, existing flatness analyses do not provide a unified account of which notion of flatness is relevant under PECL constraints, nor how adapter design reshapes it across tasks. In PECL, the relevant geometry is the sensitivity along directions that are admissible for optimization at the current task.


% We derive a sequential hierarchical PAC-Bayes bound and specialize it to frozen backbone PECL by restricting task level posteriors and priors to the adapter supported update subspaces. Under this specialization, both the sharpness dependent empirical contribution and the continual complexity depend only on task admissible update directions, rather than on arbitrary full parameter perturbations. This motivates \emph{Task-Conditional Adapter-Subspace Sharpness}, which measures sensitivity to perturbations restricted to the adapter coordinates that are trainable at task $t$, while explicitly excluding frozen components. The resulting decomposition separates within task adaptation from cross task interference accumulated over transitions, and it provides a principled explanation of why flatness conclusions can diverge when perturbations probe directions that are not admissible for optimization.



% % In this work, we develop a unified theoretical and empirical analysis of perturbation-based flatness control for LoRA-based PECL in continual learning. Theoretically, we derive a sequential hierarchical PAC-Bayes framework and specialize it to frozen-backbone PECL with adapter-supported task-level distributions. Under this specialization, both the sharpness-related empirical term and the continual complexity reduce to quantities defined on task-admissible update directions. We therefore focus on \emph{Task-Conditional Adapter-Subspace Sharpness}, defined as sensitivity to perturbations restricted to the adapter coordinates that are allowed to be updated at the current task, explicitly excluding frozen components. This restriction provides a consistent basis for analyzing and comparing flatness control in PECL, and clarifies why conclusions can diverge when perturbations are applied outside admissible directions. The resulting bound yields an interpretable decomposition that separates within-task adaptation from cross-task interference accumulated over transitions, and it explains how LoRA organization affects continual stability by reshaping the admissible directions and their overlap across tasks.


% Guided by this perspective, we study flatness control through three coupled design axes: perturbation scope (task-admissible versus full-parameter), perturbation type within admissible directions (isotropic versus curvature-aligned), and LoRA organization across tasks (how admissible directions evolve and overlap). Across class-incremental benchmarks with ViT backbones and multiple LoRA organizations, sharpness-aware optimization restricted to adapter parameters yields consistent gains and improved stability.



% \begin{figure*}[thbp]
% \centering
%     \begin{subfigure}{0.59\linewidth}
%         \centering
% \includegraphics[width=0.95\linewidth]{figures_TRML/perturbation_range2.pdf}
%         \subcaption{Perturbation Range}
%         \label{fig: pertuabtion scope}
%     \end{subfigure}
%         \begin{subfigure}{0.40\linewidth}
%         \centering
% \includegraphics[width=0.95\linewidth]{figures_TRML/triple_junction_random_theta_three_points2.pdf}
%         \subcaption{Perturbation type}
%     \end{subfigure}
%     \centering
%      \caption{Perturbation design in frozen-backbone PECL. \textbf{(a) Perturbation range:} full-parameter perturbations inject noise into both the frozen backbone and adapters, whereas task-admissible perturbations are restricted to the adapter coordinates that are trainable at the current task. \textbf{(b) Perturbation type:} schematic comparison of randomized zeroth-order sharpness $\mathrm{RS}^{(0)}_S(W,\Sigma)$, adversarial zeroth-order sharpness $\mathrm{AS}^{(0)}_S(W,\rho)$, and adversarial first-order sharpness $\mathrm{AS}^{(1)}_S(W,\rho)$, illustrating their relative strength under comparable perturbation budgets.}
% \end{figure*}







% % , whereas full-parameter perturbations that include frozen backbone directions provide no systematic benefit and often reduce continual performance.

% Our main contributions are as follows:
% \begin{itemize}
%     \item \textbf{Sequential hierarchical PAC-Bayes for continual learning.} We derive a task-wise hierarchical PAC-Bayes bound that explicitly decomposes continual complexity into within-task adaptation and across-task hyper-prior drift, yielding an accumulated drift penalty that is absent in multitask PAC-Bayes bounds with fixed hyper-priors.
%     \item \textbf{PECL specialization and task-admissible sharpness.} Under frozen-backbone PECL with adapter-supported posteriors, we prove that both the sharpness-dependent empirical term and the continual complexity depend only on task-admissible update directions, which formalizes \emph{task-admissible sharpness}  as the appropriate notion for PECL, and clarifies why full-forward perturbations advocated for full-space perturbation can be mismatched in PECL.
%     \item \textbf{Design principles for perturbations and LoRA organization.} We systematically evaluate perturbation \emph{scope} and \emph{geometry} together with LoRA organization on ViT-based continual learning benchmarks, and identify actionable criteria: reducing curvature exposure within task-admissible directions and reducing cross-task overlap of admissible directions improves PECL stability and accuracy.
% \end{itemize}








\section{Preliminaries}





% \label{sec:background}

% We first summarize the continual learning setting, the parameter efficient formulation with shared LoRA adapters, and the notions of flatness and PAC--Bayesian complexity that we use in the rest of the paper.


% \section{PAC-Bayes Bound and Sharpness}

\subsection{Classical PAC-Bayesian bound}
Let $\mathcal X$ be the input space, $\mathcal Y$ the label space, and $\mathcal Z = \mathcal X \times \mathcal Y$ the data space.
A model with parameters $W \in \mathbb R^d$ induces a predictor $f_W : \mathcal X \to \mathcal Y$.
We consider a bounded loss $\ell : \mathcal H \times \mathcal Z \to [0,1]$ and write $\ell(W,z)$ for $\ell(f_W,z)$. 
Given a distribution $\mathcal D$ over $\mathcal Z$ and a training set $S=\{z_j\}_{j=1}^m$ drawn i.i.d.\ from $\mathcal D$, the population and empirical risks are
$
L_{\mathcal D}(W)
=
\mathbb E_{z\sim\mathcal D}\ell(W,z),
\hat L_{S}(W)
=
\frac{1}{m}\sum_{j=1}^{m}\ell(W,z_{j}).
$


The PAC-Bayesian viewpoint place probability measures over the hypothesis space $\mathcal{H}$ rather than focusing on a single function: starting from a prior $P\in\mathcal M(\mathcal H)$, the learner observes data $S$ and outputs a posterior $Q(\cdot\mid P,S)\in\mathcal M(\mathcal H)$\cite{mcallesterPACBayesianModelAveraging1999}. The corresponding posterior-averaged population and empirical risks  are $
L_{\mathcal D}(Q)
=
\mathbb E_{W \sim Q(\cdot \mid P,S)}\mathbb E_{z\sim\mathcal D}\ell(W,z)
$,
$
\hat L_{S}(Q_t)
=
\mathbb E_{W \sim Q(\cdot \mid P,S)}\frac{1}{m}\sum_{j=1}^{m_t}\ell(W,z_{j}).
$
A classical PAC--Bayes bound \cite{mcallesterPACBayesianStochasticModel2003} for a single task under a fixed distribution is stated as follows.
\begin{theorem}[PAC-Bayes Generalization Bound for single task]
\label{theorem:classical}
For any loss $\ell: \mathcal{F} \times \mathcal{X} \times \mathcal{Y} \to [0,1]$, with probability at least $1-\delta$ over the choice of training set $S$:
{%
\setlength{\abovedisplayskip}{2pt}
\setlength{\belowdisplayskip}{3pt}
\setlength{\abovedisplayshortskip}{3pt}
\setlength{\belowdisplayshortskip}{3pt}
    \[\mathbb{E}_{f \sim Q}[L_{\mathcal{D}}^{\ell}(f)] \leq \mathbb{E}_{f \sim Q}[\hat{L}_S^{\ell}(f)] + \sqrt{\frac{KL(Q\|P) + \log \frac{m}{\delta}}{2(m-1)}}\]
}
\end{theorem}
\vspace{-8pt}
\subsection{Sharpness}

Flat minima are often associated with improved generalization and robustness in standard supervised learning. Motivated by PAC-Bayes analyses, sharpness-aware optimization seeks to reduce the empirical contribution in generalization bounds by controlling the loss elevation in a local neighborhood. In particular, SAM can be interpreted as minimizing an \textbf{adversarial zeroth-order sharpness} on the training set $S$ within a radius $\rho$~\citep{AWP2020,foretSharpnessAwareMinimizationEfficiently2021}:
{%
\setlength{\abovedisplayskip}{2pt}
\setlength{\belowdisplayskip}{3pt}
\setlength{\abovedisplayshortskip}{3pt}
\setlength{\belowdisplayshortskip}{3pt}
\[\mathrm{AS}^{(0)}_S(W,\rho) =\max_{ \| \epsilon \|\leq \rho } \Big[ \hat{L}_{S}(W+{\epsilon})-\hat{L}_{S}(W)\ \Big].\]
}
% \vspace{-5mm}
A complementary formulation considers randomized perturbations and minimizes the expected loss elevation under Gaussian noise, referred to as \textbf{random zeroth-order sharpness}~\citep{liRevisitingRandomWeight2024,jiang2019fantasticgeneralizationmeasures}:
{%
\setlength{\abovedisplayskip}{2pt}
\setlength{\belowdisplayskip}{3pt}
\setlength{\abovedisplayshortskip}{3pt}
\setlength{\belowdisplayshortskip}{3pt}
\[
\mathrm{RS}_S^{(0)}({W}, \Sigma)
 =
 \mathbb{E}_{\varepsilon \sim \mathcal{N}(0,\Sigma)}
 \big[
 \hat{L}_{S}({W}+\varepsilon) - \hat{L}_{S}({W})
 \big],
 \]}Typical choices for $\mathcal{N}(0,\Sigma)$ include isotropic Gaussians and structured variants such as filter-wise Gaussians~\citep{pmlr-v151-bisla22a}. Moreover, SAM can be interpreted as optimizing the mean loss under a fixed perturbation variance: for a suitable correspondence between the radius $\rho$ and the covariance $\Sigma$, the adversarial zeroth-order sharpness provides an upper bound on the random zeroth-order sharpness~\cite{mollenhoffSAMOptimalRelaxation2022}.

Beyond loss elevation, one may measure sharpness through the neighborhood gradient norm. The \textbf{adversarial first-order sharpness} is defined as the maximum gradient norm within a radius $\rho$ on the training set~\citep{Zhang_GAM}:
{%
\setlength{\abovedisplayskip}{2pt}
\setlength{\belowdisplayskip}{3pt}
\setlength{\abovedisplayshortskip}{3pt}
\setlength{\belowdisplayshortskip}{3pt}
\[\mathrm{AS}^{(1)}_S(W,\rho) = \max_{ \| \epsilon \|\leq \rho } \| \nabla_{W} \hat{L}_S(W+\epsilon)\|.\]}Compared with zeroth-order sharpness which measures  loss increase, the first-order criterion emphasizes the sharpness of the gradient itself, and is closely tied to dominant curvature directions. For any $W$ and fixed $\rho$, the adversarial first-order sharpness dominates the adversarial zeroth-order sharpness, hence minimizing $\mathrm{AS}^{(1)}_S$ also reduces $\mathrm{AS}^{(0)}_S$, making it a stronger flatness measure~\citep{Zhang_GAM}.

 







% \subsection{LoRA-Based PECL}
% \label{sec:pecl-setup}



% The accuracy at step $t$ is
% $
% \operatorname{ACC}_t = \frac{1}{t} \sum_{j=1}^{t} a_{tj},
% $ where $\{a_{tj}\}_{j\leq t}$ denotes the accuracy on task $j$ after learning task $t$,
% and we report final accuracy $\operatorname{ACC}_T^{\mathrm{cls}}$ and time-averaged accuracy $\operatorname{AAA}^{\mathrm{cls}} = \frac{1}{T}\sum_{t=1}^{T} \operatorname{ACC}_t.$ 

% We also compute the corresponding Nearest-Class-Mean Classifier (NCM) results
% % Nearest Mean of Exemplars (NME) metrics 
% $(\operatorname{ACC}_T^{\mathrm{ncm}}, \operatorname{AAA}^{\mathrm{ncm}})$ using class prototypes in the feature space, which isolates the quality and stability of representations from the classifier head.



\subsection{LoRA-Based PECL}
\label{sec:pecl-setup}

LoRA-based PECL adapts a frozen backbone by learning low-rank updates across tasks. A key design axis is how these updates are organized in parameter space, since the resulting sharing pattern determines the degree of knowledge transfer and cross-task interference. We focus on \emph{shared-subspace} PECL, where tasks interact through a common adapter mechanism, and summarize three representative organizations.

\textbf{Fully shared.}
SeqLoRA reuses a single LoRA pair $(A,B)$ for all tasks and updates it sequentially on the fixed backbone $W_0$. This maximizes parameter sharing but also couples tasks most strongly, making interference accumulate within the same low-rank directions.

\textbf{Incremental shared.}
Incremental designs expand the adapter set over time by introducing a new pair $(A_t,B_t)$ at task $t$ while keeping $W_0$ and earlier adapters fixed, and then composing multiple adapters at inference. Gated variants further select or combine adapters with task-dependent weights~\citep{Gao_2023_ICCV,10.1145/3730436.3730456,zhao-etal-2024-sapt,cheng2025exploiting,wu2025sdlorascalabledecoupledlowrank}.

\textbf{Constraint incremental shared.}
Constraint-based methods regularize the geometry of the adapter space to reduce conflicts based incremental designs. OLoRA~\citep{wang2023orthogonal} encourages orthogonality between task-specific directions to reduce parameter overlap, while gradient-based approaches align or project updates to avoid previously occupied directions~\citep{Liang_2024_CVPR,yang2024parametercollisionhinderingcontinual}.

% For completeness, some methods allocate fully disjoint adapters per task and activate them conditionally by task identification at inference. Such isolated designs are not the focus of our analysis that targets the shared-subspace.

% # regime
% where both transfer and interference are governed by how LoRA updates overlap across tasks.



% % \subsection{LoRA-Based PECL and Adapter Subspace Sharing}

% Recent work on LoRA-based PECL trains low-rank adapters across tasks over a frozen backbone\cite{}. A central design axis is the organization of these low-rank adapters within the parameter space, which governs subspace sharing, potential for knowledge transfer, and inter-task interference.
% From this viewpoint, LoRA-based PECL methods fall into two regimes: isolated per-task adaptation without subspace interaction, and shared-subspace adaptation where new tasks are trained on top of previously LoRA parameters.
% Some methods train an independent LoRA adapter for each task and only activate the corresponding adapter at inference, so there is no parameter sharing and no interaction in the LoRA subspace~\citep{Yu_2024_CVPR}. This almost removes cross-task interference, but also prevents knowledge reuse and requires explicit task selection at inference time.
% Shared-Subspace Adaptation builds a common adapter subspace over the frozen backbone, updated sequentially with incoming tasks. Three representative families illustrate how strongly tasks are coupled inside this subspace.
% % \begin{itemize}
%     % \item 

    
%     \textbf{Fully shared subspace}
%     SeqLoRA reuse a single low-rank adapter for all tasks. A LoRA pair $(A,B)$ is updated sequentially on a fixed backbone $W_0$. 
%     % Training and inference stay inside the same low-rank subspace, which maximises sharing but also maximizes tasks interference.
%     % \item 

    
%     \textbf{Incremental shared subspace.}
%     Incremental designs expand the adapter subspace over time. IncLoRA introduces a new pair $(A_t,B_t)$ at task $t$ while keeping $W_0$ and earlier adapters fixed and composes all adapters at inference.
% % , for example
% %     $
% %         W_t = W_0 + \sum_{i=1}^t B_i A_i.
% %     $
%     % The effective update still lies in a common span, so tasks interact inside a shared low-rank subspace, but the per-task decomposition provides partial separation. 
%     Gated variants further select or combine adapter with different weight~\citep{Gao_2023_ICCV,10.1145/3730436.3730456,zhao-etal-2024-sapt,cheng2025exploiting,wu2025sdlorascalabledecoupledlowrank}.
%     % \item  , aiming to reuse useful directions while avoiding unnecessary interference

    
%     \textbf{Constraint-augmented incremental sharing.}
%     Constraint-augmented methods regularise the geometry of the incremental subspace to mitigate conflicts. OLoRA~\citep{wang2023orthogonal} adds orthogonality penalties to encourage task-specific directions and reduce paramters overlap. Gradient-level approaches~\citep{Liang_2024_CVPR,yang2024parametercollisionhinderingcontinual} align or project gradients away from previously occupied directions. 
%     % These methods still share a global LoRA subspace, but seek to organise it so that shared adapter capture reusable knowledge while conflicting directions are suppressed.
% % \end{itemize}


% To model the sequential nature of knowledge accumulation across tasks in CL, we adopt a hierarchical PAC‑Bayesian formulation. In this framework, the task-level prior itself is treated as a random variable, and a hyper-posterior is defined as a probability measure $\mathcal P \in \mathcal M(\mathcal M(\mathcal H))$ over such priors~\citep{pentinaPACBayesianBoundLifelong2014,JMLR:v24:22-1254}.
% % To analyze the role of sharpness in PECL, we adopt a hierarchical PAC-Bayesian framework, where the prior $P_t$ itself as a random variable.
% % Follow hyper PAC , we treat the prior $P_t$ itself as a random variable.
% %We write $\mathcal F_t = \sigma(S_1,\dots,S_t)$ for the $\sigma$-algebra generated by the first $t$ datasets and consider the filtration $\mathcal F_0 \subseteq \mathcal F_1 \subseteq \cdots \subseteq \mathcal F_T$. 
% % At the beginning of task $t$, before observing $S_t$, the learner maintains a hyper-posterior $\mathcal P_{t-1}$ that is $\mathcal F_{t-1}$-measurable and encodes all information extracted from previous tasks. 
% Let $\mathcal F_{t-1}=\sigma(S_1,\dots,S_{t-1})$ denote the $\sigma$-algebra generated by the data observed up to task $t-1$. 
% Before observing $S_t$, the learner maintains a hyper-posterior $\mathcal P_{t-1}$ that is $\mathcal F_{t-1}$-measurable, summarizing information transferred from previous tasks. 
% % The learning process at task $t$ is then described in two stages. First, a task-level prior $P_t$ is sampled from the current hyper-posterior,
% % $
% % P_t \sim \mathcal P_{t-1}.
% % $
% % Second, after observing $S_t$, the learner updates from $P_t$ to a task-level posterior $Q_t$, and the hyper-posterior is updated to $\mathcal P_t$. 
% % We assume that these updates are absolutely continuous, so that the associated KL divergences are well defined.
% % $
% % \mathcal P_t \ll \mathcal P_{t-1},
% % \quad
% % Q_t(\cdot \mid P,S_t) \ll P(\cdot)
% % \quad \text{for } \mathcal P_{t-1}\text{-almost every } P.
% % $
% % The corresponding task-level risks under the hierarchical posterior are obtained by averaging over the sampling of priors and hypotheses,
% The learning procedure at task $t$ proceeds in two stages: (i) sample a task-level prior $P_t$ from the current hyper-posterior, $P_t\sim\mathcal P_{t-1}$; (ii) after observing $S_t$, form the task-level posterior $Q_t$ and update the hyper-posterior to $\mathcal P_t$. 
% The corresponding hierarchical risks are obtained by averaging over both the sampled prior and the posterior draw of hypotheses,
% $
% \mathcal L_t
% =
% \mathbb E_{P_t \sim \mathcal P_t}
% \,
% \mathbb E_{W_t \sim Q_t}
% \big[
% L_{\mathcal D_t}(W_t)
% \big],
% \hat{\mathcal L}_t
% =
% \mathbb E_{P_t \sim \mathcal P_t}
% \,
% \mathbb E_{W_t \sim Q_t}
% \big[
% \hat L_{S_t}(W_t)
% \big].
% $


% In PECL, a pre-trained model with frozen weights $W_0$ is adapted to sequential tasks by tuning only lightweight adapter modules. For LoRA, the adapted weights at task $t$ are given by $W_t = W_{\mathrm{frozen}} + \Delta W_t,$ where $W_{\mathrm{frozen}}$ contains the pretrained backbone and all previously learned adapters, which remain frozen at task $t$, and $\Delta W_t$ is a new low rank update confined to a shared adapter subspace $\mathcal W_\Delta$. 
% We focus on PECL methods that reuse a shared adapter subspace on a frozen backbone and distinguish three representative organizations.

% \begin{itemize}
%     \item \textbf{SeqLoRA} (fully shared subspace). A single LoRA pair $(A,B)$ is trained sequentially across tasks, yielding $W_t = W_0 + BA$ at all $t$. Both training and inference remain confined to the same low-rank subspace.
%     \item \textbf{IncLoRA} (incremental partially fused subspaces). Each task $t$ introduces a new pair $(A_t,B_t)$ while keeping $W_0$ and earlier pairs fixed:
%     $
%         W_t = W_0 + \sum_{i=1}^t B_i A_i.
%     $
%     At inference, a fixed or gated combination of adapters is used.
%     \item \textbf{OLoRA} (constraint-augmented sharing). On top of an incremental design, OLoRA~\cite{wang2023orthogonal} introduces explicit orthogonality penalties such as
%     $
%         \sum_{i<t} \|A_i^\top A_t\|_F^2
%     $
%     to encourage task-specific directions within the shared adapter subspace and reduce interference.
% \end{itemize}

% These designs share the same frozen backbone $W_0$ and differ only in how they organize and regularize the shared LoRA subspace $\mathcal W_\Delta$. Methods with completely disjoint adapters that never interact during training or inference are outside our scope.


% We consider three representative strategies for organizing this subspace: a single, fully-shared adapter (SeqLoRA); incrementally added, task-specific adapters (IncLoRA); and adapters regularized for orthogonality to minimize interference (OLoRA).


% In PECL we decompose the parameters at task $t$ as
% $$
% W_t
% =
% W_{\mathrm{frozen}} + \Delta W_t,
% \qquad
% \Delta W_t \in \mathcal W_\Delta \subseteq \mathbb R^d,
% $$
% where $W_{\mathrm{frozen}}$ contains the pretrained backbone and all previously learned adapters, which remain frozen at task $t$, and $\Delta W_t$ is a new low rank update confined to a shared adapter subspace $\mathcal W_\Delta$.

% We focus on three representative organizations of LoRA adapters on the shared backbone.
% SeqLoRA maintains a single adapter in $\mathcal W_\Delta$ and updates it sequentially across tasks.
% IncLoRA accumulates task specific adapters and combines them through a shared head on top of $W_{\mathrm{frozen}}$.
% OLoRA augments this sharing with explicit orthogonality regularization between adapter directions.
% All three methods operate in a common update subspace and differ only in how task specific components are composed inside $\mathcal W_\Delta$.






\section{PAC-Bayesian Generalization Theory with Sharpness for CL}
\label{sec:pac-pecl}

{%
\setlength{\intextsep}{8pt}        % [h] 浮动体与正文的上下间距（最关键）
\setlength{\textfloatsep}{6pt}     % [t]/[b] 时浮动体与正文间距（备用）
\setlength{\abovecaptionskip}{2pt} % 图与caption之间（上方）间距
\setlength{\belowcaptionskip}{0pt} % caption下方到正文的间距
\captionsetup{skip=4pt}            % 你已在用：图与caption的距离（更细粒度）

\begin{figure}[h]
\centering
\includegraphics[width=\linewidth]{figures_TRML/hyperbiyao3.pdf}
\caption{
Time-varying Hyperposterior in Continual Learning. \textbf{Left:} In multi-task learning, task priors are drawn from a single time-invariant hyperdistribution, so there is no sequential drift. \textbf{Right:} In continual learning, the hyperposterior is updated over time, inducing a history-dependent cumulative drift across task transitions.
% Time-varying Hyperposterior in Continual Learning. \textbf{Left:} Multi-task uses a single hyperdistribution drift term to sample task priors.
% \textbf{Right:} Continual learning updates the hyperposterior over time, yielding history-dependent cumulate hyperdistribution drift.
}
\label{fig:showPAC3}
\end{figure}
}% 结束局部分组


\subsection{From single-task bound to continual learning}
\label{sec:single_to_cl}

Theorem~\ref{theorem:classical} provides a generalization guarantee for a single task under a fixed data distribution and a prior that is fixed before observing the corresponding training set. Continual learning instead proceeds over a task sequence $t\in\{1,\dots,T\}$, where at each step a dataset $S_t=\{z_{tj}\}_{j=1}^{m_t}$ is drawn i.i.d.\ from $\mathcal D_t$ and the learner updates its hypothesis after observing $S_t$. This sequential setting implies that the inductive bias should adapt with the task history, so the task prior at time $t$ must reflect information from tasks $1,\dots,t-1$, which cannot be captured by treating $\{P_t\}$ as externally specified constants.

We model transferable inductive bias by treating the task prior as a random object and introducing a hyperposterior $\mathcal P\in\mathcal M(\mathcal M(\mathcal H))$ over task-level priors. The key motivation is that, in continual learning, the prior used at step $t$ is typically produced by a history-dependent rule that progressively incorporates information from earlier tasks, rather than being generated from a single time-invariant hyperdistribution as shown in Fig \ref{fig:showPAC3}. 
%~\citep{pentinaPACBayesianBoundLifelong2014}
% This pathwise perspective aligns with the information-budget view of online PAC-Bayes, while making explicit an additional object that must be tracked across tasks: the evolution of the prior-generation mechanism itself, beyond the within-task update from $P_t$ to $Q_t$.
% This pathwise perspective is consistent with the information-budget view of online PAC-Bayes. It additionally tracks how the prior-generation mechanism evolves across tasks, beyond the examples level update from $P_t$ to $Q_t$.
%his prior-generating rule nonstationary
Concretely, common CL mechanisms make the prior sequence nonstationary, including curvature-based regularization that updates the effective prior through Fisher or Hessian summaries~\citep{EWC2017,SI2017} and scope-allocation strategies that modify which parameters remain trainable over time~\citep{progressive2016,HAT2018,Piggyback2018,Supermasks2020,adapter2019,hu2022lora}, which naturally introduces a notion of drift between priors.
% leads to an explicit drift penalty across task transitions.


% We therefore model transferable inductive bias by treating the task prior as a random object and introducing a hyperposterior $\mathcal P\in\mathcal M(\mathcal M(\mathcal H))$ over task level priors. Figure~\ref{fig:showPAC3} highlights a key departure from classical hierarchical PAC Bayes analyses for exchangeable multi task environments~\citep{pentinaPACBayesianBoundLifelong2014}: there, a single time invariant hyperdistribution generates task priors, so the task order plays no explicit role. In continual learning, by contrast, the prior used at step $t$ is produced by a history dependent rule that is updated as tasks arrive, because the learner must revise what is transferable based on the information extracted from earlier tasks.
% This sequential view is consistent with the information budget interpretation of online PAC Bayes, but it introduces a structurally different object to be tracked across tasks: not only the within task update from $P_t$ to $Q_t$, but also the evolution of the \emph{prior generation mechanism} that maps past experience to the next task prior. The resulting decomposition separates the complexity of fitting the current task from the cost of updating the inductive bias across task transitions. Such history dependence arises broadly in continual learning, for example when importance based regularization updates the effective prior using curvature summaries~\citep{EWC2017,SI2017}, or when architectural and masking strategies modify which directions remain trainable over time~\citep{progressive2016,HAT2018,Piggyback2018,Supermasks2020,adapter2019,hu2022lora}. These mechanisms motivate modeling continual learning as a time varying hyperposterior process and make an explicit drift penalty across task transitions unavoidable.

Let $\mathcal F_{t-1}=\sigma(S_1,\dots,S_{t-1})$ denote the information available before observing task $t$. We model continual learning as a hyperposterior process $\{\mathcal P_t\}_{t=0}^T$ where $\mathcal P_{t-1}$ is $\mathcal F_{t-1}$ measurable, which yields a pathwise bound with an explicit hyperposterior drift penalty that quantifies the information budget of sequential belief updates. We define the hierarchical test and empirical risks
$
\mathcal L_t
=
\mathbb E_{P_t \sim \mathcal{P}_{t-1}}
\,
\mathbb E_{W_t \sim Q_t(\cdot \mid P_t,S_t)}
\big[
L_{\mathcal D_t}(W_t),
\big],
$
$
\hat{\mathcal L}_t
=
\mathbb E_{P_t \sim \mathcal{P}_{t-1}}
\,
\mathbb E_{W_t \sim Q_t(\cdot \mid P_t,S_t)}
\big[
\hat L_{S_t}(W_t)
\big].
$

\begin{theorem}[PAC Bayes with time varying hyperposterior]
\label{thm:pathwise_pac}
For any $\delta\in(0,1]$, with probability at least $1-\delta$ over $S_{1:T}$,
{%
\setlength{\abovedisplayskip}{2pt}
\setlength{\belowdisplayskip}{3pt}
\setlength{\abovedisplayshortskip}{3pt}
\setlength{\belowdisplayshortskip}{3pt}
\begin{equation*}
\resizebox{\linewidth}{!}{$
\begin{aligned} 
&\quad\frac{1}{T} \sum_{t=1}^T \mathcal L_t
\le  
\frac{1}{T} \sum_{t=1}^T \hat{\mathcal L}_t  + \\
&
\sqrt{
\frac{
\displaystyle
\sum_{t=1}^T
\mathbb E_{P_t\sim\mathcal P_{t-1}}
\mathrm{KL}\bigl(Q_t \,\Vert\, P_t\bigr)
+
\sum_{t=1}^T
\mathrm{KL}\bigl(\mathcal P_t\Vert\mathcal P_{t-1}\bigr)
+
\log\frac{1}{\delta}}
{2\displaystyle\sum_{t=1}^T (m_t - 1)}
}.
\end{aligned}
$}
\end{equation*}
}
\end{theorem}

% Theorem~\ref{thm:pathwise_pac} bounds the trajectory averaged hierarchical test risk by the corresponding trajectory averaged empirical risk plus a complexity penalty, where averaging over $t=1,\dots,T$ aligns with the standard continual learning protocol of summarizing performance over the full task sequence. 
Theorem~\ref{thm:pathwise_pac} bounds the trajectory-averaged hierarchical test risk by the corresponding trajectory-averaged empirical risk plus a complexity penalty, aligning with standard continual learning evaluation over a task sequence.
The complexity decomposes into a within task adaptation term $\mathbb E_{P_t\sim\mathcal P_{t-1}}\mathrm{KL}(Q_t\Vert P_t)$ and a hyperdistribution drift term $\mathrm{KL}(\mathcal P_t\Vert\mathcal P_{t-1})$. The former quantifies how much the learner must deviate from the task prior to fit $S_t$, while the latter quantifies how the task prior generation mechanism evolves as new tasks are incorporated.
% The need for the drift term is evident in continual learning because the task prior at step $t$ is typically produced by a history dependent rule rather than fixed externally, 
This drift is intrinsic to continual learning because the prior used at step $t$ is typically produced by a history-dependent rule rather than fixed externally,
for example by specifying the admissible update scope~\citep{progressive2016,HAT2018,Supermasks2020,adapter2019,hu2022lora}, or by shaping the prior geometry using curvature summaries as regularization~\citep{EWC2017,SI2017}.
Theorem~\ref{thm:pathwise_pac} also contains the classical exchangeable-task hierarchical PAC-Bayes bound~\citep{pentinaPACBayesianBoundLifelong2014} as a special case: if the hyperposterior is not a sequence change and the result reduces to a single shared hyperdistribution over task priors. 

% More generally, when $\mathcal P_t$ is generated by a tractable parametric rule, the drift admits a closed-form expression and becomes a computable proxy for how strongly the learner updates its inductive bias across task transitions.


% \paragraph{Strict containment and computable drift under common prior generation rules.}
% Theorem~\ref{thm:pathwise_pac} strictly contains classical exchangeable task hierarchical PAC Bayes as a special case. If the hyperposterior is time invariant, namely $\mathcal P_t\equiv \mathcal P$ for all $t$, then the drift term vanishes and the bound reduces to a hierarchical guarantee with a single shared hyperdistribution over task priors, consistent with the exchangeable multi task view~\citep{pentinaPACBayesianBoundLifelong2014}. In continual learning, allowing $\mathcal P_t$ to evolve introduces the additional drift budget $\sum_{t=1}^T \mathrm{KL}(\mathcal P_t\Vert\mathcal P_{t-1})$, which becomes a measurable proxy for how aggressively the learner updates its inductive bias.

% Moreover, for broad classes of prior generation rules, the drift admits an explicit upper bound in terms of computable quantities, which yields testable predictions.
% A convenient instantiation is to parameterize the hyperposterior by a distribution over prior parameters $\theta$ (for example the mean, precision, or support specification of the task prior), and to select $\mathcal P_t$ from a tractable family $\{\mathsf H(\eta)\}$ indexed by hyperparameters $\eta_t$. Then the drift becomes
% $$
% \mathrm{KL}\bigl(\mathcal P_t\Vert\mathcal P_{t-1}\bigr)
% =
% \mathrm{KL}\bigl(\mathsf H(\eta_t)\Vert \mathsf H(\eta_{t-1})\bigr),
% $$
% which is closed form for standard families.
% For instance, if $\mathsf H(\eta)$ is Gaussian over a prior mean parameter $\theta$ with fixed covariance $\Sigma_{\mathcal P}$, namely $\mathcal P_t=\mathcal N(\mu_t,\Sigma_{\mathcal P})$, then
% $$
% \mathrm{KL}\bigl(\mathcal P_t\Vert\mathcal P_{t-1}\bigr)
% =
% \frac12(\mu_t-\mu_{t-1})^\top \Sigma_{\mathcal P}^{-1}(\mu_t-\mu_{t-1}),
% $$
% so the cumulative drift is controlled by the squared step to step change of the hyper mean.

% This yields two concrete, computable regimes that connect directly to common continual learning mechanisms.
% First, under Fisher based regularization, $\mu_t$ can be taken as a curvature informed summary (for example a moving estimate of a shared center in parameter space), and the bound predicts that stronger curvature constraints reduce the drift by suppressing large changes in $\mu_t$, which is verifiable by correlating the measured drift with forgetting and with curvature statistics.
% Second, under scope allocation and masking rules, the hyperparameters can include a representation of the admissible update support, such as a binary mask or a low dimensional subspace descriptor. Choosing a hyperposterior family whose parameters encode that support yields drift terms that scale with the amount of scope change across tasks. This predicts that methods that expand or reshuffle the admissible scope more aggressively incur larger drift, while methods that reuse a stable scope incur smaller drift, a distinction that can be tested by measuring drift alongside performance along the task trajectory.




% \subsection{From single-task bound to continual learning}




% % Let $\mathcal X$ be the input space, $\mathcal Y$ the label space, and $\mathcal Z = \mathcal X \times \mathcal Y$ the data space.
% % A model with parameters $W \in \mathbb R^d$ induces a predictor $f_W : \mathcal X \to \mathcal Y$.
% % We consider a bounded loss $\ell : \mathcal H \times \mathcal Z \to [0,1]$ and write $\ell(W,z)$ for $\ell(f_W,z)$. 
% % Given a distribution $\mathcal D$ over $\mathcal Z$ and a training set $S=\{z_j\}_{j=1}^m$ drawn i.i.d.\ from $\mathcal D$, the population and empirical risks are
% % $
% % L_{\mathcal D}(W)
% % =
% % \mathbb E_{z\sim\mathcal D}\ell(W,z),
% % \hat L_{S}(W)
% % =
% % \frac{1}{m}\sum_{j=1}^{m}\ell(W,z_{j}).
% % $
% % The PAC-Bayesian viewpoint place probability measures over the hypothesis space $\mathcal{H}$ rather than focusing on a single function: starting from a prior $P\in\mathcal M(\mathcal H)$, the learner observes data $S$ and outputs a posterior $Q(\cdot\mid P,S)\in\mathcal M(\mathcal H)$. 
% % %that represents the updated belief over $\mathcal H$ rather than focusing on a single function. From a Bayesian perspective, learning corresponds to updating a belief distribution over . From a Bayesian perspective, learning corresponds to updating a belief distribution over 
% % % The learner starts from a random prior $P\in\mathcal M(\mathcal H)$ and, after learning on $S$, updates it to a posterior $Q_t(\cdot\mid P_t,S_t)\in\mathcal M(\mathcal H)$. 
% % % In this view, learning on a task involves a hypothesis $f \in \mathcal H$ equipped with a prior distribution $P \in \mathcal M(\mathcal H)$, which is updated to a posterior $Q_t(\cdot\mid P_t,S_t)\in\mathcal M(\mathcal H)$ after observing data. 
% % The corresponding posterior-averaged population and empirical risks  are
% % % The object of learning on a task has been a hypothesis $f \in \mathcal{H}$ over which one presumes a \textit{prior distribution} $P\sim \mathcal{M}(\mathcal{H})$ that is updated into a \textit{posterior} $Q\sim \mathcal{M}(\mathcal{H})$ based on observed data. We have the corresponding risk as follow
% % $
% % L_{\mathcal D}(Q)
% % =
% % \mathbb E_{W \sim Q(\cdot \mid P,S)}\mathbb E_{z\sim\mathcal D}\ell(W,z)
% % $,
% % $
% % \hat L_{S}(Q_t)
% % =
% % \mathbb E_{W \sim Q(\cdot \mid P,S)}\frac{1}{m}\sum_{j=1}^{m_t}\ell(W,z_{j}).
% % $
% % A classical PAC--Bayes bound for a single task under a fixed distribution is stated as follows.
% % \begin{theorem}[PAC-Bayes Generalization Bound]
% % \label{theorem:classical}
% % For any loss $\ell: \mathcal{F} \times \mathcal{X} \times \mathcal{Y} \to [0,1]$, with probability at least $1-\delta$ over the choice of training set $S$:
% %     \[\mathbb{E}_{f \sim Q}[L_{\mathcal{D}}^{\ell}(f)] \leq \mathbb{E}_{f \sim Q}[\hat{L}_S^{\ell}(f)] + \sqrt{\frac{KL(Q\|P) + \log \frac{m}{\delta}}{2(m-1)}}\]
% % \end{theorem}

% % The classical PAC bound \cite{mcallesterPACBayesianModelAveraging1999} is
% % $L_{\mathcal D_t}(Q_t) \leq \hat L_{S_t}(Q_t) + \sqrt{\frac{KL(Q\|P) + \log \frac{m}{\delta}}{2(m-1)}}.$



% % \vspace{-10pt}

% % This section provides a generalization guarantee that makes the role of flatness explicit under the PECL constraint.
% % We proceed in three steps.
% % First, we state a generic {pathwise} hierarchical PAC Bayes bound that controls the time averaged test risk along a task sequence for CL.
% % Second, we specialize its complexity and empirical terms to PECL by restricting the posterior support to the LoRA update subspace.
% % Third, we combine and obtain a sharpness aware bound whose relevant quantities in the adapter subspace.







% % In parameter efficient continual learning (PECL) the model at task $t$ is written as
% % $$
% % W_t \;=\; W_{\mathrm{frozen}} + \Delta W_t,
% % $$
% % where $W_{\mathrm{frozen}}$ collects the pretrained backbone and all previously learned adapters, and $\Delta W_t$ is a new low rank update confined to a shared adapter subspace $\mathcal W_\Delta \subseteq \mathbb R^d$. This section develops a PAC–Bayesian analysis that isolates the role of \emph{adapter subspace flatness} in PECL.

% % We adopt a hierarchical PAC–Bayesian view with task level priors and posteriors sampled from a hyper posterior over tasks. A pathwise PAC–Bayes inequality for continual learning with time varying hyper posteriors is provided in Appendix~A (Theorem~A.1). It bounds the task averaged test risk by an empirical term evaluated under Gaussian perturbations and two KL complexity terms on task posteriors and hyper posteriors. In what follows we specialise this general result to the PECL parameterisation and show that both flatness and complexity can be written entirely in the adapter subspace.








% % To model the sequential nature of knowledge accumulation across tasks in CL, we adopt a hierarchical PAC‑Bayesian formulation. In this framework, the task-level prior itself is treated as a random variable, and a hyper-posterior is defined as a probability measure $\mathcal P \in \mathcal M(\mathcal M(\mathcal H))$ over such priors~\citep{pentinaPACBayesianBoundLifelong2014,JMLR:v24:22-1254}.
% % % To analyze the role of sharpness in PECL, we adopt a hierarchical PAC-Bayesian framework, where the prior $P_t$ itself as a random variable.
% % % Follow hyper PAC , we treat the prior $P_t$ itself as a random variable.
% % %We write $\mathcal F_t = \sigma(S_1,\dots,S_t)$ for the $\sigma$-algebra generated by the first $t$ datasets and consider the filtration $\mathcal F_0 \subseteq \mathcal F_1 \subseteq \cdots \subseteq \mathcal F_T$. 
% % % At the beginning of task $t$, before observing $S_t$, the learner maintains a hyper-posterior $\mathcal P_{t-1}$ that is $\mathcal F_{t-1}$-measurable and encodes all information extracted from previous tasks. 
% % Let $\mathcal F_{t-1}=\sigma(S_1,\dots,S_{t-1})$ denote the $\sigma$-algebra generated by the data observed up to task $t-1$. 
% % Before observing $S_t$, the learner maintains a hyper-posterior $\mathcal P_{t-1}$ that is $\mathcal F_{t-1}$-measurable, summarizing information transferred from previous tasks. 
% % The learning process at task $t$ is then described in two stages. First, a task-level prior $P_t$ is sampled from the current hyper-posterior,
% % $
% % P_t \sim \mathcal P_{t-1}.
% % $
% % Second, after observing $S_t$, the learner updates from $P_t$ to a task-level posterior $Q_t$, and the hyper-posterior is updated to $\mathcal P_t$. 
% % We assume that these updates are absolutely continuous, so that the associated KL divergences are well defined.
% % $
% % \mathcal P_t \ll \mathcal P_{t-1},
% % \quad
% % Q_t(\cdot \mid P,S_t) \ll P(\cdot)
% % \quad \text{for } \mathcal P_{t-1}\text{-almost every } P.
% % $
% % The corresponding task-level risks under the hierarchical posterior are obtained by averaging over the sampling of priors and hypotheses,

% % To capture the sequential nature of knowledge accumulation in continual learning, we adopt a hierarchical PAC Bayes formulation in which task level priors are themselves random objects.
% % Specifically, we consider a hyperposterior $\mathcal{P} \in \mathcal{M}(\mathcal{M}(\mathcal{H}))$ over task level priors $P \in \mathcal{M}(\mathcal{H})$~\citep{pentinaPACBayesianBoundLifelong2014,JMLR:v24:22-1254}.
% % However, existing work are primarily developed in a multi task perspective: they typically assume that task specific priors are independent draws from a \emph{fixed} hyperprior, and thus do not explicitly model the sequential dependence and information accumulation that characterize continual learning.
% % In particular, this independence assumption prevents the bound from capturing how the inductive bias evolves as a consequence of earlier tasks.


% % Theorem 1 provides a sharp characterization for a single task under a fixed data distribution, continual learning differs in two structural aspects.
% % In the CL, the learner faces a sequence of tasks $t\in\{1,\dots,T\}$ rather than only one task. 
% % At task $t$, a dataset
% % $
% % S_t = \{z_{tj}\}_{j=1}^{m_t}
% % $
% % is drawn i.i.d.\ from a task distribution $\mathcal D_t$. First, data arrive as a sequence of task distributions $\{\mathcal D_t\}_{t=1}^T$ and the learner updates its hypothesis after each dataset $S_t$.
% % Second, the inductive bias is itself shaped by the task history, since the choice of a task-level prior at time $t$ should reflect what has been learned from tasks $1,\dots,t-1$.
% % Consequently, applying a classical PAC-Bayes bound independently on each task treats the priors as externally specified and leaves knowledge accumulation outside the analysis.


% Theorem~\ref{theorem:classical} provides a generalization guarantee for a single task under a fixed data distribution and a prior that is fixed before observing the corresponding training set.
% However, continual learning proceeds over a task sequence $t\in\{1,\dots,T\}$, where at each step a dataset $S_t=\{z_{tj}\}_{j=1}^{m_t}$ is drawn i.i.d.\ from $\mathcal D_t$ and the learner updates its hypothesis after observing $S_t$.
% This sequential setting implies that the inductive bias should adapt with the task history, so the task prior at time $t$ must reflect information from tasks $1,\dots,t-1$, which cannot be captured by treating $\{P_t\}$ as externally specified constants.


% We therefore model transferable inductive bias by treating the task prior as a random object and introducing a hyperposterior $\mathcal P\in\mathcal M(\mathcal M(\mathcal H))$ over task-level priors. Figure~\ref{fig:showPAC3} highlights a key departure from classical hierarchical PAC-Bayes analyses for exchangeable multi-task environments~\citep{pentinaPACBayesianBoundLifelong2014}: there a single time-invariant hyperdistribution generates task priors, so the task order plays no explicit role. In continual learning, by contrast, the prior used at step $t$ is produced by a history-dependent rule that is updated as tasks arrive, because the learner must revise what is transferable based on the information extracted from earlier tasks.


% This sequential view is consistent with the information-budget interpretation of online PAC-Bayes, but it introduces a structurally different object to be tracked across tasks: not only the within-task update from $P_t$ to $Q_t$, but also the evolution of the \emph{prior-generation mechanism} that maps past experience to the next task prior. The resulting decomposition separates the complexity of fitting the current task from the cost of updating the inductive bias across task transitions. Such history dependence arises broadly in continual learning, for example when importance-based regularization updates the effective prior using curvature summaries~\citep{EWC2017,SI2017}, or when architectural and masking strategies modify which directions remain trainable over time~\citep{progressive2016,HAT2018,Piggyback2018,Supermasks2020,adapter2019,hu2022lora}. These mechanisms motivate modeling continual learning as a time-varying hyperposterior process and make an explicit drift penalty across task transitions unavoidable.
% % #a hierarchical PAC-Bayes formulation with
% % We therefore model transferable inductive bias by treating the task prior as a random object and introducing a hyperposterior $\mathcal P\in\mathcal M(\mathcal M(\mathcal H))$ over task-level priors.
% % Crucially, to represent knowledge accumulation as an explicit sequential update, we let the hyperposterior depend on the observed history and evolve with $t$, rather than remaining time invariant as in previous setting.
% % Crucially, we let the hyperposterior evolve with the observed history to represent knowledge accumulation as an explicit sequential belief update. In CL, the effective task prior is typically a history-dependent object: importance-based regularization (EWC\cite{EWC2017}, SI\cite{SI2017}) updates the prior using Fisher and Hessian summaries of past tasks; architectural or masking methods (ProgressiveNets\cite{progressive2016}, HAT\cite{HAT2018} , Piggyback\cite{Piggyback2018}, supermasks\cite{Supermasks2020} ) change the admissible update scope $\mathcal W_{\Delta,t}$ by allocating or restricting capacity. PECL methods, including adapters\cite{adapter2019} and LoRA~\cite{hu2022lora}, follow the same principle in a structured form, since introducing new factors or reorganizing modules makes the prior over adapter coordinates depend on the task history. These mechanisms all induce a nonstationary prior generation process, motivating a time-varying hyperposterior process in sequential environments.
% % Figure~\ref{fig: showPAC3} highlights a key distinction from classical hierarchical PAC-Bayes analyses\cite{pentinaPACBayesianBoundLifelong2014} developed for exchangeable multi-task environments: there, a single time-invariant hyperdistribution is used to generate task priors, and the task order plays no explicit role. In continual learning, by contrast, the prior used at step $t$ is typically produced by a history-dependent rule that evolves as new tasks arrive, because the learner updates what is transferable based on the information extracted from tasks $1,\dots,t-1$. This motivates modeling continual learning as a hyperposterior process that is updated sequentially,
% % % Crucially, we let the hyperposterior evolve with the observed history to represent knowledge accumulation as an explicit sequential belief update. 
% % In CL, the effective task prior is typically a history-dependent object: importance-based regularization \cite{EWC2017, SI2017} updates the prior using Fisher and Hessian summaries of past tasks; architectural or masking methods \cite{progressive2016, HAT2018,Piggyback2018,Supermasks2020} change the admissible update scope by allocating or restricting capacity. These mechanisms all induce a nonstationary prior generation process, motivating a time-varying hyperposterior process in sequential environments.




% % \[
% % \mathcal L_t
% % =
% % \mathbb E_{P_t \sim \mathcal P_t}
% % \,
% % \mathbb E_{W_t \sim Q_t(\cdot \mid P_t,S_t)}
% % \big[
% % L_{\mathcal D_t}(W_t)
% % \big],
% % \qquad
% % \hat{\mathcal L}_t
% % =
% % \mathbb E_{P_t \sim \mathcal P_t}
% % \,
% % \mathbb E_{W_t \sim Q_t(\cdot \mid P_t,S_t)}
% % \big[
% % \hat L_{S_t}(W_t)
% % \big].
% % \]
% % The corresponding hierarchical risks are obtained by averaging over both the sampled prior and the posterior draw of hypotheses,




% % Theorem~\ref{theorem:classical} provides a generalization guarantee for a single task with a fixed data distribution and a fixed prior.
% % However, continual learning differs in two structural aspects.
% % First, the learner faces a sequence of tasks $t\in\{1,\dots,T\}$.
% % At task $t$, a dataset $S_t=\{z_{tj}\}_{j=1}^{m_t}$ is drawn i.i.d.\ from a task distribution $\mathcal D_t$, and the learner updates its hypothesis after each dataset.
% % Second, the inductive bias is shaped by the task history, since the prior used at time $t$ should reflect what has been learned from tasks $1,\dots,t-1$.
% % A direct application of the single-task PAC--Bayes bound to each task would therefore require task priors $\{P_t\}$ that depend on the history.
% % Such a history-dependent choice is not captured by treating $\{P_t\}$ as externally specified constants. 

% % This motivates treating the task prior itself as a random object and introducing a hierarchical PAC--Bayes formulation in which transferable inductive bias is represented by a hyperposterior $\mathcal{P} \in \mathcal{M}(\mathcal{M}(\mathcal{H}))$ over task level priors $P \in \mathcal{M}(\mathcal{H})$. In sequential learning, the learner updates its hypothesis after each task, so the hyperposterior over task priors should generally depend on the observed task history rather than remain time invariant. 
% % However, prior work~\citep{pentinaPACBayesianBoundLifelong2014,JMLR:v24:22-1254} assume a single fixed hyperposterior,and the resulting bound does not include an hyperdistribution drift penalty.






% % A natural way to model such accumulation is the hierarchical PAC-Bayes formalism, where a hyperdistribution over task-level priors represents transferable inductive bias. However, existing hierarchical analyses typically infer a single time-invariant hyperposterior from a batch of observed tasks, which is appropriate under exchangeable task environments but does not reflect sequential belief updates in continual learning.







% % treats the task priors $\{P_t\}$ as externally specified and does not model knowledge accumulation.



% % A natural way to model such accumulation is the hierarchical PAC-Bayes formalism, where a hyperdistribution over task-level priors represents transferable inductive bias.
% % However, existing hierarchical analyses typically infer a single time-invariant hyperposterior from a batch of observed tasks, which is appropriate under exchangeable task environments but does not reflect sequential belief updates in continual learning. In particular, a fixed hyperposterior yields a task-level permutation-invariant characterization and, more importantly, it does not quantify the cost of adapting the inductive bias across task transitions.


% % To model knowledge accumulation in continual learning, we adopt a hierarchical PAC-Bayes formulation and consider a hyperposterior $\mathcal{P} \in \mathcal{M}(\mathcal{M}(\mathcal{H}))$ over task level priors $P \in \mathcal{M}(\mathcal{H})$.
% % % where task-level priors are treated as random variables drawn from a  hyperdistribution $\mathcal{P} \in \mathcal{M}(\mathcal{M}(\mathcal{H}))$. 
% % In sequential learning, the learner updates its hypothesis after each task, so the hyperposterior over task priors should generally depend on the observed task history rather than remain time invariant. 
% % However, prior work~\citep{pentinaPACBayesianBoundLifelong2014,JMLR:v24:22-1254} assume a single fixed hyperposterior,and the resulting bound does not include an hyperdistribution drift penalty.
% % that tracks how the hyperdistribution changes across tasks
% % produce an explicit hyperdistribution drift term.
% % and the bound is permutation invariant at the task level and does not yield an explicit hyperdistribution drift term.
% % , which is appropriate under exchangeable task environments but does not explicitly account for the cost of updating the inductive bias across task transitions.
% % Motivated by this gap, we extend this setting 
% % by allowing a history dependent hyperposterior that evolves with $t$.
% % % which introduces a KL drift penalty between consecutive hyperposteriors and makes sequential updates explicit. 
% % %This sequential template will later be specialized to PECL, where both the complexity and sharpness dependent empirical terms collapse to the adapter subspace under a frozen backbone.
% % %rom a batch of observed tasks and upper-bounds the transfer risk on a future unseen task.
% % % We extend this setting by retain the hyper distribution formalism while allowing a history dependent hyperposterior that evolves with $t$, which introduces a drift term and makes sequential updates explicit in the bound.
% % % . Specifically, we consider a hyperdistribution $\mathcal{P} \in \mathcal{M}(\mathcal{M}(\mathcal{H}))$ over task level priors $P \in \mathcal{M}(\mathcal{H})$.
% % % However, existing analyses~\citep{pentinaPACBayesianBoundLifelong2014,JMLR:v24:22-1254} assume a single fixed hyperposterior over task priors and derives a bound on the transfer risk of a future unseen task, which do not quantify the cost of sequentially updating the hyper distribution across tasks.
% % % in terms of the empirical multi task risk plus environment level and within task KL complexities.
% % % assume a fixed hyperdistribution, and the resulting bounds do not include an explicit task-to-task drift penalty that captures sequential updates.
% % % assume a fixed hyperdistribution, and therefore cannot quantify the cost of updating the hyperposterior across tasks.
% % % which yields bounds do not provide an explicit task-to-task drift penalty to model the sequential dependence.
% % % that characterize continual learning.
% % % In particular, this independence assumption prevents the bound from capturing how the inductive bias evolves as a consequence of earlier tasks.
% % % Motivated by this gap, we retain the hyper distribution formalism while allowing a history dependent hyperposterior that evolves with $t$, which introduces a drift term and makes sequential updates explicit in the bound.
% % % retain the hyperdistribution formalism while allowing the hyperposterior to evolve over time.
% % % Motivated by this gap, we employ the hyperdistribution formalism while allowing the hyperposterior to evolve over time.
% % % This yields a pathwise bound with an explicit drift term and provides the appropriate starting point for a continual learning.
% % % Let $\mathcal{F}_{t-1}=\sigma(S_1,\dots,S_{t-1})$ be the information available before observing task $t$.
% % % We model knowledge accumulation by allowing the hyperposterior $\mathcal{P}_{t-1}$ to be $\mathcal{F}_{t-1}$-measurable and to evolve with $t$.
% % Let $\mathcal{F}_{t-1} = \sigma(S_1,\dots,S_{t-1})$ denote the information available before observing task $t$.
% % In contrast to batch hierarchical analyses that infer a single hyperposterior from all observed tasks, we model continual learning as a \emph{hyperposterior process} $\{\mathcal{P}_t\}_{t=0}^T$ where $\mathcal{P}_{t-1}$ is $\mathcal{F}_{t-1}$ measurable and may change with $t$. 
% % This pathwise viewpoint leads to an explicit hyperdistribution drift penalty, which captures the cumulative complexity of sequential belief updates and provides an appropriate starting point for continual learning.
% % At step $t$, a task-level prior $P_t$ is drawn according to $P_t \sim \mathcal{P}_{t-1}$, and after observing $S_t$ the learner outputs a task posterior $Q_t(\cdot \mid P_t,S_t)$ and updates the hyperposterior to $\mathcal{P}_t$.

% Let $\mathcal F_{t-1}=\sigma(S_1,\dots,S_{t-1})$ denote the information available before observing task $t$.
% We model continual learning as a \emph{hyperposterior process} $\{\mathcal P_t\}_{t=0}^T$ where $\mathcal P_{t-1}$ is $\mathcal F_{t-1}$-measurable, which yields a pathwise bound with an explicit hyperposterior drift penalty that quantifies the information budget of sequential belief updates. We define the hierarchical test and empirical risks $
% \mathcal L_t
% =
% \mathbb E_{P_t \sim \mathcal{P}_{t-1}}
% \,
% \mathbb E_{W_t \sim Q_t(\cdot \mid P_t,S_t)}
% \big[
% L_{\mathcal D_t}(W_t)
% \big],
% \hat{\mathcal L}_t
% =
% \mathbb E_{P_t \sim \mathcal{P}_{t-1}}
% \,
% \mathbb E_{W_t \sim Q_t(\cdot \mid P_t,S_t)}
% \big[
% \hat L_{S_t}(W_t)
% \big].
% $
% We derive the following theorem and the proof is provided in Appendix~\ref{proof:theorem all}.
% \begin{theorem}[PAC-Bayes with time-varying hyperposterior]
% \label{thm:pathwise_pac}
% For any $\delta\in(0,1]$, with probability at least $1-\delta$ over $S_{1:T}$,
% \begin{equation*}
% \resizebox{\linewidth}{!}{$
% \begin{aligned} 
% &\quad\frac{1}{T} \sum_{t=1}^T \mathcal L_t
% \le  
% \frac{1}{T} \sum_{t=1}^T \hat{\mathcal L}_t  + \\
% &
% \sqrt{
% \frac{
% \displaystyle
% \sum_{t=1}^T
% \mathbb E_{P_t\sim\mathcal P_{t-1}}
% \mathrm{KL}\bigl(Q_t \,\Vert\, P_t\bigr)
% +
% \sum_{t=1}^T
% \mathrm{KL}\bigl(\mathcal P_t\Vert\mathcal P_{t-1}\bigr)
% +
% \log\frac{1}{\delta}}
% {2\displaystyle\sum_{t=1}^T (m_t - 1)}
% }.
% \end{aligned}
% $}
% \end{equation*}
% \end{theorem}


% % The left-hand side averages the hierarchical test risks $\mathcal{L}_t$ over the task sequence, while the first term on the right-hand side averages the corresponding empirical risks $\hat{\mathcal{L}}_t$.
% % The bound is stated for the trajectory average $\frac{1}{T}\sum_t$, and our experiments report average and final step metrics.
% % IF interpret $\mathcal D_t$ as the \emph{seen-task evaluation distribution} at step $t$, Under this choice, $\mathcal L_t$ represents the expected loss of the step-$t$ predictor on the same collection of seen tasks that defines $\operatorname{ACC}_t$. And $\frac{1}{T}\sum_t \hat{\mathcal L}_t$ is the correspond training loss in the CL learning sequence.
% Theorem~\ref{thm:pathwise_pac} bounds the trajectory-averaged hierarchical test risk by the corresponding trajectory-averaged empirical risk plus a complexity penalty.
% % , where averaging over $t=1,\dots,T$ matches the standard CL practice of reporting performance across the entire task sequence.
% Averaging over $t=1,\dots,T$ aligns with the standard continual learning protocol of summarizing performance over the full task sequence.
% % The complexity decomposes into a within task adaptation term $\mathbb E_{P_t\sim\mathcal P_{t-1}}\mathrm{KL}(Q_t\Vert P_t)$ and a hyperposterior drift term $\mathrm{KL}(\mathcal P_t\Vert\mathcal P_{t-1})$, where the latter quantifies how the distribution over task priors evolves over time.
% The complexity decomposes into a within-task adaptation term $\mathbb E_{P_t\sim\mathcal P_{t-1}}\mathrm{KL}(Q_t\Vert P_t)$ and a hyperposterior drift term $\mathrm{KL}(\mathcal P_t\Vert\mathcal P_{t-1})$.
% The former quantifies how much the learner must deviate from the task prior to fit $S_t$, while the latter quantifies how the task prior generation mechanism evolves as new tasks are incorporated. 


% The need for the hyperdrift term is evident in continual learning because the task prior at step $t$ is typically produced by a history-dependent prior generation rule rather than fixed externally. 
% This rule may specifies the admissible update scope\cite{progressive2016,HAT2018,Piggyback2018,Supermasks2020,adapter2019,hu2022lora}, the initialization and geometry of the task prior\cite{EWC2017}, and the constraints of historical information\cite{farajtabarOrthogonalGradientDescent2020,SI2017}.
% As tasks arrive, the rule must be updated and is captured explicitly by the cumulative drift term $\sum_{t=1}^T \mathrm{KL}(\mathcal P_t\Vert\mathcal P_{t-1})$.

% % The need for hyperdrift term is evident in common CL paradigms.


% % Importantly, this history dependence does not require direct access to past data. It arises because $\mathcal P_{t-1}$ is $\mathcal F_{t-1}$-measurable, so $\phi_t$ may depend on past tasks only through learned parameters and summaries of the learning trajectory, such as the previous mean and curvature-based statistics.



% % Theorem~\ref{thm:pathwise_pac} allows a time varying hyperposterior process $\{\mathcal P_t\}_{t=0}^T$ and makes sequential updates explicit through the cumulative drift penalty $\sum_{t=1}^T \mathrm{KL}(\mathcal P_t\Vert\mathcal P_{t-1})$. 
% % This term has a direct operational meaning: it measures the information injected into the prior-generation mechanism as the learner incorporates new tasks. The need for such a term is evident in common CL paradigms and is essential for LoRA-based PECL. where changes in LoRA organization across tasks and the initialization of newly introduced adapter factors induce a nonstationary prior generation mechanism that cannot be faithfully summarized by a single stationary hyperposterior. 


% % The remaining terms constitute the complexity penalty: $\sum_{t=1}^T \mathbb{E}_{P_t \sim \mathcal{P}_{t-1}}\mathrm{KL}(Q_t|P_t)$ quantifies within-task adaptation from prior to posterior, $\sum_{t=1}^T \mathrm{KL}(\mathcal{P}_t|\mathcal{P}_{t-1})$ measures the drift of the hyperposterior and thus the evolution of the inductive bias across tasks, and the last term is the confidence penalty.

% % Compared with classical hierarchical PAC-Bayes bounds that use a fixed hyperprior and a single shared hyperposterior, Theorem~\ref{thm:pathwise_pac} allows the hyperposterior to vary across tasks and introduces a cumulative drift penalty that makes sequential updates explicit.

% % Compared with classical hierarchical PAC-Bayes bounds\cite{pentinaPACBayesianBoundLifelong2014} that use a fixed hyperprior and a single shared hyperposterior, 


% % In the special case where one uses a single hyperposterior, the drift term collapses to the classical hierarchical PAC-Bayes bounds\cite{pentinaPACBayesianBoundLifelong2014}.

% % Compared with classical hierarchical PAC-Bayesian analyses that use a fixed hyperprior, Theorem~\ref{thm:pathwise_pac} permits a task-independent hyperposterior and thus yields an explicit cumulative drift penalty, making sequential updates visible at the level of the bound.
% % In the special case where one uses a single hyperposterior $\mathcal P_0$, the drift term collapses to the complexity $\mathrm{KL}(\mathcal P \Vert \mathcal P_0)$, consistent with classical bound\cite{pentinaPACBayesianBoundLifelong2014}. 
% % In contrast, allowing $\mathcal P_t$ to change across tasks incurs additional drift terms that quantify the cost of sequentially adapting the hyperposterior.

% % Compared with classical hierarchical PAC-Bayesian bounds based on a fixed hyperprior, Theorem~\ref{thm:pathwise_pac} permits a time-varying hyperposterior and therefore yields an explicit drift penalty, making sequential updates visible at the level of the bound. 
% % In the special case where the hyperposterior is time invariant, the drift term disappears and the result reduces to the standard exchangeable-task formulation with a single hyperposterior, consistent with classical lifelong learning guarantees.
% % In the special case where one uses a single hyperposterior $\mathcal Q$ for all tasks after an initial update, namely $\mathcal P_1=\cdots=\mathcal P_T=\mathcal Q$ with a fixed hyperprior $\mathcal P_0$, the drift term collapses to the environment-level complexity $\mathrm{KL}(\mathcal Q \Vert \mathcal P_0)$, consistent with classical lifelong learning guarantees. In contrast, allowing $\mathcal P_t$ to change across tasks incurs additional drift terms that quantify the cost of sequentially adapting the hyperposterior.



% % hierarchical PAC Bayes bounds that draw task priors from a fixed hyperprior, Theorem~\ref{thm:pathwise_pac} allows a time varying hyperposterior and introduces an explicit drift term, which makes sequential information accumulation explicit in the bound. When the hyperposterior is time invariant, the drift term vanishes and our sequential template reduces to a standard exchangeable task environment bound with a single hyperposterior, matching the classical lifelong learning formulation.

% % In contrast of existing hierarchical PAC-Bayes bounds draw task priors from a fixed hyperprior,
% % Theorem~\ref{thm:pathwise_pac} allows a time-varying hyperposterior and introduces an explicit drift term, which is the key ingredient needed to model information accumulation in continual learning.



% % The left-hand side is the time average of the hierarchical test risks along the learning trajectory.
% % The first term on the right-hand side is the corresponding time average of posterior-averaged empirical risks.
% % The remaining terms form the complexity penalty, and Fig.~\ref{fig: showPAC3} summarizes this decomposition. The term
% % $\sum_{t=1}^T \mathbb{E}_{P_t \sim \mathcal{P}_t}\mathrm{KL}\!\big(Q_t(\cdot\mid P_t,S_t)\,\|\,P_t\big)$
% % measures how much the task posterior must deviate from its task prior to fit $S_t$.
% % It quantifies the amount of adaptation performed at each step.
% % The drift term
% % $\sum_{t=1}^T \mathrm{KL}\!\big(\mathcal{P}_t\,\|\,\mathcal{P}_{t-1}\big)$
% % captures how strongly the distribution over priors changes across tasks. This is the key pathwise ingredient that is absent under an independence assumption with a fixed hyperprior, and it explicitly reflects sequential dependence through an evolving inductive bias.
% % The final term is the standard confidence penalty.


% % Theorem~\ref{thm:pathwise_pac} upper bounds a \emph{time average of per step generalization risks} along the learning trajectory. Figure~\ref{fig:showPAC3} summarizes the structure of the bound.
% % The first term on the right hand side is the time average of the hierarchical empirical risks $\hat{\mathcal{L}}_t$, obtained by averaging the empirical loss over both the sampled prior $P_t$ and the posterior draw $W_t$.
% % The square root term is the complexity penalty.
% % It decomposes into:
% %  \textbf{Taskwise complexity} $\mathrm{KL}^{\mathrm{task}}$:
% % for each step $t$, the divergence
% % $
% % \mathbb{E}_{P_t \sim \mathcal{P}_t}\mathrm{KL}\!\bigl(Q_t(\cdot \mid P_t,S_t)\Vert P_t\bigr)
% % $
% % measures how much the task posterior must deviate from its task prior to fit $S_t$.
% % This term quantifies within task adaptation.
% %  \textbf{Hyperposterior drift} $\mathrm{KL}^{\mathrm{hyper}}$:
% % the sum
% % $
% % \sum_{t=1}^T \mathrm{KL}\!\bigl(\mathcal{P}_t \Vert \mathcal{P}_{t-1}\bigr)
% % $
% % measures how strongly the distribution over priors changes across tasks.
% % Unlike prior hierarchical PAC Bayes analyses that draw priors independently from a fixed hyperprior, this explicit drift term captures sequential dependence through a time evolving inductive bias.




% % Prior PAC Bayes analyses in lifelong and multi task learning bound averages across tasks in the full parameter space.
% % Theorem~\ref{thm:pathwise_pac} follows the same general philosophy, but it is stated in a pathwise form with a time varying hyperposterior, making sequential dependence explicit through the drift term.


% % $\sum_t \mathrm{KL}(\mathcal{P}_t\Vert\mathcal{P}_{t-1})$

% % \subsection{KL decomposition under a frozen backbone}  $\sum_t \mathbb{E}_{P_t\sim\mathcal{P}_{t-1}}\mathrm{KL}(Q_t\Vert P_t)$

% % \begin{figure}[t]
% % \centering
% % % \begin{minipage}{1\linewidth}
% % % \centering
% % \includegraphics[width=\linewidth]{figures_TRML/showPAC3.pdf}
% % % \caption{Illustration of the PAC bound.}
% % \label{fig: showPAC3}
% % % \end{minipage}
% % % \vspace{-10pt}
% % \caption{Hierarchical PAC--Bayes with a time-varying hyperposterior. At each task $t$, the hyperposterior $\mathcal{P}_{t-1}$ induces a task prior $P_t$, which is updated to a task posterior $Q_t$ after observing $S_t$. The bound controls the time-averaged test risk $\frac{1}{T}\sum_t \mathcal{L}_t$ by the time-averaged empirical risk $\frac{1}{T}\sum_t \hat{\mathcal{L}}_t$ plus two complexity components: the hyperposterior drift  and the within-task divergence, scaled by sample size and a confidence term .}
% % \end{figure}


% \begin{figure}[t]
% \centering
% % \begin{minipage}{1\linewidth}
% % \centering
% \includegraphics[width=\linewidth]{figures_TRML/hyperbiyao3.pdf}
% % \caption{Illustration of the PAC bound.}
% \label{fig: showPAC3}
% % \end{minipage}
% % \vspace{-10pt}
% \caption{Why continual learning calls for a time-varying hyperposterior. \textbf{Left:} in the exchangeable multi-task view, task priors are generated from a single time-invariant hyperdistribution, so the inductive bias is shared but does not evolve along the task order. \textbf{Right:} in continual learning, the task prior at step $t$ is history dependent and is produced by updating the hyperposterior after observing past tasks, which makes the evolution of the prior-generation mechanism explicit and motivates an additional drift penalty across task transitions.}
% \end{figure}















\subsection{From Full-Finetune to LoRA}
\label{sec:kl-decomp-pecl}





% For each task $t$, the PECL parameterisation writes
% $$
% W_t \;=\; W_{\mathrm{frozen}} + \Delta W_t,
% \qquad
% \Delta W_t \in \mathcal W_\Delta \subseteq \mathbb R^d,
% $$
% where $\mathcal W_\Delta$ is the linear subspace of admissible adapter updates. We can therefore use $\Delta W$ as coordinates on the affine set $W_{\mathrm{frozen}} + \mathcal W_\Delta$.

% Given a prior $P_t^\Delta$ and a posterior $Q_t^\Delta$ on the adapter subspace, we define the corresponding distributions on full parameters simply by translation
% \begin{equation*}
% \begin{aligned}
%         W \sim P_t^{\mathrm{full}}
% \iff
% W = W_{\mathrm{frozen}} + \Delta W,
% \quad \Delta W \sim P_t^\Delta, \\
% \quad
% W \sim Q_t^{\mathrm{full}}
% \iff
% W = W_{\mathrm{frozen}} + \Delta W,
% \quad \Delta W \sim Q_t^\Delta.
% \end{aligned}
% \end{equation*}

% Let $\mathcal W_\Delta\subseteq\mathbb R^d$ be a linear subspace and let $W+\mathcal W_\Delta$ be its affine translate. 
% Equip both sets with the trace Borel $\sigma$-algebra induced from $\mathbb R^d$. 
% $\mathcal B(\mathcal W_\Delta)\ :=\{B\cap \mathcal W_\Delta:\ B\in \mathcal B(\mathbb R^d)\},\quad
% \mathcal B(W+\mathcal W_\Delta)\ :=\{B\cap (W+\mathcal W_\Delta):\ B\in \mathcal B(\mathbb R^d)\}.$
% These coincide with the Borel $\sigma$-algebras generated by the subspace topologies on $\mathcal W_\Delta$ and $W+\mathcal W_\Delta$, respectively, so both are standard Borel spaces.
% To formalize the probabilistic relationship between the LoRA subspace and the full model space, we define the translation map:
% % \begin{align*}
% $T_W: (\mathcal W_\Delta,\mathcal B(\mathcal W_\Delta)) \to (W+\mathcal W_\Delta,\mathcal B(W+\mathcal W_\Delta)), 
% T_W(\Delta)=W+\Delta,$
% % \end{align*}
% which is a Borel isomorphism between the subspace and affine space.  The full-space prior and posterior are defined as pushforwards onto the affine subspace by
% $
% P^{\mathrm{full}}:=(T_W)_{\text{\#}} P^{\Delta},
% Q^{\mathrm{full}}:=(T_W)_{\text{\#}} Q^{\Delta},
% $
% where $P_t^{\Delta}$ and $Q_t^{\Delta}$ are probability measures on $(\mathcal{W}\Delta, \mathcal{B}(\mathcal{W}\Delta))$.
% This construction is introduced to ensure that the PAC-Bayesian complexity term in the full parameter space can be expressed entirely in the adapter subspace under the PECL constraint. 
% For any measurable $F$, the pushforward $F_{\text{\#}}\mu$ is defined by $(F_{\#}\mu)(A)=\mu(F^{-1}(A))$. 
% This pushforward construction leads to a fundamental simplification of the KL divergence between full-model distributions.
% Theorem~\ref{thm:pathwise_pac} provides a generic sequential hierarchical PAC-Bayes template. 
% To make it informative for PECL, we impose the frozen-backbone constraint and exploit the fact that all task-dependent variability is confined to the LoRA update subspace.

% Theorem~\ref{thm:pathwise_pac} is a generic sequential hierarchical template. We enforce the PECL constraint and show that both the PAC Bayes complexity and the sharpness dependent empirical term collapse to the adapter update subspace.


% Sharpness are typically defined over the full parameter space and are aligned with full fine-tuning, where future adaptation may move in arbitrary directions.
% In frozen-backbone PECL, however, the model can only change through the adapter update.
% This induces a task-dependent set of \emph{task-admissible} update directions, and only variations along these directions are actionable under continual learning.
% As a result, both the PAC-Bayes complexity and the sharpness-dependent empirical contribution should be evaluated in the geometry restricted to the task-admissible update subspace.


% \paragraph{PECL parameterization and task-admissible subspace.}
% In PECL, the parameter decomposes as
% $
% W_t = W_{\mathrm{frozen}} + \Delta W_t,
% $
% where $W_{\mathrm{frozen}}\in\mathbb{R}^d$ collects all parameters frozen while learning task $t$, and $\Delta W_t$ is the trainable increment satisfies $\Delta W \in \mathcal W_{\Delta}\subseteq\mathbb{R}^d$, with $\mathcal W_{\Delta,t}$ denoting the task-$t$ admissible update subspace induced by the chosen LoRA organization. Consequently, the PECL hypothesis class at task $t$ is the affine set $ W_{\mathrm{frozen}} +  \mathcal W_{\Delta}\subseteq\mathbb{R}^d$ and $W_t$ lies in this affine model space.
% % We denote by $\mathcal W_{\Delta}\subseteq\mathbb{R}^d$ the linear subspace of admissible increments at task $t$ induced by the chosen LoRA organization.
% The following lemma shows that, under PECL, the PAC-Bayes complexity on the induced affine model space $W_{\mathrm{frozen}}+\mathcal W_{\Delta,t}$ reduces exactly to the KL divergence defined on the task-admissible update subspace $\mathcal W_{\Delta,t}$.
% the PAC-Bayes complexity in the full parameter space reduces to the KL divergence on the LoRA update subspace.


In PECL, the task-$t$ parameter vector can be written as
$
W_t = W^{\mathrm{froz}}_t + \Delta W_t,
$
where $W^{\mathrm{froz}}_t\in\mathbb R^d$ collects all parameters frozen while learning task $t$, and $\Delta W_t$ is the trainable increment induced by the current-task LoRA coordinates. The feasible update dynamics are thus restricted to an \emph{increment subspace}
$
\Delta W_t \in \mathcal W_{\Delta,t}\subseteq \mathbb R^d,
$
where $\mathcal W_{\Delta,t}$ denotes the task-admissible linear space of \emph{increment directions} around the learned
solution (formalized in Appendix \ref{app:lora_subspace_bridge}). Consequently, the hypothesis class at task $t$ is the affine
set $W^{\mathrm{froz}}_t + \mathcal W_{\Delta,t}$, and all task-dependent randomness in the PAC--Bayes analysis is supported
on $\mathcal W_{\Delta,t}$ and then pushed forward to the affine class by translation. The proof is deferred to Appendix~\ref{proof:Lemma1}.
% We propose the following lemma to formalize the resulting translation-invariance of the KL divergence, and its proof is included in Appendix~\ref{proof:Lemma1}.
\begin{lemma}
\label{lem:kl-decomp-pecl}
Assume $Q_t^\Delta \ll P_t^\Delta$. Then $Q_t^{\mathrm{full}} \ll P_t^{\mathrm{full}}$ and
$$
\mathrm{KL}\bigl(Q_t^{\mathrm{full}} \,\Vert\, P_t^{\mathrm{full}}\bigr)
\;=\;
\mathrm{KL}\bigl(Q_t^\Delta \,\Vert\, P_t^\Delta\bigr).
$$
\end{lemma}
% Under PECL, the posterior can only deviate from the prior through the update coordinate $\Delta W_t$, so each taskwise KL term in Theorem~\ref{thm:pathwise_pac} is fully determined by the distribution of the update $\Delta W$. The effective complexity scale is governed by $r_t$ and by the amount of posterior drift within $\mathcal W_t$. In contrast, a full-space update allows posterior mass to move in more directions and the resulting KL necessarily budgets additional information for those directions.
Lemma~\ref{lem:kl-decomp-pecl} shows that, under a frozen backbone, the PAC-Bayes complexity in the full parameter space is exactly the KL divergence on the task-admissible update subspace. 
% This reduction is structural rather than a matter of dimensionality.
% In PECL, future tasks can only modify the model through directions in $\mathcal W_{\Delta}$, so any notion of continual complexity must be budgeted within the same task-admissible geometry.
% Consequently, controlling generalization under continual adaptation requires controlling how the posterior drifts \emph{within} $\mathcal W_{\Delta}$, while directions outside $\mathcal W_{\Delta}$ are neither learnable nor actionable for subsequent tasks.
% We denote by $\mathcal{W}_t \subseteq \mathbb{R}^d$ the set of admissible LoRA increments at task $t$, viewed as a linear subspace embedded in the full parameter space.
% , and let $r_t:=\dim(\mathcal{W}_t)$ denote the number of trainable LoRA parameters.
% Concretely, for a layer weight matrix $W_\ell\in\mathbb{R}^{d_o^{(\ell)}\times d_i^{(\ell)}}$, LoRA parameterizes the increment as
% $
% \Delta W_\ell \;=\; B_\ell A_\ell ,
% $
% so perturbations in the factor coordinates induce perturbations of the full increment $\Delta W_t\in \mathcal{W}_t$.
% In particular, a Gaussian posterior over the factor coordinates induces a Gaussian distribution over $\Delta W_t$ supported on $\mathcal{W}_t$.
% Under PECL, the full parameter vector is therefore restricted to the affine set
% $
% W_t^{\mathrm{fr}} + \mathcal{W}_t
% =
% \{\, W_t^{\mathrm{fr}} + \Delta W : \Delta W \in \mathcal{W}_t \,\}.
% $
% Define the translation map
% $
% T_t:\mathcal{W}_t \to \mathbb{R}^d,\quad
% T_t(\Delta W) := W_t^{\mathrm{fr}} + \Delta W.
% $
% Given a prior $P_t$ and a posterior $Q_t$ on $\mathcal{W}_t$, we lift them to full-space measures via pushforward:
% $
% P_t^{\mathrm{full}} := (T_t)_{\#} P_t,
% \quad
% Q_t^{\mathrm{full}} := (T_t)_{\#} Q_t.
% $
% Fix a task index $t$.
% In PECL, parameters are partitioned into a frozen component and a trainable LoRA update.
% Let $W_{\mathrm{frozen},t}\in\mathbb{R}^d$ denote the collection of parameters that are held fixed while learning task $t$, and let $\mathcal{W}_\Delta\subseteq\mathbb{R}^d$ be the linear subspace of admissible LoRA updates at task $t$. We emphasize that $\mathcal W_\Delta$ denotes the LoRA-induced update subspace in the full parameter space, whose dimension is determined by the number of trainable LoRA parameters rather than the rank $r$.  LoRA realizes $\mathcal W_\Delta$ through low-rank factors. For a layer weight matrix $W_\ell \in \mathbb R^{d_o^{(\ell)}\times d_i^{(\ell)}}$, LoRA introduces the update $\Delta W_\ell = B_\ell A_\ell$. Perturbations of the factor coordinates  $\theta$ therefore induce perturbations $\Delta W\in\mathcal W_\Delta$, and a Gaussian posterior on factor coordinates $\theta$ yields a Gaussian posterior on $\Delta W$ supported on $\mathcal W_\Delta$.
% The full parameter vector is restricted to the affine set
% $
% W_{\mathrm{frozen},t} + \mathcal{W}_\Delta=
% \{W_{\mathrm{frozen},t} + \Delta W : \Delta W \in \mathcal{W}_\Delta\}.
% $
% % We equip $\mathcal{W}_\Delta$ and $W_{\mathrm{frozen},t}+\mathcal{W}_\Delta$ with their Borel $\sigma$-algebras inherited from $\mathbb{R}^d$.
% Define the translation map
% $
% T_t(\Delta W) := W_{\mathrm{frozen},t} + \Delta W.
% $
% % which is a measurable bijection.
% Given a prior $P_t^\Delta$ and a posterior $Q_t^\Delta$ on $\mathcal{W}_\Delta$, we induce the corresponding full space measures by pushforward:
% $
% P_t^{\mathrm{full}} := (T_t)_{\#}P_t^\Delta,
% \quad
% Q_t^{\mathrm{full}} := (T_t)_{\#}Q_t^\Delta.
% $
% This construction makes the PECL constraint explicit at the level of probability measures:
% all randomness is confined to the LoRA update coordinates $\Delta W\in\mathcal{W}_\Delta$.
% , as established in the following lemma.
% In contrast, allowing the backbone to move introduces additional degrees of freedom and cannot decrease the full-space KL for a fixed update marginal, thereby formalizing the complexity control advantage of PECL over full fine-tuning.
% Lemma~\ref{lem:kl-decomp-pecl} implies that the frozen parameters contribute no divergence.
% Consequently, every taskwise KL term  in Theorem \ref{thm:pathwise_pac} is fully determined by the distribution of $\Delta W$. Under PECL, the induced full-space measures satisfy Lemma~\ref{lem:kl-decomp-pecl}, hence the taskwise complexity equals the KL on the update subspace.
% Therefore, holding the update marginals fixed, allowing the backbone to move can only increase the complexity term, formalizing a strict complexity control advantage of PECL relative to full fine-tuning.
% \begin{corollary}[Complexity advantage over full fine-tuning]
% \label{cor:pecl_vs_ft_kl}
% Let $Q_t^{\mathrm{ft}}$ and $P_t^{\mathrm{ft}}$ be any posterior and prior on the full parameter space used for full fine-tuning.
% Let $\pi$ be a measurable map that extracts the update component on $\mathcal W_\Delta$ (for instance, the coordinate projection induced by the decomposition into frozen and update coordinates).
% Then
% \[
% \mathrm{KL}\bigl(Q_t^{\mathrm{ft}} \,\Vert\, P_t^{\mathrm{ft}}\bigr)
% \;\ge\;
% \mathrm{KL}\bigl(\pi_{\#}Q_t^{\mathrm{ft}} \,\Vert\, \pi_{\#}P_t^{\mathrm{ft}}\bigr).
% \]
% Under PECL, the induced full-space measures satisfy Lemma~\ref{lem:kl-decomp-pecl}, hence the taskwise complexity equals the KL on the update subspace.
% Therefore, holding the update marginals fixed, allowing the backbone to move can only increase the complexity term, formalizing a strict complexity control advantage of PECL relative to full fine-tuning.
% \end{corollary}
% The full model distributions $P_t^{\mathrm{full}}$ and $Q_t^{\mathrm{full}}$ are pushforwards of $P_t^\Delta$ and $Q_t^\Delta$ under the affine map $\Delta W \mapsto W_{\mathrm{frozen}} + \Delta W$. Since this map is bijective on $W_{\mathrm{frozen}} + \mathcal W_\Delta$ and preserves Radon–Nikodym derivatives, the relative entropy is invariant. In PECL, the KL complexity is therefore completely determined by the distribution of the adapter update $\Delta W\in\mathcal W_\Delta$; the frozen backbone contributes no divergence.
% \subsection{Expected sharpness restricted to the adapter subspace}
% \label{sec:subspace-sharpness-pecl}
% We specialise the posterior to a Gaussian supported on $\mathcal W_\Delta$. For task $t$ We consider a Gaussian posterior \( Q_t^{\Delta} = \mathcal{N}(\hat{\Delta W}_t, \Sigma_t) \) on the measurable space \( (\mathcal{W}_\Delta, \mathcal{B}(\mathcal{W}_\Delta)) \), where \( \hat{\Delta W}_t \in \mathcal{W}_\Delta \) is the mean and \( \Sigma_t\) is a covariance on the adapter subspace \( \mathcal{W}_\Delta \). Sampling from \( \Delta W \sim Q_t^{\Delta} \) is equivalent to 
% $
% \Delta W = \hat{\Delta W}_t + \varepsilon, \quad \varepsilon \sim \mathcal{N}(0, \Sigma_t) \ \text{on } \mathcal{W}_\Delta.
% $ This Gaussian posterior formulation enables a decomposition of the empirical risk into nominal performance and subspace-restricted sharpness.
% Under PECL, the model can only be modified through the LoRA increment, so all task-dependent randomness is naturally supported on the task-admissible update subspace $\mathcal W_{\Delta,t}$.


We specialize the adapter posterior to a Gaussian distribution on $\mathcal W_{\Delta,t}$,
$
Q_t^\Delta = \mathcal N(\hat{\Delta W}_t,\Sigma_t),
\hat{\Delta W}_t \in \mathcal W_{\Delta,t},
$
where $\Sigma_t$ is a covariance operator defined on $\mathcal W_{\Delta,t}$.
Equivalently, sampling $W\sim Q_t^{\mathrm{full}}$ supported on $W_{\mathrm{frozen}}+\mathcal W_{\Delta,t}$ amounts to sampling
% Equivalently, sampling $W\sim Q_t^{\mathrm{full}}$ amounts to sampling
% {%
% \setlength{\abovedisplayskip}{2pt}
% \setlength{\belowdisplayskip}{4pt}
% \setlength{\abovedisplayshortskip}{3pt}
% \setlength{\belowdisplayshortskip}{4pt}
% \[
% \Delta W = \hat{\Delta W}_t + \varepsilon,
% \quad
% \varepsilon \sim \mathcal N(0,\Sigma_t) \text{ on } W_{\Delta,t}.
% \]
% }%
$\Delta W = \hat{\Delta W}_t + \varepsilon,
\quad
\varepsilon \sim \mathcal N(0,\Sigma_t) \text{ on } W_{\Delta,t}.$
and setting $W=W_{\mathrm{frozen}}+\Delta W$. We define the adapter subspace sharpness as the expected loss elevation under perturbations supported on the adapter update subspace $\mathcal W_{\Delta,t}$ for task $t$.  Lemma~\ref{lem:subspace-loss-pecl} isolates an additive loss elevation caused by perturbations supported on $\mathcal W_{\Delta,t}$ and proof is deferred to Appendix~\ref{proof:lemma2}.

% $$
% \Delta W = \hat{\Delta W}_t + \varepsilon,
% \quad
% \varepsilon \sim \mathcal N(0,\Sigma_t),
% \quad
% \varepsilon \in \mathcal W_{\Delta,t},
% $$
% and setting $W = W_{\mathrm{frozen}} + \Delta W$.
% This choice is aligned with continual learning under PECL: future tasks can only realize variations in $\mathcal W_{\Delta,t}$, hence any within-task sensitivity relevant for generalization should be measured in this task-admissible geometry.
% We specialize the posterior on $\mathcal{W}_\Delta$ to a Gaussian distribution,
% $
% Q_t^\Delta = \mathcal{N}(\hat{\Delta W}_t,\Sigma_t),
% \quad
% \hat{\Delta W}_t \in \mathcal{W}_\Delta,
% $
% where $\Sigma_t$ is a covariance operator on $\mathcal{W}_\Delta$.
% Equivalently, $\Delta W = \hat{\Delta W}_t + \varepsilon$ with $\varepsilon\sim\mathcal{N}(0,\Sigma_t)$ supported on $\mathcal{W}_\Delta$.
% The PAC-Bayes empirical contribution involves the posterior expected empirical loss.
% For PECL posteriors supported on $\mathcal W_{\Delta,t}$, it admits a decomposition into a nominal empirical loss and a sharpness-like elevation induced by admissible perturbations.

% The proof is deferred to Appendix~\ref{proof:lemma2}.
% The following lemma decomposes the posterior expected empirical loss into a nominal term plus an additive increase induced by perturbations inside $\mathcal{W}_\Delta$.

% We propose the following lemma for this decomposition, and its proof is included in Appendix~\ref{proof:lemma2}.
\begin{lemma}[Task adapter-subspace sharpness]
\label{lem:subspace-loss-pecl}
For each task $t$,
{%
\setlength{\abovedisplayskip}{2pt}
\setlength{\belowdisplayskip}{2pt}
\setlength{\abovedisplayshortskip}{2pt}
\setlength{\belowdisplayshortskip}{2pt}
\begin{equation*}
\label{eq:subspace-loss-pecl}
\begin{aligned}
&\mathbb E_{W\sim Q_t^{\mathrm{full}}}\bigl[\hat L_{S_t}(W)\bigr] \\
&= \mathbb{E}_{\varepsilon\sim\mathcal N(0,\Sigma_t)}
    \big[\hat{L}_{S_t}(W_\mathrm {frozen }+\hat{\Delta W}_t+\varepsilon)\big] \\
& =\hat L_{S_t}(W_{\mathrm{frozen}} + \hat{\Delta W}_t)
+
\mathrm{TS}_t^{\Delta}(\hat{\Delta W}_t,\Sigma_t).
\end{aligned}
\end{equation*}
}where the task adapter subspace sharpness $TS_t^{\Delta}$ is $
\mathrm{TS}_t^{\Delta}(\hat{\Delta W}_t,\Sigma_t)
= \mathbb{E}_{\varepsilon\sim\mathcal N(0,\Sigma_t)}
   \big[\hat{L}_{S_t}(W_\mathrm {frozen }+\hat{\Delta W}_t+\varepsilon)
   - \hat{L}_{S_t}(W_\mathrm {frozen }+\hat{\Delta W}_t)\big],\quad
     \varepsilon\sim\mathcal N(0,\Sigma_t) \in  \mathcal W_\Delta$.
     % \begin{equation*}
     % \begin{aligned}
     % &\mathrm{E\text{-}Sh}_t^{\Delta}(\hat{\Delta W}_t,\Sigma_t):= \\
     % &\qquad \mathbb{E}_{\varepsilon\sim\mathcal N(0,\Sigma_t)}
     %    \big[\hat{L}_{S_t}(W_\mathrm {frozen }+\hat{\Delta W}_t+\varepsilon) \\
     % & \qquad- \hat{L}_{S_t}(W_\mathrm {frozen }+\hat{\Delta W}_t)\big] \\
     % &\qquad \varepsilon\sim\mathcal N(0,\Sigma_t) \in \mathcal W_\Delta
     % \end{aligned}
     % \end{equation*}
     \end{lemma}
     % Lemma~\eqref{lem:subspace-loss-pecl} shows that, within a single task, the posterior expected empirical loss differs from the nominal empirical loss only through admissible perturbations in $\mathcal W_{\Delta,t}$.

     
     % We  refer to $\mathrm{TS}_t^\Delta(\hat{\Delta W}_t,\Sigma_t)$ as the \emph{task-admissible expected sharpness}. This notion is consistent with PECL and continual learning, since it measures sensitivity along the directions that tasks can effectively realize.


% The following lemma decomposes the posterior expected empirical loss into a nominal term plus an additive increase induced by perturbations inside $\mathcal{W}_\Delta$.
% The proof is deferred to Appendix~\ref{proof:lemma2}.
% % We propose the following lemma for this decomposition, and its proof is included in Appendix~\ref{proof:lemma2}.
% \begin{lemma}
% \label{lem:subspace-loss-pecl}
% For each task $t$,
% \begin{equation*}
% \label{eq:subspace-loss-pecl}
% \begin{aligned}
% &\mathbb E_{W\sim Q_t^{\mathrm{full}}}\bigl[\hat L_{S_t}(W)\bigr] \\
% &= \mathbb{E}_{\varepsilon\sim\mathcal N(0,\Sigma_t)}
%     \big[\hat{L}_{S_t}(W_\mathrm {frozen }+\hat{\Delta W}_t+\varepsilon)\big],\varepsilon\in\mathcal W_\Delta \\
% & =\hat L_{S_t}(W_{\mathrm{frozen}} + \hat{\Delta W}_t)
% +
% \mathrm{T-S}_t^{(0)\Delta}(\hat{\Delta W}_t,\Sigma_t),\varepsilon\in\mathcal W_\Delta.
% \end{aligned}
% \end{equation*}
% where $
% T-S_t^{\Delta}(\hat{\Delta W}_t,\Sigma_t)
% := \mathbb{E}_{\varepsilon\sim\mathcal N(0,\Sigma_t)}
%    \big[\hat{L}_{S_t}(W_\mathrm {frozen }+\hat{\Delta W}_t+\varepsilon)
%        - \hat{L}_{S_t}(W_\mathrm {frozen }+\hat{\Delta W}_t)\big],\quad
% \varepsilon\sim\mathcal N(0,\Sigma_t) \in  \mathcal W_\Delta$.
% % \begin{equation*}
% % \begin{aligned}
% % &\mathrm{E\text{-}Sh}_t^{\Delta}(\hat{\Delta W}_t,\Sigma_t):= \\
% % &\qquad \mathbb{E}_{\varepsilon\sim\mathcal N(0,\Sigma_t)}
% %    \big[\hat{L}_{S_t}(W_\mathrm {frozen }+\hat{\Delta W}_t+\varepsilon) \\
% % & \qquad- \hat{L}_{S_t}(W_\mathrm {frozen }+\hat{\Delta W}_t)\big] \\
% % &\qquad \varepsilon\sim\mathcal N(0,\Sigma_t) \in \mathcal W_\Delta
% % \end{aligned}
% % \end{equation*}
% \end{lemma}
% Lemma~\eqref{lem:subspace-loss-pecl} decomposes the posterior expected loss into the empirical loss plus an increase that depends only on perturbations inside $\mathcal W_\Delta$. 
% % The backbone is frozen and does not appear in the sharpness term. In PECL, the  flatness is captured by sharpness restricted to the adapter subspace.





% \subsection{Sharpness-aware PAC–Bayes bound for PECL}
% \label{sec:pecl-bound-main}






% We combine the hierarchical PAC-Bayes with Lemma~\ref{lem:kl-decomp-pecl} and Lemma~\ref{lem:subspace-loss-pecl} . This yields a sharpness-aware bound whose empirical and complexity terms are both confined to the adapter subspace. The proof is in Appendix~\ref{proof:theorem}.


% We combine the hierarchical PAC-Bayes framework with Lemma~\ref{lem:kl-decomp-pecl} and Lemma~\ref{lem:subspace-loss-pecl} to obtain the following sharpness-aware PAC-Bayesian bound, in which both the empirical term and the complexity term are confined to the adapter subspace. We below provide the theorem\ref{thm:pecl-bound}, and its proof is included in Appendix~\ref{proof:theorem}.

% Section~4.1 provides a pathwise hierarchical PAC Bayes bound for the time averaged test risk over a task sequence.
% Lemma~\ref{lem:kl-decomp-pecl} and Lemma~\ref{lem:subspace-loss-pecl} show that, under PECL, its complexity and empirical terms can both be expressed entirely in the LoRA update subspace.
% In the next subsection, we combine these ingredients to obtain a sharpness aware PAC Bayes bound for PECL in which all task dependent quantities are confined to $\mathcal{W}_\Delta$.

% We combine the hierarchical PAC Bayes framework with Lemma~\ref{lem:kl-decomp-pecl} and Lemma~\ref{lem:subspace-loss-pecl} to derive a sharpness aware bound whose empirical and complexity terms are both confined to the adapter subspace.
% Lemma~\ref{lem:kl-decomp-pecl} is illustrated in Fig.~\ref{fig: KL_decompose}, showing that the full space divergence reduces to $\mathrm{KL}(Q^{\Delta}\|P^{\Delta})$ under a frozen backbone, while Lemma~\ref{lem:subspace-loss-pecl} expresses the empirical term through adapter space perturbations.
% Fig.~\ref{fig: showPAC3} summarizes the resulting hierarchical construction across tasks and the corresponding decomposition of the bound.
% We now state Theorem~\ref{thm:pecl-bound}; the proof is provided in Appendix~\ref{proof:theorem}.




% Section~4.1 establishes a hierarchical PAC--Bayes bound for the time-averaged test risk along a task sequence.
Under PECL, Lemma~\ref{lem:kl-decomp-pecl} and Lemma~\ref{lem:subspace-loss-pecl} rewrite its complexity and empirical terms entirely in the task adpter subspace $\mathcal W_\Delta$, through the task KL divergence and the adaptersubspace sharpness, respectively. 
% See Fig.~\ref{fig: showPAC3} and Fig.~\ref{fig: KL_decompose} for an overview of the subspace reduction and the hierarchical construction.
Substituting these reductions into the bound yields our main result below; the proof is given in Appendix~\ref{proof:theorem}.


% \begin{figure}[t]
% \centering
% \begin{minipage}{1\linewidth}
% \centering
% \includegraphics[width=\linewidth]{figures_TRML/shiyitu668.pdf}
% % \caption{Illustration of the KL divergence decomposition and coordinate equivalence. 
% %     (1) \textbf{KL Decomposition}: In PECL, the overall model complexity, measured by $\mathrm{KL}(Q^{\mathrm{full}} \| P^{\mathrm{full}})$, is solely determined by the distribution of the updated parameters in the adapter subspace $\mathcal{W}_\Delta$, formalized as $\mathrm{KL}(Q^{\Delta} \| P^{\Delta})$, while the frozen backbone contributes no divergence. 
% %     (2) \textbf{Two equivalent views}: The optimization in the low-rank factor coordinates $(A, B)$ is theoretically isomorphic to operating directly in the update subspace $\mathcal{W}_\Delta$. This equivalence preserves geometric properties such as flatness and curvature, which allows sharpness-aware optimization to be confined to the factorized LoRA parameters.}
% % \vspace{-4pt}
% \caption{KL decomposition and coordinate equivalence in PECL. With a frozen backbone,
% $\mathrm{KL}(Q^{\mathrm{full}}\|P^{\mathrm{full}})=\mathrm{KL}(Q^{\Delta}\|P^{\Delta})$, so the complexity depends only on the adapter subspace $\mathcal W_{\Delta}$.
% LoRA factors $(A,B)$ parameterize the same update subspace, allowing sharpness control to be applied in factor space.}
% \label{fig: KL_decompose}
% \end{minipage}
% \end{figure}



% [Sharpness-aware PAC--Bayes bound for PECL]
\begin{theorem}
\label{thm:pecl-bound}
Consider a PECL model with frozen backbone $W_{\mathrm{frozen}}$ and Gaussian posteriors 
$Q_t^\Delta = \mathcal N(\hat{\Delta W}_t,\Sigma_t)$ supported on $\mathcal W_\Delta$. 
Let $P_t^\Delta$ be task priors sampled from a hyper posterior $\mathcal P_{t-1}$ over adapter 
distributions. For any $\delta\in(0,1]$, with probability at least $1-\delta$ over $S_{1:T}$,
{%
\setlength{\abovedisplayskip}{2pt}
\setlength{\belowdisplayskip}{3pt}
\setlength{\abovedisplayshortskip}{3pt}
\setlength{\belowdisplayshortskip}{3pt}
\begin{equation*}
\label{eq:pecl-main-bound}
\resizebox{\linewidth}{!}{$
\begin{aligned}
&\frac{1}{T}\sum_{t=1}^T
\mathbb E_{P_t^\Delta\sim\mathcal P_{t-1}}
\mathbb E_{W\sim Q_t^{\mathrm{full}}}
L_{\mathcal D_t}(W) \\
&\le
\frac{1}{T}\sum_{t=1}^T
\mathbb E_{P_t^\Delta\sim\mathcal P_{t-1}}
\bigl[
\hat L_{S_t}(W_{\mathrm{frozen}} + \hat{\Delta W}_t)
+
\mathrm{RS}_t^{(0),\Delta}(\hat{\Delta W}_t,\Sigma_t)
\bigr] \\
&+
\sqrt{
\frac{
\displaystyle
\sum_{t=1}^T
\mathbb E_{P_t^\Delta\sim\mathcal P_{t-1}}
\mathrm{KL}\bigl(Q_t^\Delta \,\Vert\, P_t^\Delta\bigr)
+
\sum_{t=1}^T
\mathrm{KL}\bigl(\mathcal P_t\Vert\mathcal P_{t-1}\bigr)
+
\log\frac{1}{\delta}}
{2\displaystyle\sum_{t=1}^T (m_t - 1)}
}.
\end{aligned}
$}
\end{equation*}
}
\end{theorem}

% Theorem~\ref{thm:pecl-bound} specialises the hierarchical PAC-Bayesian framework to a PECL regime with a frozen backbone and posteriors supported on the adapter subspace $\mathcal{W}_\Delta$. In this setting, the empirical term is evaluated under Gaussian perturbations \( \varepsilon \in \mathcal W_\Delta \), which gives rise to the adapter-subspace expected sharpness \( \mathrm{E\text{-}Sh}^{\Delta}_t(\hat{\Delta W_t},\Sigma_t) \), while the complexity term decomposes into adapter-level divergences \( \mathrm{KL}(Q^\Delta_t \Vert P^\Delta_t) \) and a hyper-posterior drift \( \mathrm{KL}(P_t \Vert P_{t-1}) \). Since the backbone is frozen, neither sharpness nor complexity depends on directions outside \( \mathcal W_\Delta \), so the bound singles out adapter-subspace flatness as the relevant notion of flatness in PECL and theoretically justifies focusing flatness-aware optimisation on the adapter subspace.

% {%
% \setlength{\textfloatsep}{4pt}      % figure* (top/bottom) 与正文的距离（最关键）
% \setlength{\floatsep}{4pt}          % 两个浮动体之间的距离（一般备用）
% \setlength{\abovecaptionskip}{1pt}  % caption 上方
% \setlength{\belowcaptionskip}{0pt}  % caption 下方
% \captionsetup{skip=3pt}             % 图与caption的距离（你也可继续用4pt）

% \begin{figure*}[htbp]
%     \centering
%     \begin{subfigure}[t]{0.57\linewidth}
%         \centering
% \includegraphics[width=1\linewidth]{figures_TRML/rwp_full_lora_acc20_bars_imagenetr_3result1.pdf}
%         \subcaption{Effect of perturbation scope in PECL}
%          \label{fig:perturb_scope}
%     \end{subfigure}
%     \begin{subfigure}[t]{0.4\linewidth}
%         \centering
% \includegraphics[width=1\linewidth]{figures_TRML/lossland_curv1d_combine_loss_layout2.pdf}
%        \subcaption{Loss landscapes of different perturbation type}
%         \label{fig:perturb_type_landscape}
%     \end{subfigure}
%     \vspace{-6pt}
%     \caption{\textbf{Perturbation scope and type in frozen-backbone PECL.}
%     \textbf{(a) Scope.} Final average accuracy across perturbation strengths $\rho$ under different LoRA organizations. Restricting perturbations to the task adapter subspace $\Delta W$ consistently improves over the SGD baseline, whereas applying perturbations in the full parameter space $W+\Delta W$ substantially degrades performance.
%     % Final average accuracy under different perturbation strengths $\rho$ for LoRA organizations. Restricting perturbations to the task-admissible adapter subspace $\Delta W$ consistently outperforms the SGD baseline, whereas full-parameter perturbations $W+\Delta W$ substantially degrade performance.
%     \textbf{(b) Type.} Loss curves along the dominant curvature direction.
%     The top panel evaluates loss under perturbations restricted to the LoRA coordinates $\Delta W$, while the bottom panel evaluates loss under perturbations defined over the full parameter space $W+\Delta W$.
%     % The top  evaluates the loss when perturbations are restricted to the LoRA coordinates, while the bottom evaluates the loss when perturbations are defined over all parameters.
%     }
%     \label{fig:scope_type_main}
% \end{figure*}
% }


Theorem~\ref{thm:pecl-bound} specializes the Theorem~\ref{thm:pathwise_pac} to PECL by exploiting that all task dependent randomness is confined to the task subspace $\mathcal W_\Delta$. This restriction is not merely technical.
Under a frozen backbone, $\mathcal W_{\Delta,t}$ is exactly the set of directions along which the model can be modified by future tasks.
Therefore, in continual learning the relevant notion of flatness for generalization is the effective geometry induced on $\mathcal W_{\Delta,t}$, rather than a global property of the full parameter space.
% Since subsequent adaptation is restricted to $\mathcal W_{\Delta,t}$, this term quantifies the sensitivity of the task $t$ solution to admissible variations, and thus how robust the learned solution is with respect to the only perturbations that future tasks can effectively realize.
The complexity penalty is expressed through the adapter level divergence $\mathrm{KL}(Q_t^\Delta\Vert P_t^\Delta)$ together with the hyperposterior drift, yielding a structurally aligned form of complexity control compared with full finetuning.




% The empirical risk decomposes into the nominal empirical loss and the subspace random zero-order sharpness $\mathrm{RS}_t^{(0),\Delta}(\hat{\Delta W}_t,\Sigma_t)$ induced by perturbations supported on  $\mathcal W_{\Delta,t}$.
% % The empirical contribution consists of the nominal empirical loss and a subspace expected sharpness term induced by perturbations supported on $\mathcal W_\Delta$, thereby identifying adapter subspace flatness as the notion of flatness relevant under PECL.
% The complexity penalty is expressed through the adapter level divergence $\mathrm{KL}(Q_t^\Delta\Vert P_t^\Delta)$ together with the hyperposterior drift, yielding a tighter form of complexity control than a full-fintune.
% with a frozen backbone and posteriors supported on the LoRA update subspace $\mathcal W_\Delta$.
% The empirical contribution consists of the empirical loss at $W_{\mathrm{frozen}}+\hat{\Delta W}_t$ plus the subspace sharpness $\mathrm{E\text{-}Sh}_t^\Delta(\hat{\Delta W}_t,\Sigma_t)$ induced by Gaussian perturbations within $\mathcal W_\Delta$.
% The complexity decomposes into the adapter-level divergence $\mathrm{KL}(Q_t^\Delta\Vert P_t^\Delta)$ and the hyperposterior drift $\mathrm{KL}(\mathcal P_t\Vert \mathcal P_{t-1})$, implying that tightening the bound in PECL amounts to controlling sharpness and KL entirely inside $\mathcal W_\Delta$.
% The empirical term consists of the nominal empirical loss at $W_{\mathrm{frozen}}+\hat{\Delta W}_t$ and the subspace expected sharpness induced by Gaussian perturbations within $\mathcal W_\Delta$.
% The complexity decomposes into the adapter-level divergence $\mathrm{KL}(Q_t^\Delta\Vert P_t^\Delta)$ and the hyperposterior drift.
% Therefore, tightening the bound amounts to controlling sharpness and KL entirely inside $\mathcal W_\Delta$.
% In particular, the complexity term is expressed through $\mathrm{KL}(Q_t^\Delta\Vert P_t^\Delta)$ over the LoRA update subspace, rather than through a divergence over the full parameter space.
% Since $\mathcal W_\Delta$ has far fewer trainable degrees of freedom than the full model, the resulting complexity control is correspondingly smaller than in Theorem~1, yielding a strictly tighter bound under the PECL constraint.




% Therefore, tightening the bound amounts to controlling sharpness and KL entirely inside $\mathcal W_\Delta$.





% \begin{theorem}[Sharpness-aware PAC–Bayes bound for PECL]
% \label{thm:pecl-bound}
% Consider a PECL model with frozen backbone $W_{\mathrm{frozen}}$ and Gaussian posteriors $Q_t^\Delta = \mathcal N(\hat{\Delta W}_t,\Sigma_t)$ supported on $\mathcal W_\Delta$. Let $P_t^\Delta$ be task priors sampled from a hyper posterior $\mathcal P_{t-1}$ over adapter distributions. For any $\delta\in(0,1]$, with probability at least $1-\delta$ over $S_{1:T}$,
% \begin{equation}
% \label{eq:pecl-main-bound}
% \begin{aligned}
% \frac1T\sum_{t=1}^T
% \mathbb E_{P_t^\Delta\sim\mathcal P_{t-1}}
% \mathbb E_{W\sim Q_t^{\mathrm{full}}}
% L_{\mathcal D_t}(W)
% \;\le\;
% &\;
% \frac1T\sum_{t=1}^T
% \mathbb E_{P_t^\Delta\sim\mathcal P_{t-1}}
% \Bigl[
% \hat L_{S_t}(W_{\mathrm{frozen}} + \hat{\Delta W}_t)
% +
% \mathrm{E\text{-}Sh}_t^\Delta(\hat{\Delta W}_t,\Sigma_t)
% \Bigr]
% \\[0.2em]
% &\;+\;
% \sqrt{
% \frac{
% \displaystyle \sum_{t=1}^T
% \mathbb E_{P_t^\Delta\sim\mathcal P_{t-1}}
% \mathrm{KL}\bigl(Q_t^\Delta \,\Vert\, P_t^\Delta\bigr)
% \;+\;
% \displaystyle \sum_{t=1}^T
% \mathrm{KL}(\mathcal P_t\Vert\mathcal P_{t-1})
% \;+\;
% \log\frac1\delta}
% {2\displaystyle\sum_{t=1}^T (m_t - 1)}
% }.
% \end{aligned}
% \end{equation}
% \end{theorem}

% Theorem~\ref{thm:pecl-bound} decomposes the task averaged test risk into three components.

% \begin{itemize}
% \item An empirical term consisting of the empirical loss at $W_{\mathrm{frozen}}+\hat{\Delta W}_t$ and the adapter-subspace expected sharpness $\mathrm{E\text{-}Sh}_t^\Delta(\hat{\Delta W}_t,\Sigma_t)$.
% \item A sum of adapter-level KL divergences $\sum_t \mathbb E_{P_t^\Delta}\mathrm{KL}(Q_t^\Delta\Vert P_t^\Delta)$, which measures how much each task posterior deviates from its prior inside $\mathcal W_\Delta$.
% \item A hyper-posterior drift $\sum_t \mathrm{KL}(\mathcal P_t\Vert\mathcal P_{t-1})$, which quantifies how the distribution of task priors evolves with the task sequence.
% \end{itemize}

% Curvature enters the bound only through the expected sharpness term, which is defined over Gaussian perturbations supported on $\mathcal W_\Delta$. The complexity term is also expressed entirely through KL divergences on $\mathcal W_\Delta$ and their hyper-level evolution. In a PECL regime, adapter-subspace flatness is therefore the relevant notion of flatness: tightening the bound requires controlling $\mathrm{E\text{-}Sh}_t^\Delta$ and the adapter KL divergences, while directions outside $\mathcal W_\Delta$ play no role.

% A special case is SeqLoRA, where the prior is chosen as the previous posterior, $P_t^\Delta := Q_{t-1}^\Delta$ conditioned on the realised history. In this case the hyper update is deterministic, $\mathcal P_t = \mathcal P_{t-1}$. The complexity reduces to $\sum_t \mathrm{KL}(Q_t^\Delta\Vert Q_{t-1}^\Delta)$.

% \subsubsection{What is change between full-space and lora under PECL?}







% \begin{figure*}[thbp]
%     \centering
%     \begin{subfigure}{0.3\linewidth}
%         \centering
% \includegraphics[width=0.95\linewidth]{figures_TRML/rwp_full_lora_acc20_bars_imagenetr_redraw.pdf}
%         \subcaption{Effect of perturbation scope}
%         \label{fig: pertuabtion scope}
%     \end{subfigure}
%     \begin{subfigure}{0.35\linewidth}
%         \centering
% \includegraphics[width=0.95\linewidth]{figures_TRML/lossland_curv1d_combine_loss_layout2.pdf}
%         \subcaption{Loss landscape of perturbation type}
%         \label{fig: loss landscape}
%     \end{subfigure}\hfill
%     \begin{subfigure}{0.33\linewidth}
%         \centering
% \includegraphics[width=0.95\linewidth]{figures_TRML/weight_overlap_q_13.pdf}
%         \subcaption{Similarity between $W$ and $\Delta W$}
%         \label{fig: weight similiairy}
%     \end{subfigure}
%     \caption{Empirical analysis of perturbation}
% \end{figure*}



\subsection{Discussion}
\label{sec:flatness-subspace}




% \subsection{Scope of perturbations}





% The introduction highlighted a structural tension: classical sharpness-aware arguments are often phrased in the full parameter space, whereas PECL enforces a frozen backbone and restricts learning to an adapter update subspace.
% Theorem~\ref{thm:pecl-bound} resolves this mismatch at the level of the bound.
% Under PECL, the learned predictor satisfies $\hat W_t=W_0+\hat{\Delta W}_t$ with $\hat{\Delta W}_t\in\mathcal W_{\Delta,t}$, and the perturbations entering the empirical term are supported on $\mathcal W_{\Delta,t}$ through $\varepsilon\sim\mathcal N(0,\Sigma_t)$.
% Moreover, by Lemma~\ref{lem:kl-decomp-pecl}, the task-wise complexity term depends only on the adapter-level distributions.
% Hence, directions outside $\mathcal W_{\Delta,t}$ do not enter either the sharpness term or the KL complexity under a frozen backbone.
% Consequently, the relevant comparison in PECL is determined by how $\Sigma_t$ is chosen within $\mathcal W_{\Delta,t}$ under comparable KL and drift budgets, rather than by invoking full-space perturbations that are not controlled by the PECL bound.



% Theorem~\ref{thm:pecl-bound} is formulated on an abstract trainable update subspace $\mathcal{W}_\Delta$. In PECL, this abstraction is instantiated by LoRA: although optimisation is performed in low-rank factor coordinates, the induced weight variation $\delta W$ lies in the admissible set $\mathcal{W}_\Delta$ and locally spans update directions in weight space. Therefore, both the sharpness term and the task-wise KL contribution are determined by the geometry restricted to $\mathcal{W}_\Delta$. Concretely, at task $t$ the learned predictor takes the form
% $
% \hat W_t \;=\; W_0+\hat{\Delta W}_t,
% \hat{\Delta W}_t\in\mathcal{W}_\Delta,
% $
% and admissible perturbations satisfy $\delta W\in\mathcal{W}_\Delta$.
% Perturbing parameters outside $\mathcal W_\Delta$ cannot be required to tighten the PECL generalization bound, because those directions are neither part of the admissible update distribution nor do they enter the complexity term under a frozen backbone.
% At best, full-parameter perturbations may act as an \emph{external regularizer} that changes the training dynamics, but such effects are not aligned with the PECL bound itself and may introduce instability by coupling optimization to directions that the learner is not allowed to retain at test time.
% Therefore, the appropriate theoretical comparison in PECL is not ``global flatness versus subspace flatness'', but rather which choice of $\Sigma_t$ within $\mathcal W_\Delta$ yields smaller $\mathrm{ES}_t^{0,\Delta}$ under comparable KL budgets, while controlling hyperposterior drift.

% \paragraph{Perturbations in PECL: Scope and Type.}
% The introduction raised a tension between full-space sharpness arguments and PECL, which freezes the backbone and restricts learning to an adapter subspace.
% Theorem~\ref{thm:pecl-bound} resolves this mismatch: at task $t$, $\hat W_t=W_0+\hat{\Delta W}_t$ with $\hat{\Delta W}_t\in\mathcal W_{\Delta,t}$, and the empirical perturbations are confined to $\mathcal W_{\Delta,t}$ through $\varepsilon\sim\mathcal N(0,\Sigma_t)$.
% Moreover, Lemma~\ref{lem:kl-decomp-pecl} shows that the task-wise complexity depends only on the adapter-level distributions, so directions outside $\mathcal W_{\Delta,t}$ enter neither the sharpness term nor the KL penalty under a frozen backbone.
% Hence, full-parameter perturbations are not required to tighten the PECL bound; the relevant comparison is instead which covariance geometry $\Sigma_t$ on $\mathcal W_{\Delta,t}$ yields smaller $\mathrm{RS}_t^{(0),\Delta}(\hat{\Delta W}_t,\Sigma_t)$ under comparable KL and drift budgets.

% \paragraph{Scope: Why admissible perturbations lie in $\mathcal W_{\Delta,t}$}

\vspace{-1pt}
\noindent\textbf{Scope: Why admissible perturbations lie in $\mathcal W_{\Delta,t}$.}\quad
The introduction raised a tension between ful space sharpness and the PECL setting.
In PECL, the backbone $W_{\mathrm{frozen}}$ is fixed and learning is confined to the LoRA update subspace $\mathcal W_{\Delta,t}$. Consequently, only directions in $\mathcal W_{\Delta,t}$ are reachable by future task updates and can affect learning along the training trajectory.
Theorem \ref{thm:pecl-bound} makes this restriction explicit: both the task-wise KL term and the sharpness-dependent empirical term depend only on distributions supported on $\mathcal W_{\Delta,t}$, and directions in the complement subspace contribute to neither term.
Therefore, perturbing the parameter outside $\mathcal W_{\Delta,t}$ imposes robustness requirements along unreachable directions, which cannot be compensated by admissible LoRA updates and can introduce unnecessary loss elevation and unstable optimization behavior.



% t determines the only directions along which the hypothesis can move under future tasks, and therefore the only directions that can affect continual generalization through the learning trajectory. Theorem 4.4 makes this geometric restriction explicit: both the task wise KL complexity and the sharpness dependent empirical contribution reduce exactly to distributions supported on $\mathcal W_{\Delta,t}$, so directions outside $\mathcal W_{\Delta,t}$ enter neither the empirical term nor the complexity term of the bound.
%  If one perturbs the full parameter vector with noise that has substantial mass outside $\mathcal W_{\Delta,t}$, then the training objective effectively demands robustness to directions that are unreachable by future updates. Since LoRA cannot compensate for sensitivity along the complement subspace, such perturbations introduce an irreducible loss elevation and inflate gradient variance, which manifests as conservative updates and underfitting. 


% Theorem~\ref{thm:pecl-bound} resolves this mismatch.
% At task $t$, $\hat W_t=W_0+\hat{\Delta W}_t$ with $\hat{\Delta W}_t\in\mathcal W_{\Delta,t}$, and the empirical perturbations are confined to $\mathcal W_{\Delta,t}$ through $\varepsilon\sim\mathcal N(0,\Sigma_t)$.
% Theorem~\ref{thm:pecl-bound} resolves this mismatch: at task $t$, $\hat W_t=W_0+\hat{\Delta W}_t$ with $\hat{\Delta W}_t\in\mathcal W_{\Delta,t}$, and the empirical perturbations are confined to $\mathcal W_{\Delta,t}$ through $\varepsilon\sim\mathcal N(0,\Sigma_t)$.
% Lemma~\ref{lem:kl-decomp-pecl} 



% Theorem~\ref{thm:pecl-bound} shows that the task wise complexity depends only on the adapter-level distributions, so directions outside $\mathcal W_{\Delta,t}$ enter neither the sharpness term nor the KL penalty.
% Hence, full parameter perturbations are not required to tighten the PECL bound. 
% At task $t$, only perturbations supported on $\mathcal W_{\Delta,t}$ can affect the empirical regularization term.
% The relevant comparison is instead which covariance geometry $\Sigma_t$ on $\mathcal W_{\Delta,t}$ yields smaller $\mathrm{RS}_t^{(0),\Delta}(\hat{\Delta W}_t,\Sigma_t)$,  reducing the susceptibility of the solution to curvature amplification along future admissible updates.

% \paragraph{Type: How perturbation shapes flatness}

% Scope: Why admissible perturbations lie in $\mathcal W_{\Delta,t}$
% \vspace{-1pt}
\noindent\textbf{Type: How perturbation shapes flatness.}\quad
% Within the correct scope, the bound further separates “which directions can matter” from “how strongly each admissible direction is probed”. 
%The former is controlled by $\mathcal W_{\Delta,t}$ through $C_{\Delta,t}$, while the latter is controlled by the covariance geometry $\Sigma_t$. 
% In PECL, task-admissible sharpness are restricted to the adapter update subspace.
Let $\Pi_{\Delta,t}$ be the projector onto $\mathcal W_{\Delta,t}$ and let $C_t=C(\hat W_t;S_t)$ be a task-dependent curvature operator of $\hat L_{S_t}$ at $\hat W_t$.
% (for example, the Gauss Newton or Fisher operator; if the Hessian is used, the discussion concerns the associated local quadratic form).
In PECL, task-relevant parameter updates and task-admissible sharpness are restricted to $\mathcal W_{\Delta,t}$, hence only the restricted curvature $C_{\Delta,t}=\Pi_{\Delta,t}C_t\Pi_{\Delta,t}$ and perturbations supported on $\mathcal W_{\Delta,t}$ can contribute. 
% to the sharpness dependent term
% Although $C_t$ is naturally defined on the full parameter space, PECL perturbations satisfy $\varepsilon\in\mathcal W_{\Delta,t}$, so only the restricted operator
% $
% C_{\Delta,t}=\Pi_{\Delta,t}C_t\Pi_{\Delta,t}
% $
% can affect the sharpness term.
Under a local second-order approximation and $\varepsilon\sim\mathcal N(0,\Sigma_t)$ supported on $\mathcal W_{\Delta,t}$,
{%
\setlength{\abovedisplayskip}{4pt}
\setlength{\belowdisplayskip}{3pt}
\setlength{\abovedisplayshortskip}{4pt}
\setlength{\belowdisplayshortskip}{3pt}
\[
\mathrm{TS}_t^{\Delta}(\hat{\Delta W}_t,\Sigma_t)
% =
% \mathbb E_{\varepsilon}\big[\hat L_{S_t}(\hat W_t+\varepsilon)-\hat L_{S_t}(\hat W_t)\big]
\approx
\frac12\operatorname{tr}\!\big(C_{\Delta,t}\Sigma_t\big),
\]
}where $\mathcal W_{\Delta,t}$ fixes the admissible directions, $C_{\Delta,t}$ specifies the restricted curvature on them, $\Sigma_t$ allocates perturbation variance within them, and $\operatorname{tr}(C_{\Delta,t}\Sigma_t)$ measures curvature–noise alignment inside the task admissible geometry.
% The scope $\mathcal W_{\Delta,t}$ determines which curvature modes, the curvature $C_{\Delta,t}$ controlls which directions can matter, $\Sigma_t$ specifies how variance is allocated across those admissible directions and
% $\operatorname{tr}(C_{\Delta,t}\Sigma_t)$ quantifies curvature-noise alignment.

% If the perturbation is more anisotropic and better aligned with highcurvature directions, it will facilitate escape from sharp minima\cite{pmlr-v97-zhu19e,chen2026sharpnessflatnessdecompositionframework}.


This clarifies how different optimizers act in PECL. More generally, standard sharpness criteria can be re-expressed in the task-admissible geometry by restricting perturbations to $\mathcal W_{\Delta,t}$. Standard SGD injects an implicit stochastic perturbation with an effective covariance $\Sigma_t^{\mathrm{SGD}}$ determined by minibatch noise\cite{pmlr-v97-zhu19e}. Random zeroth-order sharpness $\mathrm{RS}^{(0)}$ corresponds to explicit isotropic Gaussian noise $\Sigma=\sigma^2 I$, and measures the expected loss elevation under average-case perturbations supported on the admissible subspace. In contrast, adversarial criteria replace the average-case elevation with a worst-case concentration over the same adapter subspace. The adversarial zeroth-order sharpness $\mathrm{AS}^{(0),\Delta}$ chooses the worst-case perturbation that maximizes the local loss increase over $\Sigma=\rho\frac{\nabla L}{|\nabla L|_2}$, yielding a gradient-aligned direction\cite{foretSharpnessAwareMinimizationEfficiently2021}. The adversarial first-order sharpness $\mathrm{AS}^{(1),\Delta}$ maximizes the neighborhood gradient norm
$\Sigma=\rho\frac{\nabla|\nabla L|_2}{|\nabla|\nabla L|_2|_2}$,
which emphasizes directions directions aligned with dominant eigenmodes of $C_{\Delta,t}$~\cite{zhangGradientNormAware2023}. These definitions imply an ordering in strength within the same subspace and radius constraint~\cite{mollenhoffSAMOptimalRelaxation2022,Zhang_GAM}:
{%
\setlength{\abovedisplayskip}{3pt}
\setlength{\belowdisplayskip}{3pt}
\setlength{\abovedisplayshortskip}{1pt}
\setlength{\belowdisplayshortskip}{1pt}
\begin{equation*}
    \mathrm{AS}^{(1)}_S(W,\rho)\ \ge \mathrm{AS}^{(0)}_S(W,\rho)\ \ge \mathrm{RS}^{(0)}_S(W,\Sigma_t),
\end{equation*}
}
% $$
Accordingly, optimizing increasingly adversarial perturbation objectives typically yields solutions with smaller sharpness within task subspace. And as the perturbation becomes increasingly concentrated and aligned with the dominant high-curvature direction, it more effectively promotes escape from sharp minima and biases the trajectory toward flatter regions~\cite{pmlr-v97-zhu19e,chen2026sharpnessflatnessdecompositionframework}.
% Accordingly, optimizing these perturbation objectives leads to final solutions with different levels of flatness in the admissible subspace. It is pointd out that as the perturbation becomes increasingly concentrated and aligned with the dominant high-curvature modes of $C_{\Delta,t}$, it more effectively promotes escape from sharp minima and biases the trajectory toward flatter regions. \cite{pmlr-v97-zhu19e,chen2026sharpnessflatnessdecompositionframework}



% SGD can be locally interpreted as injecting an implicit stochastic perturbation with some effective covariance $\Sigma_t^{\mathrm{SGD}}$ determined by minibatch noise\cite{pmlr-v97-zhu19e}. 
% Random zeroth-order sharpness $\mathrm{RS}^{(0)}$ considers isotropic Gaussian noise $\Sigma=\sigma^2 I$ restricted to $\mathcal W_{\Delta,t}$, and thus measures the expected loss elevation under average-case perturbations supported on the admissible subspace.
% In contrast, adversarial criteria replace this averaging by a worst-case concentration over the same admissible set.
% Adversarial zeroth-order sharpness $\mathrm{AS}^{(0)}$ selects perturbations that maximize the loss increase, which can be represented by $\Sigma=\rho^2\frac{\nabla L}{|\nabla L|_2}$, highlighting that it is gradient-driven and hence primarily controlled by the local slope rather than curvature.
% Adversarial first-order sharpness $\mathrm{AS}^{(1)}$ further strengthens this notion by maximizing the neighborhood gradient norm, thereby favoring directions aligned with the dominant eigenmodes of $C_{\Delta,t}$; equivalently, one may view it as inducing a perturbation distribution of the form $\varepsilon\sim\mathcal N!\left(0,\rho^2\frac{\nabla|\nabla L|_2}{|\nabla|\nabla L|_2|_2}\right)$.
% Collectively, these definitions imply a strict ordering in strength:
% $$
% \mathrm{AS}^{(1)}_S(W,\rho)\ \ge\ \mathrm{AS}^{(0)}_S(W,\rho)\ \ge\ \mathrm{RS}^{(0)}_S(W,\Sigma_t),
% $$
% which makes explicit that $\mathrm{RS}^{(0)}$ is average-case in loss, $\mathrm{AS}^{(0)}$ is worst-case in loss, and $\mathrm{AS}^{(1)}$ is worst-case in gradient norm.



% Random zeroth-order sharpness $\mathrm{RS}^{(0)}$ consider Isotropic Gaussian noise $\sum =\sigma^2 I$, which corresponds to averaging the loss elevation under supported on $\mathcal W_{\Delta,t}$. In contrast, adversarial criteria replace this average-case weighting by a worst-case concentration within the admissible set.




% whereas sharpness-aware methods reshape the effective perturbation geometry.


% This also clarifies how different optimizers act in PECL. SGD can be locally viewed as injecting an implicit stochastic perturbation \cite{pmlr-v97-zhu19e}. Sharpness aware methods modify the effective perturbation geometry.
% Random zeroth-order sharpness $\mathrm{RS}^{(0)}$  consider Isotropic Gaussian noise $\sum =\sigma^2 I$, which corresponds to averaging the loss elevation under supported on $\mathcal W_{\Delta,t}$. 
% In contrast, adversarial criteria replace this average-case weighting by a worst-case concentration within the admissible set. The adversarial zeroth-order sharpness $\mathrm{AS}^{(0)}$ selects perturbations that maximize the loss increase consider$\sum=\rho^2 \frac{\nabla L}{\|\nabla L\|_2})$, and it is gradient-driven means controlled by the local slope rather than curvature. 
% The adversarial first-order sharpness $\mathrm{AS}^{(1)}$ is defined by maximizing the neighborhood gradient norm, which directions aligned with dominant eigenmodes of $C_{\Delta,t}$ and $\epsilon \sim \mathcal{N}(0,\rho^2\frac{\nabla \|\nabla L\|_2}{\|\nabla \|\nabla L\|_2\|_2})$. 



% Thus, 

% Random zeroth-order sharpness $\mathrm{RS}^{(0)}$ corresponds to averaging under $\Sigma_t$, whereas adversarial criteria concentrate mass toward locally most sensitive directions within $\mathcal W_{\Delta,t}$.
% The adversarial zeroth-order  sharpness$\mathrm{AS}^{(0)}$ is gradient-driven: it concentrates perturbations toward directions that maximize the loss change, therefore it is sensitive to the local slope rather than curvature.
% The first-order adversarial criterion $\mathrm{AS}^{(1)}$ is curvature-driven: it seeks directions that maximize the gradient which directions aligned with dominant eigenmodes of $C_{\Delta,t}$。



% While the zeroth-order $\mathrm{AS}^{(0)}$ targets the norm of the gradient, the first-order $\mathrm{AS}^{(0)}$ considers the curvature of the function.
% Thus, $\mathrm{RS}^{(0)}$, $\mathrm{AS}^{(0)}$, and $\mathrm{AS}^{(1)}$ are comparable within a unified PECL-compatible parameterization through $\Sigma_t$.

% This expression makes the roles of scope and type explicit: the admissible subspace $\mathcal W_{\Delta,t}$ determines which curvature directions can be excited at all, while the covariance geometry $\Sigma_t$ determines how strongly each admissible direction is probed.
% Moreover, $\operatorname{tr}(C_{\Delta,t}\Sigma_t)$ can be interpreted as a restricted curvature-noise alignment measure, since it aggregates curvature energy along the directions where the perturbation distribution places variance.
% A larger value indicates that the injected noise is more aligned with high-curvature modes in $\mathcal W_{\Delta,t}$, which increases the expected loss elevation and thus promotes leaving narrow basins and exploring flatter regions within the admissible update space. 
% This local view also connects flatness for generalization to forgetting.  A future admissible update $\Delta\in\mathcal W_{\Delta,t}$ induces an amplified loss change $\frac12\Delta^\top C_{\Delta,t}\Delta$ on task $t$, so reducing curvature along admissible directions reduces amplification and improves stability.
% This clarifies that scope is fixed by $\mathcal W_{\Delta,t}$, while the effect of perturbations within that scope is governed jointly by the restricted curvature $C_{\Delta,t}$ and the covariance geometry $\Sigma_t$.

% This representation clarifies that the scope of admissible variation is fixed by $\mathcal W_{\Delta,t}$, while its effect is governed jointly by the restricted curvature $C_{\Delta,t}$ and the covariance geometry $\Sigma_t$.
% This clarifies that scope is fixed by $\mathcal W_{\Delta,t}$, while the effect of perturbations within that scope is governed jointly by the restricted curvature $C_{\Delta,t}$ and the covariance geometry $\Sigma_t$.




% Consider a future admissible update $\Delta\in\mathcal W_{\Delta,t}$ performed when learning later tasks.
% A second order expansion yields a curvature amplified loss change term of the form $\frac12\Delta^\top C_{\Delta,t}\Delta$ for the task $t$ empirical risk.
% Large eigenvalues of $C_{\Delta,t}$ therefore correspond to directions within the admissible update space where small future updates can induce disproportionate loss increases on past tasks, which manifests as forgetting on the test distributions.
% Mechanisms that reduce the effective restricted curvature within $\mathcal W_{\Delta,t}$ consequently reduce this amplification and thus mitigate the degradation of continual generalization.
% Importantly, the same restricted curvature also governs the amplification mechanism underlying forgetting. Consider a future admissible update $\Delta\in\mathcal W_{\Delta,t}$. A second-order expansion yields a curvature-amplified loss change term of the form $\frac12\Delta^\top C_{\Delta,t}\Delta$. Large eigenvalues of $C_{\Delta,t}$ therefore correspond to directions within the admissible update space where even small future updates can induce disproportionate loss increases on past tasks. Mechanisms that reduce the effective curvature in $\mathcal W_{\Delta,t}$ consequently reduce this amplification, and thus systematically mitigate forgetting.


% \paragraph{Type of perturbations}
% \label{sec:type}

% In Theorem~\ref{thm:pecl-bound}, the empirical regularization is induced by Gaussian perturbations
% $\varepsilon\sim\mathcal N(0,\Sigma_t)$ supported on the adapter update subspace $\mathcal W_{\Delta,t}$.
% Accordingly, a perturbation design is fully characterized by the covariance geometry of $\Sigma_t$ on $\mathcal W_{\Delta,t}$.

% In Theorem~\ref{thm:pecl-bound}, perturbations are Gaussian $\varepsilon\sim\mathcal N(0,\Sigma_t)$ supported on $\mathcal W_{\Delta,t}$, so the perturbation type is fully specified by $\Sigma_t$ on this subspace.

% $\mathrm{AS}^{(0)}$ uses loss elevation and $\mathrm{AS}^{(1)}$ uses gradient sensitivity.


% Adversarial zeroth-order sharpness $\mathrm{AS}^{(0)}$ can be interpreted as driving the perturbation distribution toward increasingly concentrated mass on locally worst-case directions inside  $\mathcal W_{\Delta,t}$.
% , yielding a distributional surrogate of worst-direction probing within the PECL-admissible subspace.
% Adversarial first-order sharpness $\mathrm{AS}^{(1)}$ adopts the same direction-selection principle, but replaces loss elevation with gradient sensitivity, thereby emphasizing curvature through the local variation of $\|\nabla_W \hat L_{S_t}(W+\varepsilon)\|$.
% Adversarial first-order sharpness $\mathrm{AS}^{(1)}$ adopts the same direction-selection principle, but replaces loss elevation with gradient sensitivity, thereby emphasizing curvature through the local variation of $\|\nabla_W \hat L_{S_t}(W+\varepsilon)\|$ in the neighborhood.
% % , whose worst-case direction is governed by second-order effects.

% Prior work also explores intermediate choices that interpolate between  isotropic averaging and worst-direction probing by biasing $\Sigma_t$ toward a small subset of directions~\cite{pmlr-v151-bisla22a,yangRevisitingFlatnessAwareOptimization2025,NEURIPS2022_c859b99b}.
% Overall, $\mathrm{RS}^{(0)}$, $\mathrm{AS}^{(0)}$, and $\mathrm{AS}^{(1)}$ can be expressed and compared within a unified, PECL-compatible perturbation parameterization through $\Sigma_t$.

% Adversarial first-order sharpness $\mathrm{AS}^{(1)}$ retains the same directional-selection perspective but changes the measured quantity from loss elevation to gradient sensitivity and thus accentuating curvature effects. Other works are also explore intermediate choices that interpolate between purely random averaging and worst-direction probing by biasing $\Sigma_t$ toward a small subset of directions \cite{pmlr-v151-bisla22a,yangRevisitingFlatnessAwareOptimization2025,NEURIPS2022_c859b99b}. In conclusion,  $\mathrm{RS}^{(0)}$, $\mathrm{AS}^{(0)}$, and $\mathrm{AS}^{(1)}$ can all be expressed and compared within the same PECL-compatible perturbation parameterization through $\Sigma_t$.


% and the sharpness term in the bound is the subspace-restricted random sharpness
% $\mathrm{RS}^{(0),\Delta}_t(\hat{\Delta W}_t,\Sigma_t)$.
% We therefore distinguish perturbations by how $\Sigma_t$ allocates variance across directions in $\mathcal W_{\Delta,t}$.

% \paragraph{(i) Isotropic perturbations.}
% A direction-agnostic baseline is obtained by
% \[
% \Sigma_t^{\mathrm{iso}}=\sigma_t^2 I_{\mathcal W_{\Delta,t}},
% \qquad
% \varepsilon\sim\mathcal N(0,\sigma_t^2 I_{\mathcal W_{\Delta,t}}),
% \]
% which probes the average local sensitivity uniformly over all admissible adapter directions.

% \paragraph{(ii) Structured anisotropic perturbations.}
% A broad class of practically used perturbations is captured by anisotropic covariances of the form
% \[
% \Sigma_t^{\mathrm{str}}
% =
% \sum_{i=1}^{k_t}\lambda_{t,i}u_{t,i}u_{t,i}^\top,
% \qquad
% u_{t,i}\in\mathcal W_{\Delta,t},\ \lambda_{t,i}\ge 0,
% \]
% including Filter-wise\cite{pmlr-v151-bisla22a}, history-gradient scaling\cite{yangRevisitingFlatnessAwareOptimization2025} and Fisher-informed  \cite{NEURIPS2022_c859b99b} designs.
% % In particular, random zero-order sharpness correspond to choosing $\Sigma_t^{\mathrm{str}}$ with a diagonal  structure in a prescribed grouping.

% % \emph{Filter-wise noise} can be expressed by taking $\Sigma_t$ to be diagonal in a filter decomposition,
% % \[
% % \Sigma_t \;=\; \frac{\sigma^2}{n}\,
% % \mathrm{diag}\!\bigl(\|W^{1}_{:,:}\|_2^2,\|W^{2}_{:,:}\|_2^2,\dots,\|W^{m}_{:,:}\|_2^2\bigr),
% % \]
% % which assigns larger variance to groups with larger weight norms.

% % \emph{Adaptive Random Weight Perturbation (ARWP)} further rescales the covariance using accumulated gradient statistics across previous steps. In its isotropic form, ARWP corresponds to
% % \[
% % \varepsilon_{r,t}\sim\mathcal N\!\left(0,\;
% % \Sigma_t^{\mathrm{arwp}}\right),
% % \qquad
% % \Sigma_t^{\mathrm{arwp}}
% % =
% % \frac{\sigma^2 I_{\mathcal W_{\Delta,t}}}
% % {\sqrt{\,1+\eta \sum_{i=1}^{t-1}\beta^{t-i-1} g_i^2\,}},
% % \]
% % where the scalar factor in the denominator contracts the perturbation scale as the accumulated gradient magnitude increases.

% % Combining ARWP with filter-wise generation yields a diagonal (or group-diagonal) anisotropic covariance,
% % \[
% % \varepsilon_{r,t}\sim\mathcal N\!\left(0,\;\Sigma_t^{\mathrm{f\text{-}arwp}}\right),
% % \qquad
% % \Sigma_t^{\mathrm{f\text{-}arwp}}
% % =
% % \frac{\sigma^2\,\mathrm{diag}\!\bigl(\{\|w_t^{(j)}\|^2\}_{j=1}^k\bigr)}
% % {\sqrt{\,1+\eta \sum_{i=1}^{t-1}\beta^{t-i-1}\,
% % \mathrm{diag}\!\bigl(\{\|g_i^{(j)}\|^2\}_{j=1}^k\bigr)\,}},
% % \]
% % which adaptively contracts each group variance based on its own historical gradient energy.
% % These examples illustrate that many RWP-style perturbations differ only through the structure and scaling of $\Sigma_t$, while sharing the same theoretical interface via $\mathrm{RS}^{(0),\Delta}_t(\hat{\Delta W}_t,\Sigma_t)$.

% % If desired, the directions $\{u_{t,i}\}$ may also be chosen as the top eigenvectors of the restricted curvature $H_t^\Delta$ (or an approximate subspace thereof), yielding curvature-aligned anisotropy without changing the form of the bound.

% \paragraph{(iii) Direction-concentrated perturbations.}
% To isolate the effect of selecting a small set of influential directions within the admissible update set, one may consider highly concentrated covariances
% \[
% \Sigma_t^{\mathrm{con}}
% =
% \rho_t^2 v_t v_t^\top,
% \qquad
% v_t\in\mathcal W_{\Delta,t},\ \|v_t\|=1,
% \]
% which excite essentially a single direction in $\mathcal W_{\Delta,t}$.
% This remains a special case of $\mathrm{RS}^{(0),\Delta}_t(\hat{\Delta W}_t,\Sigma_t)$, but represents an extreme anisotropic regime that is often used as a distributional proxy for worst-direction probing inside the PECL-admissible subspace.

% \smallskip
% Overall, these perturbation types differ only through $\Sigma_t$ while sharing the same PAC--Bayes interface in Theorem~\ref{thm:pecl-bound}:
% isotropic choices measure average sensitivity over $\mathcal W_{\Delta,t}$,
% structured anisotropy subsumes filter-wise and adaptive RWP schemes through diagonal covariances,
% and direction-concentrated covariances isolate the role of direction selection under the PECL constraint.


% \paragraph{Subspace design}
% \label{sec:subspace-design}

% In Theorem~\ref{thm:pecl-bound}, the LoRA organization enters the bound only through the task-dependent admissible update subspace $\mathcal W_{\Delta,t}$ and the induced covariance geometry $\Sigma_t$ supported on it.
% Theorem~\ref{thm:pecl-bound} also clarifies how LoRA architecture enters PECL: a design choice induces a \emph{sequence} of admissible update subspaces $\{\mathcal W_{\Delta,t}\}_{t=1}^T$, and thereby constrains the feasible families of $(Q_t^\Delta, P_t^\Delta, \Sigma_t)$ at each task.
% From the bound, subspace design affects continual learning through three coupled mechanisms.
% First, it shapes the empirical regularization term because $\mathrm{RS}_t^{(0),\Delta}(\hat{\Delta W}_t,\Sigma_t)$ depends on the curvature restricted to $\mathcal W_{\Delta,t}$.
% Second, it controls task interference through the task-wise complexity $\mathrm{KL}(Q_t^\Delta\Vert P_t^\Delta)$: when many tasks are forced to reuse a highly overlapping subspace (as in strongly shared adapters), fitting a new task typically requires larger posterior displacement along already occupied directions, which increases $\mathrm{KL}(Q_t^\Delta\Vert P_t^\Delta)$ under comparable priors.
% By contrast, incremental and constraint-augmented designs modify $\mathcal W_{\Delta,t}$ to introduce partially separated directions, which relaxes this competition and can reduce both the KL shift required for plasticity and the induced sharpness under comparable covariance geometries.
% Third, since priors are sampled from a history-dependent hyperposterior, the organization of $\{\mathcal W_{\Delta,t}\}$ also impacts the drift penalty $\sum_t \mathrm{KL}(\mathcal P_t\Vert\mathcal P_{t-1})$: stronger subspace coupling generally forces more frequent updates of the prior-selection mechanism, while better separated subspaces can stabilize this evolution.
% \paragraph{LoRA design}
\noindent\textbf{LoRA Organization as a Structural Prior.}\quad
% For Gaussian task posteriors and priors on the adapter coordinates,
% $Q_t^\Delta=\mathcal N(\mu_{Q,t},\Sigma_{Q,t})$ and $P_t^\Delta=\mathcal N(\mu_{P,t},\Sigma_{P,t})$ with $k_t=\dim(\mathcal W_{\Delta,t})$,
% % the task-wise KL admits the closed form
% \begin{equation*}
%     \begin{aligned}
% &\mathrm{KL}(Q_t^\Delta\Vert P_t^\Delta)
% =
% \frac12\Bigg(
% (\mu_{Q,t}-\mu_{P,t})^\top\Sigma_{P,t}^{-1}(\mu_{Q,t}-\mu_{P,t}) \\
% &\quad  +\log\frac{|\Sigma_{P,t}|}{|\Sigma_{Q,t}|}
% +\operatorname{tr}\big(\Sigma_{P,t}^{-1}\Sigma_{Q,t}\big)
% -k_t
% \Bigg).
%     \end{aligned}
% \end{equation*}
To make the hyper process concrete in PECL, we represent each task-level prior by a minimal \emph{hyper-state}
$
\phi_t = \bigl(\mathcal W_{\Delta,t},\mu_{P,t},\Sigma_{P,t}\bigr),
$
where $\mathcal W_{\Delta,t}$ specifies the admissible adapter update scope and $(\mu_{P,t},\Sigma_{P,t})$ parameterize a task prior on these coordinates.
Sampling $P_t\sim\mathcal P_{t-1}$ can then be interpreted as drawing $\phi_t\sim\mathcal P_{t-1}$ and instantiating the corresponding task prior. Under this view, SeqLoRA, IncLoRA, and constraint-augmented designs correspond to different evolution rules for $\phi_t$, and thus shape both the restricted sharpness channel and the KL complexity in Theorem~\ref{thm:pathwise_pac}.


% measures how much the hyper-state generation rule is revised over time.

% determines the admissible update directions and $(\mu_{P,t},\Sigma_{P,t})$ encode the inductive bias on those coordinates. 

% Under PECL, both the sharpness-dependent empirical contribution and the KL complexities reduce to quantities supported on $\mathcal W_{\Delta,t}$, so LoRA organization affects the bound through the evolution of $\phi_t$.

% Under this view, SeqLoRA, IncLoRA, and constraint-augmented designs correspond to different evolution rules for $\phi_t$, and thus shape both the restricted sharpness and the KL complexity in Theorem~\ref{thm:pathwise_pac}.

%equivalently $\phi_t=(\mathcal W_\Delta,\mu_{Q,t-1},\Sigma_{Q,t-1})$
% SeqLoRA keeps an approximately fixed scope $\mathcal W_{\Delta,t}\equiv\mathcal W_\Delta$ and adopts the deterministic rule $P_t^\Delta=Q_{t-1}^\Delta$, in which case the complexity collapses to $\sum_t\mathrm{KL}(Q_t^\Delta\Vert Q_{t-1}^\Delta)$. 



SeqLoRA keeps an approximately fixed scope $\mathcal W_{\Delta,t}\equiv\mathcal W_\Delta$ and, under the deterministic choice $P_t^\Delta=Q_{t-1}^\Delta$, the complexity reduces to $\sum_t\mathrm{KL}(Q_t^\Delta\Vert Q_{t-1}^\Delta)$.
% This shared scope limit admissible directions for new task and amplify cross-task interference for previous task.
A fixed scope concentrates successive updates in the same geometry, which tends to induce stronger direction competition across tasks.
IncLoRA relaxes the fixed-scope assumption by expanding $\mathcal W_{\Delta,t}$ over time, which redistribute adaptation across additional degrees of freedom and reduced subspace overwrite, while the increased flexibility may also admit more directions with projected curvature.
% . At the same time, a larger freedom may admit more directions with projected curvature.
% This provides extra admissible directions and reduce the influence of previous task direction. 
Constraint-augmented designs further control how new degrees of freedom are allocated and initialized by imposing structure on $\mathcal W_{\Delta,t}$, for example via orthogonality regularization~\cite{wang2023orthogonal}, gradient projection~\citep{Liang_2024_CVPR,yang2024parametercollisionhinderingcontinual}.
It reduce repeated overwrite and stabilize the evolution of the admissible geometry.

% reduce harmful overlap with previously important modes, thereby suppressing persistent sharp contributions within the admissible scope.


% suppress persistent sharp contributions within the admissible scope and can reduce posterior displacement relative to the prior.


% concentrates both curvature and KL accumulation in a single geometry, which can amplify cross-task interference .

% This shared scope forces all tasks to traverse the same restricted geometry, so persistent admissible sharp directioncan repeatedly affect subsequent updates and amplify cross-task interference.
% Sharing a single geometry concentrates both curvature weighting and
% KL accumulation in a single geometry, which can amplify
% cross-task interference when high curvature direction of previous task remain admissible for future task.
% Sharing a single geometry concentrates curvature exposure and task KL accumulation, which can exacerbate cross-task interference when restricted high-curvature modes persist across tasks. 
% Sharing a single geometry concentrates both curvature weighting and
% KL accumulation in a single geometry, which can amplify
% cross-task interference when high curvature direction of previous task remain admissible for future task. Alternatively, one may allow the prior-generation mechanism to adapt across taskst.
% IncLoRA relaxes the fixed-scope assumption by expanding $\mathcal W_{\Delta,t}$ over time, which can avoid forcing task-specific directions to be represented in the same shared coordinates, reduce mismatch between $Q_t^\Delta$ and $P_t^\Delta$ when accommodating novel directions, thereby lowering $\mathrm{KL}(Q_t^\Delta\Vert P_t^\Delta)$. and mitigate repeated excitation of a few persistent sharp modes.
 

% It can also reduce repeated excitation of a few persistent sharp direction by distributing updates across a larger admissible space.

% IncLoRA can be interpreted as a relaxation of the SeqLoRA assumption $\mathcal W_{\Delta,t}\equiv\mathcal W_\Delta$. Instead of reusing a fixed admissible scope, IncLoRA lets the scope expand over time.  
% This expansion can reduce the need to absorb task mismatch within a single fixed geometry, potentially lowering the posterior displacement required to fit novel directions and mitigating repeated excitation of a small set of persistent sharp modes inside the admissible space.
% Constraint-augmented design further constrains how the new degrees of freedom are introduced, aiming to reduce harmful overlap with history-induced important directions. In the hyper-state $\phi_t=(\mathcal W_{\Delta,t},\mu_{P,t},\Sigma_{P,t})$, this corresponds to imposing structural constraints on the evolution of $\mathcal W_{\Delta,t}$ and the initialization $(\mu_{P,t},\Sigma_{P,t})$, for example via parameter-level orthogonality regularization\cite{wang2023orthogonal}, gradient projection~\citep{Liang_2024_CVPR,yang2024parametercollisionhinderingcontinual} or inherit the previous posterior to initialization. These constraints suppress persistent sharp contributions within the admissible scope and can reduce posterior displacement relative to the prior.

% thereby suppressing persistent sharp contributions within the admissible scope and, when aligned with task relatedness, reducing posterior displacement relative to the history-induced prior while stabilizing the evolution of $\phi_t$.



 % IncLoRA expands the scope by increasing $k_t=\dim(\mathcal W_{\Delta,t})$, which can alleviate repeated excitation of a small set of sharp shared modes and reduce the posterior shift needed to fit novel directions, at the price of an explicit drift cost because $\mathcal W_{\Delta,t}$ itself becomes time-varying. Constraint-augmented designs (e.g., OLoRA) reshape $\mathcal W_{\Delta,t}$ to reduce harmful overlap with previously important modes, thereby suppressing persistent sharp contributions within the admissible scope and, when aligned with task relatedness, reducing posterior displacement relative to the history-induced prior while stabilizing the evolution of $\phi_t$.

% IncLoRA induces an expanding sequence $\{\mathcal W_{\Delta,t}\}$ with increasing $k_t$, which can redistribute task adaptation across additional admissible directions.

% Constraint-augmented designs such as OLoRA shape ${\mathcal W_{\Delta,t}}$ to reduce overlap across tasks, which in turn can lower the restricted curvature coupling in $\operatorname{tr}(C_{\Delta,t}\Sigma_t)$ and reduce the mean-shift.

% SeqLoRA, IncLoRA, and constraint-augmented designs can be viewed as three rules for how the hyper-state $\phi_t=(\mathcal W_{\Delta,t},\mu_{P,t},\Sigma_{P,t})$ evolves across tasks, thereby shaping both the sharpness and the KL complexity in Theorem~\ref{thm:pathwise_pac}.


% SeqLoRA keeps an approximately fixed scope $\mathcal W_{\Delta,t}\equiv\mathcal W_\Delta$, so persistent admissible high-curvature modes can repeatedly contribute to $\operatorname{tr}(C_{\Delta,t}\Sigma_t)$, while task mismatch must be absorbed within the same coordinates, typically increasing $\mathbb E_{P_t\sim\mathcal P_{t-1}}\mathrm{KL}(Q_t^\Delta\Vert P_t^\Delta)$. This mismatch can be paid either by revising the prior generator, which increases the cumulative hyperdrift $\sum_t\mathrm{KL}(\mathcal P_t\Vert\mathcal P_{t-1})$, or by adopting the deterministic special case $P_t^\Delta:=Q_{t-1}^\Delta$, which sets $\phi_t=(\mathcal W_\Delta,\mu_{Q,t-1},\Sigma_{Q,t-1})$ and reduces the complexity to $\sum_t\mathrm{KL}(Q_t^\Delta\Vert Q_{t-1}^\Delta)$. 


% SeqLoRA keeps an approximately fixed scope $\mathcal W_{\Delta,t}\equiv\mathcal W_\Delta$, so persistent admissible high-curvature modes can repeatedly contribute to the restricted coupling $\operatorname{tr}(C_{\Delta,t}\Sigma_t)$, while task mismatch must be absorbed within the same coordinates, often increasing $\mathbb E_{P_t\sim\mathcal P_{t-1}}\mathrm{KL}(Q_t^\Delta\Vert P_t^\Delta)$. If the prior generator is revised, the cumulative hyperdrift.


% where $\mathcal W_{\Delta,t}$ is the admissible update scope, and $(\mu_{P,t},\Sigma_{P,t})$ parameterize a Gaussian prior on the corresponding coordinates,
% $
% P_t(\cdot\mid \phi_t)=\mathcal N(\mu_{P,t},\Sigma_{P,t})
% $ supported on $\mathcal W_{\Delta,t}$.
% The hyperposterior $\mathcal P_{t-1}$ is then a distribution over such hyper-states, and sampling $P_t\sim\mathcal P_{t-1}$ can be understood as first drawing $\phi_t\sim\mathcal P_{t-1}$ and then instantiating the corresponding task prior $P_t(\cdot\mid\phi_t)$.
% In Theorem~\ref{thm:pecl-bound}, LoRA organization affects the bound through two coupled channels: a sharpness-dependent empirical term and a KL-based complexity term. The empirical term is governed by the admissible update subspaces ${\mathcal W_{\Delta,t}}$ and the perturbation  supported on them.
%, while the complexity term is governed by task-level distributions on $\mathcal W_{\Delta,t}$
% In Theorem~\ref{thm:pecl-bound}, architectural choices in LoRA-based PECL affect the bound through the induced sequence of admissible update subspaces $\{\mathcal W_{\Delta,t}\}_{t=1}^T$ and the covariance geometry $\Sigma_t$ supported on them.

% The subspace design determines which curvature directions are accessible through $C_{\Delta,t}$ and how strongly they are probed through the spectrum and range of $\Sigma_t$.
% Let $\Pi_{\Delta,t}$ be the orthogonal projector, and let $C_t$ be a task-dependent curvature operator (e.g., a Hessian, Gauss-Newton, or Fisher information matrix) evaluated at $W_0+\hat{\Delta W}_t$. Under a second-order Taylor approximation of $\hat L_{S_t}$ around $W_0+\hat{\Delta W}_t$ and for $\varepsilon\sim\mathcal N(0,\Sigma_t)$ supported on $\mathcal W_{\Delta,t}$, the expected loss elevation satisfies
% $$
% \mathrm{RS}_t^{(0),\Delta}(\hat{\Delta W}_t,\Sigma_t)
% \approx
% \frac12\operatorname{tr}\big(H_{\Delta,t}\Sigma_t\big),
% F_{\Delta,t}=\Pi_{\Delta,t}H_t\Pi_{\Delta,t},
% $$
% which shows that subspace design governs the empirical term by (i) restricting the observable curvature through $\Pi_{\Delta,t}$ and (ii) modulating how strongly each admissible direction is probed through the spectrum of $\Sigma_t$.
% Let $\Pi_{\Delta,t}$ denote the orthogonal projector onto the admissible update subspace $\mathcal W_{\Delta,t}$, and let $C_t=C(\hat W_t;S_t)$ be a task-dependent curvature operator (e.g., Hessian, Gauss--Newton, or Fisher) of the empirical loss evaluated at $\hat W_t=W_0+\hat{\Delta W}_t$. Although $C_t$ is naturally defined on the full parameter space, the PECL perturbations satisfy $\varepsilon\in\mathcal W_{\Delta,t}$, so only the restriction of curvature to admissible directions enters the sharpness term. Under a second-order local approximation of $\hat L_{S_t}$ around $\hat W_t$ and for $\varepsilon\sim\mathcal N(0,\Sigma_t)$ supported on $\mathcal W_{\Delta,t}$, the expected loss elevation admits
% $$
% \mathrm{RS}_t^{(0),\Delta}(\hat{\Delta W}_t,\Sigma_t)
% \approx
% \frac12\operatorname{tr}\!\big(C_{\Delta,t}\Sigma_t\big),
% C_{\Delta,t}=\Pi_{\Delta,t} C_t \Pi_{\Delta,t}.
% $$
% Thus, subspace design governs the empirical contribution through two coupled effects: it determines which curvature directions are accessible via the restriction $C_{\Delta,t}$, and it controls how strongly these admissible directions are probed through the spectrum of $\Sigma_t$.
% The same design simultaneously governs the complexity term through the geometry of task-level distributions on $\mathcal W_{\Delta,t}$.
% For Gaussian task posteriors and priors on the adapter coordinates,
% $Q_t^\Delta=\mathcal N(\mu_{Q,t},\Sigma_{Q,t})$ and $P_t^\Delta=\mathcal N(\mu_{P,t},\Sigma_{P,t})$ with $k_t=\dim(\mathcal W_{\Delta,t})$,
% % the task-wise KL admits the closed form
% \begin{equation*}
%     \begin{aligned}
% &\mathrm{KL}(Q_t^\Delta\Vert P_t^\Delta)
% =
% \frac12\Bigg(
% (\mu_{Q,t}-\mu_{P,t})^\top\Sigma_{P,t}^{-1}(\mu_{Q,t}-\mu_{P,t}) \\
% &\quad  +\log\frac{|\Sigma_{P,t}|}{|\Sigma_{Q,t}|}
% +\operatorname{tr}\big(\Sigma_{P,t}^{-1}\Sigma_{Q,t}\big)
% -k_t
% \Bigg).
%     \end{aligned}
% \end{equation*}
% $$
% \mathrm{KL}(Q_t^\Delta\Vert P_t^\Delta)
% =
% \frac12\Bigg(
% \log\frac{|\Sigma_{P,t}|}{|\Sigma_{Q,t}|}
% -k_t
% +\operatorname{tr}\big(\Sigma_{P,t}^{-1}\Sigma_{Q,t}\big)
% +
% (\mu_{Q,t}-\mu_{P,t})^\top\Sigma_{P,t}^{-1}(\mu_{Q,t}-\mu_{P,t})
% \Bigg).
% $$
% This expression makes explicit how LoRA organization modulates the bound through the mean drift$(\mu_{Q,t}-\mu_{P,t})$, the mismatch between uncertainty shapes $(\Sigma_{Q,t},\Sigma_{P,t})$ and the dimension $k_t$.
% In continual learning, these task wise effects further interact with the hyperposterior drift $\sum_t \mathrm{KL}(\mathcal P_t\Vert \mathcal P_{t-1})$, because the evolution of prior selection is history-dependent and sensitive to how strongly successive tasks couple through the evolving sequence $\{\mathcal W_{\Delta,t}\}$.
% since stronger coupling within a heavily reused subspace typically forces larger and more frequent updates of the prior-selection mechanism.


% SeqLoRA corresponds to a fixed admissible subspace across tasks. This concentrates both curvature weighting and KL accumulation in a single geometry, which can amplify cross-task interference when high curvature direction remain admissible.
 % and posterior displacement in the same geometry across tasks, increasing the likelihood that high curvature directions induced by earlier tasks remain accessible and can be repeatedly amplified by future admissible updates.
% From this viewpoint, SeqLoRA corresponds to a fixed admissible subspace
% , so both $\operatorname{tr}(C_{\Delta,t}\Sigma_t)$ and the KL increments accumulate within a fixed  fixed $\mathcal W_{\Delta,t}$ across tasks. 
%geometry across tasks.
% which maximizes reuse but concentrates both curvature probing and posterior displacement in the same geometry.
% IncLoRA expands $\mathcal W_{\Delta,t}$ over time, increasing $k_t$ and allocating additional directions that can reduce competition under comparable KL and drift budgets.

%under comparable uncertainty scales.
% Constraint-augmented designs such as OLoRA shape ${\mathcal W_{\Delta,t}}$ to reduce overlap across tasks, which in turn can lower the restricted curvature coupling in $\operatorname{tr}(C_{\Delta,t}\Sigma_t)$ and reduce the mean-shift.


% can be interpreted as shaping $\{\mathcal W_{\Delta,t}\}$ and the induced priors so that successive task updates have reduced overlap in admissible directions, thereby lowering both the curvature coupling $\operatorname{tr}(C_{\Delta,t}\Sigma_t)$ and the mean-shift component of $\mathrm{KL}(Q_t^\Delta\Vert P_t^\Delta)$.

% As a complementary remark, classical CL methods such as EWC\cite{EWC2017}, SI\cite{pmlr-v70-zenke17a}, and MAS\cite{Aljundi_2018_ECCV} can be interpreted as KL-shaping mechanisms on adapter coordinates. They anchor the prior mean at a history-dependent reference and set the prior precision using importance or curvature estimates, which acts by penalizing the mean-shift component of $\mathrm{KL}(Q_t^\Delta\Vert P_t^\Delta)$ along important directions of previous tasks.

% along important directions in $\mathcal W_{\Delta,t}$. This is complementary to sharpness-aware designs, which primarily act through the empirical term by controlling the perturbation geometry $\Sigma_t$.



% can be interpreted within our bound as imposing an implicit constraint on the task wise KL term through the choice of the Gaussian prior on adapter coordinates. Concretely, these methods assign an importance score to parameters learned from past tasks, and then discourage subsequent updates along directions deemed important, by anchoring the prior mean at a history dependent reference solution and shaping the prior precision according to an importance or curvature estimate. For example a Fisher based estimate in EWC, an online path based importance in SI, or function sensitivity based importance in MAS. Under the Gaussian KL closed form in Theorem~\ref{thm:pecl-bound}, this design primarily acts by tightening the mean shift component of $\mathrm{KL}(Q_t^\Delta\Vert P_t^\Delta)$ within $\mathcal W_{\Delta,t}$, thereby explicitly limiting admissible drift that would otherwise amplify past task loss. In this sense, our framework subsumes these consolidation regularizers as KL shaping mechanisms, complementary to sharpness aware designs that act through the empirical term via the perturbation covariance $\Sigma_t$.
% As a complementary remark, classical consolidation regularizers such as EWC, SI, and MAS can be cast in the same KL framework by an implicit choice of the task prior $P_t^\Delta$ in the Gaussian family. Concretely, if one sets $P_t^\Delta=\mathcal N(\mu_{P,t},\Sigma_{P,t})$ with $\mu_{P,t}$ centered at a history dependent reference solution and with a precision $\Sigma_{P,t}^{-1}$ aligned to an importance or curvature estimate on $\mathcal W_{\Delta,t}$ (e.g., projected Fisher or parameter-importance weights), then the KL closed form implies that the dominant $\mu_{Q,t}$ dependent term is the mean shift, which recovers the familiar Fisher weighted quadratic penalties up to constants independent of $\mu_{Q,t}$. 
% Under our PECL bound, such consolidation mechanisms therefore primarily shape the KL budget through the prior geometry on $\mathcal W_{\Delta,t}$, constraining the magnitude of admissible updates along important directions.

% Under our PECL bound, such consolidation mechanisms therefore primarily shape the KL budget through the prior geometry on $\mathcal W_{\Delta,t}$, whereas sharpness-aware designs primarily act through the perturbation covariance $\Sigma_t$ to reduce the sharpness-dependent empirical contribution.












\section{Empirical}

\label{sec:results-local-global}





\subsection{Experimental Setup}



% \begin{table*}[htpb]
% \centering
% \caption{
% Class incremental performance, weight-space flatness and feature-space robustness on ImageNet-R ($T=20$, $R=16$). 
% }
% \label{tab:table_all}
% \setlength{\tabcolsep}{3pt}
% \renewcommand{\arraystretch}{0.95}
% \resizebox{\linewidth}{!}{
% \begin{tabular}{l l cccc cc cc}
% \toprule
% \multirow{2}{*}{\textbf{Method}} & \multirow{2}{*}{\textbf{Opt.}} &
% \multicolumn{4}{c}{\textbf{CL Performance} ($\uparrow$)} &
% \multicolumn{2}{c}{\textbf{Weight Flatness} ($\downarrow$)} &
% \multicolumn{2}{c}{\textbf{Feature Robustness}} \\
% \cmidrule(lr){3-6}\cmidrule(lr){7-8}\cmidrule(lr){9-10}
% & &
% $\operatorname{ACC}_{20}^{\mathrm{cls}}$ &
% $\operatorname{AAA}^{\mathrm{cls}}$ &
% $\operatorname{ACC}_{20}^{\mathrm{nme}}$ &
% $\operatorname{AAA}^{\mathrm{nme}}$ &
% Sh$^{(0)}(\rho)$ &
% $\mathrm{tr}(F)$ &
% L2 drift ($\downarrow$) &
% $\mathrm{tr}(E)$ \\
% \midrule



\paragraph{Datasets and task protocol}
We evaluate PECL methods on vision benchmarks derived from ImageNet. ImageNet-R~\cite{ImagenetR_Hendrycks_2021_ICCV} serves as our main class-incremental benchmark and is split into $T=20$ tasks over $200$ classes following prior work~\cite{liangInfLoRAInterferenceFreeLowRank2024a}.  To assess robustness under corruptions and perturbations, we further use Tiny ImageNet-C and Tiny ImageNet-P~\cite{hendrycks2018benchmarking}. 
Tiny ImageNet-C applies common corruptions to Tiny ImageNet validation images, while Tiny ImageNet-P provides short perturbation sequences per image.



% \begin{table*}[htpb]
% \centering
% \caption{
% Class incremental performance, weight-space flatness and feature-space robustness on ImageNet-R ($T=20$, $R=16$). 
% }
% \label{tab:table_all}
% \setlength{\tabcolsep}{4pt}
% \renewcommand{\arraystretch}{0.95}
% \resizebox{\linewidth}{!}{
% \begin{tabular}{l l cccc cc cc}
% \toprule
% \multirow{2}{*}{\textbf{Method}} & \multirow{2}{*}{\textbf{Opt.}} &
% \multicolumn{4}{c}{\textbf{CL Performance} ($\uparrow$)} &
% \multicolumn{2}{c}{\textbf{Weight Flatness} ($\downarrow$)} &
% \multicolumn{2}{c}{\textbf{Feature Robustness}} \\
% \cmidrule(lr){3-6}\cmidrule(lr){7-8}\cmidrule(lr){9-10}
% & &
% $\operatorname{ACC}_{20}^{\mathrm{cls}}$ &
% $\operatorname{AAA}^{\mathrm{cls}}$ &
% $\operatorname{ACC}_{20}^{\mathrm{nme}}$ &
% $\operatorname{AAA}^{\mathrm{nme}}$ &
% Sh$^{(0)}(\rho)$ &
% $\mathrm{tr}(F)$ &
% L2 drift ($\downarrow$) &
% $\mathrm{tr}(E)$ \\
% \midrule
% \multirow{1}{*}{\textbf{PerTaskLP}}
% & SGD
% & $55.22^{\pm 0.31}$ & $59.36^{\pm 0.19}$
% & $62.96^{\pm 0.12}$ & $66.58^{\pm 0.37}$
% & 0.072 & 1015.07
% & 0.00 & 0.82 \\
% \midrule
% \multirow{3}{*}{\textbf{SeqFT}}
% & SGD
% & $50.72^{\pm 2.20}$ & $59.89^{\pm 1.32}$
% & $69.49^{\pm 0.16}$ & $75.10^{\pm 0.57}$
% & 0.582 & 29703.66
% & 29.57 & 0.58 \\
% & SAM
% & $56.07^{\pm 2.60}$ & $66.47^{\pm 1.53}$
% & $72.36^{\pm 2.46}$ & $78.31^{\pm 0.50}$
% & 0.259 & 5881.32
% & 17.04 & 0.74 \\
% & GAM
% & $59.78^{\pm 0.74}$ & $67.31^{\pm 1.18}$
% & $73.74^{\pm 0.88}$ & $78.32^{\pm 0.94}$
% & 0.177 & 3273.08
% & 13.80 & 0.78 \\
% \midrule
% \multirow{3}{*}{\textbf{SeqLoRA}}
% & SGD
% & $63.14^{\pm 2.17}$ & $69.07^{\pm 1.33}$
% & $72.82^{\pm 1.27}$ & $76.62^{\pm 0.68}$
% & 0.142 & 1401.15
% & 16.58 & 0.77 \\
% & SAM
% & $67.04^{\pm 0.71}$ & $71.42^{\pm 0.71}$
% & $74.53^{\pm 0.60}$ & $77.62^{\pm 0.43}$
% & 0.120 & 821.45
% & 14.82 & 0.81 \\
% & GAM
% & $73.31^{\pm 0.52}$ & $76.80^{\pm 0.56}$
% & $75.74^{\pm 0.15}$ & $78.86^{\pm 0.53}$
% & 0.060 & 240.91
% & 11.28 & 0.78 \\
% \midrule
% \multirow{3}{*}{\textbf{IncLoRA}}
% & SGD
% & $65.77^{\pm 2.84}$ & $71.29^{\pm 0.78}$
% & $75.41^{\pm 0.26}$ & $78.37^{\pm 0.26}$
% & 0.076 & 836.91
% & 16.87 & 0.67 \\
% & SAM
% & $68.86^{\pm 0.92}$ & $72.59^{\pm 0.39}$
% & $75.95^{\pm 0.07}$ & $78.68^{\pm 0.25}$
% & 0.053 & 526.49
% & 15.14 & 0.70 \\
% & GAM
% & $72.86^{\pm 0.23}$ & $76.52^{\pm 0.48}$
% & $76.72^{\pm 0.54}$ & $79.42^{\pm 0.40}$
% & 0.035 & 192.25
% & 12.15 & 0.73 \\
% \midrule
% \multirow{3}{*}{\textbf{OLoRA}}
% & SGD
% & $66.64^{\pm 0.44}$ & $71.41^{\pm 0.47}$
% & $75.10^{\pm 0.41}$ & $78.65^{\pm 0.67}$
% & 0.078 & 831.23
% & 17.23 & 0.68 \\
% & SAM
% & $68.53^{\pm 0.78}$ & $72.75^{\pm 0.76}$
% & $75.73^{\pm 0.32}$ & $78.70^{\pm 0.67}$
% & 0.066 & 628.12
% & 15.17 & 0.70 \\
% & GAM
% & $72.44^{\pm 0.36}$ & $76.05^{\pm 0.52}$
% & $76.56^{\pm 0.53}$ & $79.16^{\pm 0.43}$
% & 0.036 & 198.30
% & 12.16 & 0.74 \\
% % \midrule
% % \multirow{1}{*}{\textbf{l2p}}
% % & Adam
% % & $71.35^{\pm 0.05}$ & $74.84^{\pm 0.73}$
% % & -- & --
% % & -- & --
% % & -- & -- \\
% % \midrule
% % \multirow{1}{*}{\textbf{dualprompt}}
% % & Adam
% % & $71.59^{\pm 0.68}$ & $74.80^{\pm 1.18}$
% % & -- & --
% % & -- & --
% % & -- & -- \\
% \bottomrule
% \end{tabular}
% }
% \end{table*}

% \paragraph{Models}
\noindent\textbf{Models}\quad
% All experiments use a ViT-B/16 backbone~\cite{dosovitskiy2020image} pre-trained on ImageNet-21K. 
% % To probe the intrinsic representation capacity of the backbone, we include a PerTask Linear Probe baseline that trains an independent linear classifier per task. 
% For PECL, we attach LoRA adapters to each transformer block and study three adapter  representative organizations: SeqLoRA, IncLoRA, and OLoRA\cite{wang2023orthogonal}.
% % SeqLoRA uses a single adapter that is updated across all tasks. IncLoRA incrementally adds task-specific adapters on top of a shared backbone. OLoRA augments this incremental design with orthogonality-promoting constraints. 
% % All methods share the same backbone and training budget so that differences can be attributed to adapter geometry and flatness optimization.
% For comparison with perturbations scope in PECL, we use Flat-LoRA \cite{li2025flatloralowrankadaptationflat} that applies $\mathrm{RS}_S^{(0)}({W}, \Sigma)$ to all model parameters instead of restricting updates to the adapters. Unless stated otherwise, perturbations and optimizer steps are applied only to the active LoRA parameters of the current task, while the backbone and previously learned adapters are excluded from both perturbation and update
% To evaluate different perturbation type, we integrate several sharpness-based methods. 
% SAM~\cite{foretSharpnessAwareMinimizationEfficiently2021} optimize $\mathrm{AS}^{(0)}_S(W,\rho)$. RWP~\cite{pmlr-v151-bisla22a} apply $\mathrm{RS}_S^{(0)}({W}, \Sigma)$ and GAM~\cite{zhangGradientNormAware2023} optimize $\mathrm{AS}^{(1)}_S(W,\rho)$. In the PECLthese optimizers act only on the trainable adapter parameters. 
% \paragraph{Models.}
All experiments use a ViT-B/16 backbone~\cite{dosovitskiy2020image} pre-trained on ImageNet-21K.
We attach LoRA adapters to each transformer block and consider three organizations: SeqLoRA, IncLoRA, and OLoRA~\cite{wang2023orthogonal}.
Unless stated otherwise, PECL updates and perturbations are applied only to
to the trainable LoRA coordinates of the current task, while the backbone and past adapters remain frozen.

%; the induced weight perturbation is $\Delta W(\theta_\Delta+\xi)-\Delta W(\theta_\Delta)$.

% the trainable LoRA parameters

To study perturbation {scope}, we include Flat-LoRA~\cite{li2025flatloralowrankadaptationflat} for PECL, which defines $\mathrm{RS}^{(0)}_S(W,\Sigma)$ perturbations in the full parameter space.
% To study perturbation {type}, we compare SAM~\cite{foretSharpnessAwareMinimizationEfficiently2021} ($\mathrm{AS}^{(0)}_S(W,\rho)$), RWP~\cite{pmlr-v151-bisla22a} ($\mathrm{RS}^{(0)}_S(W,\Sigma)$), and GAM~\cite{zhangGradientNormAware2023} ($\mathrm{AS}^{(1)}_S(W,\rho)$), all acting on the trainable adapters in PECL.
To study perturbation \emph{type}, we compare SAM~\cite{foretSharpnessAwareMinimizationEfficiently2021}, RWP~\cite{pmlr-v151-bisla22a}, and GAM~\cite{zhangGradientNormAware2023}, which optimize different sharpness objectives: SAM minimizes adversarial zeroth-order sharpness $\mathrm{AS}^{(0)}_S(W,\rho)$, RWP targets random zeroth-order sharpness $\mathrm{RS}^{(0)}_S(W,\Sigma)$, and GAM targets adversarial first-order sharpness $\mathrm{AS}^{(1)}_S(W,\rho)$. For a controlled comparison in PECL, all methods apply their perturbations only to the trainable LoRA coordinates.
% $\theta_\Delta$; the induced weight perturbation is $\Delta W(\theta_\Delta+\xi)-\Delta W(\theta_\Delta)$.
%all methods apply their perturbations only to the trainable adapters.
Unless otherwise stated, we fix the perturbation radius to $\rho=0.05$, following the standard empirical setting used in SAM, and use a dedicated radius sweep only when diagnosing the scope effect.
% Unless otherwise specified, we fix $\rho=0.05$, to ensure a controlled comparison where differences are attributable to scope/type rather than tuning of perturbation strength.


% \paragraph{Metrics}
\noindent\textbf{Metrics}\quad
% We adopt standard continual-learning metrics based on the task accuracy matrix $a_{tj}$, where $a_{tj}$ denotes the accuracy on task $j$ after learning task $t$. The accuracy at step $t$ is
% $
% \operatorname{ACC}_t = \frac{1}{t} \sum_{j=1}^{t} a_{tj},
% $
% and we report final accuracy $\operatorname{ACC}_T^{\mathrm{cls}}$, time-averaged accuracy $\operatorname{AAA}^{\mathrm{cls}} = \frac{1}{T}\sum_{t=1}^{T} \operatorname{ACC}_t.$
% and backward transfer $\operatorname{BWT}^{\mathrm{cls}}$.  , \operatorname{BWT}^{\mathrm{nme}}
% In addition, we compute the corresponding Nearest-Class-Mean Classifier (NCM) results
% % Nearest Mean of Exemplars (NME) metrics 
% $(\operatorname{ACC}_T^{\mathrm{ncm}}, \operatorname{AAA}^{\mathrm{ncm}})$ using class prototypes in the feature space, which isolates the quality and stability of representations from the classifier head.
% We adopt standard continual-learning metrics accuracy on overall task  $\operatorname{ACC}_t^{\mathrm{cls}}$, and also  compute 
% Nearest Mean of Exemplars (NME) metrics 
% $\operatorname{ACC}_t^{\mathrm{ncm}}$ to check performance.
% % time-averaged accuracy $\operatorname{AAA}^{\mathrm{cls}} = \frac{1}{T}\sum_{t=1}^{T} \operatorname{ACC}_t$.
% % with the corresponding Nearest-Class-Mean Classifier (NCM) results.
% % In addition of standard continual-learning metrics, we compute 
% % % Nearest Mean of Exemplars (NME) metrics 
% % $(\operatorname{ACC}_T^{\mathrm{ncm}}, \operatorname{AAA}^{\mathrm{ncm}})$ using class prototypes in the feature space, which isolates the quality and stability of representations from the classifier head.
% To characterize flatness, we approximate the maximum eigenvalue and the trace of the Hessiean matric and Fisher information using Lanczos methods~\cite{ghorbaniInvestigationNeuralNet2019,foretSharpnessAwareMinimizationEfficiently2021}.  
We adopt standard continual-learning metrics based on the task accuracy matrix $a_{tj}$, where $a_{tj}$ denotes the accuracy on task $j$ after learning task $t$. The accuracy at step $t$ is
$
\operatorname{ACC}_t = \frac{1}{t} \sum_{j=1}^{t} a_{tj}
$.
We report class-incremental performance using the task-wise classification accuracy $\operatorname{ACC}_t^{\mathrm{cls}}$.
In addition, we evaluate representation quality via Nearest Mean of Exemplars (NME), denoted by $\operatorname{ACC}_t^{\mathrm{nme}}$.
% This metric reduces the influence of the classifier head and provides a complementary view of how stable and transferable the learned representations are across tasks.
To quantify flatness, we estimate the largest eigenvalue $\lambda_{\max}$ and the trace $\mathrm{tr}(\cdot)$ of both the Hessian and the Fisher information matrix using Lanczos~\cite{ghorbaniInvestigationNeuralNet2019,foretSharpnessAwareMinimizationEfficiently2021}.



% To compare the overlap between subspact in the training process, we evalute  the distance of subspace $
% d_P\big(\Delta W_{t},\Delta W_{t-1}\big)
% $ using Projection metric \cite{hammGrassmannDiscriminantAnalysis2008}. 



% use Lanczos methods~\cite{ghorbaniInvestigationNeuralNet2019,foretSharpnessAwareMinimizationEfficiently2021} to compute the top eigenvalues of Hessian and Hutchinson’s method [4, 5, 65] to compute the Hessian trace.
% We further measure the $\ell_2$ drift of class prototypes across tasks and the spectrum of the empirical feature matrix. 
% Together, these quantities provide complementary views on sharpness and feature stability over the task sequence.
% We use power iteration [66] to compute the top eigenvalues of Hessian and Hutchinson’s method [4, 5, 65] to compute the Hessian trace.





% \subsection{Main Reuslt}


% \paragraph{CL Benchmark}
% Table~\ref{table:q1_ci_with_weight_flatness_reorg} reports, on ImageNet-R, the final and time-averaged accuracies for both classifier heads and NCM evaluation, together with weight-space flatness proxies (adversarial zero-order sharpness $\mathrm{AS}^{(0)}_S(W,\rho)$, largest eigenvalue of curvature $\lambda_{\max} C$ and trace of curvature $tr(C)$) . 

% \paragraph{Robustness Evaluation}
% We next  robustness under corruptions and perturbations on Tiny ImageNet-C . Fig.~\ref{fig:robustness_curves_comparison} shows the trajectories of $\operatorname{ACC}_t$ over the task sequence. 

% \paragraph{NLP benchmarks} 
% We further conduct continual task learning experiments on two standard NLP benchmarks under multiple task orders, and report the final accuracy $ACC_T$ in Table~\ref{tab:nlp_accT_orders}.





% \begin{table*}[t]
% \centering
% \caption{
% Class incremental performance, weight-space flatness and feature-space robustness on ImageNet-R ($T=20$, $R=16$).
% }
% \label{tab:table_all}

% % 上下空白略微压缩
% % \vspace{-0.1em}

% % 列间距和行高压缩一点
% \setlength{\tabcolsep}{2pt}      % 原来是 3pt
% \renewcommand{\arraystretch}{0.85} % 原来是 0.95
% \small                            % 字体略小一号

% % 如果不想整体缩放，可以去掉 resizebox
% \resizebox{0.97\textwidth}{!}{   % 比原来略缩小一点，顺带减小高度
% \begin{tabular}{l l cccc ccc c}
% \toprule
% \multirow{2}{*}{\textbf{Method}} & \multirow{2}{*}{\textbf{Opt.}} &
% \multicolumn{4}{c}{\textbf{CL Performance} ($\uparrow$)} &
% \multicolumn{3}{c}{\textbf{Weight Flatness} ($\downarrow$)} &
% \multicolumn{1}{c}{\textbf{Feature Robustness}} \\
% \cmidrule(lr){3-6}\cmidrule(lr){7-9}\cmidrule(lr){10-10}
% & &
% $\operatorname{ACC}_{20}^{\mathrm{cls}}$ &
% $\operatorname{AAA}^{\mathrm{cls}}$ &
% $\operatorname{ACC}_{20}^{\mathrm{ncm}}$ &
% $\operatorname{AAA}^{\mathrm{ncm}}$ &
% Sh$^{(0)}(\rho)$ 
% & $\lambda_{\max}(F)$ &
% $tr(F)$ &
% L2 drift ($\downarrow$) \\
% \midrule
% \multirow{1}{*}{\textbf{PerTaskLP}}
% & SGD
% & $55.22^{\pm 0.31}$ & $59.36^{\pm 0.19}$
% & $62.96^{\pm 0.12}$ & $66.58^{\pm 0.37}$
% & 0.072  & 8.47  & 1015.07  
% & 0.00 \\
% \midrule
% \multirow{3}{*}{\textbf{SeqFT}}
% & SGD
% & $50.72^{\pm 2.20}$ & $59.89^{\pm 1.32}$
% & $69.49^{\pm 0.16}$ & $75.10^{\pm 0.57}$
% & 0.582   & 538.79               & 29703.66
% & 29.57 \\
% & + $\mathrm{AS}^{0}_S(W,\rho)$
% & $56.07^{\pm 2.60}$ & $66.47^{\pm 1.53}$
% & $72.36^{\pm 2.46}$ & $78.31^{\pm 0.50}$
% & 0.259 & 75.43                & 5881.32
% & 17.04  \\
% & + $\mathrm{AS}^{(1)}_S(W,\rho)$
% & $\mathbf{59.78}^{\pm 0.74}$ & \textbf{$\mathbf{67.31}^{\pm 1.18}$}
% & $\mathbf{73.74}^{\pm 0.88}$ & \textbf{$\mathbf{78.32}^{\pm 0.94}$}
% & 0.177 & 32.00                & 3273.08      
% & 13.80  \\
% \midrule
% \multirow{3}{*}{\textbf{SeqLoRA}}
% & SGD
% & $63.14^{\pm 2.17}$ & $69.07^{\pm 1.33}$
% & $72.82^{\pm 1.27}$ & $76.62^{\pm 0.68}$
% & 0.142 & 35.51                & 1401.15
% & 16.58  \\
% & + $\mathrm{AS}^{(1)}_S(W,\rho)$
% & $67.04^{\pm 0.71}$ & $71.42^{\pm 0.71}$
% & $74.53^{\pm 0.60}$ & $77.62^{\pm 0.43}$
% & 0.120 & 21.39                & 821.45
% & 14.82  \\
% &  + $\mathrm{AS}^{(1)}_S(W,\rho)$
% & $\mathbf{73.31}^{\pm 0.52}$ & \textbf{$\mathbf{76.80}^{\pm 0.56}$}
% & $\mathbf{75.74}^{\pm 0.15}$ & \textbf{$\mathbf{78.86}^{\pm 0.53}$}
% & 0.060 & 3.72                 & 240.91 
% & 11.28  \\
% \midrule
% \multirow{3}{*}{\textbf{IncLoRA}}
% & SGD
% & $65.77^{\pm 2.84}$ & $71.29^{\pm 0.78}$
% & $75.41^{\pm 0.26}$ & $78.37^{\pm 0.26}$
% & 0.076 & 8.89                 & 836.91    
% & 16.87  \\
% & + $\mathrm{AS}^{(0)}_S(W,\rho)$
% & $68.86^{\pm 0.92}$ & $72.59^{\pm 0.39}$
% & $75.95^{\pm 0.07}$ & $78.68^{\pm 0.25}$
% & 0.053 & 4.35                 & 526.49     
% & 15.14  \\
% &  + $\mathrm{AS}^{(1)}_S(W,\rho)$
% % & $72.86^{\pm 0.23}$ & \textbf{$76.52^{\pm 0.48}$}
% % & $76.72^{\pm 0.54}$ & \textbf{$79.42^{\pm 0.40}$}
% & $\mathbf{72.86}^{\pm 0.23}$ & $\mathbf{76.52}^{\pm 0.48}$
% & $\mathbf{76.72}^{\pm 0.54}$ & $\mathbf{79.42}^{\pm 0.40}$ 
% & 0.035 & 2.09                 & 192.25 
% & 12.15  \\
% \midrule
% \multirow{3}{*}{\textbf{OLoRA}}
% & SGD
% & $66.64^{\pm 0.44}$ & $71.41^{\pm 0.47}$
% & $75.10^{\pm 0.41}$ & $78.65^{\pm 0.67}$
% & 0.078 & 9.00                 & 831.23
% & 17.23  \\
% & + $\mathrm{AS}^{(0)}_S(W,\rho)$
% & $68.53^{\pm 0.78}$ & $72.75^{\pm 0.76}$
% & $75.73^{\pm 0.32}$ & $78.70^{\pm 0.67}$
% & 0.066 & 7.41                 & 628.12   
% & 15.17 \\
% & + $\mathrm{AS}^{(1)}_S(W,\rho)$
% % & $72.44^{\pm 0.36}$ & \textbf{$76.05^{\pm 0.52}$}
% % & $76.56^{\pm 0.53}$ & \textbf{$79.16^{\pm 0.43}$}
% & $\mathbf{72.44}^{\pm 0.36}$ & $\mathbf{76.05}^{\pm 0.52}$
% & $\mathbf{76.56}^{\pm 0.53}$ & $\mathbf{79.16}^{\pm 0.43}$
% & 0.036& 2.08                 & 198.30  
% & 12.16  \\
% % \midrule
% % \multirow{1}{*}{\textbf{l2p}}
% % & Adam
% % & $71.35^{\pm 0.05}$ & $74.84^{\pm 0.73}$
% % & -- & --
% % & -- & --
% % & -- & -- \\
% % \midrule
% % \multirow{1}{*}{\textbf{dualprompt}}
% % & Adam
% % & $71.59^{\pm 0.68}$ & $74.80^{\pm 1.18}$
% % & -- & --
% % & -- & --
% % & -- & -- \\
% \bottomrule
% \end{tabular}
% }
% \end{table*}


\subsection{Result Analysis}

% \paragraph{Scope of perturbations}
% We examine whether flatness should be promoted locally in the adapter subspace or globally in the full parameter space in PECL by compare $RS^0_S(W,\Sigma_t)$ to all model parameters with restricts to the LoRA branch only as FlatLoRA setting\cite{li2025flatloralowrankadaptationflat}.


% {%
% \setlength{\textfloatsep}{4pt}      % figure* (top/bottom) 与正文的距离（最关键）
% \setlength{\floatsep}{4pt}          % 两个浮动体之间的距离（一般备用）
% \setlength{\abovecaptionskip}{1pt}  % caption 上方
% \setlength{\belowcaptionskip}{0pt}  % caption 下方
% \captionsetup{skip=3pt}             % 图与caption的距离（你也可继续用4pt）

% \begin{figure}[htbp]
%     \centering
% \includegraphics[width=1\linewidth]{figures_TRML/rwp_full_lora_acc20_bars_imagenetr_3result1.pdf}
% \label{fig:perturb_scope}
%     \vspace{-8pt}
%     \caption{Effect of perturbation scope in PECL. Final average accuracy across perturbation strengths $\rho$ under different LoRA organizations.
%     % Restricting perturbations to the task adapter subspace $\Delta W$
%     %
%     % Restricting perturbations to be supported on the task-admissible subspace $\mathcal W_{\Delta,t}$consistently improves over the SGD baseline, whereas applying perturbations in the ambient parameter space $\mathbb R^d$ (perturbing both the frozen backbone and the adapters) substantially degrades performance.
%     %
%     % applying perturbations in the full parameter space $W+\Delta W$ substantially degrades performance.
%     % Final average accuracy under different perturbation strengths $\rho$ for LoRA organizations. Restricting perturbations to the task-admissible adapter subspace $\Delta W$ consistently outperforms the SGD baseline, whereas full-parameter perturbations $W+\Delta W$ substantially degrade performance.
%     % \textbf{(b) Type.} Loss curves along the dominant curvature direction.
%     % The top panel evaluates loss under perturbations restricted to the LoRA coordinates $\Delta W$, while the bottom panel evaluates loss under perturbations defined over the full parameter space $W+\Delta W$.
%     % The top  evaluates the loss when perturbations are restricted to the LoRA coordinates, while the bottom evaluates the loss when perturbations are defined over all parameters.
%     }
%     \label{fig: effect of scope}
% \end{figure}
% }
% \vspace{-4pt}




% \begin{table}[t]
% \centering
% \caption{$\operatorname{AAA}^{\mathrm{cls}}$ under Full vs LoRA-only perturbations.}
% \label{tab:lora_or_full_cls}
% \footnotesize
% \setlength{\tabcolsep}{3pt}
% \renewcommand{\arraystretch}{0.88}
% \begin{tabular}{@{}c c c c c@{}}
% \toprule
% \textbf{Scope} & \textbf{Strength} & \textbf{SeqLoRA} & \textbf{IncLoRA} & \textbf{OLoRA} \\
% \midrule
% Base & --   & 69.07 & 71.29 & 71.41 \\
% \midrule
% \multirow{3}{*}{Full}
%  & 0.05 & -3.50 & -2.00 & -2.62 \\
%  & 0.10 & -4.82 & -4.74 & -5.53 \\
%  & 0.15 & -6.72 & -6.84 & -6.49 \\
% \midrule
% \multirow{3}{*}{LoRA}
%  & 0.05 & 0.27 & 0.79 & 0.12 \\
%  & 0.10 & 0.70 & 0.86 & 0.28 \\
%  & 0.15 & 0.41 & 0.77 & 0.27 \\
% \bottomrule
% \end{tabular}
% \end{table}


% \begin{table}[t]
% \centering
% \caption{$\operatorname{AAA}^{\mathrm{cls}}$ under global (Full) versus subspace (LoRA-only) perturbations.}
% \label{tab:lora_or_full_cls}
% \setlength{\tabcolsep}{4pt}
% \renewcommand{\arraystretch}{1}
% \resizebox{\linewidth}{!}{
% \begin{tabular}{c c c c c}
% \toprule
% \textbf{Scope} & \textbf{Strength} & \textbf{SeqLoRA} & \textbf{IncLoRA} & \textbf{OLoRA} \\
% \midrule
% Base & -   & 69.07 & 71.29 & 71.41 \\
% \midrule
% \multirow{3}{*}{Full}
%  & $0.05$ & -3.50 & -2.00 & -2.62 \\
%  & $0.10$ & -4.82 & -4.74 & -5.53 \\
%  & $0.15$ & -6.72 & -6.84 & -6.49 \\
% \midrule
% \multirow{3}{*}{LoRA}
%  & $0.05$ & 0.27 & 0.79 & 0.12 \\
%  & $0.10$ & 0.70 & 0.86 & 0.28 \\
%  & $0.15$ & 0.41 & 0.77 & 0.27 \\
% \bottomrule
% \end{tabular}}
% \end{table}



% \begin{table}[t]
% \centering
% \caption{$\operatorname{AAA}^{\mathrm{cls}}$ under global (Full) versus subspace (LoRA-only) perturbations.}
% \label{tab:lora_or_full_cls}
% \setlength{\tabcolsep}{4pt}
% \renewcommand{\arraystretch}{1}
% \resizebox{\linewidth}{!}{
% \begin{tabular}{c c c c c}
% \toprule
% \multicolumn{2}{c}{\textbf{Opt.}} &
% \textbf{SeqLoRA} &
% \textbf{IncLoRA} &
% \textbf{OLoRA} \\
% \cmidrule(lr){1-2} \cmidrule(lr){3-3} \cmidrule(lr){4-4} \cmidrule(lr){5-5}
% \textbf{Range} & \textbf{Strength} &
% $\operatorname{AAA}^{\mathrm{cls}}$ &
% $\operatorname{AAA}^{\mathrm{cls}}$ &
% $\operatorname{AAA}^{\mathrm{cls}}$ \\
% \midrule
% Base & $0$   & 69.07 & 71.29 & 71.41 \\
% \midrule
% \multirow{3}{*}{Full}
%  & $0.05$ & -3.50 & -2.00 & -2.62 \\
%  & $0.10$ & -4.82 & -4.74 & -5.53 \\
%  & $0.15$ & -6.72 & -6.84 & -6.49 \\
% \midrule
% \multirow{3}{*}{LoRA}
%  & $0.05$ & 0.27 & 0.79 & 0.12 \\
%  & $0.10$ & 0.70 & 0.86 & 0.28 \\
%  & $0.15$ & 0.41 & 0.77 & 0.27 \\
% \bottomrule
% \end{tabular}}
% \end{table}



% \begin{table}[t]
% \centering
% \caption{$\operatorname{AAA}^{\mathrm{cls}}$ and $\operatorname{AAA}^{\mathrm{ncm}}$ under global (full model) versus subspace (LoRA only).}
% \label{tab:lora_or_full}
% \setlength{\tabcolsep}{3pt}
% \renewcommand{\arraystretch}{1}
% \resizebox{1\linewidth}{!}{
% \begin{tabular}{c c cc cc cc}
% \toprule
% \multicolumn{2}{c}{\textbf{Opt.}} &
% \multicolumn{2}{c}{\textbf{SeqLoRA}} &
% \multicolumn{2}{c}{\textbf{IncLoRA}} &
% \multicolumn{2}{c}{\textbf{OLoRA}} \\
% \cmidrule(lr){1-2} \cmidrule(lr){3-4} \cmidrule(lr){5-6} \cmidrule(lr){7-8}
% \textbf{Range} & \textbf{Strength} &
% $\operatorname{AAA}^{\mathrm{cls}}$ & $\operatorname{AAA}^{\mathrm{ncm}}$ &
% $\operatorname{AAA}^{\mathrm{cls}}$ & $\operatorname{AAA}^{\mathrm{ncm}}$ &
% $\operatorname{AAA}^{\mathrm{cls}}$ & $\operatorname{AAA}^{\mathrm{ncm}}$ \\
% \midrule
% \multirow{1}{*}{Base} & $0$ &
% 69.07 & 76.62 & 71.29 & 78.29 & 71.41 & 78.47 \\
% \midrule
% \multirow{3}{*}{Full} & $0.05$ & -3.50 & -1.54 & -2.00 & -0.38 & -2.62 & -0.80 \\
%                      & $0.10$ & -4.82 & -1.29 & -4.74 & -0.74 & -5.53 & -1.18 \\
%                      & $0.15$ & -6.72 & -1.69 & -6.84 & -0.98 & -6.49 & -1.29 \\
% \midrule
% \multirow{3}{*}{LoRA} & $0.05$ & 0.27 & 0.02 & 0.79 & 0.51 & 0.12 & 0.16 \\
%                      & $0.10$ & 0.70 & 0.29 & 0.86 & 0.53 & 0.28 & 0.12 \\
%                      & $0.15$ & 0.41 & 0.08 & 0.77 & 0.53 & 0.27 & 0.11 \\

% \bottomrule
% \end{tabular}}
% \end{table}



% \begin{table}[t]
% \centering
% \caption{$\operatorname{AAA}^{\mathrm{cls}}$ and $\operatorname{AAA}^{\mathrm{ncm}}$ under full model versus LoRA only.}
% \label{tab:lora——or_full}
% \setlength{\tabcolsep}{3pt}
% \renewcommand{\arraystretch}{1}
% \resizebox{1\linewidth}{!}{
% \begin{tabular}{c c cc cc cc}
% \toprule
% \multicolumn{2}{c}{\textbf{Opt.}} &
% \multicolumn{2}{c}{\textbf{SeqLoRA}} &
% \multicolumn{2}{c}{\textbf{IncLoRA}} &
% \multicolumn{2}{c}{\textbf{OLoRA}} \\
% \cmidrule(lr){1-2} \cmidrule(lr){3-4} \cmidrule(lr){5-6} \cmidrule(lr){7-8}
% \textbf{Strength} & \textbf{Range} &
% $\operatorname{AAA}^{\mathrm{cls}}$ & $\operatorname{AAA}^{\mathrm{ncm}}$ &
% $\operatorname{AAA}^{\mathrm{cls}}$ & $\operatorname{AAA}^{\mathrm{ncm}}$ &
% $\operatorname{AAA}^{\mathrm{cls}}$ & $\operatorname{AAA}^{\mathrm{ncm}}$ \\
% \midrule
% \multirow{1}{*}{$0$} &Base &
% 69.07 & 76.62 & 71.29 & 78.29 & 71.41 & 78.47 \\
% \midrule
% % 组 1：\rho_1
% \multirow{2}{*}{$0.05$}
% &Full 
% &-3.5  & -1.54 & -2    & -0.38 & -2.62 & -0.8  \\ 
% & LoRA &0.27  & 0.02  & 0.79  & 0.51  & 0.12  & 0.16  \\ 
% \midrule
% % 组 2：\rho_2
% \multirow{2}{*}{$0.10$} & Full &-4.82 & -1.29 & -4.74 & -0.74 & -5.53 & -1.18 \\
% & LoRA &0.7   & 0.29  & 0.86  & 0.53  & 0.28  & 0.12  \\
% \midrule
% % 组 3：\rho_3
% \multirow{2}{*}{$0.15$} & Full &-6.72 & -1.69 & -6.84 & -0.98 & -6.49 & -1.29 \\
% & LoRA &0.41  & 0.08  & 0.77  & 0.53  & 0.27  & 0.11 \\
% \bottomrule
% \end{tabular}}
% \end{table}


{%
\setlength{\textfloatsep}{0pt}      % 页顶/页底 table 与正文的距离（最关键）
\setlength{\intextsep}{2pt}         % 若 table 被放在正文中时的上下距离（备用）
\setlength{\floatsep}{1pt}          % 两个浮动体之间的距离（备用）
\setlength{\abovecaptionskip}{2pt}  % caption 与表格之间
\setlength{\belowcaptionskip}{4pt}  % caption 下方到正文
% \captionsetup{skip=2pt}           % 若你用 caption 包，可用它更细调
\begin{table}[t]
\centering
\caption{$\operatorname{AAA}^{\mathrm{cls}}$ under Full vs LoRA-only perturbations.}
\label{tab:lora_or_full_cls—show}
\footnotesize
\setlength{\tabcolsep}{3pt}
\renewcommand{\arraystretch}{0.88}
\begin{tabular}{@{}c c c c c@{}}
\toprule
\textbf{Scope} & \textbf{Strength} & \textbf{SeqLoRA} & \textbf{IncLoRA} & \textbf{OLoRA} \\
\midrule
% Base & --   & 69.07 & 71.29 & 71.41 \\
% \midrule
\multirow{3}{*}{Full}
 & 0.05 & -3.50 & -2.00 & -2.62 \\
 & 0.10 & -4.82 & -4.74 & -5.53 \\
 & 0.15 & -6.72 & -6.84 & -6.49 \\
\midrule
\multirow{3}{*}{LoRA}
 & 0.05 & 0.27 & 0.79 & 0.12 \\
 & 0.10 & 0.70 & 0.86 & 0.28 \\
 & 0.15 & 0.41 & 0.77 & 0.27 \\
\bottomrule
\end{tabular}
\end{table}
}% end local group



\noindent\textbf{Scope of perturbations}\quad
To study perturbation scope in PECL, we compare $\mathrm{RS}^{(0)}_S(W,\Sigma)$ under two scopes: perturbations restricted to the task-admissible adapter coordinates versus perturbations defined over the full parameter space.
% with perturbations restricted to the task adapters versus perturbations in the full parameter space.
% We examine whether flatness in PECL should be promoted within the task-admissible adapter subspace or in the full parameter space. 
% Specifically, we compare $\mathrm{RS}^{(0)}_S(W,\Sigma)$ under two perturbation scopes: 
% (i) \textbf{adapter-scope}, where perturbations are restricted to the trainable LoRA parameters of the current task, and 
% (ii) \textbf{full-scope} (Flat-LoRA~\cite{li2025flatloralowrankadaptationflat}), where perturbations are defined over all model parameters rather than being restricted to adapters.
% The experiment is repeated for the three LoRA organizations and different perturbation strengths and results are shown in Table\ref{tab:lora——or_full}.
% As shown in Table~\ref{tab:lora——or_full}, the perturbation {scope} is the dominant factor in frozen-backbone PECL.
%supported on the task-admissible update subspace $\mathcal W_{\Delta,t}$ (i.e., 
As shown in Table~\ref{tab:lora_or_full_cls—show}, restricting perturbations to the trainable adapters consistently improves average accuracy over the SGD baseline, whereas applying perturbations that include the frozen backbone substantially degrades performance, with the strongest drop observed for SeqLoRA.
% when perturbations are restricted to the trainable adapters, we consistently obtain higher average accuracy than the SGD baseline. In contrast, applying perturbations including the frozen backbone parameters substantially degrades performance, with a particularly pronounced drop for SeqLoRA.

{%
% 1) table* 与上下正文的距离（页顶/页底浮动体）
\setlength{\textfloatsep}{2pt}      % 可在 4pt--10pt 之间调
\setlength{\floatsep}{2pt}          % 多个浮动体之间（备用）
\setlength{\dbltextfloatsep}{2pt}   % 某些模板对 table* 用这个（双栏更关键）

\begin{table*}[h]
\centering
\caption{
Class incremental performance and weight space flatness metrics on ImageNet-R ($T=20$, $R=16$). 
The table reports standard class incremental metrics  together with the sharpness and empirical Hessian and Fisher statistics restricted in the weight space.
}
\label{table:q1_ci_with_weight_flatness_reorg}
\setlength{\tabcolsep}{2.5pt}       % 列间距更小
\renewcommand{\arraystretch}{0.92}  % 行高更紧凑（可选）
\resizebox{\linewidth}{!}{
\scriptsize % 整体字号更小（可改为 \footnotesize 或 \tiny）
\begin{tabular}{@{}p{1.15cm} p{1.8cm} c c c c c c c c@{}}
\toprule
\multirow{2}{*}{\textbf{Method}} & \multirow{2}{*}{\textbf{Opt.}} &
\multicolumn{2}{c}{\textbf{CL Performance ($\uparrow$)}} &
\multicolumn{2}{c}{\textbf{Sharpness ($\downarrow$)}} &
\multicolumn{4}{c}{\textbf{Curvature ($\downarrow$)}} 
% \multicolumn{2}{c}{\textbf{Emp Fisher ($\downarrow$)}}
\\
\cmidrule(lr){3-4}\cmidrule(lr){5-6}\cmidrule(lr){7-10}
 & & $\operatorname{ACC}_{20}^{\mathrm{cls}}$  %&$\operatorname{ACC}^{\mathrm{cls}}$ 
&$\operatorname{ACC}_{20}^{\mathrm{nme}}$ 
 & AS$^{(0)}(\rho)$  & AS$^{(1)}(\rho)$
 & $\lambda_{\max}(H)$ & $\mathrm{tr}(H)$  
 & $\lambda_{\max}(F)$ & $\mathrm{tr}(F)$   \\
\midrule
% \multirow{1}{*}{\textbf{PerTaskLP}}
% & SGD   & $55.22^{\pm 0.31}$ & $59.36^{\pm 0.19}$ & $-8.35^{\pm 0.75}$ & 0.072   & 0.072 &     &   &    &   \\
% \midrule
% \multirow{2}{*}{\textbf{SeqFT}}
% & SGD   &$50.72^{\pm 2.20}$ & $59.89^{\pm 1.32}$ & $-38.94^{\pm 1.87}$ &   &  &   &  &   & \\
% & SAM   &$56.07^{\pm 2.60}$ & $66.47^{\pm 1.53}$ & $-33.04^{\pm 4.71}$&   &  &   &  &   &  \\
%  &GAM  &$59.78^{\pm 0.74}$ & $67.31^{\pm 1.18}$ & $-32.70^{\pm 0.26}$ &   &  &   &  &   &   \\
% \midrule
\multirow{4}{*}{{SeqLoRA}}
& SGD    &$63.14^{\pm 2.17}$ & $72.82^{\pm 1.27}$  & 0.142  & 0.114 & 47.37 & 1502.55 & 35.51 & 1401.15 \\\
& + $\mathrm{RS}^{0}_S(W,\rho)$  &$64.25^{\pm 0.87}$ & $73.01^{\pm 1.11}$ &0.121 &0.108 &47.29 &1444.96 &30.35 &1338.16\\
& + $\mathrm{AS}^{0}_S(W,\rho)$    &$67.04^{\pm 0.71}$ & $74.53^{\pm 0.60}$ & 0.120  & 0.093 & 24.41   & 996.59   & 21.39  & 821.45   \\
 & + $\mathrm{AS}^{(1)}_S(W,\rho)$   &$73.31^{\pm 0.52}$ & $75.74^{\pm 0.15}$  & 0.060  &0.058  & 4.81    & 339.49   & 3.72   & 240.91   \\
\midrule
\multirow{4}{*}{{IncLoRA}}
& SGD    &$65.77^{\pm 2.84}$  &$75.41^{\pm 0.26}$ & 0.076    & 0.062 & 15.46 & 969.34  & 8.89  & 836.91 \\
& + $\mathrm{RS}^{0}_S(W,\rho)$ & $67.42^{\pm 1.23}$ & $75.79^{\pm 0.42}$  &0.075 &0.073 &7.54 &917.35 &8.82 &774.56\\
& + $\mathrm{AS}^{(0)}_S(W,\rho)$    &$68.86^{\pm 0.92}$ & $75.95^{\pm 0.07}$  & 0.053   & 0.052 & 3.72    & 712.88   & 4.35   & 526.49   \\
 &+ $\mathrm{AS}^{(1)}_S(W,\rho)$  &$72.86^{\pm 0.23}$ & $76.72^{\pm 0.54}$  & 0.035   & 0.035 & 2.69    & 281.22   & 2.09   & 192.25   \\
\midrule
\multirow{4}{*}{{OLoRA}}
& SGD    &$66.64^{\pm 0.44}$ & $75.10^{\pm 0.41}$  & 0.078   & 0.058 & 6.94    & 938.93   & 9.00   & 831.23   \\
& + $\mathrm{RS}^{0}_S(W,\rho)$  & $67.53^{\pm 1.04}$ & $75.71^{\pm 0.27}$  &0.076 &0.073 &7.18 &918.33 &8.97 &813.62\\
& + $\mathrm{AS}^{(0)}_S(W,\rho)$     &$68.53^{\pm 0.78}$ & $75.73^{\pm 0.32}$  & 0.066   & 0.048 & 4.40    & 775.40   & 7.41   & 628.12   \\
 & + $\mathrm{AS}^{(1)}_S(W,\rho)$  &$72.44^{\pm 0.36}$ & $76.56^{\pm 0.53}$  & 0.036   & 0.036 & 2.22    & 261.47   & 2.08   & 198.30\\
% \midrule
% \multirow{1}{*}{\textbf{l2p}}
% & Adam   & $71.35^{\pm 0.05}$ & $74.84^{\pm 0.73}$ & $-6.34^{\pm 0.50}$  &-  &- &- &- &- &- \\
% \midrule
% \multirow{1}{*}{\textbf{dualprompt}}
% & Adam   & $71.59^{\pm 0.68}$ & $74.80^{\pm 1.18}$ & $-3.38^{\pm 0.45}$ &-  &- &- &- &- &- \\
\bottomrule
\end{tabular}
}
\end{table*}
}


% To further explore correlation between the learned adapter update $\Delta W$ and the frozen backbone weight $W$ in PECL, 
To understand why the perturbation scope is decisive, we further analyze how the learned adapter update $\Delta W_t$ interacts with the frozen backbone weight $W$ and with the task-conditioned curvature.
We adopt the feature amplification factor~\cite{hu2022lora}
$
A_t =
\frac{\left\lVert \Delta W_{t}\right\rVert_F}
{\left\lVert U_{\Delta,t} W  V_{\Delta,t}\right\rVert_F}
$
which measures how much of task-specific directions are amplified by $\Delta W$.
Here $(U_{\Delta,t},V_{\Delta,t})$ are the rank-$r$ left and right singular vector matrices of $\Delta W $
% left and right singular vectors of $\Delta W_t$ from its SVD~\cite{hu2022lora}.
and $\Delta U_t \Delta U_t^\top W \Delta V_t\Delta V_t^\top$ corresponds to the projection of $W$ onto the subspace spanned by $\Delta W_t$.
% Larger $A_t$ suggests that the task relies on feature directions that are important but not emphasized by $W$, and that are selectively boosted by $\Delta W_t$.
Similarly, to assess whether the dominant task-conditioned curvature
concentrates more on the adapter-induced subspace or the backbone reference subspace, we define the curvature projection ratio 
$
\text{C}_t=\frac{\sum_{i=1}^k \lambda_{t, i}\left\|U_{\Delta, t}^{\top} v_{t, i} V_{\Delta, t}\right\|_F^2}{\sum_{i=1}^k \lambda_{t, i}\left\|U_{W, t}^{\top} v_{t, i} V_{W, t}\right\|_F^2},
$
where ${(\lambda_{t,i},v_{t,i})}_{i=1}^k$ are the top-$k$ curvature eigenpairs of the full model at task $t$. The highe values of $C_t$ indicate that the dominant curvature modes align more with the same adapter-induced subspace spanned by  $(U_{\Delta,t},V_{\Delta,t})$ than with the backbone top $r$ singular vectors of W.


% Across the task sequence, $A_t$ shows a clear increasing trend with task index, indicating that the learned adapter updates $\Delta W_t$ progressively strengthen task-specific directions. In parallel, $C_t$ also increases over tasks and remains above the neutral reference level $C_t\approx 1$, suggesting that the dominant task-conditioned curvature becomes increasingly concentrated in the adapter-induced subspace rather than in the backbone reference subspace.

Fig.~\ref{fig:evolute} shows consistent trends across the task sequence: $A_t$ increases with the task index, indicating that $\Delta W_t$ progressively strengthens task-specific directions. $C_t$ also increases over tasks and remains above the neutral reference level $C_t\approx 1$, suggesting that the dominant task-conditioned curvature becomes increasingly concentrated in the adapter-induced subspace rather than in the backbone reference subspace.


Taken together, these diagnostics provide a mechanism-level explanation for Table~\ref{tab:lora_or_full_cls—show}: both adaptation and the curvature that governs sensitivity are increasingly localized within the LoRA-admissible coordinatesin CL, so restricting perturbations to the adapter scope matches the controllable geometry, whereas full-scope perturbations inject variations into directions that are outside the admissible update subspace and thus harm continual adaptation.

% Restricting perturbations to the task-admissible adapter subspace $\Delta W$ consistently improves average accuracy over the SGD baseline across both SeqLoRA and IncLoRA. In contrast, applying perturbations in the full parameter space $W+\Delta W$ is substantially worse, with a particularly severe degradation for SeqLoRA. 
%in the ambient parameter space $\mathbb R^d$ (

% Moreover, we sweep the perturbation strength $\rho$ and observe that the gain from increasing $\rho$ is not monotonic even under $\Delta W$ restricted perturbations: performance typically saturates or peaks at a moderate radius $\rho$, indicating an effective strength range. This behavior is consistent with prior  results that  report a method-dependent optimal range~\citep{foretSharpnessAwareMinimizationEfficiently2021,zhangGradientNormAware2023}, and our controlled $\rho$-sweep provides a direct verification of this trend in the PECL setting.

% Moreover, the benefit of increasing perturbation strength is not monotonic even under $\Delta W$-restricted perturbations; performance saturates or peaks at a moderate $\rho$, suggesting an effective strength range.

% The result in   Fig\ref{fig: pertuabtion scope} shows global perturbations are  harmful in the PECL regime. 
% Applying $\mathrm{RS}_S^{(0)}({W}, \Sigma)$ to the full model substantially degrades on perturbation ranges, often performing worse than the plain SGD baseline. 
% In contrast, restricting $\mathrm{RS}_S^{(0)}({W}, \Sigma)$ to the LoRA subspace yields small but consistent improvements over SGD.
% % These findings are aligned with the structural assumptions in our theory. 
% % In PECL, the posterior covariance is confined to the adapter subspace $W_\Delta$, while the backbone and past adapters are treated as deterministic and are not supposed to carry posterior uncertainty. 
% Global perturbations inject noise disrupt the stability of previously learned tasks. 
% By comparison, local perturbations confined to $W_\Delta$ respect the PECL structure: they reduce the subspace expected sharpness that appears in the PAC Bayesian bound without disturbing the frozen components. 
% This explains why Flat LoRA style full perturbations are not suitable in the PECL.



% in which $\Delta W_t$ acts, rather than with the backbone space.


% indicate that dominant curvature modes of the current task are more concentrated in the adapter-induced subspace than in the backbone subspace.


% modes are more aligned with the adapter subspace or with the backbone subspace, we define the curvature projection preference as
% $
% \text {CurvateProj}_t=\frac{\sum_{i=1}^k \lambda_{t, i}\left\|U_{\Delta, t}^{\top} V_{t, i} V_{\Delta, t}\right\|_F^2}{\sum_{i=1}^k \lambda_{t, i}\left\|U_{W, t}^{\top} V_{t, i} V_{W, t}\right\|_F^2+\varepsilon},
% $
% where larger values indicate that high-curvature directions of the current task concentrate more in the adapter-induced subspace than in the backbone subspace.


% Let $(U_{\Delta,t},V_{\Delta,t})$ denote the rank-$r$ adapter-induced subspace (from the SVD of $\Delta W_t$), and let $(U_{W,t},V_{W,t})$ denote the rank-$r$ backbone subspace (from the SVD of $W$).


% how strongly the task-induced directions encoded by $\Delta W_t$ are amplified relative to the backbone. Here $U_t$ and $V_t$ are the top-$r$ left and right singular vectors of $\Delta W_t$ from its SVD, and $U_tU_t^\top W V_tV_t^\top$ corresponds to the projection of $W$ onto the subspace spanned by $\Delta W_t$~\cite{hu2022lora}.





% to measures how much of task-specific directions are amplified by $\Delta W$, where $U$ and $V$  are the left- and right-singular matrices of the SVD decomposition and $UU^TWVV^TV$ gives the “projection” of $W$ onto the subspace spanned by $\Delta W$. \cite{hu2022lora}.  

% Similiar, to assess whether the dominant task-conditioned curvature are more aligned with the adapter subspace or with the backbone subspace, we compute  the top-$k$ curvature eigenpairs ${(\lambda_{t,i},v_{t,i})}{i=1}^k$ (with $\lambda{t,i}\ge 0$). the curvature projection ratio is
% $ r^\Delta_{i}
% =\frac{\sum_{i=1}^k \lambda_{i}\left\lVert U_{\Delta}^{\top} V_{i} V_{\Delta}\right\rVert_F^2} {\sum_{i=1}^k\lambda_{i}\left\lVert U_{W}^{\top} V_{i} V_{W}\right\rVert_F^2}$ over the top-$k$ curvature eigenpairs ${(\lambda_{i},v_{i})}_{i=1}^k$ (with $\lambda_{i}\ge 0$).




% Flat minima are often associated with improved generalization and robustness in standard supervised learning. 
% % From PAC-Bayes generalization analyses, SAM try to optimize Zeroth-order sharpness,which is the worst-case loss within a defined neighborhood, defined over training set $S$.
% % since they reduce the sensitivity of the empirical loss to small perturbations in parameter space. 
% % Sharpness-aware minimization was originally proposed to improve generalization by guiding optimization toward flat minima in standard supervised learning
% From PAC-Bayes generalization analyses, SAM try to reduce the empirical loss term by minimizes  adversarial zeroth order sharpness over trainging set $S$ within a radius $\rho$~\citep{AWP2020,foretSharpnessAwareMinimizationEfficiently2021}: 
% $$\mathrm{AS}^{0}_S(W,\rho) :=\max_{{\epsilon}\in B(W,\rho)} \Big[\ \hat{L}_{S}(W+{\epsilon})-\hat{L}_{S}(W)\ \Big].$$

% The other minimizes the expected Bayes objective with random  zeroth order sharpness~\citep{liRevisitingRandomWeight2024,jiang2019fantasticgeneralizationmeasures}:
% $$\mathrm{ES}^0_t(\hat{W}_t, \Sigma_t)
%  :=
%  \mathbb{E}_{\varepsilon \sim \mathcal{N}(0,\Sigma_t)}
%  \big[
%  \hat{L}_{S_t}(\hat{W}_t+\varepsilon) - \hat{L}_{S_t}(\hat{W}_t)
%  \big],$$
% Typical choices for $\mathcal{N}(0,\Sigma_t)$ include an isotropic Gaussian and a filter-wise Gaussian~\citep{pmlr-v151-bisla22a}. This definition is exactly the term we used before. Actually it is pointed out the SAM  as a relaxation of the Bayes objective where the expected negative-loss is replaced by the optimal convex lower bound~\cite{mollenhoffSAMOptimalRelaxation2022}.


 %distributional






% Flat minima are often associated with improved generalization and robustness, since they reduce the sensitivity of the empirical loss to small perturbations in parameter space. 
% We formalize flatness through sensitivity to parameter perturbations and introduce representative sharpness definitions widely used in the literature.






% The adversarial measures the maximal loss increase within a radius $\rho$~\citep{AWP2020,foretSharpnessAwareMinimizationEfficiently2021}: 
% $\mathrm{Sh}^{(0)}_S(W,\rho) :=\max_{{\epsilon}\in B(W,\rho)} \Big[\ \hat{L}_{S}(W+{\epsilon})-\hat{L}_{S}(W)\ \Big].$
% A %distributional 
% stochastic variant averages the loss increase under random perturbations~\citep{liRevisitingRandomWeight2024,jiang2019fantasticgeneralizationmeasures}:
% $\mathrm{E\text{-}Sh}_t(\hat{W}_t, \Sigma_t)
%  :=
%  \mathbb{E}_{\varepsilon \sim \mathcal{N}(0,\Sigma_t)}
%  \big[
%  \hat{L}_{S_t}(\hat{W}_t+\varepsilon) - \hat{L}_{S_t}(\hat{W}_t)
%  \big],$
% Typical choices for $\mathcal{N}(0,\Sigma_t)$ include an isotropic Gaussian and a filter-wise Gaussian~\citep{pmlr-v151-bisla22a}. 
% The first order sharpness controls the maximal gradient norm in a neighborhood~\citep{Zhang_GAM}, which is consistent with the zeroth order definition:
% $\mathrm{Sh}^{(1)}_S(W,\rho) := \rho\cdot \max_{W'\in B(W,\rho)} \big\|\nabla_{W} \hat{L}_S(W')\big\|_2,$

% Near a local minimum $\hat W$ with $\nabla \hat L_S(\hat W) \approx 0$, a second order expansion gives
% $
% \hat L_S(\hat W + \varepsilon) - \hat L_S(\hat W)
% \approx
% \tfrac12\,\varepsilon^\top H(\hat W)\,\varepsilon.
% $
% The sharpness introduced above reduce to different ways of aggregating the quadratic form $\epsilon^{\top}H(\hat W)\epsilon$: Sh$(0)$ approximates its maximal value over a small norm ball, $\mathrm{E\text{-}Sh}$ approximates its expectation under the perturbation distribution. These identifications justify the use of spectral proxies such as  $\mathrm{tr}\big(H(\hat{W})\big)$, as well as their Fisher surrogates $\mathrm{tr}\big(F(\hat{W})\big)$, where the Fisher coincides with the expected Hessian under standard regularity conditions~\citep{dinhSharpMinimaCan2017,caoTradeoffFlatnessOptimization2024,tsuzukuNormalizedFlatMinima2020,miMakeSharpnessAwareMinimization2022}. 


% If $\varepsilon$ is sampled from $\mathcal N(0,\Sigma)$ inside $\mathcal W_\Delta$, then $\mathrm{E\text{-}Sh}_\Delta(\hat W,\Sigma)$ can be approximated by a quadratic form involving the restriction of $H(\hat W)$ or its Fisher surrogate.

% In this sense, $\mathrm{E\text{-}Sh}_\Delta$ is a curvature dependent notion of flatness that only probes directions along the trainable adapter subspace, which is the relevant geometry in PECL.

% \subsection{A Minimal PAC--Bayes View of Sequential Learning}
% \label{sec:pac-minimal}

% PAC--Bayesian analysis places a prior $P$ and a posterior $Q$ on the hypothesis space and controls generalization through a combination of empirical risk and a KL divergence between $Q$ and $P$.
% Classical PAC-Bayesian results show that, with high probability over the training sample, the expected test risk of a stochastic predictor drawn from $Q$ is upper bounded by its empirical risk plus a complexity term that increases with $\mathrm{KL}(Q\Vert P)$ and decreases with the sample size. Intuitively, good generalization arises when the posterior both fits the data and remains close to the prior.





% In the single task case, for a bounded loss and any prior $P$ and posterior $Q$ with $Q\ll P$, a classical PAC--Bayesian inequality yields, with probability at least $1-\delta$ over the sample $S$,
% $$
% \mathbb E_{W\sim Q} L_{\mathcal D}(W)
% \le
% \mathbb E_{W\sim Q}\hat L_S(W)
% +
% \sqrt{
% \frac{\mathrm{KL}(Q\Vert P) + \log(1/\delta)}
% {2(m-1)}
% }.
% $$
% The right hand side consists of an empirical term and a complexity term that scales with the posterior divergence and the sample size.
% When $Q$ is a Gaussian distribution $Q = \mathcal N(\hat W,\Sigma)$, the empirical term $\mathbb E_{W\sim Q}\hat L_S(W)$ can be written as the loss at $\hat W$ plus an expected sharpness term.

% In Section~\ref{sec:pac-pecl} we build on a task level extension of this inequality with a time varying hyper posterior over task priors.
% This yields a continual learning bound whose complexity decomposes into task wise KL divergences and a hyper posterior drift across tasks.
% We then specialise this bound to the PECL parameterisation above and show that both the empirical term and the complexity term can be expressed entirely in the adapter subspace $\mathcal W_\Delta$.

% \vspace{-15pt}









% \subsection{Discussion}
% \label{sec:geometry-lora-pecl}


% To relate $\mathrm{E\text{-}Sh}_t^\Delta$ to measurable quantities, we decompose the model into a representation map and a linear classifier. Let $\phi(x;W)$ be the representation produced by the frozen backbone together with adapters, and let $W_{\mathrm{cls}}$ be the linear classifier. The predictive distribution is
% $
% p(y\mid x;W,W_{\mathrm{cls}})=
% \mathrm{softmax}\bigl(W_{\mathrm{cls}}^\top \phi(x;W)\bigr),
% $
% A small parameter perturbation induces the first-order feature shift
% $
% \Delta \phi(x)\approx J_\phi(x;W)\delta W$,
% $
% J_\phi(x;W)=\frac{\partial \phi(x;W)}{\partial W}.
% $
% Define the local feature-space Fisher
% $
% E_\phi(x)=
% \mathbb{E}_{y\sim p(\cdot\mid x;W,W_{\mathrm{cls}})}
% \Big[
% \nabla_\phi \log p(y\mid x;W,W_{\mathrm{cls}})
% \nabla_\phi \log p(y\mid x;W,W_{\mathrm{cls}})^\top
% \Big].
% $
% By the chain rule, the parameter-space Fisher satisfies
% $
% F_W(x)=J_\phi(x;W)^\top E_\phi(x),J_\phi(x;W),
% $
% and a standard local expansion yields, 



% Finally, 
% this perspective clarifies the role of flatness-aware optimizers in LoRA-based PECL.
% SAM reduces the worst-case loss increase within an $\ell_2$ neighborhood in $\mathcal{W}_\Delta$ and thus targets the sharpness in the subspace. RWP targets a stochastic relaxation aligned with $\mathrm{E\text{-}Sh}$ through the perturbation covariance $\Sigma_t$, while GAM controls a first-order neighborhood relaxation. 

% By reducing curvature within $\mathcal{W}_\Delta$, these methods decrease $\mathrm{E\text{-}Sh}_t^\Delta$ and make low empirical loss compatible with posterior choices that remain closer to their priors, improving the empirical–complexity tradeoff in Theorem~\ref{thm:pecl-bound}. Moreover, since $\Delta\phi(x)\approx J_\phi(x;W)\delta W$, smaller curvature along $\mathcal{W}_\Delta$ also implies reduced sensitivity of representations to adapter updates, providing a mechanistic link between subspace flatness and stability in PECL.


% \subsection{Comparing flatness optimizers}
% \label{sec:results-optimizers}
{%
\setlength{\intextsep}{4pt}
\setlength{\textfloatsep}{4pt}
\setlength{\abovecaptionskip}{6pt}
\setlength{\belowcaptionskip}{4pt}
\captionsetup{skip=10pt}
\begin{figure}[t]
\centering
\includegraphics[width=0.8\linewidth]{figures_TRML/weight_overlap_q_amp_full_proj3.pdf}
\caption{Evolution of the amplification factor $A_t$ and curvature projection ratio $C_t$ across tasks in SeqLoRA. The top panel shows $A_t$, measuring how strongly the adapter update $\Delta W_t$ amplifies directions relative to the backbone restricted to the same rank-$r$ subspace. The bottom panel shows $C_t$, comparing the concentration of dominant task-conditioned curvature modes in the adapter-induced subspace versus the backbone subspace. The weight matrices are taken from the  11th layer of ViT-B/16.
}
\label{fig:evolute}
\end{figure}
}


% \paragraph{Type of perturbations }
\noindent\textbf{Type of perturbations }\quad
% To isolate the effect of different flatness optimizers, we apply $\mathrm{AS}_S^{(0)}({W}, \Sigma)$,$\mathrm{RS}_S^{(0)}({W}, \Sigma)$ and $\mathrm{AS}^{(1)}_S(W,\rho)$ applied only to the adapter parameters, with SGD as a baseline. 
To isolate the effect of perturbation \emph{type} under PECL constraints, we apply three representative flatness optimizers to the {trainable LoRA parameters only}: 
$\mathrm{RS}^{(0)}_S(W,\Sigma)$, $\mathrm{AS}^{(0)}_S(W,\rho)$, and $\mathrm{AS}^{(1)}_S(W,\rho)$, 
% randomized zeroth-order perturbations $\mathrm{RS}^{(0)}_S(W,\Sigma)$, adversarial zeroth-order perturbations $\mathrm{AS}^{(0)}_S(W,\rho)$ , and adversarial first-order perturbations $\mathrm{AS}^{(1)}_S(W,\rho)$, 
with SGD as the baseline.
% We compare perturbation \emph{types} under a matched perturbation budget: adversarial methods share the same radius $\rho$, and randomized perturbations use a calibrated $\Sigma$ with comparable magnitude.
% With all other settings fixed, the observed differences reflect the method-specific perturbation geometry rather than strength selection.
% The comparison uses the same ImageNet-R PECL setting as in Section~\ref{sec:results-main-inr} and reports $\operatorname{AAA}^{\mathrm{cls}}$, $\operatorname{AAA}^{\mathrm{nme}}$, and a subset of sharpness metrics.
% The results in Table\ref{tab:aaa_nme_lora_org}  and Fig\ref{fig:robustness_curves_comparison}exhibit a clear ordering. This trend is consistent with the curves in in Fig\ref{fig:robustness_curves_comparison} and with the sharpness results in Table\ref{tab:table_all}. Besides, we also polt the lossland of these method in \ref{fig: loss landscape}.
Fig.~\ref{fig:robustness_imagenet_cp} shows a consistent ordering across tasks on ImageNet-C and ImageNet-P: both $\mathrm{AS}^{(0)}$ and $\mathrm{AS}^{(1)}$ dominate SGD throughout the sequence, and the gap becomes more visible in later tasks, indicating improved stability under continual adaptation.
% This trend holds across LoRA organizations, suggesting that the benefit is not an artifact of a particular adapter structure.
% This pattern indicates that promoting flatness in the LoRA subspace does not trade clean accuracy for robustness but instead enhances generalization in a way that transfers to distribution shift. 
% $\mathrm{RS}_S^{(0)}({W}, \Sigma)$ only slightly above  SGD. 
% $\mathrm{AS}^{(0)}_S(W,\rho)$ achieve intermediate gains. 
% $\mathrm{AS}^{(1)}_S(W,\rho)$ achieves the highest $\operatorname{AAA}^{\mathrm{cls}}$ and $\operatorname{AAA}^{\mathrm{ncm}}$. 
The same ordering is reflected by the final class-incremental performance and weight-space geometry in Table~\ref{table:q1_ci_with_weight_flatness_reorg}. Across LoRA configurations, moving from SGD to $\mathrm{RS}^{(0)}$, $\mathrm{AS}^{(0)}$, and then $\mathrm{AS}^{(1)}$ yields progressively higher $\operatorname{ACC}_{20}^{\mathrm{cls}}$ and $\operatorname{ACC}_{20}^{\mathrm{nme}}$, while simultaneously reducing sharpness diagnostics and projected curvature statistics, including $\lambda{\max}(H)$ and $\operatorname{tr}(H)$, and similarly for Fisher.
% reducing both sharpness indices and the projected curvature statistics (e.g., $\lambda_{\max}(H)$ and $\mathrm{tr}(H)$, and similarly for Fisher).
Overall, perturbation types that more strongly suppress adapter-subspace sharpness achieve better continual performance, consistent with the sharpness-dependent empirical term in Theorem~\ref{thm:pecl-bound}.
% In general, optimizers that suppress adapter-subspace sharpness more strongly also achieve better continual performance, which matches the dependence of the empirical term on task-admissible sharpness in Theorem~\ref{thm:pecl-bound}.



Fig.~\ref{fig:perturb_type_landscape} and Fig.~\ref{fig:evolute} further clarify \emph{how} different perturbation geometries translate into different full-model solutions even when updates are restricted to LoRA.
Within a LoRA-restricted slice (top of Fig.~\ref{fig:perturb_type_landscape}), the loss profiles can appear similarly flat across perturbation types. However, evaluating the same model under full-parameter perturbations (bottom of Fig.~\ref{fig:perturb_type_landscape}) separates the curves, revealing distinct sensitivity of the effective model $W_0+\Delta W$.
In particular, compared with SGD, $\mathrm{AS}^{(1)}$ yields a larger and steadily increasing curvature localization ratio $C_t$, together with a noticeably flatter loss profile under full-parameter evaluation. This indicates that $\mathrm{AS}^{(1)}$ not only concentrates dominant task-conditioned curvature more strongly into the adapter-induced subspace, but also drives optimization toward flatter task-conditioned solutions in the effective model $W_0+\Delta W$. Consequently, $\mathrm{AS}^{(1)}$ learns updates that are both more task-specific and better regularized in the full-model geometry, which is consistent with its consistently improved continual-learning performance.
% In particular, compared with SGD, $\mathrm{AS}^{(1)}$ yields a larger and increasing curvature localization ratio $C_t$ together with a flatter full-parameter loss profile, indicating that it both concentrates dominant task-conditioned curvature more strongly into the adapter-induced subspace and optimizes toward a flatter task-conditioned solution.
% Consequently,  $\mathrm{AS}^{(1)}$ learns updates that are simultaneously more task-specific and better regularized in the full-model geometry, which aligns with its consistently improved  performance.
These results show that, under a fixed PECL scope, the choice of perturbation \emph{type} controls the perturbation geometry within the LoRA coordinates, which in turn determines the resulting $\Delta W_t$ and reshapes the \emph{full-model} landscape of $W_0+\Delta W_t$.

{%
% {%
%
\setlength{\intextsep}{4pt}
\setlength{\textfloatsep}{4pt}
\setlength{\abovecaptionskip}{6pt}
\setlength{\belowcaptionskip}{2pt}
\captionsetup{skip=10pt}
\begin{figure}[h]
\centering
\includegraphics[width=0.8\linewidth]{figures_TRML/hebing_robutness2.pdf}
\vspace{-6pt}
\caption{\textbf{Robustness of perturbation types under CL} Classification average accuracy $\operatorname{ACC}_t^{\mathrm{cls}}$ over tasks evaluated under distribution shift on ImageNet-C  and ImageNet-P.
Across tasks and LoRA organizations, both $\mathrm{AS}^{(0)}$ and $\mathrm{AS}^{(1)}$ consistently outperform SGD, and the gap widens in later tasks.}
\label{fig:robustness_imagenet_cp}
\end{figure}
}%


% Figures~\ref{fig:perturb_type_landscape} and \ref{fig:evolute}  further clarifies the mechanism. Compared with SGD and  , GAM yields systematically larger $C_t$ and a visibly flatter loss profile under full-parameter evaluation, indicating that it not only localizes dominant curvature more strongly into the adapter-induced subspace but also optimizes toward a flatter task-conditioned solution. Consequently, GAM learns updates that are both more task-specific and better regularized in geometry, which aligns with its consistently improved continual-learning performance. This explains why methods can look comparably flat within the adapter coordinates yet yield markedly different accuracy: although perturbations and updates are confined to LoRA, they produce different $\Delta W$ solutions whose interaction with the frozen backbone changes the \emph{full-model} geometry for the current task.

% Consequently, adapter-only flatness control can still translate into systematic differences in the robustness and generalization of the overall network under distribution shift.



% When the loss is visualized along the dominant curvature direction while perturbations are restricted to LoRA coordinates (top panel), the landscapes appear similarly flat across sharpness methods, so a LoRA-plane view alone can understate differences between perturbation types. However, when the same solutions are evaluated under full-parameter perturbations (bottom panel), the curves separate more clearly, revealing distinct sensitivity profiles in the effective model $W_0+\Delta W$. This explains why methods can look comparably flat within the adapter coordinates yet yield markedly different accuracy: although perturbations and updates are confined to LoRA, they produce different $\Delta W$ solutions whose interaction with the frozen backbone changes the \emph{full-model} geometry for the current task. Consequently, adapter-only flatness control can still translate into systematic differences in the robustness and generalization of the overall network under distribution shift.






% This ordering reveals a parallel hierarchy on the geometric side: optimisers that more aggressively suppress subspace sharpness also achieve better continual learning performance. In other words, the trend appears simultaneously in the flatness diagnostics and in the CI metrics. This  relationship is precisely what Theorem~\ref{thm:pecl-bound} predicts: for a fixed LoRA geometry and training budget, methods that reduce the adapter-subspace sharpness tighten the empirical term of the PAC–Bayesian bound.

% \vspace{-5pt}
% \begin{table}[thpb]
% \centering
% \caption{The $\operatorname{AAA}^{\mathrm{cls}}$ and $\operatorname{AAA}^{\mathrm{ncm}}$ performance across LoRA organizations under different flatness optimizers. 
% % Rows below SGD report the change relative to the SGD baseline.
% }
% \label{tab:aaa_nme_lora_org}
% % \setlength{\tabcolsep}{5pt}
% % \resizebox{\linewidth}{!}
% \setlength{\tabcolsep}{3pt}
% \renewcommand{\arraystretch}{1} % 可选：稍微压缩行距
% \resizebox{1\linewidth}{!}{
% \begin{tabular}{l cc cc cc}
% \toprule
% \multirow{3}{*}{\textbf{Opt.}} &
% \multicolumn{2}{c}{\textbf{SeqLoRA}} &
% \multicolumn{2}{c}{\textbf{IncLoRA}} &
% \multicolumn{2}{c}{\textbf{OLoRA}} \\
% \cmidrule(lr){2-3} \cmidrule(lr){4-5} \cmidrule(lr){6-7}
% % & $\overline{\text{ACC}}_{20}$ (CL) & $\overline{\text{ACC}}_{20}$ (NME) & $\overline{\text{ACC}}_{20}$ (CL) & $\overline{\text{ACC}}_{20}$ (NME) & $\overline{\text{ACC}}_{20}$ (CL) & $\overline{\text{ACC}}_{20}$ (NME) \\
% % \cmidrule(lr)
% % &  CL &  NME & CL &  NME &  CL &  NME \\
% % & $\overline{\text{ACC}}_{20}$ & $\overline{\text{ACC}}_{20}$  & $\overline{\text{ACC}}_{20}$  & $\overline{\text{ACC}}_{20}$  & $\overline{\text{ACC}}_{20}$  & $\overline{\text{ACC}}_{20}$  \\
% & $\operatorname{AAA}^{\mathrm{cls}}$ &  $\operatorname{AAA}^{\mathrm{ncm}}$ & $\operatorname{AAA}^{\mathrm{cls}}$ & $\operatorname{AAA}^{\mathrm{ncm}}$ & $\operatorname{AAA}^{\mathrm{cls}}$ & $\operatorname{AAA}^{\mathrm{ncm}}$ \\
% \midrule
% % $\overline{\text{ACC}}_{20}$
% SGD  & 69.07 & 76.62 & 71.29 & 78.29 & 71.41 & 78.47 \\
% \midrule
% $\mathrm{RS}_S^0({W}, \Sigma_t)$   &0.70  & 0.29  & 0.86  & 0.53  & 0.28  & 0.12   \\
% $\mathrm{AS}^{0}_S(W,\rho)$   & 2.35  & 1.00  & 1.29  & 0.40  & 1.08  & 0.22  \\
% % $\Delta$ CFLAT & 2.35  & 1.04  & 2.67  & 1.26  & 1.07  & 0.12  \\
% $\mathrm{AS}^{(1)}_S(W,\rho)$  
% % & 7.73  & 2.24  & 5.23  & 1.13  & 4.64  & 0.68
% & \textbf{7.73}  & \textbf{2.24}  & \textbf{5.23}  & \textbf{1.13}  & \textbf{4.64}  & \textbf{0.68}


% \\
% \bottomrule
% \end{tabular}}
% \end{table}















% The induced feature shift admits the first-order approximation
% \[
% \Delta \phi(x)
% :=
% \phi(x;W+\delta W)-\phi(x;W)
% \approx
% J_W(x)\,\delta W,
% \qquad
% J_W(x):=\frac{\partial \phi(x;W)}{\partial W}.
% \]
% Keeping $W_{\mathrm{cls}}$ fixed, the resulting change of the predictive distribution can be quantified by a local second-order expansion of the KL divergence in feature space,
% \[
% \mathrm{KL}\!\left(
% p(\cdot \mid \phi(x)+\Delta \phi(x); W_{\mathrm{cls}})
% \;\middle\|\;
% p(\cdot \mid \phi(x); W_{\mathrm{cls}})
% \right)
% \approx
% \frac12\,\Delta \phi(x)^\top E_\phi(x)\,\Delta \phi(x),
% \]
% where the feature-space Fisher (local feature matrix) is

%     \resizebox{\linewidth}{!}{$
% E_\phi(x)
% :=
% \mathbb{E}_{\tilde y \sim p(\cdot \mid x;W,W_{\mathrm{cls}})}
% \Big[
% s_\phi(x,\tilde y)\,s_\phi(x,\tilde y)^\top
% \Big],
% \qquad
% s_\phi(x,y):=\nabla_\phi \log p(y \mid x;W,W_{\mathrm{cls}}).
%     $}

% \[

% \]





% To understand how flatness in $\mathcal{W}_\Delta$ influences model behaviour, we decompose the network into a shared
% representation mapping and a linear classifier.
% Let $\phi(x;W)$ be the representation produced by the frozen backbone together with adapters, and let
% $W_{\mathrm{cls}}$ be the linear classifier.
% The predictive distribution is
% $
% p(y\mid \phi;W_{\mathrm{cls}})
% =
% \mathrm{softmax}\!\bigl(W_{\mathrm{cls}}^\top \phi(x;W)\bigr)
% $ and the negative log-likelihood loss is
% $
% \ell(x,y;W,W_{\mathrm{cls}}):=-\log p(y\mid \phi;W_{\mathrm{cls}}).
% $


% Focus on the representation model and consider a small parameter perturbation $\delta W$. The induced feature shift admits the first-order approximation
% $
% \Delta \phi(x)
% \approx
% J_\phi(x;W)\delta W,
% \quad
% J_\phi(x;W):=\frac{\partial \phi(x;W)}{\partial W}.
% $
% The local Fisher (feature-space curvature) at input $x$ is
% $
% E_\phi(x)
% =
% \mathbb{E}_{y\sim p(\cdot\mid x;W,W_{\mathrm{cls}})}
% \Big[
% \nabla_\phi \log p\bigl(y\mid \phi;W_{\mathrm{cls}}\bigr)
% \nabla_\phi \log p\bigl(y\mid \phi;W_{\mathrm{cls}}\bigr)^\top
% \Big].
% $
% By the chain rule,
% % $
% % \nabla_W \log p(y\mid x;W,W_{\mathrm{cls}})=
% % J_\phi(x;W)^\top
% % \nabla_\phi \log p\bigl(y\mid \phi;W_{\mathrm{cls}}\bigr).
% % $
% % Consequently, 
% the Fisher information in parameter space at $x$ is
% $
% F_W(x)
% :=
% \mathbb{E}_{y\sim p(\cdot\mid x;W,W_{\mathrm{cls}})}
% \Big[
% \nabla_W \log p(y\mid x;W,W_{\mathrm{cls}})
% \nabla_W \log p(y\mid x;W,W_{\mathrm{cls}})^\top
% \Big]=
% J_\phi(x;W)^\top E_\phi(x)J_\phi(x;W).
% $
% Using the standard local KL expansion, for sufficiently small $\delta W$,
% $
% \mathrm{KL}\Big(
% p(\cdot\mid x;W+\delta W,W_{\mathrm{cls}})
% \|
% p(\cdot\mid x;W,W_{\mathrm{cls}})
% \Big)
% \approx
% \frac{1}{2}\delta W^\top F_W(x)\delta W.
% $
% In practice, we approximate the curvature by the empirical Fisher computed from per-sample gradients on a dataset $S=\{(x_i,y_i)\}$:
% $$
% F^{\mathrm{emp}}
% :=
% \frac{1}{|S|}\sum_{i=1}^{|S|} g_i g_i^\top,
% \quad
% g_i:=\nabla_W \ell(x_i,y_i;W,W_{\mathrm{cls}}).
% $$


% % The change of the predictive distribution in feature space can be quantified as
% % $
% % \Delta \phi(x)^\top E_\phi(x) \Delta \phi(x),
% % $
% % where {$
% % E_\phi(x)
% % =
% % \mathbb{E}_{y \sim p(\cdot \mid x;W,W_{\mathrm{cls}})}
% % \Big[
% % \nabla_\phi \log p(y \mid x)
% % \nabla_\phi \log p(y \mid x)^\top
% % \Big]
% % $} is the local feature matrix.
% % For a linear classifier followed by a softmax, there is an analytic formulation
% % $  E_\phi(x)
% % =
% % W_{\mathrm{cls}}
% %  \Big(
% %  \mathrm{Diag}(p) - p p^\top
% %  \Big)
% %  W_{\mathrm{cls}}^\top.
% %  $
% % By the chain rule, the corresponding parameter-space quadratic approximation is
% % $$
% % \mathrm{KL}\big(
% % p(\cdot \mid x; W + \delta W)
% % \big\|
% % p(\cdot \mid x; W)
% % \big)
% % \approx
% % \delta W^\top F_W(x) \delta W
% % $$
% % where
% % $
% % F_W(x)
% % = \mathbb{E}_{y \sim p(\cdot \mid x;W,W_{\mathrm{cls}})} 
% % \Big(
% % \frac{\partial \log p(y \mid x)}{\partial W}
% % \Big)^\top
% % \Big(\frac{\partial \log p(y \mid x)}{\partial W} \Big)
% % $
% % is the Fisher information matrix for $W$ at $x$. 
% % We approximate this curvature by the empirical Fisher constructed from per-sample gradients,
% % \[
% % F^{\mathrm{emp}}
% % :=
% % \frac{1}{|S|}\sum_{i=1}^{|S|} g_i g_i^\top,
% % \quad
% % g_i=\nabla_W \ell(x_i,y_i;W,W_{\mathrm{cls}}).
% % \]
% % This $F^{\mathrm{emp}}$ serves as a local curvature proxy along $\mathcal W_\Delta$ for connecting subspace perturbation sensitivity to the sharpness term in Theorem~\ref{thm:pecl-bound}.
% % Averaging over a dataset $S$ yields the empirical Fisher
% % matrix
% % $
% % F_\mathrm{emp}
% % =
% % \frac{1}{|S|}
% % \sum_{(x_i,y_i) \in S} \nabla_W\ell
% % \nabla_W\ell^T,
% % $
% % which serves as the curvature measure that controls the flatness term in the PAC-Bayesian bound.


% Theorem~\ref{thm:pecl-bound} is stated for an abstract trainable update subspace $\mathcal W_\Delta$.
% In PECL, this abstraction is instantiated by LoRA: although optimisation is performed in low-rank factor coordinates, the induced weight variation $\delta W$ lies in the admissible set $\mathcal W_\Delta$ and locally spans update directions in weight space.
% Therefore, both the sharpness term $\mathrm{E\text{-}Sh}_t^\Delta$ and the task-wise KL contribution are determined by the geometry restricted to $\mathcal W_\Delta$.
% Accordingly, the action of the empirical Fisher $F^{\mathrm{emp}}$ on $\delta W\in\mathcal W_\Delta$ serves as a concrete local curvature proxy for the perturbation sensitivity encoded by $\mathrm{E\text{-}Sh}_t^\Delta$.

% This viewpoint yields a direct interpretation of flatness-aware optimisation in LoRA-based PECL.
% Applying SAM  to LoRA factors penalizes high curvature directions within $\mathcal W_\Delta$, thereby reducing $\mathrm{E\text{-}Sh}_t^\Delta$.
% Under Gaussian posteriors with controlled covariance, improved subspace flatness also allows low empirical loss with posteriors closer to their priors, improving the empirical--complexity tradeoff in Theorem~\ref{thm:pecl-bound}.
% Finally, by the same chain-rule relation, smaller curvature along $\mathcal W_\Delta$ suggests local representation stability, as $\phi(x;W)$ becomes less sensitive to adapter updates.


% Theorem~\ref{thm:pecl-bound} is stated on an abstract trainable update subspace $\mathcal W_\Delta$.
% In PECL this abstraction is realised concretely by LoRA: although optimisation is performed in low-rank factor coordinates, the induced weight perturbations lie in the admissible update set, and under a local first-order approximation the resulting variations span the corresponding update directions in weight space.
% Consequently, both the sharpness term $\mathrm{E\text{-}Sh}_t^\Delta$ and the taskwise KL term are governed by geometry restricted to $\mathcal W_\Delta$, rather than by curvature in the full parameter space. When the perturbations are restricted to the admissible update subspace $\mathcal W_\Delta$, the action of $F^{\mathrm{emp}}$ on vectors $\delta W\in\mathcal W_\Delta$ provides a concrete local curvature proxy for linking subspace perturbation sensitivity to the sharpness term in Theorem~\ref{thm:pecl-bound}.


% For the negative log-likelihood loss $
% p(y\mid x;W,W_{\mathrm{cls}})$, a standard local KL surrogate yields a quadratic approximation in terms of Fisher curvature,
% $$
% \mathrm{KL}\bigl(P_{x,y}(\theta+d) \| P_{x,y}(\theta)\bigr)
% =
% \frac12\,d^\top F(x) d+O(d^3),
% $$
% where $F=\mathbb E_{y\sim p(\cdot\mid x)}\left[\nabla \log p(y \mid x, \theta) \nabla \log p(y \mid x, \theta)^{\top}\right]$.
% We only care the representation paramter and ignore the linear classifier. Let $\phi(x;W)$ be the representation produced by the frozen backbone together with adapters, and let $W_{\mathrm{cls}}$ be the linear classifier defining $p(y\mid x)$.  By the chain rule, the gradient at $W_{cls}$ is $\nabla_\phi \log p$ and the corresponding gradient at $W$ is $\frac{\partial \phi(x;W)}{\partial W}
% \Bigr)$.
% For a small perturbation $\Delta W$, the change in predictive distribution satisfies the quadratic approximation
% $$
% \mathrm{KL}\bigl(p(\cdot\mid x;W+\Delta W)\,\Vert\,p(\cdot\mid x;W)\bigr)
% \approx
% \frac12\Delta W^\top F_W(x)\Delta W.
% $$ where a Fisher
% $
% F_W(x)
% =
% \Bigl(\frac{\partial \phi(x;W)}{\partial W}\Bigr)^\top
% E_\phi(x)
% \frac{\partial \phi(x;W)}{\partial W}.
% $Under PECL, admissible perturbations satisfy $\Delta W\in\mathcal W_\Delta$; we therefore interpret the relevant curvature through the projected Fisher restricted to $\mathcal W_\Delta$. Specificly, the empirical Fisher is
% $
% F^{\mathrm{emp}}_{\Delta W}:=\frac{1}{|S|} \sum_{(x, y) \in S} \nabla \log p(y \mid x, \theta) \nabla \log p(y \mid x, \theta)^{\top}.
% $






% Let $\phi(x;W)$ be the representation produced by the frozen backbone together with adapters, and let $W_{\mathrm{cls}}$ be the linear classifier defining $p(y\mid x)$.
% For the negative log-likelihood loss $
% p(y\mid x;W,W_{\mathrm{cls}})$, a standard local KL surrogate yields a quadratic approximation in terms of Fisher curvature,
% $$
% \mathrm{KL}\bigl(p(\cdot\mid x;W+\Delta W)\,\Vert\,p(\cdot\mid x;W)\bigr)
% =
% \frac12\,\Delta W^\top F_W(x)\Delta W+O(\Delta W^3).
% $$
% where $F=\mathrm{E}\left[\nabla \log p(y \mid x, \theta) \nabla \log p(y \mid x, \theta)^{\top}\right]$
% The local feature matrix
% $
% E_\phi(x)
% =
% \mathbb E_{y\sim p(\cdot\mid x)}
% \bigl[
% \nabla_\phi \log p\,
% \nabla_\phi \log p^\top
% \bigr]
% $
% induces a Fisher
% $
% F_W(x)
% =
% \Bigl(\frac{\partial \phi(x;W)}{\partial W}\Bigr)^\top
% E_\phi(x)
% \frac{\partial \phi(x;W)}{\partial W}.
% $
% Under PECL, admissible perturbations satisfy $\Delta W\in\mathcal W_\Delta$; we therefore interpret the relevant curvature through the projected Fisher restricted to $\mathcal W_\Delta$. Specificly, the empirical Fisher is
% $
% F^{\mathrm{emp}}_{\Delta W}:=\frac{1}{|S|} \sum_{(x, y) \in S} \nabla \log p(y \mid x, \theta) \nabla \log p(y \mid x, \theta)^{\top}.
% $


% This yields an interpretation of flatness-aware optimisation in LoRA-based PECL.
% Methods such as SAM applied to LoRA factors discourage solutions with large curvature in $\mathcal W_\Delta$, which is consistent with reducing perturbation sensitivity measured by $\mathrm{E\text{-}Sh}_t^\Delta$.
% Moreover, for common Gaussian constructions with controlled posterior covariance, improved subspace flatness can make low empirical loss compatible with posteriors that remain closer to their priors, improving the empirical--complexity tradeoff in Theorem~\ref{thm:pecl-bound}.
% Through the same chain-rule decomposition, reduced curvature along $\mathcal W_\Delta$ also suggests a local notion of representation stability, in the sense that $\phi(x;W)$ becomes less sensitive to adapter updates.


% Theorem~\ref{thm:pecl-bound} is stated on an abstract trainable update subspace $\mathcal W_\Delta$. 
% In PECL this abstraction is realized concretely by LoRA: although optimization is performed in low-rank factor coordinates, the induced weight perturbations span exactly the admissible update directions, so any posterior over LoRA factors induces a posterior over $\Delta W\in\mathcal W_\Delta$. 
% Consequently, both the sharpness term $\mathrm{E\text{-}Sh}_t^\Delta$ and the taskwise KL term are governed by geometry restricted to $\mathcal W_\Delta$, rather than by curvature in the full parameter space.

% % To connect $\mathrm{E\text{-}Sh}_t^\Delta$ to measurable curvature, we use the Fisher decomposition. 
% Let $\phi(x;W)$ be the representation produced by the frozen backbone together with adapters, and let $W_{\mathrm{cls}}$ be the linear classifier defining $p(y\mid x)$. The predictive distribution is
% $
% p(y\mid x;W,W_{\mathrm{cls}})
% =
% \mathrm{softmax}\bigl(W_{\mathrm{cls}}^\top \phi(x;W)\bigr).
% $
% The local feature matrix
% $
% E_\phi(x)
% =
% \mathbb E_{y\sim p(\cdot\mid x)}
% \bigl[
% \nabla_\phi \log p\,
% \nabla_\phi \log p^\top
% \bigr]
% $
% induces a Fisher
% $
% F_W(x)
% =
% \Bigl(\frac{\partial \phi(x;W)}{\partial W}\Bigr)^\top
% E_\phi(x)
% \frac{\partial \phi(x;W)}{\partial W}
% $ for parameters $W$ at $x$,
% so that small weight perturbations satisfy the local quadratic approximation
% $
% \mathrm{KL}\bigl(p(\cdot\mid x;W+\Delta W)\,\Vert\,p(\cdot\mid x;W)\bigr)
% \approx
% \frac12\,\Delta W^\top F_W(x)\Delta W.
% $
% Under PECL, admissible perturbations are restricted to $\Delta W\in\mathcal W_\Delta$, hence the curvature that controls $\mathrm{E\text{-}Sh}_t^\Delta$ is precisely the action of $F_W(x)$ along $\mathcal W_\Delta$.

% This yields a direct interpretation of flatness-aware optimization in LoRA-based PECL. 
% Methods such as SAM or GAM applied to LoRA factors primarily suppress large curvature along $\mathcal W_\Delta$, which reduces the perturbation sensitivity measured by $\mathrm{E\text{-}Sh}_t^\Delta$. 
% At the same time, flatter geometry in $\mathcal W_\Delta$ makes low empirical loss compatible with posteriors that remain closer to their priors, improving the empirical–complexity tradeoff in Theorem~\ref{thm:pecl-bound}. 
% Through the same chain-rule decomposition, reduced curvature along $\mathcal W_\Delta$ also implies that the learned representation $\phi(x;W)$ is less sensitive to adapter updates, providing a feature-space view of why subspace flatness correlates with improved stability in continual learning.


% \subsection{Discussion of the bound in LoRA-based PECL}
% \label{sec:bound-discussion}

% Theorem~\ref{thm:pecl-bound} shows that, under PECL, all task-dependent contributions in the PAC-Bayes guarantee are confined to the trainable update subspace $\mathcal W_\Delta$. Concretely, the empirical part depends on the subspace expected sharpness $\mathrm{E\text{-}Sh}_t^\Delta$, and the complexity depends on the adapter-level KL term $\mathrm{KL}(Q_t^\Delta\Vert P_t^\Delta)$ together with the hyperposterior drift $\mathrm{KL}(\mathcal P_t\Vert \mathcal P_{t-1})$.

% % \paragraph{Notation used in this section.}
% For each task $t$, the frozen-backbone PECL solution is evaluated at
% $
% \hat W_t := W_{\mathrm{frozen}} + \hat{\Delta W}_t,
% \quad
% \hat{\Delta W}_t \in \mathcal W_\Delta,
% $
% where $\mathcal W_\Delta\subseteq\mathbb R^d$ is the linear subspace of admissible LoRA updates in weight space.
% Let $\Pi_\Delta:\mathbb R^d\to\mathcal W_\Delta$ denote the orthogonal projection onto $\mathcal W_\Delta$ (equivalently, $\Pi_\Delta = U_\Delta U_\Delta^\top$ for any orthonormal basis $U_\Delta$ of $\mathcal W_\Delta$).
% The Gaussian perturbation defining subspace sharpness is
% $
% \varepsilon \sim \mathcal N(0,\Sigma_t),
% \ \ \mathrm{supp}(\varepsilon)\subseteq \mathcal W_\Delta,
% $
% where $\Sigma_t$ is a covariance operator acting on $\mathcal W_\Delta$.

% We write $\lambda_{\max}(A)$ for the largest eigenvalue of a symmetric matrix $A$, and $\mathrm{tr}(A)$ for its trace.
% Finally, $\theta$ denotes a vector of trainable adapter coordinates (either a coordinate system for $\mathcal W_\Delta$ or the stacked LoRA factor parameters in a local chart).

% % \paragraph{LoRA induces the update subspace.}
% Although optimization is carried out in LoRA coordinates, admissible weight-space perturbations are restricted to $\mathcal W_\Delta$ because the backbone is frozen.
% Therefore, the perturbations that enter $\mathrm{E\text{-}Sh}_t^\Delta$ can be specified directly as $\varepsilon\in\mathcal W_\Delta$ in weight space, or equivalently via a covariance in adapter coordinates $\theta$ that induces the same distribution over $\mathcal W_\Delta$.
% This is precisely the sense in which Theorem~\ref{thm:pecl-bound} makes LoRA geometry explicit: it charges only uncertainty and curvature along admissible update directions.

% % \paragraph{Projected Fisher as a curvature proxy for subspace sharpness.}
% To relate $\mathrm{E\text{-}Sh}_t^\Delta$ to measurable curvature, we use a standard local KL quadratic approximation.
% Let $p(\cdot\mid x;W)$ be the predictive distribution at parameters $W$, and define the (per-sample) Fisher information matrix at $W$ by
% \[
% F_W(x)
% :=
% \mathbb E_{y\sim p(\cdot\mid x;W)}
% \bigl[
% \nabla_W \log p(y\mid x;W)\,
% \nabla_W \log p(y\mid x;W)^\top
% \bigr]
% \in\mathbb R^{d\times d}.
% \]
% We evaluate curvature at the PECL solution $\hat W_t$ and restrict it to $\mathcal W_\Delta$ through projection. Let
% \[
% F_{\Delta,t}(x)
% :=
% \Pi_\Delta^\top F_{\hat W_t}(x)\,\Pi_\Delta,
% \qquad
% F_{\Delta,t}
% :=
% \frac{1}{m_t}\sum_{x\in S_t} F_{\Delta,t}(x).
% \]
% The following proposition summarizes the proxy used in our experiments.

% \begin{proposition}[Projected Fisher proxy for subspace expected sharpness]
% \label{prop:fisher-proxy}
% Assume the local approximation
% $
% \mathrm{KL}\bigl(p(\cdot\mid x;\hat W_t+\delta)\,\Vert\,p(\cdot\mid x;\hat W_t)\bigr)
% \approx
% \frac12\,\delta^\top F_{\hat W_t}(x)\,\delta
% $
% holds for small $\delta\in\mathbb R^d$.
% If $\varepsilon\sim\mathcal N(0,\Sigma_t)$ with $\mathrm{supp}(\varepsilon)\subseteq\mathcal W_\Delta$, then
% \[
% \mathrm{E\text{-}Sh}_t^\Delta(\hat{\Delta W}_t,\Sigma_t)
% \approx
% \frac12\,\mathrm{tr}\!\bigl(F_{\Delta,t}\Sigma_t\bigr)
% \le
% \frac12\,\lambda_{\max}(F_{\Delta,t})\,\mathrm{tr}(\Sigma_t).
% \]
% \end{proposition}

% Proposition~\ref{prop:fisher-proxy} motivates reporting $\lambda_{\max}(F_{\Delta,t})$ and $\mathrm{tr}(F_{\Delta,t})$ as curvature surrogates in the LoRA subspace, since these quantities control the magnitude of the sharpness term for a given posterior covariance $\Sigma_t$.

% \paragraph{Implications for SAM and a KL-based interpretation of EWC.}
% Theorem~\ref{thm:pecl-bound} separates an empirical contribution (which includes $\mathrm{E\text{-}Sh}_t^\Delta$) from a complexity contribution (which includes $\mathrm{KL}(Q_t^\Delta\Vert P_t^\Delta)$ and $\mathrm{KL}(\mathcal P_t\Vert\mathcal P_{t-1})$).
% Sharpness-aware optimizers such as SAM act directly on the empirical side by discouraging solutions with large subspace curvature, thereby reducing $\mathrm{E\text{-}Sh}_t^\Delta$ as reflected by smaller projected Fisher spectra in Proposition~\ref{prop:fisher-proxy}.
% While the KL terms depend on the chosen prior and posterior, improved subspace flatness can make low empirical loss compatible with posterior choices that remain closer to their priors, improving the empirical–complexity tradeoff for common Gaussian constructions.

% Classical continual learning regularizers can be expressed in the same KL language by specifying $P_t^\Delta$ appropriately.
% For EWC, consider trainable adapter coordinates $\theta$ and define an \emph{adapter-space} Gaussian prior centered at the previous solution,
% \[
% \tilde P_t(\theta)
% =
% \mathcal N\!\bigl(\theta_{t-1},(\gamma F_{\Delta,t-1})^{-1}\bigr),
% \]
% where $\gamma>0$ is a regularization strength and $F_{\Delta,t-1}$ is an estimated projected Fisher on the trainable coordinates from the previous task.
% Approximating the within-task posterior by $\tilde Q_t(\theta)=\mathcal N(\theta_t,\Sigma_Q)$, the dominant $\theta_t$-dependent component of $\mathrm{KL}(\tilde Q_t\Vert \tilde P_t)$ satisfies
% \[
% \mathrm{KL}(\tilde Q_t\Vert \tilde P_t)
% =
% \frac{\gamma}{2}\,(\theta_t-\theta_{t-1})^\top F_{\Delta,t-1}(\theta_t-\theta_{t-1})
% \;+\; \text{const.},
% \]
% recovering the EWC quadratic penalty up to additive constants that do not depend on $\theta_t$.
% Under this view, EWC primarily constrains the taskwise complexity term, whereas sharpness-aware methods primarily reduce the sharpness-dependent empirical contribution.
% Both effects are governed by the same LoRA-induced geometry, namely the projected curvature and posterior uncertainty inside $\mathcal W_\Delta$.


% % \subsection{Discussion of the bound}
% % Theorem~4.4 implies that under PECL all task-dependent terms are confined to the trainable LoRA update subspace $\mathcal W_\Delta$ (at task $t$ evaluated at $\hat W_t:=W_{\mathrm{frozen}}+\hat{\Delta W}_t$). 
% % As a result, the bound charges only uncertainty and curvature along admissible LoRA update directions.

% % To relate $\mathrm{E\text{-}Sh}_t^\Delta$ to measurable curvature, consider the local approximation around $\hat W_t$:
% % \[
% % \mathrm{KL}\bigl(p(\cdot\mid x;\hat W_t+\delta)\,\Vert\,p(\cdot\mid x;\hat W_t)\bigr)\approx \tfrac12\,\delta^\top F_{\hat W_t}(x)\,\delta.
% % \]
% % Restricting $\delta\in\mathcal W_\Delta$ corresponds to applying the orthogonal projection onto $\mathcal W_\Delta$, yielding an empirical projected Fisher $F_{\Delta,t}$ on task $t$.
% % For $\varepsilon\sim\mathcal N(0,\Sigma_t)$ supported on $\mathcal W_\Delta$,
% % \[
% % \mathrm{E\text{-}Sh}_t^\Delta \approx \tfrac12\mathrm{tr}(F_{\Delta,t}\Sigma_t)\le \tfrac12\lambda_{\max}(F_{\Delta,t})\,\mathrm{tr}(\Sigma_t).
% % \]

% % Within the same decomposition, sharpness-aware methods (e.g., SAM) primarily reduce the empirical sharpness term, while EWC can be viewed as choosing a Fisher-shaped \emph{adapter-space} prior $\tilde P_t=\mathcal N(\theta_{t-1},(\lambda F_{\Delta,t-1})^{-1})$ for trainable coordinates $\theta$, so that the dominant part of $\mathrm{KL}(\tilde Q_t\Vert \tilde P_t)$ recovers the usual quadratic penalty.


% % % \subsection{Discussion of the bound}
% % % \label{sec:geometry-lora-pecl}

% % % Theorem~\ref{thm:pecl-bound} shows that, under PECL, all task-dependent terms are confined to the LoRA update subspace $\mathcal W_\Delta$.
% % % % Both the sharpness-dependent empirical contribution $\mathrm{E\text{-}Sh}_t^\Delta$ and the within-task complexity $\mathrm{KL}(Q_t^\Delta\Vert P_t^\Delta)$ are defined on $\mathcal W_\Delta$, rather than on the full parameter space.
% % % As a result, the bound charges only curvature and uncertainty along admissible LoRA update directions, which clarifies why subspace-restricted flatness is the relevant notion of flatness in PECL.

% % % In practice, $\mathcal W_\Delta$ corresponds to the weight-space directions induced by perturbing LoRA factors.
% % % To connect $\mathrm{E\text{-}Sh}_t^\Delta$ to measurable quantities, a local quadratic approximation yields
% % % $$
% % % \mathrm{KL}\bigl(p(\cdot\mid x;W+\Delta W)\,\Vert\,p(\cdot\mid x;W)\bigr)
% % % \approx
% % % \frac12\,\Delta W^\top F_W(x)\Delta W,
% % % $$
% % % so restricting $\Delta W\in\mathcal W_\Delta$ amounts to projecting Fisher curvature onto $\mathcal W_\Delta$.
% % % Let $F_{\Delta,t}$ denote the empirical Fisher restricted to $\mathcal W_\Delta$ on task $t$.
% % % If $\varepsilon\sim\mathcal N(0,\Sigma_t)$ with $\varepsilon\in\mathcal W_\Delta$, then
% % % $$
% % % \mathrm{E\text{-}Sh}_t^\Delta
% % % \approx
% % % \frac12\,\mathrm{tr}\!\bigl(F_{\Delta,t}\Sigma_t\bigr)
% % % \le
% % % \frac12\,\lambda_{\max}(F_{\Delta,t})\,\mathrm{tr}(\Sigma_t),
% % % $$
% % % which motivates $\lambda_{\max}(F_{\Delta,t})$ and $\mathrm{tr}(F_{\Delta,t})$ as curvature proxies in the LoRA subspace.

% % % This perspective also clarifies how different continual learning strategies interact with the bound.
% % % Sharpness-aware updates such as SAM act most directly on the empirical contribution by reducing sensitivity to subspace perturbations, thereby lowering $\mathrm{E\text{-}Sh}_t^\Delta$, and can indirectly ease taskwise KL control by making low loss compatible with more prior-aligned posteriors.
% % % In contrast, EWC admits a KL interpretation that primarily targets the complexity term by constructing a Fisher-shaped Gaussian prior,
% % % $$
% % % P_t(\theta)=\mathcal N\!\bigl(\theta_{t-1},(\lambda F_{t-1})^{-1}\bigr),
% % % \qquad
% % % Q_t(\theta)=\mathcal N(\theta_t,\Sigma_Q),
% % % $$
% % % for which the dominant $\theta_t$-dependent part satisfies
% % % $
% % % \mathrm{KL}(Q_t\Vert P_t)
% % % \approx
% % % \frac{\lambda}{2}\,(\theta_t-\theta_{t-1})^\top F_{t-1}(\theta_t-\theta_{t-1})
% % % \quad+\ \text{const.}
% % % $
% % % Thus, within the same PAC--Bayes decomposition, SAM primarily tightens the sharpness-related empirical term, whereas EWC primarily constrains the KL complexity, and under PECL both effects are mediated by the LoRA-induced geometry $\mathcal W_\Delta$.



% % % % \subsection{From Subspace Sharpness to Fisher Curvature}
% % % % \label{sec:geometry-lora-pecl}

% % % % % The bound in Theorem~\ref{thm:pecl-bound} is formulated on an abstract update subspace $\mathcal W_\Delta$. In practice, this subspace is realised by LoRA adapters, and the sharpness term can be expressed in terms of feature-space matrices.

% % % % Theorem~\ref{thm:pecl-bound} is stated for an abstract LoRA update subspace $\mathcal W_\Delta$.
% % % % To connect the bound to measurable quantities, we relate the subspace expected sharpness to local curvature surrogates derived from the Fisher information.

% % % % Let $\phi(x;W)\in\mathbb R^D$ be the representation produced by the backbone and adapters, and let $W_{\mathrm{cls}}\in\mathbb R^{D\times K}$ be the linear classifier. The predictive distribution is
% % % % $
% % % % p(y\mid x;W,W_{\mathrm{cls}})
% % % % =
% % % % \mathrm{softmax}\bigl(W_{\mathrm{cls}}^\top \phi(x;W)\bigr).
% % % % $
% % % % The Fisher information matrix for parameters $W$ at $x$ is
% % % % $
% % % % F_W(x)=\Bigl(
% % % % \frac{\partial \phi(x;W)}{\partial W}
% % % % \Bigr)^\top
% % % % E_\phi(x)
% % % % \frac{\partial \phi(x;W)}{\partial W},
% % % % $ where
% % % % $
% % % % E_\phi(x)
% % % % :=
% % % % \mathbb E_{y\sim p(\cdot\mid x;W,W_{\mathrm{cls}})}
% % % % \bigl[
% % % % \nabla_\phi \log p
% % % % \nabla_\phi \log p^\top
% % % % \bigr]
% % % % $ is the local feature matrix at $x$ \cite{EFC2024}.
% % % % For a small perturbation $\Delta W$, the change in predictive distribution satisfies the quadratic approximation
% % % % $
% % % % \mathrm{KL}\bigl(p(\cdot\mid x;W+\Delta W)\,\Vert\,p(\cdot\mid x;W)\bigr)
% % % % \approx
% % % % \frac12\Delta W^\top F_W(x)\Delta W.
% % % % $


% % % % Under PECL, admissible perturbations are restricted to $\Delta W\in\mathcal W_\Delta$.
% % % % Let $U_\Delta$ be an orthonormal basis of $\mathcal W_\Delta$ and define the projected Fisher
% % % % $
% % % % F_\Delta(x) := U_\Delta^\top F_W(x)\,U_\Delta.
% % % % $
% % % % Equivalently, when parameterizing $\mathcal W_\Delta$ via LoRA factors, the same curvature is obtained by pulling back $F_W(x)$ to factor coordinates, which formalizes the coordinate equivalence between factor space and the update subspace.
% % % % This projection links the subspace expected sharpness to Fisher curvature restricted to $\mathcal W_\Delta$.
% % % % Let $\varepsilon\sim\mathcal N(0,\Sigma_t)$ be supported on $\mathcal W_\Delta$, and let $F_{\Delta,t}$ denote the empirical average of $F_\Delta(x)$ over task-$t$ samples.
% % % % Under the local approximation,
% % % % $
% % % % \mathrm{E\text{-}Sh}_t^\Delta(\hat{\Delta W}_t,\Sigma_t)
% % % % \approx
% % % % \frac{1}{2}\,\mathrm{tr}\bigl(F_{\Delta,t}\Sigma_t\bigr)
% % % % \le
% % % % \frac{1}{2}\,\lambda_{\max}(F_{\Delta,t})\,\mathrm{tr}(\Sigma_t),
% % % % $.
% % % % This motivates reporting $\lambda_{\max}(F_{\Delta,t})$ and $\mathrm{tr}(F_{\Delta,t})$ as curvature surrogates in the LoRA subspace.
% % % % In this view, sharpness-aware optimization influences the empirical contribution in Theorem~\ref{thm:pecl-bound} primarily through $\mathrm{E\text{-}Sh}_t^\Delta$, whereas the PECL constraint limits the complexity term by confining the KL divergences to $\mathcal W_\Delta$.

% Classical continual learning regularizers can be interpreted in the same KL-based language.
% Consider EWC, which penalizes parameter drift via a Fisher-weighted quadratic form.
% Model the prior as a Gaussian centered at the previous solution with precision proportional to an estimated Fisher matrix,
% $
% P_t(\theta)=\mathcal N\!\bigl(\theta_{t-1},(\lambda F_{t-1})^{-1}\bigr),
% $
% and approximate the within-task posterior by a Gaussian $Q_t(\theta)=\mathcal N(\theta_t,\Sigma_Q)$.
% Then the taskwise complexity term admits the closed-form decomposition
% \[
% \mathrm{KL}(Q_t \,\Vert\, P_t)
% =
% \frac{\lambda}{2}\,(\theta_t-\theta_{t-1})^\top F_{t-1}(\theta_t-\theta_{t-1})
% +
% \frac12\Bigl(
% \lambda\,\mathrm{tr}(F_{t-1}\Sigma_Q)
% -
% \log\det(\lambda F_{t-1}\Sigma_Q)
% -
% d
% \Bigr),
% \]
% where $d$ is the dimension of $\theta$.
% When $\Sigma_Q$ is treated as fixed or when one focuses on the dominant term that depends on $\theta_t$, this recovers the EWC quadratic penalty up to constants.
% Under this view, EWC indirectly constrains the complexity term, whereas sharpness-aware methods act more directly on the empirical part via $\mathrm{E\text{-}Sh}_t^\Delta$.
% This provides a unified interpretation of different continual learning strategies within the PAC--Bayes decomposition into empirical and complexity contributions.




% % % % % LoRA realises $\mathcal W_\Delta$ by low rank factors. For a layer with weight matrix $W_\ell \in \mathbb R^{d_o^{(\ell)}\times d_i^{(\ell)}}$, a LoRA module introduces an update
% % % % % $
% % % % % \Delta W_\ell = B_\ell A_\ell.
% % % % % $
% % % % % % Let $(\widehat A_\ell,\widehat B_\ell)$ be the current factors and $(\delta A_\ell,\delta B_\ell)$ small perturbations. 
% % % % % Collecting all factor perturbations into a vector $\theta$ and  $\theta$ span exactly the update subspace. 
% % % % % A Gaussian posterior on the factor coordinates induces a Gaussian posterior on $\Delta W$ supported on $\mathcal W_\Delta$.
% % % % % Restricting to admissible updates $\Delta W\in\mathcal W_\Delta$ yields the curvature that governs the subspace expected sharpness $\mathrm{E\text{-}Sh}_t^\Delta$.
% % % % % linearising gives
% % % % % $
% % % % % \mathrm{vec}(\Delta W)
% % % % % \;\approx\;
% % % % % J_{\mathrm{LoRA}}\,\theta,
% % % % % $
% % % % % where $J_{\mathrm{LoRA}}$ is the Jacobian from factor coordinates to full parameters. The image of $J_{\mathrm{LoRA}}$ coincides with $\mathcal W_\Delta$, so variations of $\theta$ span exactly the update subspace.


% % % % % Theorem~\ref{thm:pecl-bound} shows that, under PECL, both the empirical and complexity contributions are governed by quantities restricted to the LoRA update subspace $\mathcal W_\Delta$.
% % % % % Let $\varepsilon\sim\mathcal N(0,\Sigma_t)$ supported on $\mathcal W_\Delta$, and let $F_{\Delta,t}$ denote the Fisher matrix restricted to $\mathcal W_\Delta$.
% % % % % Under the standard local approximation, the perturbation-induced increase admits the proxy
% % % % % $
% % % % % \mathrm{E\text{-}Sh}_t^\Delta(\hat{\Delta W}_t,\Sigma_t)
% % % % % \approx
% % % % % \frac{1}{2}\,\mathrm{tr}\bigl(F_{\Delta,t}\Sigma_t\bigr)
% % % % % \le
% % % % % \frac{1}{2}\,\lambda_{\max}(F_{\Delta,t})\,\mathrm{tr}(\Sigma_t).
% % % % % $
% % % % % This relation motivates reporting $\lambda_{\max}(F)$ and $\mathrm{tr}(F)$ as curvature surrogates in the LoRA subspace.
% % % % % In this view, flatness-aware updates primarily target the empirical part through $\mathrm{E\text{-}Sh}_t^\Delta$, while PECL structurally limits the taskwise complexity through subspace-restricted KL terms.

% % % % % Classical continual learning regularizers can be interpreted in the same KL-based language.
% % % % % EWC penalizes parameter drift using a Fisher-weighted quadratic form, which corresponds to a local Gaussian approximation of a KL divergence. If one models the prior as a Gaussian centered at the previous solution with precision proportional to an estimated Fisher matrix,
% % % % % $
% % % % % P_t(\theta)=\mathcal N(\theta_{t-1},(\lambda F_{t-1})^{-1}),
% % % % % $
% % % % % and approximates the within-task posterior by a concentrated distribution around $\theta_t$, then
% % % % % $
% % % % % \mathrm{KL}\bigl(Q_t \,\Vert\, P_t\bigr)
% % % % % \approx
% % % % % \frac{\lambda}{2}\,(\theta_t-\theta_{t-1})^\top F_{t-1}(\theta_t-\theta_{t-1}),
% % % % % $
% % % % % which recovers the EWC penalty up to constants. Under this view, EWC indirectly constrains the complexity term in Theorem~\ref{thm:pathwise_pac}, while sharpness-aware methods act more directly on the empirical part via $\mathrm{E\text{-}Sh}_t^\Delta$. This provides a unified interpretation of different continual learning strategies within the PAC-Bayes decomposition into empirical and complexity contributions.

% % % % % Classical continual learning regularizers can be interpreted through the same KL-based language.
% EWC penalizes parameter drift via a Fisher-weighted quadratic term, which corresponds to a local Gaussian approximation of a KL divergence.
% Concretely, if one models the prior as a Gaussian centered at the previous solution with precision proportional to an estimated Fisher matrix, $P_t(\theta)=\mathcal N(\theta_{t-1},(\lambda F_{t-1})^{-1})$, and approximates the within-task posterior by a concentrated distribution around $\theta_t$, then the taskwise complexity term admits the surrogate
% $$
% \mathrm{KL}\bigl(Q_t \,\Vert\, P_t\bigr)
% \;\approx\;
% \frac{\lambda}{2}\,(\theta_t-\theta_{t-1})^\top F_{t-1}(\theta_t-\theta_{t-1}),
% $$
% which recovers the EWC penalty up to constants.
% % % % % From this perspective, EWC can be seen as indirectly controlling the KL complexity term in Theorem~\ref{thm:pathwise_pac}, whereas sharpness-aware optimizers act more directly on the empirical component through the expected sharpness term.
% % % % % This unified viewpoint helps explain why different CL strategies can be compared within a single PAC-Bayes decomposition into empirical and complexity contributions.

% % % % % This correspondence implies that flatness-aware optimisation applied to LoRA factors, such as SAM, directly reduces the dominant eigenvalues of the Fisher matrix restricted to $\mathcal W_\Delta$. For a fixed posterior covariance this contraction decreases $\mathrm{E\text{-}Sh}_t^\Delta$ and allows the adapter posterior to remain closer to its prior, which jointly tightens the sharpness term and the KL terms in Theorem~\ref{thm:pecl-bound}. Through the chain rule the same spectral contraction implies that the representation $\phi(x;W)$ becomes less sensitive to LoRA updates.

% % % % % Class prototypes drift less, empirical feature matrices exhibit larger effective rank and lower anisotropy, and NME-based evaluation becomes more stable. The empirical trends we observe in AAA, NME, prototype drift, CKA and feature-spectrum metrics are therefore consistent with the mechanism predicted by the bound: adapter-subspace flatness simultaneously improves PAC–Bayesian complexity and feature stability in PECL.



% \paragraph{Subspace design}

% \paragraph{Interaction with LoRA organization.}

% {%
% \setlength{\intextsep}{4pt}        % [h] 浮动体与正文的上下间距（最关键）
% \setlength{\textfloatsep}{4pt}     % [t]/[b] 时浮动体与正文间距（备用）
% \setlength{\abovecaptionskip}{1pt} % 图与caption之间（上方）间距
% \setlength{\belowcaptionskip}{0pt} % caption下方到正文的间距
% \captionsetup{skip=3pt}            % 你已在用：图与caption的距离（更细粒度）
% \begin{figure*}[thbp]
%     \centering
%     % ImageNet-C: cls / nme
%     % \begin{minipage}{0.495\linewidth}
%     %     \centering
%     %     \includegraphics[width=0.85\linewidth]{figures_TRML/AAA_curve_Imagenet-C_summary_inc10_rank16_ponit2.pdf}
%     %     % \subcaption{ImageNet-C: $\operatorname{ACC}_t^{\mathrm{cls}}$}
%     %     \label{fig:aaa_imagenet_c}
%     % \end{minipage}\hfill
%     \begin{subfigure}{0.45\linewidth}
%         \centering
%         \includegraphics[width=01\linewidth]{figures_TRML/AAA_curve_Imagenet-C_summary_inc10_rank16_dot2.pdf}
%         \subcaption{{Tiny ImageNet-C:} $\operatorname{ACC}_t^{\mathrm{cls}}$ under corruption robustness.}
%         \label{fig:nme_imagenet_c}
%     \end{subfigure}\hfill
%     % ImageNet-P: cls / nme
%     % \begin{minipage}{0.495\linewidth}
%     %     \centering
%     %     \includegraphics[width=0.85\linewidth]{figures_TRML/AAA_curve_Imagenet-P_summary_inc10_rank16_point2.pdf}
%     %     % \subcaption{ImageNet-P: $\operatorname{ACC}_t^{\mathrm{cls}}$}
%     %     \label{fig:aaa_imagenet_p}
%     % \end{minipage}\hfill
%     \begin{subfigure}{0.45\linewidth}
%         \centering
%         \includegraphics[width=1\linewidth]{figures_TRML/AAA_curve_Imagenet-P_summary_inc10_rank16_dot2.pdf}
%         \subcaption{{Tiny ImageNet-P:} $\operatorname{ACC}_t^{\mathrm{cls}}$ under perturbation robustness.}
%         \label{fig:nme_imagenet_p}
%     \end{subfigure}
% \vspace{-6pt}
%     \caption{ 
% Robustness in LoRA-based CL.
% We report the full task-wise classification accuracy trajectories $\operatorname{ACC}_t^{\mathrm{cls}}$ over the task sequence on Tiny ImageNet-C and Tiny ImageNet-P.
% Across LoRA organizations, adapter-subspace sharpness regularization via $\mathrm{AS}^{(0)}_S(W,\rho)$ and $\mathrm{AS}^{(1)}_S(W,\rho)$ consistently outperforms SGD, indicating that encouraging flatter minima within the trainable adapters improves robustness of the learned representations.
% The corresponding NME result are reported in the appendix.
% %         The full trajectories
% % of$\operatorname{ACC}_t$ on ImageNet-C and ImageNet-P. 
% %         % The four panels report performance after each task for the same class incremental methods and optimization configurations as in Figure~\ref{fig:combined_curves_imageR}. 
% %          $\mathrm{AS}^{(0)}_S(W,\rho)$and $\mathrm{AS}^{(1)}_S(W,\rho)$ consistently outperform SGD on both benchmark, which indicates that promoting flat minima in the adapter improves robustness of the learned representations.
%          %   and $\operatorname{ACC}_t^{\mathrm{cls}}$
%     }
%     \label{fig:robustness_curves_comparison}
% \end{figure*}
% }

% {%
% \setlength{\intextsep}{4pt}
% \setlength{\textfloatsep}{4pt}
% \setlength{\abovecaptionskip}{6pt}
% \setlength{\belowcaptionskip}{4pt}
% \captionsetup{skip=10pt}
% \begin{figure}[t]
% \centering
% \includegraphics[width=0.8\linewidth]{figures_TRML/AAA_curve_Imagenet-C_summary_inc10_rank16_dot2.pdf}
% \caption{$\operatorname{ACC}_t^{\mathrm{cls}}$ under corruption on Tiny ImageNet-C.}
% \label{fig:robustness_imagenet_c}
% \end{figure}
% }%

% {%
% \setlength{\intextsep}{4pt}
% \setlength{\textfloatsep}{4pt}
% \setlength{\abovecaptionskip}{6pt}
% \setlength{\belowcaptionskip}{4pt}
% \captionsetup{skip=10pt}
% \begin{figure}[t]
% \centering
% \includegraphics[width=0.8\linewidth]{figures_TRML/AAA_curve_Imagenet-P_summary_inc10_rank16_dot2.pdf}
% \caption{$\operatorname{ACC}_t^{\mathrm{cls}}$ under perturbation on Tiny ImageNet-P.}
% \label{fig:robustness_imagenet_p}
% \end{figure}
% }




% {%
% \setlength{\textfloatsep}{4pt}
% \setlength{\abovecaptionskip}{1pt}
% \setlength{\belowcaptionskip}{0pt}
% \captionsetup{skip=3pt}

% \begin{figure*}[t]
% \centering
% \begin{subfigure}[t]{0.95\linewidth}
%     \centering
%     \includegraphics[width=\linewidth]{figures_TRML/AAA_curve_Imagenet-C_summary_inc10_rank16_dot2.pdf}
%     \subcaption{{Tiny ImageNet-C:} $\operatorname{ACC}_t^{\mathrm{cls}}$ under corruption robustness.}
%     \label{fig:nme_imagenet_c}
% \end{subfigure}

% \vspace{2pt} % 这里可调小到 0pt 或 -2pt
% \begin{subfigure}[t]{0.95\linewidth}
%     \centering
%     \includegraphics[width=\linewidth]{figures_TRML/AAA_curve_Imagenet-P_summary_inc10_rank16_dot2.pdf}
%     \subcaption{{Tiny ImageNet-P:} $\operatorname{ACC}_t^{\mathrm{cls}}$ under perturbation robustness.}
%     \label{fig:nme_imagenet_p}
% \end{subfigure}

% \vspace{-4pt}
% \caption{
% Robustness in LoRA-based CL.
% We report the full task-wise classification accuracy trajectories $\operatorname{ACC}_t^{\mathrm{cls}}$ over the task sequence on Tiny ImageNet-C and Tiny ImageNet-P.
% Across LoRA organizations, adapter-subspace sharpness regularization via $\mathrm{AS}^{(0)}_S(W,\rho)$ and $\mathrm{AS}^{(1)}_S(W,\rho)$ consistently outperforms SGD, indicating that encouraging flatter minima within the trainable adapters improves robustness of the learned representations.
% The corresponding NME result are reported in the appendix.
% }
% \label{fig:robustness_curves_comparison}
% \end{figure*}
% }%





\noindent\textbf{Interaction with LoRA organization.}\quad
We further investigate how the effect of perturbation {type} varies with the LoRA organization. 
As shown in Table~\ref{table:q1_ci_with_weight_flatness_reorg} and  Fig.~\ref{fig:robustness_imagenet_cp}, the ordering of perturbation types is stable across LoRA organizations, but the performance gains vary in magnitude, indicating that perturbation geometry interacts with the organization-defined admissible update directions.
% As show in Table~\ref{table:q1_ci_with_weight_flatness_reorg} and Fig.~\ref{fig:robustness_curves_comparison},the ranking of perturbation types is consistent across organizations, but the gain varies in magnitude, suggesting an interaction between perturbation geometry and the organization-induced admissible update directions.
In particular, SeqLoRA remains more robust on ImageNet-C/P, whereas IncLoRA and OLoRA exhibit larger accuracy drops under shift.
% Across all organizations, the relative ordering of perturbation types is preserved, yet the magnitude of the gains differs, indicating that perturbation geometry interacts with the organization-induced set of admissible update directions.
% Empirically, the LoRA organizations behave differently on ImageNet-R and on the robustness benchmarks ImageNet-C and ImageNet-P: SeqLoRA over performance  IncLoRA and OLoRA under distribution shift.
% The perfomance of different LoRA organization are diffetent in the  ImageNet-R and  robustness benchmark ImageNet-C and ImageNet-P. SeqLoRA shows the smallest accuracy degradation across tasks, whereas IncLoRA and OLoRA exhibit larger drops underrobustness. 
% This pattern is consistent with our PECL perspective.
Since the sharpness term is defined on the task adapter subspace, it is inherently sensitive to how different LoRA organization shapes the admissible scope $\mathcal W_{\Delta,t}$ across tasks. When $\mathcal W_{\Delta,t}$ is larger or expands over time, perturbations can probe a broader set of admissible directions, including more directions with non-negligible projected curvature.

% In contrast, SeqLoRA maintains an approximately fixed, shared adapter subspace, which limits the directions that are simultaneously admissible for both perturbation and optimization, and is therefore associated with smaller degradation under shift.

% This observation aligns with our PECL viewpoint: the sharpness-relevant empirical contribution is defined on the task-admissible adapter subspace, and is therefore sensitive to how the organization shapes the admissible scope $\mathcal W_{\Delta,t}$ across tasks.
% When $\mathcal W_{\Delta,t}$ is larger or evolves over time, perturbations can probe a broader set of admissible directions.
% In contrast, SeqLoRA keeps an approximately fixed and shared adapter subspace, which restricts the directions that are simultaneously admissible for perturbation and optimization, and is associated with smaller degradation under shift.




% \begin{figure}[thbp]
% \centering
% \includegraphics[width=1\linewidth]{figures_TRML/AAA_curve_Imagenet-C_summary_inc10_rank16_changeColor.pdf}
% \label{fig:nme_imagenet_c}
%  \caption{ 
% The full trajectories
% of$\operatorname{ACC}_t$ on ImageNet-C. 
% % $\mathrm{AS}^{(0)}_S(W,\rho)$and $\mathrm{AS}^{(1)}_S(W,\rho)$ consistently outperform SGD on both benchmark, which indicates that promoting flat minima in the adapter improves robustness of the learned representations. 
% }
% \label{fig:robustness_curves_comparison}
% \end{figure}

% \begin{figure}[thbp]
% \centering
% \includegraphics[width=1\linewidth]{figures_TRML/AAA_curve_Imagenet-C_summary_inc10_rank16_changeColor.pdf}
% \label{fig:nme_imagenet_c}
%  \caption{ 
% The full trajectories
% of$\operatorname{ACC}_t$ on ImageNet-C. 
% % $\mathrm{AS}^{(0)}_S(W,\rho)$and $\mathrm{AS}^{(1)}_S(W,\rho)$ consistently outperform SGD on both benchmark, which indicates that promoting flat minima in the adapter improves robustness of the learned representations. 
% }
% \label{fig:robustness_curves_comparison}
% \end{figure}




%%%%%%%%%%%%%%%%%%%%%%%%%
% \paragraph{Collaboration between LoRA design and perturbation selection}


% \paragraph{Take Away}



% \section{Results and Analysis}
% \label{sec:results}






% \begin{table*}[htpb]
% \centering
% \caption{
% Class incremental performance, weight-space flatness and feature-space robustness on ImageNet-R ($T=20$, $R=16$). 
% }
% \label{tab:table_all}
% \setlength{\tabcolsep}{4pt}
% \renewcommand{\arraystretch}{0.95}
% \resizebox{\linewidth}{!}{
% \begin{tabular}{l l cccc cc cc}
% \toprule
% \multirow{2}{*}{\textbf{Method}} & \multirow{2}{*}{\textbf{Opt.}} &
% \multicolumn{4}{c}{\textbf{CL Performance} ($\uparrow$)} &
% \multicolumn{2}{c}{\textbf{Weight Flatness} ($\downarrow$)} &
% \multicolumn{2}{c}{\textbf{Feature Robustness}} \\
% \cmidrule(lr){3-6}\cmidrule(lr){7-8}\cmidrule(lr){9-10}
% & &
% $\operatorname{ACC}_{20}^{\mathrm{cls}}$ &
% $\operatorname{AAA}^{\mathrm{cls}}$ &
% $\operatorname{ACC}_{20}^{\mathrm{nme}}$ &
% $\operatorname{AAA}^{\mathrm{nme}}$ &
% Sh$^{(0)}(\rho)$ &
% $\mathrm{tr}(F)$ &
% L2 drift ($\downarrow$) &
% $\mathrm{tr}(E)$ \\
% \midrule
% \multirow{1}{*}{\textbf{PerTaskLP}}
% & SGD
% & $55.22^{\pm 0.31}$ & $59.36^{\pm 0.19}$
% & $62.96^{\pm 0.12}$ & $66.58^{\pm 0.37}$
% & 0.072 & 1015.07
% & 0.00 & 0.82 \\
% \midrule
% \multirow{3}{*}{\textbf{SeqFT}}
% & SGD
% & $50.72^{\pm 2.20}$ & $59.89^{\pm 1.32}$
% & $69.49^{\pm 0.16}$ & $75.10^{\pm 0.57}$
% & 0.582 & 29703.66
% & 29.57 & 0.58 \\
% & SAM
% & $56.07^{\pm 2.60}$ & $66.47^{\pm 1.53}$
% & $72.36^{\pm 2.46}$ & $78.31^{\pm 0.50}$
% & 0.259 & 5881.32
% & 17.04 & 0.74 \\
% & GAM
% & $59.78^{\pm 0.74}$ & $67.31^{\pm 1.18}$
% & $73.74^{\pm 0.88}$ & $78.32^{\pm 0.94}$
% & 0.177 & 3273.08
% & 13.80 & 0.78 \\
% \midrule
% \multirow{3}{*}{\textbf{SeqLoRA}}
% & SGD
% & $63.14^{\pm 2.17}$ & $69.07^{\pm 1.33}$
% & $72.82^{\pm 1.27}$ & $76.62^{\pm 0.68}$
% & 0.142 & 1401.15
% & 16.58 & 0.77 \\
% & SAM
% & $67.04^{\pm 0.71}$ & $71.42^{\pm 0.71}$
% & $74.53^{\pm 0.60}$ & $77.62^{\pm 0.43}$
% & 0.120 & 821.45
% & 14.82 & 0.81 \\
% & GAM
% & $73.31^{\pm 0.52}$ & $76.80^{\pm 0.56}$
% & $75.74^{\pm 0.15}$ & $78.86^{\pm 0.53}$
% & 0.060 & 240.91
% & 11.28 & 0.78 \\
% \midrule
% \multirow{3}{*}{\textbf{IncLoRA}}
% & SGD
% & $65.77^{\pm 2.84}$ & $71.29^{\pm 0.78}$
% & $75.41^{\pm 0.26}$ & $78.37^{\pm 0.26}$
% & 0.076 & 836.91
% & 16.87 & 0.67 \\
% & SAM
% & $68.86^{\pm 0.92}$ & $72.59^{\pm 0.39}$
% & $75.95^{\pm 0.07}$ & $78.68^{\pm 0.25}$
% & 0.053 & 526.49
% & 15.14 & 0.70 \\
% & GAM
% & $72.86^{\pm 0.23}$ & $76.52^{\pm 0.48}$
% & $76.72^{\pm 0.54}$ & $79.42^{\pm 0.40}$
% & 0.035 & 192.25
% & 12.15 & 0.73 \\
% \midrule
% \multirow{3}{*}{\textbf{OLoRA}}
% & SGD
% & $66.64^{\pm 0.44}$ & $71.41^{\pm 0.47}$
% & $75.10^{\pm 0.41}$ & $78.65^{\pm 0.67}$
% & 0.078 & 831.23
% & 17.23 & 0.68 \\
% & SAM
% & $68.53^{\pm 0.78}$ & $72.75^{\pm 0.76}$
% & $75.73^{\pm 0.32}$ & $78.70^{\pm 0.67}$
% & 0.066 & 628.12
% & 15.17 & 0.70 \\
% & GAM
% & $72.44^{\pm 0.36}$ & $76.05^{\pm 0.52}$
% & $76.56^{\pm 0.53}$ & $79.16^{\pm 0.43}$
% & 0.036 & 198.30
% & 12.16 & 0.74 \\
% % \midrule
% % \multirow{1}{*}{\textbf{l2p}}
% % & Adam
% % & $71.35^{\pm 0.05}$ & $74.84^{\pm 0.73}$
% % & -- & --
% % & -- & --
% % & -- & -- \\
% % \midrule
% % \multirow{1}{*}{\textbf{dualprompt}}
% % & Adam
% % & $71.59^{\pm 0.68}$ & $74.80^{\pm 1.18}$
% % & -- & --
% % & -- & --
% % & -- & -- \\
% \bottomrule
% \end{tabular}
% }
% \end{table*}


% \subsection{Main results on ImageNet-R}
% \label{sec:results-main-inr}


 
% Table~\ref{tab:table_all} summarises final accuracy $\operatorname{ACC}_{20}^{\mathrm{cls}}$, averaged accuracy $\operatorname{AAA}^{\mathrm{cls}}$, with the corresponding NME metrics together with representative sharpness proxies in the parameter spac and feature space diagnostics. Figure~\ref{fig:combined_curves_imageR} reports classifier head accuracy $\operatorname{ACC}_t^{\mathrm{cls}}$ and NME based accuracy $\operatorname{ACC}_t^{\mathrm{nme}}$ over the task sequence on ImageNet-R.

% Figure~\ref{fig:combined_curves_imageR} complements this summary with the full trajectories of $\operatorname{ACC}_t^{\mathrm{cls}}$ and $\operatorname{ACC}_t^{\mathrm{nme}}$ over the task sequence.

% To check the design of the effect of lora. The figure \ref{fig: weight similiairy} show distance in the trainging process between $\Delta W_t$ and $\Delta W_{t-1}$ and feature amplification factor $
% A_t =
% \frac{\left\lVert \Delta W^{\mathrm{eff}}_{t}\right\rVert_F}
% {\left\lVert U_{t}^{\top} W V_{t}\right\rVert_F+\varepsilon}
% $ which represent how  “strong” do we need to amplify  task-specific directions in order for the model adaptation to work well.  There are  feature directions need to be amplified in order to achieve our reported accuracy for the downstream specific task. It show seqlora have a larege overlap in the traing,  inclora and olora reduce a lot by apply new admissable space. The amp factor show the same trace that the inclora and olora learned more task specific direction compare seqlora.


% SeqLoRA becomes the most robust configuration, while IncLoRA and OLoRA suffer larger accuracy drops. This is consistent with the dependence of the effective sharpness term on the spectrum of $F_W$ restricted to the adapter subspace. Larger update subspaces admit more directions along which curvature and shift induced drift can accumulate, whereas the compact SeqLoRA subspace constrains both $\Sigma_t$ and $F_\Delta$ to a smaller region and therefore exposes fewer unstable directions.

% The relative behaviour of different LoRA organizations changes in the presence of strong corruptions. On
% ImageNet-C and ImageNet-P, SeqLoRA becomes the most robust configuration. This is consistent with the dependence of the effective sharpness term on the . Larger update subspaces admit more directions along
% which curvature and shift induced drift can accumulate, whereas the compact SeqLoRA subspace constrains to a smaller region and therefore exposes fewer unstable directions.


% Across all three LoRA organizations, applying SAM or GAM to the adapter parameters improves class-incremental performance compared with SGD under the same training budget. Both final and time-averaged accuracies increase for the classifier head and for NCM evaluation, that is, $\operatorname{ACC}_{20}^{\mathrm{cls}}$, $\operatorname{AAA}^{\mathrm{cls}}$, $\operatorname{ACC}_{20}^{\mathrm{ncm}}$ and $\operatorname{AAA}^{\mathrm{ncm}}$ all improve, with GAM typically achieving the strongest gains. 
% The trajectories in Figure~\ref{fig:combined_curves_imageR} show that this advantage is in all process: SAM and GAM maintain higher accuracy throughout the sequence. The  alignment between the classifier-head and NME curves indicates that these gains reflect more stable shared representations.
% These performance gains are tightly aligned with our flatness and feature metrics. For each LoRA organization, moving from SGD to SAM and then to GAM monotonically decreases adversarial sharpness $\mathrm{Sh}^{(0)}(\rho)$ as well as the largest eigenvalue and trace of the empirical Fisher, while simultaneously reducing prototype drift, and yielding a more isotropic feature matrix. 
% \paragraph{Take awat}

% Taken together, these trends support the central prediction of our PAC-Bayesian analysis that reducing expected sharpness within the adapter subspace is sufficient to tighten the generalization guarantee and improve the stability–plasticity trade-off in PECL.

% \begin{figure}[hpbt]
%     \centering
%     \begin{minipage}{0.49\linewidth}
%         \centering
%         \includegraphics[width=\linewidth]{figures_TRML/AAA_curve_Imagenet-R_summary_inc10_rank16—point2.pdf}
%         % \subcaption{Classifier head accuracy $\operatorname{ACC}_t^{\mathrm{cls}}$ on ImageNet-R}
%         \label{fig:curves_imageR_1_AAA}
%     \end{minipage}
%      \begin{minipage}{0.49\linewidth}
%         \centering
%         \includegraphics[width=\linewidth]{figures_TRML/NME_curve_Imagenet-R_summary_inc10_rank16_point2.pdf}
%         % \subcaption{Prototype based accuracy $\operatorname{ACC}_t^{\mathrm{nme}}$ on ImageNet-R}
%         \label{fig:curves_imageR_2_NME}
%     \end{minipage}
%     \caption{
%     The full trajectories
% of $\operatorname{ACC}_t^{\mathrm{cls}}$ and $\operatorname{ACC}_t^{\mathrm{nme}}$ on ImageNet-R. 
%     % Summary curves of classifier head accuracy $\operatorname{ACC}_t^{\mathrm{cls}}$ and NME based accuracy $\operatorname{ACC}_t^{\mathrm{nme}}$ on ImageNet-R. 
%     % For all three LoRA organizations, flatness promoting optimizers such as SAM and GAM consistently outperform SGD, and their advantage becomes more pronounced on later tasks, which indicates more stable representations and reduced forgetting.
%     }
%     \label{fig:combined_curves_imageR}
% \end{figure}








% tables \ref{table:q1_ci_with_weight_flatness_reorg} and \ref{tab:q1_lp_feat_reorg} remain as in Section~\ref{sec: flat minima}. We only change how they are discussed here.



% \subsection{Robustness on ImageNet-C and ImageNet-P}
% \label{sec:results-robustness}




% % $\operatorname{ACC}_t^{\mathrm{cls}}$ and 

% Across a wide range of corruption types and perturbation sequences, sharpness-based LoRA training consistently improves performance: SAM and GAM achieve higher $\operatorname{ACC}_t^{\mathrm{cls}}$ and $\operatorname{ACC}_t^{\mathrm{ncm}}$ than SGD. This pattern indicates that promoting flatness in the LoRA subspace does not trade clean accuracy for robustness but instead enhances generalization in a way that transfers to distribution shift. 




% \begin{figure}[htbp]
%     \centering
%     % ImageNet-C: cls / nme
%     \begin{minipage}{0.48\linewidth}
%         \centering
%         \includegraphics[width=\linewidth]{figures_TRML/AAA_curve_Imagenet-C_summary_inc10_rank16_ponit2.pdf}
%         % \subcaption{ImageNet-C: $\operatorname{ACC}_t^{\mathrm{cls}}$}
%         \label{fig:aaa_imagenet_c}
%     \end{minipage}\hfill
%     % ImageNet-P: cls / nme
%     \begin{minipage}{0.48\linewidth}
%         \centering
%         \includegraphics[width=\linewidth]{figures_TRML/AAA_curve_Imagenet-P_summary_inc10_rank16_point2.pdf}
%         % \subcaption{ImageNet-P: $\operatorname{ACC}_t^{\mathrm{cls}}$}
%         \label{fig:aaa_imagenet_p}
%     \end{minipage}\hfill
    
%     \caption{
%         The NME-based ($\operatorname{ACC}_t^{\mathrm{nme}}$) on ImageNet-C and ImageNet-P. 
%         % The four panels report performance after each task for the same class incremental methods and optimization configurations as in Figure~\ref{fig:combined_curves_imageR}. 
%          SAM and GAM consistently outperform SGD on both corruption robustness (ImageNet-C) and perturbation robustness (ImageNet-P), which indicates that promoting flat minima in the adapter improves robustness of the learned representations.
%     }
%     \label{fig:robustness_curves_comparison}
% \end{figure}


% We next evaluate robustness under corruptions and perturbations on Tiny ImageNet-C and Tiny ImageNet-P.
%  Figure~\ref{fig:robustness_curves_comparison} complements this summary with the full trajectories of $\operatorname{ACC}_t^{\mathrm{cls}}$ and $\operatorname{ACC}_t^{\mathrm{nme}}$ over the task sequence.
% % For each LoRA organization and optimizer, we measure performance using $\operatorname{ACC}_t^{\mathrm{cls}}$ and $\operatorname{ACC}_t^{\mathrm{nme}}$.

% Across a wide range of corruption types and perturbation sequences, flatness oriented LoRA training consistently improves robustness. 
% SAM and GAM achieve higher corrupted $\operatorname{ACC}_t^{\mathrm{cls}}$ and $\operatorname{ACC}_t^{\mathrm{nme}}$ than SGD, and the ordering observed on clean ImageNet-R is preserved on ImageNet-C and ImageNet-P. 
% The gains under corruption are strongly correlated with the gains on clean data: configurations that achieve higher clean AAA also yield higher corrupted AAA. 
% This pattern indicates that LoRA subspace flatness does not trade clean accuracy for robustness. 
% Instead, it leads to a global enhancement of generalization that extends to distribution shifts, which is consistent with the PAC Bayesian view that smaller expected sharpness tightens the error bound uniformly over the data distribution.





% \subsection{Result on NLP Benchmark}

% \begin{table}[thb]
% \centering
% \caption{Summary of the results on two standard CL benchmarks with T5-large model. $ACC_T$ is reported.} % 标题不改：如你原本有 caption 就把它放这里
% \label{tab:nlp_accT_orders} % 如你原本有 label 就沿用
% % 列间距和行高压缩
% \setlength{\tabcolsep}{3pt}
% \renewcommand{\arraystretch}{0.9}
% \small
% \resizebox{0.98\linewidth}{!}{
% \begin{tabular}{l cccc cccc}
% \toprule
% \multirow{2}{*}{} &
% \multicolumn{4}{c}{\textbf{Standard CL Benchmark}} &
% \multicolumn{4}{c}{\textbf{Large Number of Tasks}} \\
% \cmidrule(lr){2-5}\cmidrule(lr){6-9}
% \textbf{Method} &
% \textbf{Order1} & \textbf{Order2} & \textbf{Order3} & \textbf{AVG} &
% \textbf{Order4} & \textbf{Order5} & \textbf{Order6} & \textbf{AVG} \\
% \midrule
% SeqLoRA     & 61.66          & 59.99          & 72.22          & 64.62          & 55.85          & 40.65          & 27.05          & 41.18         \\
% +$\mathrm{RS}_S^{(0)}({W}, \Sigma)$ & \textbf{63.33} & \textbf{63.16} & \textbf{72.95} & \textbf{66.48} & \textbf{56.20} & \textbf{56.25} & \textbf{12.77} & \textbf{41.74}   \\
% \midrule
% IncLoRA     &\textbf{71.18}          & \textbf{71.56}          & 72.87          & 71.87          & 58.13          & 56.42          & 39.26          & 51.27            \\
% +$\mathrm{RS}_S^{(0)}({W}, \Sigma)$ & 70.20          & 71.18          & \textbf{75.78} & \textbf{72.39}          & \textbf{67.5}  & \textbf{56.54}          & \textbf{66.38}          & \textbf{63.47} \\
% \midrule
% OLoRA       & 76.41          & 76.93          & 73.69          & 75.68          & 66.73          & 68.62          &\textbf{71.51}          & 68.95         \\
% +$\mathrm{RS}_S^{(0)}({W}, \Sigma)$   & \textbf{78.31} & \textbf{77.96} & \textbf{75.95} & \textbf{77.41} & \textbf{71.08} & \textbf{69.05} & {70.08} & \textbf{70.07}   \\
% \bottomrule
% \end{tabular}
% }
% \end{table}

% We further conduct continual task learning experiments on two standard NLP benchmarks under multiple task orders, and report the final accuracy $ACC_T$. Table~\ref{tab:nlp_accT_orders} shows that SAM consistently improves $ACC_T$ over the SGD baselines across most configurations. The gains are more pronounced in the long-sequence regime, suggesting that promoting flatness within the LoRA subspace remains effective when inter-task interference accumulates. Although a few order-specific deviations are observed, the overall trend indicates that SAM is a generally beneficial for LoRA-based CL methods in NLP.









\subsection{Ablation study on task length and adapter rank}
\label{sec:results-ablation}


{%
\setlength{\intextsep}{4pt}        % [h] 浮动体与正文的上下间距（最关键）
\setlength{\textfloatsep}{4pt}     % [t]/[b] 时浮动体与正文间距（备用）
\setlength{\abovecaptionskip}{1pt} % 图与caption之间（上方）间距
\setlength{\belowcaptionskip}{0pt} % caption下方到正文的间距
\captionsetup{skip=3pt}            % 你已在用：图与caption的距离（更细粒度）
\begin{figure}[h]
    \centering
\includegraphics[width=1\linewidth]{figures_TRML/hebing_ablation1.pdf}
\vspace{-10pt}
    \caption{ Sensitivity to Adapter Capacity and Sequence Difficulty
    % Ablation on adapter rank and task sequence length on ImageNet-R.
% We report the final NME-based accuracy $\operatorname{ACC}_T^{\mathrm{nme}}$ while varying (a) the adapter rank $r$ and (b) the task length $T$.
    }
    \label{fig:ablation}
\end{figure}
}


Ablation on adapter rank and task sequence length on ImageNet-R.
We study how adapter capacity and sequence difficulty affect PECL by varying the LoRA rank $r$ and the task sequence length $T$, and report the final NME accuracy $\operatorname{ACC}_T^{\mathrm{nme}}$ in Fig.~\ref{fig:ablation}.
Across all settings, the relative ordering is stable: flatness-aware optimizers consistently outperform SGD.
Accuracy improves from small $r$ but quickly saturates at moderate ranks. This supports our theoretical view that PECL primarily requires a low-dimensional task-admissible update subspace to capture task-specific directions, so enlarging $r$ beyond this intrinsic dimension adds capacity with limited benefit while leaving the bound-relevant geometry largely unchanged.

% We further vary the adapter rank $r$ and the task sequence length $T$ on ImageNet-R and report $\operatorname{ACC}_t^{\mathrm{nme}}$ in Fig.~\ref{fig:ablation}.
% Across all values of $T$ and $r$, the relative ordering remains stable: flatness-aware optimizers consistently outperform SGD.

% {%
% \setlength{\intextsep}{4pt}        % [h] 浮动体与正文的上下间距（最关键）
% \setlength{\textfloatsep}{4pt}     % [t]/[b] 时浮动体与正文间距（备用）
% \setlength{\abovecaptionskip}{1pt} % 图与caption之间（上方）间距
% \setlength{\belowcaptionskip}{0pt} % caption下方到正文的间距
% \captionsetup{skip=3pt}            % 你已在用：图与caption的距离（更细粒度）
% \begin{figure}[thbp]
%     \centering
%     % ImageNet-C: cls / nme
%     % \begin{minipage}{0.495\linewidth}
%     %     \centering
%     %     \includegraphics[width=0.85\linewidth]{figures_TRML/AAA_curve_Imagenet-C_summary_inc10_rank16_ponit2.pdf}
%     %     % \subcaption{ImageNet-C: $\operatorname{ACC}_t^{\mathrm{cls}}$}
%     %     \label{fig:aaa_imagenet_c}
%     % \end{minipage}\hfill
%     \begin{subfigure}{0.45\linewidth}
%         \centering
%         \includegraphics[width=0.8\linewidth]{figures_TRML/rank_NME_T_by_method_opt_imagenetr2.pdf}
%         % \subcaption{ Adapter rank.}
%     \end{subfigure}\hfill
%     % ImageNet-P: cls / nme
%     % \begin{minipage}{0.495\linewidth}
%     %     \centering
%     %     \includegraphics[width=0.85\linewidth]{figures_TRML/AAA_curve_Imagenet-P_summary_inc10_rank16_point2.pdf}
%     %     % \subcaption{ImageNet-P: $\operatorname{ACC}_t^{\mathrm{cls}}$}
%     %     \label{fig:aaa_imagenet_p}
%     % \end{minipage}\hfill
%     \begin{subfigure}{0.45\linewidth}
%         \centering
%         \includegraphics[width=0.8\linewidth]{figures_TRML/length_NME_T_by_method_opt_imagenetr2.pdf}
%         % \subcaption{Task length}
%     \end{subfigure}
% \vspace{-8pt}
%     \caption{Ablation on adapter rank and task sequence length on ImageNet-R.
% We report the final NME-based accuracy $\operatorname{ACC}_T^{\mathrm{nme}}$ while varying (a) the adapter rank $r$ and (b) the task length $T$.
%          %   and $\operatorname{ACC}_t^{\mathrm{cls}}$
%     }
%     \label{fig:robustness_curves_comparison}
% \end{figure}
% }





% We further vary the task sequence length and the adapter rank on ImageNet-R and report $\operatorname{AAA}^{\mathrm{ncm}}$ in Fig.~\ref{fig:ablation}. Across all values of $T$ and $r$, the relative behavior observed remains stable: flatness aware optimizers consistently outperform SGD.

% Finally, we study the robustness of the above conclusions with respect to the task sequence length $T$ and the LoRA rank $r$. 
% We repeat the ImageNet-R experiments for shorter and longer sequences and for several ranks $r$ while keeping all other hyperparameters fixed. 
% The detailed $\operatorname{ACC}_t^{\mathrm{cls}}$ and $\operatorname{ACC}_t^{\mathrm{nme}}$ curves are reported in Appendix~\ref{app:ablation-length-rank}.

% \begin{figure}[htbp]
%     \centering
%         \includegraphics[width=0.6\linewidth]{figures_TRML/rank_ACC_T_by_method_opt_imagenetr_changeColor.pdf}
%         % \subcaption{Different Length}
%         % \label{fig:curves_imageR_length}
%     % \begin{minipage}{0.495\linewidth}
%     %     \centering
%     %     \includegraphics[width=\linewidth]{figures_TRML/rank_NME_T_avg_by_method_opt_imagenetr.pdf}
%     %     % \subcaption{Different Rank}
%     %     % \label{fig:curves_imageR_rank}
%     % \end{minipage}
%     \caption{
%         Effect of adapter rank on ImageNet-R.
%     }
%     \label{fig:ablation}
% \end{figure}




% \begin{figure}[htbp]
%     \centering
%     \begin{minipage}{0.495\linewidth}
%         \centering
%         \includegraphics[width=\linewidth]{figures_TRML/length_NME_T_avg_by_method_opt_imagenetr.pdf}
%         % \subcaption{Different Length}
%         % \label{fig:curves_imageR_length}
%     \end{minipage}%
%     \begin{minipage}{0.495\linewidth}
%         \centering
%         \includegraphics[width=\linewidth]{figures_TRML/rank_NME_T_avg_by_method_opt_imagenetr.pdf}
%         % \subcaption{Different Rank}
%         % \label{fig:curves_imageR_rank}
%     \end{minipage}
%     \caption{
%         Ablation of task sequence length and adapter rank on ImageNet-R.
%     }
%     \label{fig:ablation}
% \end{figure}

% The qualitative picture remains stable. 
% For all tested values of $T$ and $r$, applying SAM or GAM in the adapter subspace consistently improves $\operatorname{AAA}^{\mathrm{cls}}$ and $\operatorname{AAA}^{\mathrm{nme}}$ over SGD, and the accuracy gap tends to be more pronounced on later tasks.




% Changing $T$ or $r$ affects the absolute level of performance and the magnitude of forgetting, as expected, but does not alter the relative ordering among optimizers. 
% Feature space diagnostics follow the same pattern: when flatness is promoted in $W_\Delta$, prototype drift is reduced and the empirical feature matrix exhibits higher effective rank across all ablation settings. 

% These results indicate that the empirical support for the sharpness aware PAC Bayesian analysis is robust to reasonable variations in task length and adapter capacity, and does not depend on a narrow choice of hyperparameters.



% \section{Discussion, Limitations, and Related Work}



% \section{Conclusion and Limitations}
\section{Related Work}


\subsection{Flat Minima}
Prior work has explored flat minima in CL through sharpness-aware optimization~\cite{pmlr-v151-bisla22a,foretSharpnessAwareMinimizationEfficiently2021,zhangGradientNormAware2023}, mode connectivity~\cite{izmailov2019averagingweightsleadswider,mirzadeh2020linearmodeconnectivitymultitask}, and improved representation learning~\cite{hu2022doesselfsupervisedpretrainingperform,ramasesh2022effect}.
These methods are mainly developed for full model training and do not account for the adapter subspace constraints\cite{yangRevisitingFlatnessAwareOptimization2025}. Whether sharpness within the adapter subspace alone can ensure generalization in LoRA-based CL remains unclear, motivating further investigation.


% % Flat minima are widely believed to promote generalization and robustness by reducing the sensitivity of the loss to parameter perturbations~\citep{FlatMinima1997,wu2017understandinggeneralizationdeeplearning,chaudhariEntropySGDBiasingGradient2019,NEURIPS2020_518a38cc,bahri2022sharpnessawareminimizationimproveslanguage,NEURIPS2021_995f5e03,JMLRAnEmpiricalInvestigatio,yangRevisitingFlatnessAwareOptimization2025}. Sharpness-aware optimization methods such as SAM and its variants explicitly minimize adversarial or stochastic sharpness and have been shown to improve both accuracy and robustness in full finetuning regimes~\citep{foretSharpnessAwareMinimizationEfficiently2021,pmlr-v151-bisla22a,Zhang_GAM}. 

% In continual learning, flatness has been studied through mode connectivity and representation smoothing~\citep{izmailov2019averagingweightsleadswider,mirzadeh2020linearmodeconnectivitymultitask,ramasesh2022effect,hu2022doesselfsupervisedpretrainingperform}, but these works operate on the full parameter space and do not distinguish between backbone and adapter directions. 

% More recently, Flat-LoRA-type approaches apply perturbations to the entire parameter space when training LoRA adapters, arguing that subspace flatness may be misaligned with full-space flatness. However, in PECL with a strictly frozen backbone, full-model perturbations conflict with the design principle that old knowledge is preserved in $W_0$. Existing analyses therefore leave open a central question that motivates our work: in LoRA-based PECL, is it sufficient to control flatness within the shared adapter subspace, or must one still reason in the full parameter space to explain forgetting and generalization?

\subsection{PAC-Bayesian Theory }

The PAC-Bayesian framework provides a principled foundation for analyzing generalization by balancing empirical risk and posterior complexity~\cite{mcallesterPACBayesianStochasticModel2003,pentinaPACBayesianBoundLifelong2014,germainPACBayesianTheoryMeets2016,lotfiPACBayesCompressionBounds2022}.
Online PAC-Bayes~\citep{NEURIPS2022_onlineBayes,NEURIPS2024_xinrui} allows time-varying, data-dependent priors, but it is formulated at the level of individual examples rather than tasks and does not explicitly model task-level inheritance in continual learning.
Hierarchical PAC-Bayes bounds for meta learning~\citep{pentinaPACBayesianBoundLifelong2014,JMLR:v24:22-1254} introduce hyperposteriors over task priors, but they assume a single stationary hyperposterior and therefore do not model the sequential, history-dependent evolution of the inductive bias that underlies forgetting in continual learning.
% a fixed hyperposterprior to consider multi-task risk. rather than forgetting in CL and omit the nonstationary accumulation of information across tasks.
% However, continual learning is inherently history dependent, and this is amplified in LoRA-based PECL where new adapters are introduced on top of the current model with initializations that can depend on past training.
% This motivates our pathwise hierarchical PAC-Bayes analysis, which models task priors as history-measurable draws from a time-varying hyperposterior.
In continual learning, the inductive bias must adapt to the observed task history, and this nonstationarity is particularly pronounced in LoRA-based PECL, where the update scope and the initialization of newly introduced adapters can change with training.

% Motivated by this mismatch, we develop a pathwise hierarchical PAC-Bayes analysis that treats the task prior at step $t$ as a history-measurable draw from a time-varying hyperposterior process, yielding an explicit hyperposterior drift penalty that quantifies the information budget of sequential belief updates.


% The PAC-Bayesian framework analyzes generalization by trading off empirical risk and posterior complexity~\cite{mcallesterPACBayesianStochasticModel2003,pentinaPACBayesianBoundLifelong2014,germainPACBayesianTheoryMeets2016,lotfiPACBayesCompressionBounds2022}.
% Online PAC-Bayes~\citep{NEURIPS2022_onlineBayes,NEURIPS2024_xinrui} allows time-varying, data-dependent priors and thus captures sequential dependence in the data stream. But it is not tailored to represent the task-level inheritance mechanism that drives continual learning.
% Hierarchical PAC-Bayes bounds for meta-learning~\citep{pentinaPACBayesianBoundLifelong2014,JMLR:v24:22-1254} introduce hyperposteriors over task-specific priors to formalize transfer, yet they assume a fixed hyperprior and therefore obscure the information accumulation and nonstationarity across tasks.
% However, continual learning is intrinsically history dependent: each task is trained on top of previously learned parameters, so the effective task priors evolve over time rather than arising as independent stationary draws.
% This effect is especially pronounced in LoRA-based PECL, where new adapter components are added on top of the current model and their initialization can explicitly depend on past training\cite{}.
% Motivated by this mismatch, we develop a pathwise hierarchical PAC-Bayes analysis that treats task priors as history-measurable draws from a time-varying hyperposterior.

% The PAC-Bayesian framework provides a principled foundation for analyzing generalization by balancing empirical risk and posterior complexity~\cite{mcallesterPACBayesianStochasticModel2003,pentinaPACBayesianBoundLifelong2014,germainPACBayesianTheoryMeets2016,lotfiPACBayesCompressionBounds2022}.
% Online PAC-Bayesian work~\citep{NEURIPS2022_onlineBayes,NEURIPS2024_xinrui}introduces time-varying, data-dependent priors with guarantees that depend on cumulative data, but it can't excatly the  inheriter learning process.
% Hierarchical PAC-Bayesian bounds~\citep{pentinaPACBayesianBoundLifelong2014,JMLR:v24:22-1254} learn a hyper-posterior over task-specific priors to express the meat learing process.
% But they assume a fixed hyper-prior and therefore do not model the sequential dependence and information accumulation that characterize continual learning. 
% However,Continual learning is intrinsically history-dependent: each task is trained on top of the parameters obtained from previous tasks, so the effective task priors evolve over time rather than arising as independent stationary draws from a fixed hyperprior.
% This dependence is further amplified in LoRA-based PECL, where both the adapter organization and initialization rely on past model.
% This motivates our pathwise hierarchical PAC-Bayes analysis, which models task priors as history-measurable random variables drawn from a time-varying hyperposterior.

% In parallel, online PAC-Bayesian work~\citep{NEURIPS2022_onlineBayes,NEURIPS2024_xinrui}introduces time-varying, data-dependent priors with guarantees that depend on cumulative data, yet it is formulated at the level of individual examples rather than tasks and does not consider task-conditioned posteriors.
% For meta-learning, hierarchical PAC-Bayes bounds introduce a hyper-prior over task-specific priors and derive complexity terms that couple task-level posteriors through the hyper posterior~\citep{pentinaPACBayesianBoundLifelong2014,JMLR:v24:22-1254}. In these formulations, all task priors are assumed to be independent draws from a fixed hyper-prior, which is natural for multi-task transfer but does not encode the sequential dependence and information accumulation that characterize continual learning.
% Parallel work on online PAC-Bayes develops bounds with time-varying, data-dependent priors that are measurable with respect to a filtration of past observations~\citep{NEURIPS2022_onlineBayes,NEURIPS2024_xinrui}. These results capture sequential adaptation at the level of individual examples and have inspired practical algorithms, yet they are not formulated at the task level and remain geometry agnostic: they do not distinguish between different parameter subspaces or relate sharpness to the PAC-Bayesian complexity.
% Recent studies connect PAC-Bayes and flat minima by interpreting sharpness as a surrogate for posterior sensitivity~\citep{NIPS2017_10ce03a1,jiang2019fantasticgeneralizationmeasures,mllenhoff2023sam}, which yields bounds for flat solutions in full finetuning. 


% However, these analyses treat all parameters symmetrically and do not address adapter-based PECL, where the is constrained to an update subspace. Moreover, continual learning is intrinsically history-dependent: each task is trained on top of the parameters obtained from previous tasks, so the effective task priors evolve over time rather than arising as independent stationary draws from a fixed hyperprior. 

% This dependence is further amplified in LoRA-based PECL, where both the adapter organization and the initialization of newly introduced task-specific factors are determined by the current model state inherited from the past. This motivates our pathwise hierarchical PAC-Bayes analysis, which models task priors as history-measurable random variables drawn from a time-varying hyperposterior

% However, these analyses treat all parameters symmetrically and do not address adapter-based PECL, where the posterior is constrained to live in a subspace. Continual training is inherently history-dependent: each task is learned on top of the parameters produced by previous tasks, so the induced task priors evolve nonstationarily rather than as independent stationary draws from a fixed hyperprior. In practice,  the adapter organization and the initialization of new task also depend on the previous result.
% This motivates our pathwise hierarchical PAC-Bayes analysis, which treats task priors as history-measurable random variables drawn from a time-varying hyperposterior and introduces an explicit drift penalty that quantifies the accumulation of inductive-bias changes across tasks.

% These gaps are consequential in PECL. Beyond restricting the posterior to an adapter subspace, practical LoRA based continual learning often change the adapter organization and inilitzation  across tasks.
% Also, task priors evolve in a history-dependent and nonstationary manner,  each new task is trained on top of the parameters learned from previous tasks, rather than arising as independent stationary draws from a fixed hyperprior.
% %over tasks and introduces newly activated low rank factors that must be initialized.
% %, for example via  history conditioned reuse or Kaiming style
% % This induces a task dependent and nonstationary mechanism for generating task priors, rather than a single stationary draw from a fixed hyperprior. at the task level, the explicit cost of such sequential prior evolution.
% This motivates our pathwise hierarchical PAC Bayes analysis, which models task priors as history measurable random variables drawn from a time varying hyperposterior and yields an explicit hyperposterior drift term that quantifies the accumulation of inductive bias changes across tasks.



% \paragraph{Our position relative to prior PAC-Bayesian work.}


% \paragraph{Our position relative to prior PAC-Bayesian work.}
% Compared to the hierarchical bound of \citet{pentinaPACBayesianBoundLifelong2014}, our analysis makes two structural changes that are essential for LoRA-based PECL. First, we introduce a task-level PAC-Bayesian bound with a time-varying hyper-posterior that is adapted to the task filtration. This removes the independence assumption on task priors and yields a complexity term that decomposes into (i) task-wise KL divergences and (ii) a hyper-posterior drift term that explicitly measures how much the prior over adapters moves across tasks. Ignoring this modification would prevent us from isolating cross-task interference as a separate, quantifiable source of complexity in continual learning.

% Second, we specialize the bound to PECL with a frozen backbone and a shared LoRA subspace. By restricting the posterior support to the adapter subspace $\mathcal W_\Delta$, we show that the full-parameter KL and empirical terms can be rewritten in terms of (i) a KL over adapter updates and (ii) an expected sharpness term that only depends on perturbations within $\mathcal W_\Delta$. This subspace-aware formulation is absent from previous PAC-Bayesian work and is crucial for our main question. Without it, one cannot formally distinguish between global and adapter-subspace flatness, nor explain why LoRA-only sharpness-aware optimization can be beneficial in PECL while full-model perturbations may be systematically harmful. Our experiments later validate these theoretical insights by contrasting local and global perturbations under identical PECL protocols.







\section{Conclusion}

We develop a sequential hierarchical PAC-Bayes framework for frozen-backbone LoRA-based PECL, showing that both the sharpness-dependent empirical term and the KL complexity are naturally supported on the task-admissible adapter update subspace $\mathcal W_{\Delta,t}$. This identifies \emph{adapter-subspace sharpness} as the relevant notion of flatness in PECL.
% We formalize perturbation design in LoRA-based PECL via a sequential hierarchical PAC-Bayes bound, showing that both sharpness and complexity are naturally supported on the task-admissible adapter subspace.
% Our core contribution is a unified perturbation-based perspective for LoRA-based CL, supported by a sequential hierarchical PAC-Bayes analysis. The resulting bound shows that both the sharpness-dependent empirical term and the KL complexity are naturally supported on the current task-admissible adapter update subspace, establishing adapter-subspace sharpness as the relevant notion of flatness in CL.
In LoRA-Based CL, sharpness should be defined on the task-admissible LoRA subspace, while the perturbation radius has an effective operating range and the perturbation type modulates the strength of subspace-aligned flatness control and its associated performance gains.

% The key takeaway is simple: perturbations should be defined and optimized within the trainable LoRA subspace, while the perturbation type primarily modulates the strength of flatness control and the attainable gains; LoRA organization further shapes these gains by altering the admissible subspace and its exposed curvature.

% We study sharpness through the geometry of LoRA update subspaces. Using a sequential hierarchical PAC-Bayes analysis under a frozen backbone, we derive a bound whose KL complexity and sharpness term both reduce to adapter-subspace quantities, identifying subspace sharpness as the relevant notion of flatness in PECL. 
% % We further relate this sharpness to a LoRA-restricted Fisher and feature stability, clarifying stability-plasticity trade-offs across LoRA organizations.
% % Our theory applies to PECL and does not cover partial backbone fine-tuning. Extending the analysis to hybrid update regimes, and characterizing the interaction between global updates and subspace geometry, is left for future work. 
% Within scope, the results suggest simple guidelines: enforce flatness in the trainable adapter subspace, and view adapter organization as a structural prior shaping curvature and interference.


% We study PECL through the geometry of LoRA update subspaces. Within a sequential hierarchical PAC-Bayes framework, we derive a generalization bound for PECL under a frozen backbone and adapter-supported posteriors. The bound reduces both the KL complexity and the sharpness-dependent empirical term to adapter-subspace quantities, thereby identifying subspace sharpness as the relevant notion of flatness in PECL. Our analysis further connects subspace sharpness to the spectrum of a LoRA-restricted Fisher matrix and to feature stability, providing a principled lens for understanding stability--plasticity trade-offs across LoRA organizations.

% Our theory is specific to the strict PECL regime with a fixed backbone and does not cover settings with partial or periodic backbone fine-tuning. Extending the analysis to hybrid update regimes, and characterizing how global updates interact with subspace geometry, is an important direction for future work. Within its scope, our results suggest practical guidelines: apply flatness control to the trainable adapter subspace rather than through full-space perturbations, interpret adapter organization as a structural prior that shapes curvature and interference.



% This work revisits PECL through the geometry of LoRA adapter subspaces. Within a hierarchical PAC Bayesian framework, we derive a generalization bound.
% % in which the continual learning risk is controlled by an empirical term under Gaussian perturbations, the divergence, and the hyper posterior drift across tasks. 
% By restricting updates to the LoRA manifold $W_\Delta$ and working in factor coordinates, the analysis identifies subspace sharpness as the relevant notion of flatness and links it to the spectrum of a LoRA restricted Fisher matrix and to the stability of feature representations.


% The framework is aligned with a PECL regime and does not cover partial finetuning of the backbone. Extending the analysis to such settings is a direction for future work and may reveal additional interactions between subspace geometry and global updates. Even within its current scope, the results suggest several practical guidelines for LoRA based PECL: flatness control should primarily act in the trainable adapter subspace, the organisation of adapters can be viewed as a structural prior that shapes the Fisher spectrum and stability plasticity trade offs, and subspace sharpness together with Fisher based diagnostics provides a useful lens for interpreting and comparing PECL methods.

\bibliography{example_paper}
\bibliographystyle{icml2025}

\clearpage      % 换页，避免和正文混排
\onecolumn      % 从这里开始改为单栏
\appendix       % 如果你需要正式标记为附录，可加这一句

\section{Appendix Proof}



\subsection{Proof of theorem \ref{thm:pathwise_pac}}
\label{proof:theorem all}

\begin{theorem}[Pathwise PAC--Bayes with time-varying hyper-posterior]
\label{the: full}
For any $\delta\in(0,1]$, with probability at least $1-\delta$ over $S_{1:T}$,
\begin{equation}
\begin{aligned} 
& \frac{1}{T} \sum_{t=1}^T \underbrace{\mathbb{E}_{P_t \sim \mathcal{P}_t} \mathbb{E}_{W \sim Q_t(\cdot \mid P_t, S_t)} L_{\mathcal{D}_t}(W)}_{\text{test risk at step } t} 
\le \ \frac{1}{T} \sum_{t=1}^T \underbrace{\mathbb{E}_{P_t \sim \mathcal{P}_t} \mathbb{E}_{W \sim Q_t(\cdot \mid P_t, S_t)} \hat{L}_{S_t}(W)}_{\text{empirical risk at step } t} \\ 
& + \sqrt{\frac{ \sum_{t=1}^T \mathbb{E}_{P_t \sim \mathcal{P}_t} \mathrm{KL}\big(Q_t(\cdot \mid P_t, S_t) \,\|\, P_t\big) + \sum_{t=1}^T \mathrm{KL}(\mathcal{P}_t \,\|\, \mathcal{P}_{t-1}) + \log\frac{1}{\delta} }{2 \sum_{t=1}^T (m_t - 1)}}. 
\end{aligned} 
\end{equation}
\end{theorem}


\begin{proof}
    


% \subsection*{A.1\quad Two-level variational inequalities}
Let $B_T:=\sum_{t=1}^T (m_t-1)$ and define $\alpha_t:=(m_t-1)/B_T$ so that $\sum_{t=1}^T\alpha_t=1$.
% \paragraph{Goal quantity.}
Define the step-$t$ generalization gap
\[
\Delta_t
:=\mathbb E_{P_t\sim\mathcal P_{t-1}}\,
\mathbb E_{W\sim Q_t(\cdot\mid P_t,S_t)}
\big(L_{\mathcal D_t}(W)-\hat L_{S_t}(W)\big),
\]
and the size-weighted gap $\bar\Delta:=\sum_{t=1}^T \alpha_t \Delta_t$.
Fix $\lambda>0$ and set
\(
g_t(W):=\lambda\alpha_t\big(L_{\mathcal D_t}(W)-\hat L_{S_t}(W)\big).
\)

\paragraph{Step 1 Model level.}
For any fixed $P$ and $S_t$,
\begin{equation}
\label{eq:dv-w}
\mathbb E_{W\sim Q_t(\cdot\mid P,S_t)} g_t(W)
\;\le\;
\mathrm{KL}\!\big(Q_t(\cdot\mid P,S_t)\,\|\,P\big)
+\log\mathbb E_{W\sim P}\!\big[e^{g_t(W)}\big].
\end{equation}
Taking $\mathbb E_{S_t}$ and using the standard Hoeffding--MGF bound for bounded losses,
\begin{equation}
\label{eq:mgf-mean-gap}
\mathbb E_{S_t}\log\mathbb E_{W\sim P} e^{g_t(W)}
\;\le\; \frac{(\lambda\alpha_t)^2}{8\,(m_t-1)}.
\end{equation}
Therefore,
\begin{equation}
\label{eq:model-step}
\mathbb E_{S_t}\mathbb E_{W\sim Q_t(\cdot\mid P,S_t)} g_t(W)
\;\le\;
\mathbb E_{S_t}\mathrm{KL}\!\big(Q_t(\cdot\mid P,S_t)\,\|\,P\big)
+\frac{(\lambda\alpha_t)^2}{8\,(m_t-1)}.
\end{equation}
 % (change of measure in $P$-space)
\paragraph{Step 2 Hyper level.} 
Define
\[
h_t(P):=\mathbb E_{S_t}\mathbb E_{W\sim Q_t(\cdot\mid P,S_t)} g_t(W)
-\mathbb E_{S_t}\mathrm{KL}\!\big(Q_t(\cdot\mid P,S_t)\,\|\,P\big).
\]
By \eqref{eq:model-step}, $h_t(P)\le \frac{(\lambda\alpha_t)^2}{8\,(m_t-1)}$ for all $P$.
Applying DV to change the measure from $\mathcal P_{t-1}$ to $\mathcal P_t$ gives
\begin{equation}
\label{eq:dv-p}
\log\mathbb E_{P\sim\mathcal P_t}\!\big[e^{h_t(P)}\big]
\;\le\; \mathrm{KL}(\mathcal P_t\|\mathcal P_{t-1})
+\log\mathbb E_{P\sim\mathcal P_{t-1}}\!\big[e^{h_t(P)}\big]
\;\le\; \mathrm{KL}(\mathcal P_t\|\mathcal P_{t-1})
+\frac{(\lambda\alpha_t)^2}{8\,(m_t-1)}.
\end{equation}
Using $\log\mathbb E[e^X]\ge \mathbb E[X]$ yields
\begin{equation}
\label{eq:p-expect}
\mathbb E_{P_t\sim\mathcal P_t} h_t(P_t)
\;\le\; \mathrm{KL}(\mathcal P_t\|\mathcal P_{t-1})
+\frac{(\lambda\alpha_t)^2}{8\,(m_t-1)}.
\end{equation}
Unfolding $h_t(P)$ and interchanging expectations,
\begin{equation}
\label{eq:two-term}
\mathbb E_{S_t}\Big[
\lambda\alpha_t\Delta_t
-\alpha_t\,\mathbb E_{P_t\sim\mathcal P_{t-1}}\mathrm{KL}\!\big(Q_t(\cdot\mid P_t,S_t)\,\|\,P_t\big)
\Big]
\;\le\;
\mathrm{KL}(\mathcal P_t\|\mathcal P_{t-1})
+\frac{(\lambda\alpha_t)^2}{8\,(m_t-1)}.
\end{equation}

% \subsection*{A.2\quad Conditional MGF and supermartingale}
\paragraph{Step 3 Martingale}
From \eqref{eq:two-term} and the tower property, one obtains the conditional exponential bound
\begin{equation}
\label{eq:cond-mgf}
\mathbb E\!\left[
\exp\!\Big(
\lambda\alpha_t\Delta_t
-\alpha_t\,\mathbb E_{P_t\sim\mathcal P_{t-1}}\mathrm{KL}(Q_t\|P_t)
-\mathrm{KL}(\mathcal P_t\|\mathcal P_{t-1})
-\frac{(\lambda\alpha_t)^2}{8\,(m_t-1)}
\Big)\,\Big|\,\mathcal F_{t-1}\right]\;\le\;1.
\end{equation}
Define $X_t$ to be the exponential increment inside \eqref{eq:cond-mgf} and
\[
\mathcal{M}_0:=1,\qquad
\mathcal{M}_t:=\prod_{s=1}^t X_s.
\]
Then $(\mathcal{M}_t)_{t\ge0}$ is a nonnegative supermartingale, so
$\mathbb E[\mathcal{M}_T]\le 1$.
By Markov's inequality, with probability at least $1-\delta$,
\begin{equation}
\label{eq:hpp}
\sum_{t=1}^T \lambda\alpha_t\Delta_t
\;\le\;
\sum_{t=1}^T \alpha_t\,\mathbb E_{P_t\sim\mathcal P_{t-1}}\mathrm{KL}(Q_t\|P_t)
\;+\; \sum_{t=1}^T \mathrm{KL}(\mathcal P_t\|\mathcal P_{t-1})
\;+\; \frac{\lambda^2}{8}\sum_{t=1}^T \frac{\alpha_t^2}{(m_t-1)}
\;+\; \log\tfrac1\delta.
\end{equation}

% \subsection*{A.3\quad Size-weighted self-normalization and optimization}

By $\alpha_t=(m_t-1)/B_T$,
\[
\sum_{t=1}^T \frac{\alpha_t^2}{(m_t-1)}
=\sum_{t=1}^T \frac{(m_t-1)^2/B_T^2}{(m_t-1)}
=\frac{1}{B_T}.
\]
Hence, \eqref{eq:hpp} reads, for all $\lambda>0$ and with probability at least $1-\delta$,
\begin{equation}
\label{eq:gap-preopt}
\bar\Delta
\;\le\;
\frac{1}{\lambda}\Big(
\sum_{t=1}^T \alpha_t\,\mathbb E_{P_t\sim\mathcal P_{t-1}}\mathrm{KL}(Q_t\|P_t)
+ \sum_{t=1}^T \mathrm{KL}(\mathcal P_t\|\mathcal P_{t-1})
+ \log\tfrac1\delta\Big)
\;+\;
\frac{\lambda}{8B_T}.
\end{equation}
Optimizing the right-hand side over $\lambda>0$ (take $\lambda^\star=\sqrt{8B_T\,A}$ with
$A:=\sum_{t=1}^T \alpha_t\,\mathbb E_{P_t\sim\mathcal P_{t-1}}\mathrm{KL}(Q_t\|P_t)
+ \sum_{t=1}^T \mathrm{KL}(\mathcal P_t\|\mathcal P_{t-1}) + \log\tfrac1\delta$) yields
\begin{equation}
\label{eq:weighted-final}
\bar\Delta
\;\le\;
\sqrt{\frac{A}{2B_T}}
\;\le\;
\sqrt{\frac{
\sum_{t=1}^T \mathbb E_{P_t\sim\mathcal P_{t-1}}\mathrm{KL}(Q_t\|P_t)
+ \sum_{t=1}^T \mathrm{KL}(\mathcal P_t\|\mathcal P_{t-1})
+ \log\tfrac1\delta}{2\,\sum_{t=1}^T (m_t-1)}}.
\end{equation}
The last inequality uses $\sum_t \alpha_t x_t \le \sum_t x_t$ for nonnegative $(x_t)$.

\paragraph{Step 4 From weighted to the statement of Theorem~\ref{the: full}.}
The theorem controls the unweighted average $\frac1T\sum_{t=1}^T \Delta_t$ and uses the smoothed empirical risk on the right-hand side.
When $m_t$ are homogeneous, $\alpha_t\equiv 1/T$ and $\bar\Delta=\frac1T\sum_t\Delta_t$.
For heterogeneous $m_t$, we keep the same right-hand side and upper-bound the numerator $\sum_t \alpha_t \mathbb E_{P_t}\mathrm{KL}(Q_t\|P_t)$ by $\sum_t \mathbb E_{P_t}\mathrm{KL}(Q_t\|P_t)$, which yields exactly the denominator $2\sum_t(m_t-1)$ in the theorem.

% \subsection*{A.4\quad Gaussian smoothing (LoRA) corollary}

If $Q_t(\cdot\mid P_t,S_t)=\mathcal N(\hat W_t,\Sigma_t)$ (e.g., supported on an adapter subspace), then
\[
\mathbb E_{W\sim Q_t(\cdot\mid P_t,S_t)}\hat L_{S_t}(W)
=\mathbb E_{\varepsilon\sim\mathcal N(0,\Sigma_t)}\hat L_{S_t}(\hat W_t+\varepsilon),
\]
which converts the empirical term into the smoothed empirical loss (expected sharpness) used in the main text; substituting this identity into \eqref{eq:weighted-final} yields the empirical term in Theorem~\ref{the: full}.

Theorem~A.1 provides a tight complexity penalty of the form.
Using $\sqrt{x+y+z}\le \sqrt{x}+\sqrt{y}+\sqrt{z}$ for $x,y,z\ge 0$ yields the decomposed bound, which is the form stated in Theorem~\ref{thm:pathwise_pac} for interpretability.

\end{proof}



\subsection{Proof of Lemma~\ref{lem:kl-decomp-pecl}}
\label{proof:Lemma1}

Let $\mathcal W_\Delta\subseteq\mathbb R^d$ be a linear subspace and let $W+\mathcal W_\Delta$ be its affine translate. 
Equip both sets with the trace Borel $\sigma$-algebra induced from $\mathbb R^d$. 
$\mathcal B(\mathcal W_\Delta)\ :=\{B\cap \mathcal W_\Delta:\ B\in \mathcal B(\mathbb R^d)\},\quad
\mathcal B(W+\mathcal W_\Delta)\ :=\{B\cap (W+\mathcal W_\Delta):\ B\in \mathcal B(\mathbb R^d)\}.$
These coincide with the Borel $\sigma$-algebras generated by the subspace topologies on $\mathcal W_\Delta$ and $W+\mathcal W_\Delta$, respectively, so both are standard Borel spaces.
To formalize the probabilistic relationship between the LoRA subspace and the full model space, we define the translation map:
% \begin{align*}
$T_W: (\mathcal W_\Delta,\mathcal B(\mathcal W_\Delta)) \to (W+\mathcal W_\Delta,\mathcal B(W+\mathcal W_\Delta)), 
T_W(\Delta)=W+\Delta,$
% \end{align*}
which is a Borel isomorphism between the subspace and affine space.  The full-space prior and posterior are defined as pushforwards onto the affine subspace by
$
P^{\mathrm{full}}:=(T_W)_{\text{\#}} P^{\Delta},
Q^{\mathrm{full}}:=(T_W)_{\text{\#}} Q^{\Delta},
$
where $P_t^{\Delta}$ and $Q_t^{\Delta}$ are probability measures on $(\mathcal{W}\Delta, \mathcal{B}(\mathcal{W}\Delta))$.

% \begin{lemma}[KL decomposition with a frozen backbone]
% \label{equation: subspace-full}
% Let $W\in\mathbb R^d$ be fixed and let $P^{\Delta},Q^{\Delta}$ be probability measures on $(\mathcal W_\Delta,\mathcal B(\mathcal W_\Delta))$ with $Q^{\Delta}\ll P^{\Delta}$. Define their pushforwards onto the affine subspace by
% \[
% P^{\mathrm{full}}:=(T_W)_{\text{\#}} P^{\Delta},\qquad
% Q^{\mathrm{full}}:=(T_W)_{\text{\#}} Q^{\Delta}.
% \]
% Then $Q^{\mathrm{full}}\ll P^{\mathrm{full}}$ and
% \[
% \mathrm{KL}\left(Q^{\mathrm{full}}\middle|P^{\mathrm{full}}\right)
% =\mathrm{KL}\left(Q^{\Delta}\middle|P^{\Delta}\right).
% \]
% Equivalently, identifying ${W}\times\mathcal W_\Delta$ with $W+\mathcal W_\Delta$ via $\Phi$, one may write
% $P^{\mathrm{full}}=\Phi_{\text{\#}}(\delta_W\otimes P^{\Delta})$ and
% $Q^{\mathrm{full}}=\Phi_{\text{\#}}(\delta_W\otimes Q^{\Delta})$, and the same equality holds.
% \end{lemma}

\begin{lemma}[KL in PECL is adapter-subspace KL]
Assume $Q_t^\Delta \ll P_t^\Delta$. Then $Q_t^{\mathrm{full}} \ll P_t^{\mathrm{full}}$ and
$$
\mathrm{KL}\bigl(Q_t^{\mathrm{full}} \,\Vert\, P_t^{\mathrm{full}}\bigr)
\;=\;
\mathrm{KL}\bigl(Q_t^\Delta \,\Vert\, P_t^\Delta\bigr).
$$
\end{lemma}

\begin{proof}
  Absolute continuity is preserved by pushforward under measurable maps, hence $Q^{\mathrm{full}}\ll P^{\mathrm{full}}$. Since $T_W$ is a Borel isomorphism, by the Radon-Nikodym chain rule
 \[
 \frac{dQ^{\mathrm{full}}}{dP^{\mathrm{full}}}(T_W\Delta)=\frac{dQ^{\Delta}}{dP^{\Delta}}(\Delta)\quad\text{for }P^{\Delta}\text{-a.e.\ }\Delta.
 \]

\begin{equation*}
     \begin{aligned}
 \mathrm{KL}\left(Q^{\mathrm{full}}\middle|P^{\mathrm{full}}\right)
 &= \int_{W+\mathcal W_\Delta}
 \log\Big(\frac{dQ^{\mathrm{full}}}{dP^{\mathrm{full}}}(y)\Big)
 Q^{\mathrm{full}}(dy)\\
 &= \int_{\mathcal W_\Delta}
 \log\Big(\frac{dQ^{\mathrm{full}}}{dP^{\mathrm{full}}}\big(T_W\Delta\big)\Big)
 Q^{\Delta}(d\Delta) \quad \text{(} Q^{\mathrm{full}}=(T_W)_{\text{\#}}Q^{\Delta}\text{)} 
 \\
 &= \int_{\mathcal W_\Delta}
 \log\Big(\frac{dQ^{\Delta}}{dP^{\Delta}}(\Delta)\Big)
 Q^{\Delta}(d\Delta) \\
 &= \mathrm{KL}\left(Q^{\Delta}\middle|P^{\Delta}\right).
 \end{aligned}
\end{equation*} 
\end{proof}

% \section{ Proof of Lemma~\ref{lemma: subspace-full}}
% \begin{lemma}[Sharpness decomposition in PECL]
% \label{lemma: Sharpness in PECL}
% Let $W_0\in\mathbb R^d$ be fixed and let the trainable update lie in the adapter subspace $\mathcal W_\Delta\subseteq\mathbb R^d$. Consider a Gaussian posterior on $(\mathcal W_\Delta,\mathcal B(\mathcal W_\Delta))$,
%  $
%  Q_t^{\Delta}=\mathcal N(\hat{\Delta W}_t,\Sigma_t),
%  $
%  with mean $\hat{\Delta W}_t\in\mathcal W_\Delta$ and covariance operator $\Sigma_t\succeq 0$ acting on $\mathcal W_\Delta$ (viewed in $\mathbb R^d$, this covariance is supported in $\mathcal W_\Delta$). Define the full-space posterior on ${W_0}\times\mathcal W_\Delta$ by
%  $
%  Q_t^{\mathrm{full}}:=\delta_{W_0}\otimes Q_t^{\Delta},
%  $
% and identify ${W_0}\times\mathcal W_\Delta$ with the affine subspace $W_0+\mathcal W_\Delta$ via $(w,\Delta)\mapsto w+\Delta$. Let $\hat L_{S_t}:\mathbb R^d\to\mathbb R$ be measurable and integrable under $Q_t^{\mathrm{full}}$. Then
%  \[
%  \mathbb{E}_{W\sim Q_t^{\mathrm{full}}}\big[\hat{L}_{S_t}(W)\big]
%  = \mathbb{E}_{\varepsilon\sim\mathcal N(0,\Sigma_t)}
%  \big[\hat{L}_{S_t}(W_0+\hat{\Delta W}_t+\varepsilon)\big]
%  = \hat{L}_{S_t}(W_0+\hat{\Delta W}_t) +\mathrm{E\text{-}Sh}_t^{\Delta}(\hat{\Delta W}_t,\Sigma_t) \]
%    where the subspace expected sharpness is
%    \[
%    \mathrm{E\text{-}Sh}_t^{\Delta}(\hat{\Delta W}_t,\Sigma_t)
%    = \mathbb{E}_{\varepsilon\sim\mathcal N(0,\Sigma_t)}
%    \big[\hat{L}_{S_t}(W_0+\hat{\Delta W}_t+\varepsilon)
%   - \hat{L}_{S_t}(W_0+\hat{\Delta W}_t)\big].
%      \]
% \end{lemma}
% \begin{proof}
%     Sampling $W\sim Q_t^{\mathrm{full}}$ is equivalent to sampling $\varepsilon\sim\mathcal N(0,\Sigma_t)$ on $\mathcal W_\Delta$ and setting $W=W_0+\hat{\Delta W}_t+\varepsilon$. Taking expectations and adding-subtracting $\hat L_{S_t}(W_0+\hat{\Delta W}_t)$ inside the expectation gives the decomposition. 
% \end{proof}

\subsection{Proof of Lemma~\ref{lem:subspace-loss-pecl}}
\label{proof:lemma2}

% \begin{equation*}
% \label{eq:subspace-loss-pecl}
% \begin{aligned}
% &\mathbb E_{W\sim Q_t^{\mathrm{full}}}\bigl[\hat L_{S_t}(W)\bigr] \\
% &= \mathbb{E}_{\varepsilon\sim\mathcal N(0,\Sigma_t)}
%     \big[\hat{L}_{S_t}(W_\mathrm {frozen }+\hat{\Delta W}_t+\varepsilon)\big],\varepsilon\in\mathcal W_\Delta \\
% & =\hat L_{S_t}(W_{\mathrm{frozen}} + \hat{\Delta W}_t)
% +
% \mathrm{E\text{-}Sh}_t^\Delta(\hat{\Delta W}_t,\Sigma_t),\varepsilon\in\mathcal W_\Delta.
% \end{aligned}
% \end{equation*}
We specialise the posterior to a Gaussian supported on $\mathcal W_\Delta$. For task $t$ We consider a Gaussian posterior \( Q_t^{\Delta} = \mathcal{N}(\hat{\Delta W}_t, \Sigma_t) \) on the measurable space \( (\mathcal{W}_\Delta, \mathcal{B}(\mathcal{W}_\Delta)) \), where \( \hat{\Delta W}_t \in \mathcal{W}_\Delta \) is the mean and \( \Sigma_t\) is a covariance on the adapter subspace \( \mathcal{W}_\Delta \). Sampling from \( \Delta W \sim Q_t^{\Delta} \) is equivalent to 
$
\Delta W = \hat{\Delta W}_t + \varepsilon, \quad \varepsilon \sim \mathcal{N}(0, \Sigma_t) \ \text{on } \mathcal{W}_\Delta.
$ 

\begin{lemma}
    The posterior expected empirical loss decomposes as
\begin{equation*}
    \mathbb{E}_{W\sim Q_t^{\mathrm{full}}}\big[\hat{L}_{S_t}(W)\big]
= \mathbb{E}_{\varepsilon\sim\mathcal N(0,\Sigma_t)}
    \big[\hat{L}_{S_t}(W_\mathrm {frozen }+\hat{\Delta W}_t+\varepsilon)\big]
= \hat{L}_{S_t}(W_\mathrm {frozen }+\hat{\Delta W}_t)
  + \mathrm{E\text{-}Sh}_t^{\Delta}(\hat{\Delta W}_t,\Sigma_t),
\end{equation*}
where the expected sharpness is
$$
\mathrm{E\text{-}Sh}_t^{\Delta}(\hat{\Delta W}_t,\Sigma_t)
:= \mathbb{E}_{\varepsilon\sim\mathcal N(0,\Sigma_t)}
   \big[\hat{L}_{S_t}(W_\mathrm {frozen }+\hat{\Delta W}_t+\varepsilon)
       - \hat{L}_{S_t}(W_\mathrm {frozen }+\hat{\Delta W}_t)\big],\quad
\varepsilon\sim\mathcal N(0,\Sigma_t)\ \text{on } \mathcal W_\Delta$$
\end{lemma}


\begin{proof}
Recall that the adapter posterior is a Gaussian supported on the update subspace,
$
Q_t^\Delta = \mathcal N(\hat{\Delta W}_t,\Sigma_t)
$
on $(\mathcal W_\Delta,\mathcal B(\mathcal W_\Delta))$.
Hence we can represent a draw $\Delta W \sim Q_t^\Delta$ as
$
\Delta W = \hat{\Delta W}_t + \varepsilon
$
with $\varepsilon \sim \mathcal N(0,\Sigma_t)$ on $\mathcal W_\Delta$.

Under the PECL constraint, the full parameter is obtained by translating the adapter update by the frozen part,
$
W = W_{\mathrm{frozen}} + \Delta W.
$
Therefore, for the empirical risk,
\begin{align*}
\mathbb E_{W \sim Q_t^{\mathrm{full}}}\big[\hat L_{S_t}(W)\big]
&= \mathbb E_{\Delta W \sim Q_t^\Delta}\big[\hat L_{S_t}(W_{\mathrm{frozen}}+\Delta W)\big] \\
&= \mathbb E_{\varepsilon \sim \mathcal N(0,\Sigma_t)}
\big[\hat L_{S_t}(W_{\mathrm{frozen}}+\hat{\Delta W}_t+\varepsilon)\big].
\end{align*}
Finally, add and subtract the mean point loss to obtain
\begin{align*}
\mathbb E_{\varepsilon \sim \mathcal N(0,\Sigma_t)}
\big[\hat L_{S_t}(W_{\mathrm{frozen}}+\hat{\Delta W}_t+\varepsilon)\big]
&=
\hat L_{S_t}(W_{\mathrm{frozen}}+\hat{\Delta W}_t) \\
&\quad +
\mathbb E_{\varepsilon \sim \mathcal N(0,\Sigma_t)}
\Big[
\hat L_{S_t}(W_{\mathrm{frozen}}+\hat{\Delta W}_t+\varepsilon)
-
\hat L_{S_t}(W_{\mathrm{frozen}}+\hat{\Delta W}_t)
\Big].
\end{align*}
The second term is exactly the expected Sharpness on the subspace
\[
\mathrm{E\text{-}Sh}_t^\Delta(\hat{\Delta W}_t,\Sigma_t)
:=
\mathbb E_{\varepsilon \sim \mathcal N(0,\Sigma_t)}
\Big[
\hat L_{S_t}(W_{\mathrm{frozen}}+\hat{\Delta W}_t+\varepsilon)
-
\hat L_{S_t}(W_{\mathrm{frozen}}+\hat{\Delta W}_t)
\Big],
\]
completes the proof.
\end{proof}



\subsection{Proof of Theorem \ref{thm:pecl-bound}}
\label{proof:theorem}

% \subsection{Proof of }





\begin{theorem}[Sharpness-aware PAC--Bayes bound for PECL]
% \label{thm:pecl-bound}

For any $\delta\in(0,1]$, with probability at least $1-\delta$ over $S_1,\ldots,S_T$,
\begin{equation}\label{eq:pecl-main}
\begin{aligned}
&\frac{1}{T}\sum_{t=1}^T
\mathbb{E}_{P_t\sim\mathcal P_{t-1}}\mathbb{E}_{W\sim Q_t}L_{\mathcal{D}_t}(W)
\ \le\
\ \frac{1}{T}\sum_{t=1}^T
\mathbb{E}_{P_t\sim\mathcal P_{t-1}}
\mathbb{E}_{\varepsilon \sim \mathcal{N}(0, \Sigma_t)}[\hat{L}_{S_t}(W_0 + \hat{\Delta W}_t + \varepsilon)
\\
&\quad +\sqrt{\frac{
\displaystyle \sum_{t=1}^T \mathbb{E}_{P_t^\Delta\sim\mathcal P_{t-1}}\!\Big[\mathrm{KL}\!\big(Q_t^\Delta\|P_t^\Delta\big)\Big]
\;+\;\displaystyle \sum_{t=1}^T \mathrm{KL}\!\big(\mathcal P_t\|\mathcal P_{t-1}\big)
\;+\;\log\!\frac{1}{\delta}
}{2\,\displaystyle\sum_{t=1}^T (m_t-1)}}.
\end{aligned}
\end{equation}
\end{theorem}



\begin{proof}
The bound follows by specializing Theorem \ref{the: full} to the PECL setting through two key decompositions:
By Lemma~\ref{lem:kl-decomp-pecl}, the full-space KL terms reduce to their adapter-space counterparts:
\[
\mathrm{KL}(Q_t^{\mathrm{full}} \| P_t^{\mathrm{full}}) = \mathrm{KL}(Q_t^\Delta \| P_t^\Delta),
\quad
\mathrm{KL}(\mathcal{P}_t^{\mathrm{full}} \| \mathcal{P}_{t-1}^{\mathrm{full}}) = \mathrm{KL}(\mathcal{P}_t^\Delta \| \mathcal{P}_{t-1}^\Delta).
\]

% \textbf{Step 2: Sharpness decomposition via Lemma 6.}  
Under the Gaussian posterior $Q_t^\Delta = \mathcal{N}(\Delta\hat{W}_t, \Sigma_t)$, Lemma \ref{lem:subspace-loss-pecl} gives:
\[
\mathbb{E}_{W \sim Q_t^{\mathrm{full}}}[\hat{L}_{S_t}(W)] 
= \mathrm{E}_{\varepsilon \sim \mathcal{N}(0, \Sigma_t)}{L}_{S_t}(W_{t-1}+\Delta\hat{W}_t + \varepsilon)
 = \hat{L}_{S_t}(W_{t-1} + \Delta\hat{W}_t) + \mathrm{E\text{-}Sh}_t^\Delta(\Delta\hat{W}_t, \Sigma_t),
\]
where the perturbations are confined to the adapter subspace.

% \textbf{Step 3: Substitution into general bound.}
Substituting these decompositions into Lemma \ref{the: full} yields the final PECL-specific bound, where all complexity and sharpness terms are now restricted to the trainable LoRA subspace $\mathcal{W}_\Delta$.
% Substituting these decompositions into Theorem 2 yields the final PECL-specific bound, where all complexity and sharpness terms are now restricted to the trainable LoRA subspace $\mathcal{W}_\Delta$.

\end{proof}

\section{Bridging the Abstract Update Subspace and the LoRA Implementation}
\label{app:lora_subspace_bridge}

This appendix clarifies how the abstract admissible update subspace
$\mathcal W_{\Delta,t}$ used in the PAC--Bayesian analysis relates to the concrete LoRA
parameterization implemented in practice.
The main point is conceptual: throughout the paper, $\mathcal W_{\Delta,t}$ denotes a
\emph{linear subspace of the ambient weight-increment space} $\mathbb R^d$ with
$d=\dim(\mathrm{vec}(\Delta W))$, rather than the rank-$r$ factor coordinate space of $(A,B)$.
This distinction is essential because the LoRA mapping $\theta_\Delta=(A,B)\mapsto \Delta W=BA$
is bilinear, hence it does not define a global linear set of reachable increments in weight space.
Our bound is local and only requires a linear model of admissible \emph{directions} around the
learned solution, which is precisely what the Jacobian provides.

\paragraph{Step 1: Abstract object in the bound: a linear increment subspace in $\Delta W$-space.}
In the main text, PECL is modeled as restricting task-$t$ updates to a linear subspace
$\mathcal W_{\Delta,t}\subseteq\mathbb R^d$, leading to the affine model class
\[
W_{\mathrm{froz}}+\mathcal W_{\Delta,t}.
\]
All priors, posteriors, and perturbation distributions appearing in
Theorem~\ref{thm:pathwise_pac} are defined on $(\mathcal W_{\Delta,t},\mathcal B(\mathcal W_{\Delta,t}))$,
and are pushed forward to the affine class by translation.
Concretely, letting $T_{W_{\mathrm{froz}}}(\Delta)=W_{\mathrm{froz}}+\Delta$, we work with
\[
P_t^{\mathrm{full}} := (T_{W_{\mathrm{froz}}})_{\#} P_t^{\Delta},
\qquad
Q_t^{\mathrm{full}} := (T_{W_{\mathrm{froz}}})_{\#} Q_t^{\Delta},
\]
where $P_t^\Delta$ and $Q_t^\Delta$ are measures supported on $\mathcal W_{\Delta,t}$.
Lemma~\ref{lem:kl-decomp-pecl} then shows that the complexity term satisfies
\[
\mathrm{KL}\!\left(Q_t^{\mathrm{full}}\Vert P_t^{\mathrm{full}}\right)
=
\mathrm{KL}\!\left(Q_t^{\Delta}\Vert P_t^{\Delta}\right),
\]
so in PECL the KL contribution is entirely determined by the adapter-subspace distributions.

\paragraph{Step 2: Implemented object: LoRA coordinates and a local linearization in increment space.}
For each adapted weight block (or layer) $\ell$, LoRA introduces trainable coordinates
$\theta_{\Delta,\ell}=(A_\ell,B_\ell)$ and defines the effective increment
\[
\Delta W_\ell(\theta_{\Delta,\ell}) = B_\ell A_\ell ,
\qquad
\theta_{\Delta,\ell}\in\mathbb R^{r\times n_\ell}\times\mathbb R^{m_\ell\times r}.
\]
Let $\hat\theta_{\Delta,t}$ denote the learned LoRA coordinates after training on task $t$,
and consider a small coordinate perturbation $\delta\theta_\Delta$ around $\hat\theta_{\Delta,t}$.
Define the mapping into the ambient increment vector space by
\[
g(\theta_\Delta) := \mathrm{vec}(\Delta W(\theta_\Delta))\in\mathbb R^d.
\]
A first-order expansion yields
\begin{equation}
\label{eq:lora_local_linearization}
g(\hat\theta_{\Delta,t}+\delta\theta_\Delta)
=
g(\hat\theta_{\Delta,t})
+
J_{\mathrm{LoRA}}(\hat\theta_{\Delta,t})\,\delta\theta_\Delta
+
R(\delta\theta_\Delta),
\end{equation}
where $J_{\mathrm{LoRA}}(\hat\theta_{\Delta,t})$ is the Jacobian of $g$ evaluated at $\hat\theta_{\Delta,t}$,
and the remainder satisfies the standard second-order bound
\[
\|R(\delta\theta_\Delta)\|_2 = O(\|\delta\theta_\Delta\|_2^2)
\quad\text{as}\quad
\|\delta\theta_\Delta\|_2\to 0.
\]
Equation~\eqref{eq:lora_local_linearization} makes explicit that LoRA induces a linear space of admissible
\emph{increment directions} through the Jacobian image at the learned coordinates:
\begin{equation}
\label{eq:WDelta_Jacobian_image_new}
\mathcal W_{\Delta}(\hat\theta_{\Delta,t})
:=
\mathrm{Im}\!\big(J_{\mathrm{LoRA}}(\hat\theta_{\Delta,t})\big)
\subseteq \mathbb R^d .
\end{equation}
This object is a \emph{local tangent subspace in increment space}, not a global reachable set.

\paragraph{Step 3: Identification used by the theory: $\mathcal W_{\Delta,t}$ as a local tangent model.}
The abstract subspace $\mathcal W_{\Delta,t}$ in Theorem~\ref{thm:pathwise_pac} should be interpreted as
a local linear model of admissible update \emph{directions} around the task-$t$ solution in the ambient
increment space. Accordingly, we identify
\begin{equation}
\label{eq:WDelta_t_identification}
\mathcal W_{\Delta,t}
\equiv
\mathcal W_{\Delta}(\hat\theta_{\Delta,t})
=
\mathrm{Im}\!\big(J_{\mathrm{LoRA}}(\hat\theta_{\Delta,t})\big),
\end{equation}
and define the PAC--Bayesian distributions directly on this linear subspace, as required by
Lemma~\ref{lem:kl-decomp-pecl} and Theorem~\ref{thm:pathwise_pac}.

\paragraph{Step 4: How coordinate-space perturbations induce subspace perturbations in the bound.}
In implementation, perturbations are naturally applied in the LoRA coordinates.
Under the local linearization~\eqref{eq:lora_local_linearization}, a coordinate perturbation
$\delta\theta_\Delta$ induces an increment perturbation
\[
\delta\Delta := g(\hat\theta_{\Delta,t}+\delta\theta_\Delta) - g(\hat\theta_{\Delta,t})
\approx
J_{\mathrm{LoRA}}(\hat\theta_{\Delta,t})\,\delta\theta_\Delta,
\]
whose support lies in $\mathrm{Im}(J_{\mathrm{LoRA}}(\hat\theta_{\Delta,t}))=\mathcal W_{\Delta,t}$.
For instance, if $\delta\theta_\Delta\sim\mathcal N(0,\Sigma_{\theta,t})$ and the perturbation radius
is chosen such that the second-order remainder $R(\delta\theta_\Delta)$ is negligible with high probability,
then the induced increment perturbation is approximately Gaussian in the ambient space,
\begin{equation}
\label{eq:Sigma_pushforward}
\delta\Delta \approx \mathcal N(0,\Sigma_t),
\qquad
\Sigma_t := J_{\mathrm{LoRA}}(\hat\theta_{\Delta,t})\,\Sigma_{\theta,t}\,
J_{\mathrm{LoRA}}(\hat\theta_{\Delta,t})^\top,
\end{equation}
and is supported on $\mathcal W_{\Delta,t}$.
This provides an implementation-level meaning of the perturbation geometry $\Sigma_t$ in the bound:
it is the covariance of the \emph{induced} increment perturbation in the ambient space,
obtained (to first order) by pushing forward coordinate noise through the Jacobian.

\paragraph{Takeaway.}
The bound is formulated on $\mathcal W_{\Delta,t}\subset\mathbb R^d$ because it is a learning-theoretic
statement about perturbations and posterior complexity \emph{in weight-increment space} under PECL constraints.
LoRA is one concrete mechanism that realizes such a constraint.
Due to bilinearity, LoRA does not yield a global linear increment class, but it does induce a local
tangent subspace of admissible directions around the learned solution.
Equation~\eqref{eq:WDelta_t_identification} explains why the abstract update scope $\mathcal W_{\Delta,t}$
can be consistently interpreted and instantiated from the LoRA implementation while keeping the KL analysis
of Lemma~\ref{lem:kl-decomp-pecl} intact.


% \section{Bridging the Abstract Update Subspace and the LoRA Implementation}
% \label{app:lora_subspace_bridge}

% This appendix clarifies the relation between the abstract admissible update subspace
% $\mathcal W_{\Delta,t}$ used in the PAC--Bayesian analysis and the concrete LoRA parameterization
% implemented in practice. The key point is that LoRA does not define a global linear set of reachable
% increments in weight space, since $\Delta W(\theta_\Delta)=BA$ is bilinear. However, the bound terms
% are local and rely on perturbations around the learned solution. In this local regime, LoRA induces a
% linear space of admissible increment \emph{directions} through the Jacobian of the mapping
% $\theta_\Delta\mapsto \Delta W(\theta_\Delta)$.

% % \paragraph{Abstract object: admissible increment subspace.}
% In the main text we represent PECL as restricting task-$t$ updates to a linear subspace
% $\mathcal W_{\Delta,t}\subseteq\mathbb R^d$, leading to the affine model class
% $W_{\mathrm{froz}}+\mathcal W_{\Delta,t}$.
% All posteriors, priors, and perturbation distributions in Theorem~\ref{thm:pathwise_pac}
% are defined on $\mathcal W_{\Delta,t}$ and then pushed forward to the affine class by translation.

% % \paragraph{Implemented object: LoRA coordinates and induced increments.}
% For each adapted weight block (or layer) $\ell$, LoRA introduces trainable coordinates
% $\theta_{\Delta,\ell}=(A_\ell,B_\ell)$ and defines the effective increment
% \[
% \Delta W_\ell(\theta_{\Delta,\ell}) = B_\ell A_\ell .
% \]
% Let $\hat\theta_{\Delta,t}$ denote the learned LoRA coordinates after training on task $t$,
% and consider a small coordinate perturbation $\delta\theta_\Delta$ around $\hat\theta_{\Delta,t}$.
% By first-order expansion, the induced increment variation in the ambient parameter space satisfies
% \[
% \mathrm{vec}\!\big(\delta(\Delta W)\big)
% \approx
% J_{\mathrm{LoRA}}(\hat\theta_{\Delta,t})\,\delta\theta_\Delta,
% \]
% where $J_{\mathrm{LoRA}}(\hat\theta_{\Delta,t})$ is the Jacobian of the mapping
% $\theta_\Delta\mapsto \mathrm{vec}(\Delta W(\theta_\Delta))$ evaluated at $\hat\theta_{\Delta,t}$.
% Therefore, LoRA induces a linear space of admissible increment \emph{directions} given by the Jacobian image
% \begin{equation}
% \label{eq:WDelta_Jacobian_image}
% \mathcal W_{\Delta}(\hat\theta_{\Delta,t})
% =
% \mathrm{Im}\!\big(J_{\mathrm{LoRA}}(\hat\theta_{\Delta,t})\big)
% \subseteq \mathbb R^d .
% \end{equation}

% % \paragraph{Identification with the abstract subspace.}
% The abstract subspace $\mathcal W_{\Delta,t}$ in Theorem~\ref{thm:pathwise_pac} should be interpreted as a
% local linear model of admissible update \emph{directions} around the task-$t$ solution. Concretely, we
% identify it with the LoRA-induced Jacobian image at the learned coordinates:
% \begin{equation}
% \label{eq:WDelta_t_approx}
% \mathcal W_{\Delta,t}
% \approx
% \mathrm{Im}\!\big(J_{\mathrm{LoRA}}(\hat\theta_{\Delta,t})\big)
% =
% \mathcal W_{\Delta}(\hat\theta_{\Delta,t}) .
% \end{equation}
% Under this identification, any perturbation distribution defined in LoRA coordinate space and applied
% locally around $\hat\theta_{\Delta,t}$ induces a perturbation distribution supported on
% $\mathcal W_{\Delta,t}$ in the ambient parameter space.
% The identification~\eqref{eq:WDelta_t_approx} provides a concrete implementation-level meaning of the
% ``admissible update scope'' in the bound: $\mathcal W_{\Delta,t}$ is the local space of weight-space
% directions that can be induced by infinitesimal changes of the trainable LoRA coordinates around the
% task-$t$ solution. This clarifies why the bound is formulated on $\mathcal W_{\Delta,t}$ while training is
% performed in the factor coordinates $\theta_\Delta$.
% % which is precisely the support assumption used
% % by the bound.

% % \paragraph{SeqLoRA vs IncLoRA: fixed coordinates versus expanding degrees of freedom.}
% % Equation~\eqref{eq:WDelta_t_approx} highlights that the evolution of $\mathcal W_{\Delta,t}$ depends on how
% % the LoRA organization evolves across tasks.

% % \begin{itemize}
% % \item \textbf{SeqLoRA (fixed coordinates, evolving tangent geometry).}
% % SeqLoRA keeps the set of trainable coordinates fixed across tasks (same adapter slots and ranks).
% % Thus, the coordinate domain of $\theta_\Delta$ does not change over $t$.
% % Nevertheless, the induced admissible direction space
% % $\mathcal W_{\Delta,t}\approx \mathrm{Im}(J_{\mathrm{LoRA}}(\hat\theta_{\Delta,t}))$
% % may still evolve because the Jacobian is evaluated at different learned coordinates $\hat\theta_{\Delta,t}$.
% % In this sense, SeqLoRA corresponds to an \emph{evolving tangent geometry under fixed coordinates}.

% % \item \textbf{IncLoRA (expanded coordinates, evolving tangent geometry).}
% % IncLoRA expands the trainable degrees of freedom over tasks, for instance by adding new branches,
% % increasing effective rank, or enlarging the set of adapted blocks.
% % As a result, the coordinate dimension of $\theta_\Delta$ changes with $t$, and the Jacobian
% % $J_{\mathrm{LoRA}}(\hat\theta_{\Delta,t})$ is defined on an expanding coordinate space.
% % Consequently, $\mathcal W_{\Delta,t}$ evolves both because it is evaluated at a changing point
% % $\hat\theta_{\Delta,t}$ and because the underlying degrees of freedom are expanded.
% % This corresponds to an \emph{evolving tangent geometry under expanding coordinates}.
% % \end{itemize}













\section{Implementation Details}


Unless stated otherwise, we use SGD as the base optimizer, train each session for $20$ epochs and SAM radius $\rho = 0.05$. LoRA based variants are trained with learning rate $0.01$ without momentum or weight decay. The same training configuration is used across SeqLoRA, IncLoRA, and OLoRA, so that differences in performance can be attributed to the geometry of the adapter subspace and the choice of flatness optimizer rather than to changes in training budget.


\section{Distinguishing the Abstract Admissible Subspace and the Implemented LoRA Parameterization}

    
\section{Supplementary results}

\subsection{CORRELATION BETWEEN W AND $\Delta$ W}






\begin{table}[thpb]
\centering
\caption{Performance under different scope.}
\label{tab:lora——or_full}
\setlength{\tabcolsep}{3pt}
\renewcommand{\arraystretch}{1}
\resizebox{1\linewidth}{!}{
\begin{tabular}{c c cc cc cc}
\toprule
\multicolumn{2}{c}{\textbf{Opt.}} &
\multicolumn{2}{c}{\textbf{SeqLoRA}} &
\multicolumn{2}{c}{\textbf{IncLoRA}} &
\multicolumn{2}{c}{\textbf{OLoRA}} \\
\cmidrule(lr){1-2} \cmidrule(lr){3-4} \cmidrule(lr){5-6} \cmidrule(lr){7-8}
\textbf{Strength} & \textbf{Scope} &
$\operatorname{AAA}^{\mathrm{cls}}$ & $\operatorname{AAA}^{\mathrm{ncm}}$ &
$\operatorname{AAA}^{\mathrm{cls}}$ & $\operatorname{AAA}^{\mathrm{ncm}}$ &
$\operatorname{AAA}^{\mathrm{cls}}$ & $\operatorname{AAA}^{\mathrm{ncm}}$ \\
\midrule
\multirow{1}{*}{$0$} & &
69.07 & 76.62 & 71.29 & 78.29 & 71.41 & 78.47 \\
\midrule
% 组 1：\rho_1
\multirow{2}{*}{$0.05$}
& $W+\Delta W$ 
&-3.5  & -1.54 & -2    & -0.38 & -2.62 & -0.8  \\ 
& ${W}_{\Delta}$ &0.27  & 0.02  & 0.79  & 0.51  & 0.12  & 0.16  \\ 
\midrule
% 组 2：\rho_2
\multirow{2}{*}{$0.10$} & $W+\Delta W$ &-4.82 & -1.29 & -4.74 & -0.74 & -5.53 & -1.18 \\
& ${W}_{\Delta}$ &0.7   & 0.29  & 0.86  & 0.53  & 0.28  & 0.12  \\
\midrule
% 组 3：\rho_3
\multirow{2}{*}{$0.15$} & $W+\Delta W$ &-6.72 & -1.69 & -6.84 & -0.98 & -6.49 & -1.29 \\
& $\Delta W$ &0.41  & 0.08  & 0.77  & 0.53  & 0.27  & 0.11 \\
\bottomrule
\end{tabular}}
\end{table}


\subsection{Detailed results on ImageNet-R}


% \begin{table}[htpb]
% \centering
% \caption{
% Class incremental performance and weight space flatness metrics on ImageNet-R ($T=20$, $R=16$). 
% The table reports standard class incremental metrics  together with three sharpness indices and empirical Hessian and Fisher statistics in the weight space.
% }
% \label{table:q1_ci_with_weight_flatness_reorg}
% \setlength{\tabcolsep}{4pt}
% \resizebox{\linewidth}{!}{
% \begin{tabular}{l l ccc ccccc c}
% \toprule
% \multirow{2}{*}{\textbf{Method}} & \multirow{2}{*}{\textbf{Opt.}} &
% \multicolumn{3}{c}{\textbf{CL ($\uparrow$)}} &
% \multicolumn{2}{c}{\textbf{Sharpness ($\downarrow$)}} &
% \multicolumn{2}{c}{\textbf{Emp Hessian ($\downarrow$)}} &
% \multicolumn{2}{c}{\textbf{Emp Fisher ($\downarrow$)}} \\
% \cmidrule(lr){3-5}\cmidrule(lr){6-7}\cmidrule(lr){8-9}\cmidrule(lr){10-11}
%  & & $\operatorname{ACC}_{20}^{\mathrm{cls}}$  &$\operatorname{AAA}^{\mathrm{cls}}$   & $\operatorname{BWT}^{\mathrm{cls}}$
%  & Sh$^{(0)}(\rho)$  & Sh$^{(1)}(\rho)$
%  & $\lambda_{\max}(H)$ & $\mathrm{tr}(H)$  
%  & $\lambda_{\max}(F)$ & $\mathrm{tr}(F)$   \\
% \midrule
% \multirow{1}{*}{\textbf{PerTaskLP}}
% & SGD   & $55.22^{\pm 0.31}$ & $59.36^{\pm 0.19}$ & $-8.35^{\pm 0.75}$ & 0.072   & 0.072 & 5.31    & 1182.49  & 8.47   & 1015.07  \\
% \midrule
% \multirow{2}{*}{\textbf{SeqFT}}
% & SGD   &$50.72^{\pm 2.20}$ & $59.89^{\pm 1.32}$ & $-38.94^{\pm 1.87}$ & 0.582  & 0.517 & 4692.68 & 25621.40 & 538.79 & 29703.66 \\
% & SAM   &$56.07^{\pm 2.60}$ & $66.47^{\pm 1.53}$ & $-33.04^{\pm 4.71}$ & 0.259  & 0.252 & 474.25  & 19747.59 & 75.43  & 5881.32  \\
%  &GAM  &$59.78^{\pm 0.74}$ & $67.31^{\pm 1.18}$ & $-32.70^{\pm 0.26}$ & 0.177  & 0.176 & 877.97  & 18879.15 & 32.00  & 3273.08  \\
% \midrule
% \multirow{2}{*}{\textbf{SeqLoRA}}
% & SGD    &$63.14^{\pm 2.17}$ & $69.07^{\pm 1.33}$ & $-11.91^{\pm 1.82}$ & 0.142  & 0.114 & 47.37 & 1502.55 & 35.51 & 1401.15 \\
% & SAM    &$67.04^{\pm 0.71}$ & $71.42^{\pm 0.71}$ & $-9.26^{\pm 0.66}$ & 0.120  & 0.093 & 24.41   & 996.59   & 21.39  & 821.45   \\
%  &GAM  &$73.31^{\pm 0.52}$ & $76.80^{\pm 0.56}$ & $-9.37^{\pm 0.61}$ & 0.060  &0.058  & 4.81    & 339.49   & 3.72   & 240.91   \\
% \midrule
% \multirow{2}{*}{\textbf{IncLoRA}}
% & SGD    &$65.77^{\pm 2.84}$ & $71.29^{\pm 0.78}$ & $-4.47^{\pm 2.85}$ & 0.076    & 0.062 & 15.46 & 969.34  & 8.89  & 836.91 \\
% & SAM    &$68.86^{\pm 0.92}$ & $72.59^{\pm 0.39}$ & $-3.51^{\pm 0.50}$ & 0.053   & 0.052 & 3.72    & 712.88   & 4.35   & 526.49   \\
%  &GAM  &$72.86^{\pm 0.23}$ & $76.52^{\pm 0.48}$ & $-6.41^{\pm 0.48}$ & 0.035   & 0.035 & 2.69    & 281.22   & 2.09   & 192.25   \\
% \midrule
% \multirow{2}{*}{\textbf{OLoRA}}
% & SGD    &$66.64^{\pm 0.44}$ & $71.41^{\pm 0.47}$ & $-3.52^{\pm 0.43}$ & 0.078   & 0.058 & 6.94    & 938.93   & 9.00   & 831.23   \\
% & SAM    &$68.53^{\pm 0.78}$ & $72.75^{\pm 0.76}$ & $-3.16^{\pm 0.68}$ & 0.066   & 0.048 & 4.40    & 775.40   & 7.41   & 628.12   \\
%  &GAM &$72.44^{\pm 0.36}$ & $76.05^{\pm 0.52}$ & $-7.12^{\pm 0.88}$ & 0.036   & 0.036 & 2.22    & 261.47   & 2.08   & 198.30\\
% \midrule
% \multirow{1}{*}{\textbf{l2p}}
% & Adam   & $71.35^{\pm 0.05}$ & $74.84^{\pm 0.73}$ & $-6.34^{\pm 0.50}$  &-  &- &- &- &- &- \\
% \midrule
% \multirow{1}{*}{\textbf{dualprompt}}
% & Adam   & $71.59^{\pm 0.68}$ & $74.80^{\pm 1.18}$ & $-3.38^{\pm 0.45}$ &-  &- &- &- &- &- \\
% \bottomrule
% \end{tabular}
% }
% \end{table}




\begin{table}[htpb]
\centering
\caption{
NME based performance and feature space diagnostics on ImageNet-R ($T=20$, $R=16$). 
The table reports NME accuracy metrics together with prototype drift, CKA similarity to the initial representation, and spectral properties of the empirical feature matrix.
}
\label{tab:q1_lp_feat_reorg}
\setlength{\tabcolsep}{4pt}
\renewcommand{\arraystretch}{0.95} % 可选：稍微压缩行距
\resizebox{0.9\linewidth}{!}{
\begin{tabular}{l l ccc  cc cc}
\toprule
\multirow{2}{*}{\textbf{Method}} & \multirow{2}{*}{\textbf{Opt.}} &
\multicolumn{3}{c}{\textbf{NME} ($\uparrow$)} &
\multicolumn{2}{c}{\textbf{Feature Stability}} &
\multicolumn{2}{c}{\textbf{Feature Matrix}} \\
\cmidrule(lr){3-5} \cmidrule(lr){6-7} \cmidrule(lr){8-9}
&  & $ \operatorname{ACC}_{20}^{\mathrm{nme}}$  & $\operatorname{AAA}^{\mathrm{nme}}$  & $\operatorname{BWT}^{\mathrm{nme}}$ 
 & L2 drift & CKA$_{0}$
 & $\mathrm{tr}(E)$ & eff \\ 
\midrule
\multirow{1}{*}{\textbf{PerTaskLP}}
  & SGD &$62.96^{\pm0.12}$ & $66.58^{\pm0.37}$ & $-7.66^{\pm1.27}$
     & 0        & 1        & 0.82 & 128.11  \\
\midrule
\multirow{2}{*}{\textbf{SeqFT}}
 & SGD     &$69.49^{\pm0.16}$ & $75.10^{\pm0.57}$ & $-16.11^{\pm0.23}$ & 29.57 & 0.78 & 0.58 & 93.53  \\
 & SAM     &$72.36^{\pm2.46}$ & $78.31^{\pm0.50}$ & $-14.04^{\pm3.30}$ & 17.04 & 0.86 & 0.74 & 127.28 \\
 &GAM  &$73.74^{\pm0.88}$ & $78.32^{\pm0.94}$ & $-13.36^{\pm0.20}$ & 13.80 & 0.85 & 0.78 & 127.87 \\
\midrule
\multirow{2}{*}{\textbf{SeqLoRA}}
 & SGD     &$72.82^{\pm1.27}$ & $76.62^{\pm0.68}$ & $-6.75^{\pm1.29}$  & 16.58 & 0.73 & 0.77 & 125.50 \\
 & SAM     &$74.53^{\pm0.60}$ & $77.62^{\pm0.43}$ & $-5.54^{\pm0.72}$  & 14.82 & 0.73 & 0.81 & 135.38 \\
 &GAM   &$75.74^{\pm0.15}$ & $78.86^{\pm0.53}$ & $-6.19^{\pm0.33}$  & 11.28 & 0.78 & 0.78 & 133.96 \\
\midrule
\multirow{2}{*}{\textbf{IncLoRA}}
 & SGD      &$75.41^{\pm0.26}$ & $78.37^{\pm0.26}$ & $-4.32^{\pm0.59}$   & 16.87 & 0.69 & 0.67 & 121.18 \\
 & SAM     &$75.95^{\pm0.07}$ & $78.68^{\pm0.25}$ & $-3.87^{\pm0.03}$ & 15.14 & 0.68 & 0.70 & 127.16 \\
&GAM  &$76.72^{\pm0.54}$ & $79.42^{\pm0.40}$ & $-3.63^{\pm0.67}$  & 12.15 & 0.75 & 0.73 & 119.68 \\
\midrule
\multirow{2}{*}{\textbf{OLoRA}}
 & SGD     &$75.10^{\pm0.41}$ & $78.65^{\pm0.67}$ & $-3.90^{\pm0.54}$ & 17.23 & 0.69 & 0.68 & 120.89 \\
 & SAM     &$75.73^{\pm0.32}$ & $78.70^{\pm0.67}$ & $-3.47^{\pm0.54}$ & 15.17 & 0.72 & 0.70 & 123.76 \\
 &GAM  &$76.56^{\pm0.53}$ & $79.16^{\pm0.43}$ & $-3.37^{\pm0.68}$  & 12.16 & 0.75 & 0.74 & 119.36\\
\bottomrule
\end{tabular}
}
\end{table}

\begin{figure}[thbp]
    \centering
    \begin{minipage}{0.495\textwidth}
        \centering
        \includegraphics[width=0.9\linewidth]{figures_TRML/AAA_curve_Imagenet-R_summary_inc10_rank16—point2.pdf}
        % {figures_TRML/AAA_curve_Imagenet-R_summary_inc10_rank16_replot2.pdf}
        % \subcaption{Accuracy$\operatorname{ACC}_t^{\mathrm{cls}}$ on ImageNet-R}
        \label{fig:curves_imageR_1_AAA}
    \end{minipage}
    \begin{minipage}{0.495\textwidth}
        \centering
        \includegraphics[width=0.9\linewidth]{figures_TRML/NME_curve_Imagenet-R_summary_inc10_rank16_point2.pdf}
        % {figures_TRML/NME_curve_Imagenet-R_summary_inc10_rank16_replot2.pdf}
        % \subcaption{NME-based accuracy $\operatorname{ACC}_t^{\mathrm{nme}}$ on ImageNet-R}
        \label{fig:curves_imageR_2_NME}
    \end{minipage}
    \caption{
    % Summary curves of accuracy with the classifier head $\operatorname{ACC}_t^{\mathrm{cls}}$ and prototype-based accuracy $\operatorname{ACC}_t^{\mathrm{nme}}$ on ImageNet-R ($T=20, r=16$). 
    %     The curves report performance after each task for different class incremental methods and optimization configurations. 
        Summary curves of classifier-head accuracy $\operatorname{ACC}_t^{\mathrm{cls}}$ and prototype-based accuracy $\operatorname{ACC}_t^{\mathrm{nme}}$ on ImageNet-R ($T=20, r=16$) over the task sequence, comparing different class-incremental methods and optimization configurations.
        Flatness-promoting optimizers such as SAM and GAM consistently outperform SGD, and their advantage becomes more pronounced on later tasks, which suggests more stable representations and improved performance.
        % Flatness promotion optimizers (SAM and GAM) consistently outperform SGD, and their advantage tends to increase over later tasks, indicating more stable representations and better performance.
    }
    \label{fig:combined_curves_imageR}
\end{figure}




%%%%%%%%%%%%%%%%%%%%%%%%%%%%%%%%%%%%%%%%%%%%%%%%%%%%%%%%%%%%%%%%%%%%%%%%%%%%%%%
%%%%%%%%%%%%%%%%%%%%%%%%%%%%%%%%%%%%%%%%%%%%%%%%%%%%%%%%%%%%%%%%%%%%%%%%%%%%%%%



% \textbf{EFM analyse}

% \begin{figure}[htbp]
% \centering
% \begin{minipage}{0.9\linewidth}
% \centering
% \includegraphics[width=0.9\linewidth]{figures_TRML/finetune_42_gam-sam-sgd_efm_spectrum_compare.pdf}
% % \subcaption{Eigenvalue Spectrum in Sequential Fine-tuning (SeqFT)}
% \label{fig:EFM_seqFT}
% \end{minipage}
% \caption{
% Comparison of EFM eigenvalue spectra in the feature space across incremental learning methods. The horizontal axis represents eigenvalue rank (sorted by magnitude), while the vertical axis shows eigenvalue magnitude. Only the top 200 eigenvalues are non-zero, with remaining dimensions at zero. SAM consistently produces flatter, heavier-tailed spectra compared to SGD across both configurations, indicating reduced energy concentration in the leading principal directions. This spectral profile suggests that SAM learns more isotropic representations with higher effective dimensionality.
% }
% \label{fig:EFM_spectra_comparison}
% \end{figure}



\subsection{Detailed Resuts on ImageNet-C and ImageNet-P}


\begin{figure*}[tbp]
    \centering
    % ImageNet-C: cls / nme
    \begin{minipage}{0.495\linewidth}
        \centering
        \includegraphics[width=0.85\linewidth]{figures_TRML/AAA_curve_Imagenet-C_summary_inc10_rank16_ponit2.pdf}
        % \subcaption{ImageNet-C: $\operatorname{ACC}_t^{\mathrm{cls}}$}
        \label{fig:aaa_imagenet_c}
    \end{minipage}\hfill
    \begin{minipage}{0.495\linewidth}
        \centering
        \includegraphics[width=0.85\linewidth]{figures_TRML/NME_curve_Imagenet-C_summary_inc10_rank16_point2.pdf}
        % \subcaption{ImageNet-C: $\operatorname{ACC}_t^{\mathrm{nme}}$}
        \label{fig:nme_imagenet_c}
    \end{minipage}\hfill
    % ImageNet-P: cls / nme
    \begin{minipage}{0.495\linewidth}
        \centering
        \includegraphics[width=0.85\linewidth]{figures_TRML/AAA_curve_Imagenet-P_summary_inc10_rank16_point2.pdf}
        % \subcaption{ImageNet-P: $\operatorname{ACC}_t^{\mathrm{cls}}$}
        \label{fig:aaa_imagenet_p}
    \end{minipage}\hfill
    \begin{minipage}{0.495\linewidth}
        \centering
        \includegraphics[width=0.85\linewidth]{figures_TRML/NME_curve_Imagenet-P_summary_inc10_rank16_point2.pdf}
        % \subcaption{ImageNet-P: $\operatorname{ACC}_t^{\mathrm{nme}}$}
        \label{fig:nme_imagenet_p}
    \end{minipage}
    \caption{ 
        The full trajectories
of$\operatorname{ACC}_t^{\mathrm{nme}}$ on ImageNet-C and ImageNet-P. 
        % The four panels report performance after each task for the same class incremental methods and optimization configurations as in Figure~\ref{fig:combined_curves_imageR}. 
         SAM and GAM consistently outperform SGD on both benchmark, which indicates that promoting flat minima in the adapter improves robustness of the learned representations.
         %   and $\operatorname{ACC}_t^{\mathrm{cls}}$
    }
    % \label{fig:robustness_curves_comparison}
\end{figure*}

% \subsection{Ablation Study}


\subsection{Computational Cost Analysis}

\begin{table}[thpb]
\centering
\caption{Average per‑task training time (seconds) across methods and optimizers on ImageNet‑R}
\label{tab:lora_opt_metric}
\setlength{\tabcolsep}{4pt}
\renewcommand{\arraystretch}{1}
\resizebox{\linewidth}{!}{
\begin{tabular}{l ccccc}
\toprule
\textbf{Method} & \textbf{SGD} & \textbf{RWP} & \textbf{SAM} & \textbf{GAM} & \textbf{C-FLAT} \\
\midrule
SeqLoRA & 282.72 ± 61.34  & 244.95 ± 6.38  & 472.29 ± 83.70  & 670.20 ± 217.92 & 532.62 ± 35.93   \\
IncLoRA & 310.54 ± 141.00 & 658.34 ± 83.90 & 777.91 ± 239.29 & 1247.56 ± 5.05  & 1189.40 ± 109.93 \\
OLoRA   & 412.99 ± 146.84 & 644.19 ± 55.05 & 628.44 ± 181.29 & 836.99 ± 327.65 & 990.88 ± 360.79  \\
\bottomrule
\end{tabular}
}
\end{table}









\end{document}


% This document was modified from the file originally made available by
% Pat Langley and Andrea Danyluk for ICML-2K. This version was created
% by Iain Murray in 2018, and modified by Alexandre Bouchard in
% 2019 and 2021 and by Csaba Szepesvari, Gang Niu and Sivan Sabato in 2022.
% Modified again in 2023 and 2024 by Sivan Sabato and Jonathan Scarlett.
% Previous contributors include Dan Roy, Lise Getoor and Tobias
% Scheffer, which was slightly modified from the 2010 version by
% Thorsten Joachims & Johannes Fuernkranz, slightly modified from the
% 2009 version by Kiri Wagstaff and Sam Roweis's 2008 version, which is
% slightly modified from Prasad Tadepalli's 2007 version which is a
% lightly changed version of the previous year's version by Andrew
% Moore, which was in turn edited from those of Kristian Kersting and
% Codrina Lauth. Alex Smola contributed to the algorithmic style files.
